\documentclass[11pt]{article}

\usepackage[english]{babel}
\usepackage[utf8]{inputenc}
\usepackage{amsmath,amssymb}
\usepackage[margin=1in]{geometry}
\usepackage[affil-it]{authblk}
\usepackage{booktabs}
\usepackage{enumitem}
\usepackage{hyperref}
\hypersetup{
	colorlinks = true,
	linkcolor = blue,
	citecolor = blue,
	urlcolor = gray,
}

\newcommand{\bs}[1]{\boldsymbol{#1}}
\newcommand{\mbf}[1]{\mathbf{#1}}
\newcommand{\mc}[1]{\mathcal{#1}}
\newcommand{\dd}{\,\mathrm{d}}
\newcommand{\Div}{\operatorname{Div}}
\newcommand{\Grad}{\nabla_{\!X}}
\newcommand{\grad}{\nabla_{\!x}}
\newcommand{\tr}{\operatorname{tr}}
\newcommand{\refdomain}{\Omega_0^s}
\newcommand{\curdomain}{\Omega_t}
\newcommand{\test}[1]{\delta #1}
\newcommand{\rate}[1]{\dot{#1}}
\newcommand{\Dt}[2]{\frac{D_{#1} #2}{Dt}}


\title{Finite-element implementation and verification of a multicomponent reacting porous-media theory on the solid reference configuration}

\author[1,2,3]{John T. Foster}
\author[1]{Syed Talha Tirmizi}
\author[1]{Emmanuel Ikpesu}
\author[1,3]{Marc Hesse}

\affil[1]{Hildebrand Department of Petroleum and Geosystems Engineering, The University of Texas at Austin}
\affil[2]{Department of Aerospace Engineering and Engineering Mechanics, The University of Texas at Austin}
\affil[3]{Oden Institute for Computational Engineering and Sciences, The University of Texas at Austin}

\date{\today}

\begin{document}
\sloppy
\maketitle

\begin{abstract}
	This companion manuscript records the finite-element form, implementation
	interfaces, and verification plan for the multicomponent reacting porous-media
	theory developed in the companion theory paper.  The computational body is
	the reference configuration of the solid skeleton, and every residual is
	assembled per unit solid reference volume.  The immediate purpose is to define
	a closed first implementation problem and to separate it from the broader
	theory: residuals, constitutive material interfaces,
	automatic-differentiation boundaries, validation problems, and unresolved
	thermodynamic derivatives are identified before MOOSE kernels are written.
\end{abstract}

\section{Scope and manuscript contract}
\label{sec:scope}

This paper is the implementation and verification companion to the
multicomponent reacting porous-media theory paper.  Its source of truth is the
strong-form theory in that manuscript, especially the equations pulled back to
the solid skeleton \cite{foster2026reactingporous}.  The formulation inherits
the reacting-mixture setting of Drumheller and Bedford
\cite{drumheller1980thermomechanical,drumheller2000} and the finite-deformation
poroelastic structure that motivates the nonlinear Biot coefficient
\cite{foster2025}.  This document does not rederive the variational theory and
does not claim that any MOOSE object has already been implemented.  Instead, it
defines the finite-element residuals, material-property interfaces, validation
matrix, and unresolved thermodynamic derivatives needed before implementation in
MOOSE \cite{harbour2025moose}.

The computational body is the reference configuration of the solid skeleton,
\(\refdomain\).  The solid motion
\[
	\mbf x = \bs\chi_s(\mbf X,t), \qquad
	\mbf F = \Grad \bs\chi_s, \qquad
	J = \det \mbf F
\]
maps \(\refdomain\) to the current mixture domain \(\curdomain\).  All storage
terms, flux divergences, stress divergences, and source terms in this paper are
assembled per unit solid reference volume unless explicitly stated otherwise.
Spatial gradients appear only inside constitutive closures that are immediately
pulled back by \(\mbf F^{-T}\).

The first implementation target is a fully implicit finite-element solve on
\(\refdomain\) for displacement, phase/component storage variables, conversion
potentials or equivalent thermodynamic variables, phase partition variables, and
the Lagrange multipliers required by the constrained thermodynamics.  The exact
primitive variable set remains a design choice; the residual and algebraic
definitions below are written in terms of physical storage and thermodynamic
variables so that later agents can map them to MOOSE variables without changing
the mathematics or changing the reference measure of the residuals.

\subsection{Closure requirement and first closed problem}
\label{sec:closure_requirement}

For the purpose of implementation, the system is closed only when each residual
term in Section~\ref{sec:reference_solid_weak_forms} is supplied by either a
finite-element unknown, a boundary condition, or an automatic-differentiation
material property.  The general multicomponent theory is not closed by this
scaffold alone; the following list is the contract that a closed implementation
must satisfy:
\begin{itemize}[leftmargin=*]
	\item solid-skeleton kinematics: \(\mbf F\), \(J\), \(\mbf F^{-1}\), and
	the maps between spatial and reference gradients;
	\item storage variables: \(M_f^\alpha\) and \(M_s\), including the volume
	and mass-fraction constraints that define \(\phi_\xi\) and
	\(\eta_\xi^\alpha\);
	\item thermodynamic forces: \(\bar p_f\), \(\mu_\xi^\alpha\),
	\(\bar p_E\), and \(B\), with explicitly recorded held-fixed variables for
	each derivative;
	\item flux and stress closures: \(\mbf w_f\), \(\mbf W_f\),
	\(\widetilde{\bs\kappa}_f\), \(\widetilde{\mbf K}_f\),
	\(\mbf P^{\prime\prime}\), and the pressure stress contribution
	\(B\bar p_EJ\mbf F^{-T}\);
	\item conversion closures: \(\dot r_{(m)}\), \(\dot c_\xi^\alpha\), and
	any affinity or activity model used by the reaction network.
	\item phase kinematics for mobile phases: solid-frame advection of
	phase-attached scalar fields through \(\mbf W_f\).
\end{itemize}
If any item in this list is not provided, the corresponding kernel would encode
an implicit constitutive assumption and the implementation would not yet be a
closed discretization of the theory.

The first closed review problem is deliberately narrower than the full theory
but retains the full conversion structure needed for verification.  It is
isothermal, uses one mobile fluid phase and one solid skeleton phase, prescribes
a single reaction-rate law, and solves on \(\refdomain\) for displacement,
solid-reference component storage variables, phase partition variables, and the
conversion potentials required by the affinity law.  A thermodynamic material
maps those variables and \(\mbf F\) to
\(\phi_\xi\), \(\eta_\xi^\alpha\), \(\bar p_f\), \(\mu_\xi^\alpha\),
\(\bar p_E\), \(B\), and \(\dot c_\xi^\alpha\).  The first kernels therefore
include the component storage residuals, the solid storage residual, the
conversion-potential evolution equations, the reference-flux closure, and the
quasi-static skeleton momentum residual.  Additional
multifluid phase equilibrium, thermal effects, and plasticity are planned
extensions.

\subsection{Out of scope for this scaffold}
\label{sec:out_of_scope}

The following items are deliberately excluded from this initial scaffold:
\begin{itemize}[leftmargin=*]
	\item Implemented C++ kernels, materials, actions, input decks, or tests.
	\item A final choice between direct component-mass variables, pressure-composition variables, or reduced saturation-composition variables.
	\item Stabilization, upwinding, limiters, or finite-volume alternatives.
	\item Full benchmark data or production-scale reservoir examples.
\end{itemize}

Those decisions should be added only after the weak forms, thermodynamic
closures, and automatic-differentiation dependencies have been checked against
the theory manuscript.

\section{Algebraic and kinematic closure on the solid frame}
\label{sec:algebraic_kinematic_closure}

The weak forms in Section~\ref{sec:reference_solid_weak_forms} are not a
closed finite-element system unless the partition constraints, thermodynamic
multiplier relations, transfer-potential evolution laws, and phase kinematics
are solved or supplied as material equations.  This section records those
relations in the variables needed by a solid-reference implementation.

\subsection{Partition and storage equations}
\label{sec:partition_storage_equations}

For each phase \(\xi\) and component \(\alpha\), the component partial density,
phase partial density, intrinsic density, volume fraction, and mass fraction
are tied by
\begin{subequations}
	\label{eq:implementation_partition_relations}
	\begin{align}
		\rho_\xi^\alpha
		 & =\phi_\xi\bar\rho_\xi\eta_\xi^\alpha,
		\label{eq:implementation_component_density_partition} \\
		\rho_\xi
		 & =\phi_\xi\bar\rho_\xi,
		\label{eq:implementation_phase_density_partition} \\
		\sum_\xi\phi_\xi
		 & =1,
		\label{eq:implementation_volume_partition} \\
		\sum_\alpha\eta_\xi^\alpha
		 & =1 .
		\label{eq:implementation_mass_fraction_partition}
	\end{align}
\end{subequations}
The phase-attached material mass identity gives the pointwise storage relation
\begin{equation}
	\rho_{\xi0}^\alpha+\mc C_\xi^\alpha
	=
	J_\xi\phi_\xi\bar\rho_\xi\eta_\xi^\alpha .
	\label{eq:implementation_phase_material_storage}
\end{equation}
For a fluid phase \(f\), the solid-reference storage that enters the component
balance is \(J\phi_f\bar\rho_f\eta_f^\alpha\), where \(J\) is the solid
skeleton Jacobian.  Equation~\eqref{eq:implementation_phase_material_storage}
is still needed because the current source rate is
\begin{equation}
	\dot c_\xi^\alpha
	=
	\frac{\dot{\mc C}_\xi^\alpha}{J_\xi}
	=
	\sum_m\nu_{\xi (m)}^\alpha\dot r_{(m)} .
	\label{eq:implementation_conversion_rate_constraint}
\end{equation}
Thus \(J_\xi\) is part of the phase-attached definition of
\(\dot{\mc C}_\xi^\alpha\).  The solid-reference implementation keeps this
distinction explicit.  The input deck registers a finite phase set \(\mc P\)
and selects its reference solid phase \(s\in\mc P\).  The weak forms are pulled
back with that solid skeleton map, while every registered mobile phase
\(a\in\mc P\setminus\{s\}\) may carry phase-history deformation gradients and
Jacobians when a closure needs phase-attached material rates.

\subsection{Fluid material derivatives on the solid frame}
\label{sec:fluid_material_derivatives_solid_frame}

For each registered mobile phase \(a\in\mc P\setminus\{s\}\), the
solid-frame convective velocity is reconstructed from the same reference
relative mass flux that enters the component balance:
\begin{equation}
	\mbf v_a
	=
	\mbf v_s+\frac{1}{J_s\rho_a}\mbf F_s\mbf W_a,
	\qquad
	\mbf c_a
	=
	\mbf F_s^{-1}\left(\mbf v_a-\mbf v_s\right)
	=
	\frac{1}{J_s\rho_a}\mbf W_a .
	\label{eq:implementation_phase_velocity_reconstruction}
\end{equation}
Here \(\rho_a=\phi_a\bar\rho_a\) is the current bulk phase density per unit
current mixture volume; using the intrinsic density \(\bar\rho_a\) alone would
omit the phase volume fraction and give an incorrect material derivative.
Therefore any phase-attached scalar or tensor field \(q_a\) stored on the solid
mesh has
\begin{equation}
	\frac{D_a q_a}{Dt}
	=
	\left.\frac{\partial q_a}{\partial t}\right|_{\mbf X}
	+\mbf c_a\cdot\Grad q_a .
	\label{eq:fluid_material_derivative_on_solid_frame}
\end{equation}
The phase-history kinematics use each mobile phase velocity gradient in the
current configuration,
\begin{equation}
	\mbf l_a
	=
	\grad \mbf v_a
	=
	\Grad\mbf v_a\,\mbf F_s^{-1},
	\label{eq:implementation_phase_velocity_gradient}
\end{equation}
and evolve on the solid mesh according to
\begin{subequations}
	\label{eq:implementation_phase_history_kinematics}
	\begin{align}
		\left.\frac{\partial \mbf F_a}{\partial t}\right|_{\mbf X}
		+\mbf c_a\cdot\Grad\mbf F_a
		 & =
		\mbf l_a\mbf F_a,
		\label{eq:implementation_phase_deformation_gradient_history} \\
		\left.\frac{\partial J_a}{\partial t}\right|_{\mbf X}
		+\mbf c_a\cdot\Grad J_a
		 & =
		J_a\tr\mbf l_a .
		\label{eq:implementation_phase_jacobian_history}
	\end{align}
\end{subequations}
Here \(\mbf F_a\) and \(J_a\) are phase-history kinematics associated with the
phase labels, stored on the solid mesh.  They are distinct from
\(\mbf F_s\) and \(J_s\), which define the reference configuration, all
Grad/Div operators, boundary measures, and Piola transformations in the weak
forms.  Input decks must initialize a history-free phase with
\(\mbf F_a=\mbf I\) and \(J_a=1\), or provide another mutually consistent
history.  If a phase is inactive, the implementation sets both
\(\mbf c_a=\mbf 0\) and \(\mbf l_a=\mbf 0\), thereby preserving the stored
history until a documented closure reactivates that phase with consistent data.

\subsection{Thermodynamic multiplier and partition relations}
\label{sec:thermodynamic_partition_relations}

The general EOS input is a user-supplied Helmholtz free-energy density per unit
current phase volume,
\(\mc A_a=\mc A_a(\rho_a^1,\ldots,\rho_a^N,\theta,\ldots)\), where the
\(\rho_a^\alpha\) are constituent mass densities per unit current phase
volume.  The parsed material differentiates this single potential rather than
accepting pressure or chemical potentials as independent constitutive inputs:
\begin{subequations}
\begin{align}
 \mu_a^\alpha &=
 \left.\frac{\partial\mc A_a}{\partial\rho_a^\alpha}\right|_{\rho_a^{\beta\ne\alpha},\theta,\ldots},
 \label{eq:implementation_eos_chemical_potential}\\
 \bar p_a &= \sum_\alpha\rho_a^\alpha\mu_a^\alpha-\mc A_a,
 \label{eq:implementation_eos_pressure}\\
 s_a &= -\left.\frac{\partial\mc A_a}{\partial\theta}\right|_{\rho_a^1,\ldots,\rho_a^N,\ldots}.
 \label{eq:implementation_eos_entropy}
\end{align}
\end{subequations}
Second derivatives of \(\mc A_a\) provide the composition and thermal
thermodynamic Jacobian.  Thus pressure, chemical potentials, entropy, and their
Newton couplings remain mutually consistent for every admissible user
expression.  For a single density argument, this Euler relation reduces to
\begin{equation}
	\bar p_\xi
	=
	\bar\rho_\xi^2
	\left.\frac{\partial\psi_\xi}{\partial\bar\rho_\xi}\right|_{\mbf F_\xi,\eta_\xi^\beta,\theta}.
	\label{eq:implementation_phase_pressure}
\end{equation}
The equivalent pore pressure is
\begin{equation}
	\bar p_E
	=
	-\lambda
	=
	\sum_{f\in\mc F}S_f\bar p_f-\gamma .
	\label{eq:implementation_equivalent_pressure_relation}
\end{equation}
The individual capillary relation
\(\lambda+\bar p_f-\gamma_f=0\) and the pressure-difference relation
\(\bar p_f-\bar p_g=\gamma_f-\gamma_g\) are useful diagnostics, but they need
not be separate residual equations if the selected primitive variables enforce
\(\bar p_E=-\lambda\) and the capillary closure directly.
The mass-fraction Euler--Lagrange equations provide the composition partition
relation on the constrained tangent space:
\begin{align}
		\rho_\xi
		\left(
		\left.
		\frac{\partial\psi_\xi}{\partial\eta_\xi^\beta}
		\right|_{\mbf F_\xi,\bar\rho_\xi,\eta_\xi^{\kappa\neq\beta},\theta}
		-
		\left.
		\frac{\partial\psi_\xi}{\partial\eta_\xi^N}
		\right|_{\mbf F_\xi,\bar\rho_\xi,\eta_\xi^{\kappa\neq N},\theta}
		\right)
		 & =
		\frac{\pi_\xi^\beta}{\eta_\xi^\beta}
		-\frac{\pi_\xi^N}{\eta_\xi^N},
		\quad \beta=1,\ldots,N-1.
	\label{eq:implementation_pi_difference_system}
\end{align}
These algebraic equations are not optional postprocessing quantities.  They are
part of the nonlinear system that defines \(\phi_\xi\),
\(\eta_\xi^\alpha\), \(\pi_\xi^\alpha\), \(\lambda\), \(\bar p_f\), and
\(\bar p_E\) from the chosen primitive variables.  If all component mass
fractions are solved directly, the unprojected equation with the multiplier
\(\Pi_\xi\) may be used instead, but the projected system is the minimal
independent implementation form.

\subsection{Transfer-potential evolution on the solid frame}
\label{sec:conversion_potential_solid_frame}

Variation of the phase-attached conversion variables first gives a relation
containing the conversion-work potential \(L_\xi^\alpha\).  Subtracting its
reference-component form gives
\begin{equation}
	L_\xi^\alpha-L_\xi^N
	=
	\frac{\pi_\xi^\alpha}
	{\phi_\xi\bar\rho_\xi\eta_\xi^\alpha}
	-
	\frac{\pi_\xi^N}
	{\phi_\xi\bar\rho_\xi\eta_\xi^N}
	=
	\hat\mu_\xi^\alpha-\hat\mu_\xi^N,
	\qquad
	\alpha\ne N .
	\label{eq:implementation_transfer_work_component_difference}
\end{equation}
The neutral potentials are AD material outputs implementing the theory
manuscript's constituent-partial-density derivative
\(\hat\mu_\xi^\alpha
=\partial(\rho_\xi\psi_\xi+\phi_\xi\omega_\xi^+)/
\partial\rho_\xi^\alpha\) with the specified held-fixed fields.  They satisfy
the theory Euler identity
\(\hat\mu_\xi^\alpha
=\psi_\xi+\pi_\xi^\alpha/
(\phi_\xi\bar\rho_\xi\eta_\xi^\alpha)\); neither relation is an additional
finite-element residual.
Choose one skeleton-attached solid phase \(s_\star\) as the transfer reference
and impose
\(\psi_{s_\star}+L_{s_\star}^N=\hat\mu_{s_\star}^N\).
The only evolution equation for the retained transfer potential is then
\begin{align}
	\left.\frac{\partial\tau}{\partial t}\right|_{\mbf X}
	=
	\frac{1}{2}\mbf v_s\cdot\mbf v_s
	+\hat\mu_{s_\star}^N
	-\psi_{s_\star}
	-\frac{\pi_{s_\star}^N}
	{\phi_{s_\star}\bar\rho_{s_\star}\eta_{s_\star}^N}.
	\label{eq:implementation_tau_solid_frame}
\end{align}
Once \(\tau\) is known, the \(P-1\) non-reference phase levels are algebraic.
Using \eqref{eq:fluid_material_derivative_on_solid_frame}, a mobile fluid phase
\(f\ne s_\star\) has
\begin{align}
	L_f^N
	 & =
	\left.\frac{\partial\tau}{\partial t}\right|_{\mbf X}
	+\frac{\mbf W_f}{J\rho_f}\cdot\Grad\tau
	-\frac{1}{2}\mbf v_f\cdot\mbf v_f
	+\frac{\pi_f^N}
	{\phi_f\bar\rho_f\eta_f^N},
	\notag \\
	 & =
	\hat\mu_f^N-\psi_f
	+\left.\frac{\partial\tau}{\partial t}\right|_{\mbf X}
	+\frac{\mbf W_f}{J\rho_f}\cdot\Grad\tau
	-\frac{1}{2}\mbf v_f\cdot\mbf v_f,
	\label{eq:implementation_fluid_transfer_work_recovery}
\end{align}
while another skeleton-attached solid phase \(s\ne s_\star\) has
\begin{align}
	L_s^N
	 & =
	\left.\frac{\partial\tau}{\partial t}\right|_{\mbf X}
	-\frac{1}{2}\mbf v_s\cdot\mbf v_s
	+\frac{\pi_s^N}
	{\phi_s\bar\rho_s\eta_s^N},
	\notag \\
	 & =
	\hat\mu_s^N-\psi_s
	+\left.\frac{\partial\tau}{\partial t}\right|_{\mbf X}
	-\frac{1}{2}\mbf v_s\cdot\mbf v_s.
	\label{eq:implementation_solid_transfer_work_recovery}
\end{align}
These equalities are the direct algebraic closure
\(L_\xi^\alpha=\hat\mu_\xi^\alpha-\psi_\xi
+D_\xi\tau/Dt-\mbf v_\xi\cdot\mbf v_\xi/2\), written on the solid
reference frame.  Thus \(L_\xi^\alpha\) is recovered from the constitutive
potential \(\hat\mu_\xi^\alpha\), the retained field \(\tau\), and the phase
kinematics; it is not an independent constitutive function.

The traditional electrochemical affinity of mechanism \(m\) is the
stoichiometric projection
\begin{equation}
	\mc A_{(m)}
	=
	-\sum_\xi\sum_\alpha
	\mu_\xi^\alpha
	\nu_{\xi (m)}^\alpha .
	\label{eq:implementation_affinity_projection}
\end{equation}
The coefficient supplied by the variational conversion power is related to
this affinity by
\begin{equation}
	-\sum_\xi\sum_\alpha
	\left(
	\psi_\xi+L_\xi^\alpha+z_\xi^\alpha\varphi
	\right)
	\nu_{\xi (m)}^\alpha
	=
	\mc A_{(m)}
	-
	\sum_\xi
	\left(
	\frac{D_\xi\tau}{Dt}
	-\frac{1}{2}\mbf v_\xi\cdot\mbf v_\xi
	\right)
	\sum_\alpha\nu_{\xi (m)}^\alpha .
	\label{eq:implementation_generalized_conversion_affinity_relation}
\end{equation}
The correction vanishes for chemistry that conserves mass separately within
each phase and at full reversible equilibrium.  It must be retained for phase
change, dissolution, precipitation, and pore-filling mechanisms.
Equations~\eqref{eq:implementation_tau_solid_frame}--
\eqref{eq:implementation_generalized_conversion_affinity_relation} are required whenever the
implementation retains \(\tau\) or uses affinity-based kinetics.  For the first
closed problem, \(\tau\) is the single solved transfer field on the solid mesh;
the \(P-1\) levels \(L_\xi^N\) are algebraic variables or AD material
properties, and \eqref{eq:implementation_transfer_work_component_difference}
reconstructs the other component levels.  Their residuals are given in
Section~\ref{sec:conversion_potential_weak_form} so that the affinity and
the generalized conversion coefficient are differentiated consistently by
MOOSE AD.

\section{Weak forms on the solid reference configuration}
\label{sec:reference_solid_weak_forms}

Let \(\test{\mbf u}\), \(\test q_f^\alpha\), and \(\test q_s\) denote
finite-element test functions on \(\refdomain\).  Boundary subsets with
prescribed traction or flux are denoted by \(\Gamma_{0t}\) and
\(\Gamma_{0W}\).  The reference outward unit normal is \(\mbf N\).

\subsection{Fluid component balance}
\label{sec:fluid_component_weak_form}

For component \(\alpha\) in fluid phase \(f\), define the solid-reference
storage
\[
	M_f^\alpha = J\phi_f\bar\rho_f\eta_f^\alpha .
\]
The strong form on the solid trajectory is
\[
	\left.\frac{\partial M_f^\alpha}{\partial t}\right|_{\mbf X}
	+\Div\left(\eta_f^\alpha\mbf W_f\right)
	=
	J\sum_m \nu_{f(m)}^\alpha \dot r_{(m)} .
\]
The weak form uses the same reference measure; no current-volume source term is
inserted without the explicit pull-back factor \(J\).
The Galerkin residual is
\begin{align}
	R_f^\alpha
	=&
	\int_{\refdomain}
	\test q_f^\alpha
	\left[
		\left.\frac{\partial M_f^\alpha}{\partial t}\right|_{\mbf X}
		-
		J\sum_m \nu_{f(m)}^\alpha \dot r_{(m)}
	\right]\dd V
	\nonumber \\
	&-
	\int_{\refdomain}
	\Grad \test q_f^\alpha \cdot
	\left(\eta_f^\alpha\mbf W_f\right)\dd V
	+
	\int_{\Gamma_{0W}}
	\test q_f^\alpha
	\eta_f^\alpha\bar W_f
	\dd A .
	\label{eq:fluid_component_residual}
\end{align}
Here \(\bar W_f=\mbf W_f\cdot\mbf N\) is the prescribed outward reference mass
flux of the phase before multiplication by \(\eta_f^\alpha\).  The sign
convention makes outward flux positive.

\subsection{Solid component or solid phase balance}
\label{sec:solid_component_weak_form}

For a one-component solid phase \(s\), the storage variable is
\[
	M_s = J\phi_s\bar\rho_s .
\]
Because each solid phase is attached to the skeleton, there is no relative-flux
term:
\begin{equation}
	R_s =
	\int_{\refdomain}
	\test q_s
	\left[
		\left.\frac{\partial M_s}{\partial t}\right|_{\mbf X}
		-
		J\sum_m \nu_{s(m)}^s \dot r_{(m)}
	\right]\dd V .
	\label{eq:solid_component_residual}
\end{equation}
If several solid components share a solid phase, this residual must be promoted
to a component balance with the appropriate mass fractions and conversion
sources.

\subsection{Fluid phase kinematic residuals}
\label{sec:phase_kinematic_weak_forms}

No mobile-phase deformation-gradient or Jacobian residual is part of the first
solid-reference finite-element system.  Phase-attached scalar derivatives are
evaluated with \eqref{eq:fluid_material_derivative_on_solid_frame}, using the
same \(\mbf W_f\) that enters the component balance.  If a later closure
requires \(J_f\) or the full tensor \(\mbf F_f\), their evolution should be
introduced as auxiliary/history kinematics and tied back to the phase map at
that point.

\subsection{Reference flux closure}
\label{sec:reference_flux_closure}

The reference relative mass flux entering
\eqref{eq:fluid_component_residual} has the form
\begin{align}
	\mbf W_f
	= & \;
	\rho_f\widetilde{\mbf K}_f
	\left[
		\rho_f\mbf F^T\left(\mbf g-\mbf a_f\right)
		-\rho_f
		\left(
			\Grad\psi_f
			+
			\mathfrak{s}_f\Grad\theta_{\mc F}
		\right)
		-\phi_f\Grad
		\left(
			\bar p_E+\gamma_f
		\right)
		\right.
		\notag \\
	& \phantom{\;\rho_f\widetilde{\mbf K}_f}
		\left.
		+
		\mbf F^T
		\left[
			\grad\cdot
			\left(
				\phi_f\mbf E\otimes\mbf d_f
				\right)
		\right]
		\right.
		\notag \\
	& \phantom{\;\rho_f\widetilde{\mbf K}_f}
		\left.
		+
		\sum_\alpha\dot c_f^\alpha
		\left(
			\Grad \tau
			-\mbf F^T\mbf v_s
		\right)
	\right],
	\label{eq:reference_flux_closure}
\end{align}
where \(\rho_f=\phi_f\bar\rho_f\) and
\begin{equation}
	\widetilde{\mbf K}_f
	=
	J\mbf F^{-1}
	\left(
			\phi_f^2\mu_f^{\mathrm{visc}}\bs\kappa_f^{-1}
		+\sum_\alpha\dot c_f^\alpha\mbf I
	\right)^{-1}
	\mbf F^{-T}.
	\label{eq:reference_modified_permeability}
\end{equation}
This is the Piola transform of the current generalized Darcy mass flux.  The
tensor \(\widetilde{\mbf K}_f\) is a pulled-back
conversion-corrected mobility, not a permeability: its current-configuration
inverse contains the net conversion phase-mass source
\(\sum_\alpha\dot c_f^\alpha\).  The inverse exists only when none of the
eigenvalues
	\(\phi_f^2\mu_f^{\mathrm{visc}}/\kappa_{f,i}+\sum_\alpha\dot c_f^\alpha\) vanishes.
Positive definiteness is the stronger positive-resistance condition
\(\sum_\alpha\dot c_f^\alpha>
	-\phi_f^2\mu_f^{\mathrm{visc}}/\kappa_{f,\max}\); it is sufficient, but not necessary, for
algebraic elimination.  These source terms and the \(\Grad\tau\) insertion
force are part of the conversion-momentum structure and must not be dropped
when reactions or phase transformations produce appreciable mass in a mobile
phase.  The additional
\(-\rho_f(\Grad\psi_f+\mathfrak{s}_f\Grad\theta_{\mc F})\) term is the
reversible bulk-transport contribution required by the theory paper's complete
fluid--skeleton interaction closure; it is not part of the dissipative drag
tensor.

\subsection{Transfer-potential and transfer-work residuals}
\label{sec:conversion_potential_weak_form}

Let \(\test t\) be the scalar test function for the single transfer potential
\(\tau\).  The Galerkin residual associated with
\eqref{eq:implementation_tau_solid_frame} is
\begin{align}
	R_\tau
	=&
	\int_{\refdomain}
	\test t
	\left[
		\left.\frac{\partial\tau}{\partial t}\right|_{\mbf X}
		-\frac{1}{2}\mbf v_s\cdot\mbf v_s
		-\hat\mu_{s_\star}^N
		+\psi_{s_\star}
		+\frac{\pi_{s_\star}^N}
		{\phi_{s_\star}\bar\rho_{s_\star}\eta_{s_\star}^N}
	\right]\dd V .
	\label{eq:tau_evolution_residual}
\end{align}
For every non-reference phase \(\xi\ne s_\star\), let
\(\test\ell_\xi\) test its algebraic reference-component transfer-work level.
Using \eqref{eq:fluid_material_derivative_on_solid_frame} for a mobile phase,
the recovery residual is
\begin{align}
	R_{L_f^N}
	=&
	\int_{\refdomain}
	\test\ell_f
	\left[
		L_f^N
		-\left.\frac{\partial\tau}{\partial t}\right|_{\mbf X}
		-\frac{\mbf W_f}{J\rho_f}\cdot\Grad\tau
		+\frac{1}{2}\mbf v_f\cdot\mbf v_f
		-\frac{\pi_f^N}
		{\phi_f\bar\rho_f\eta_f^N}
	\right]\dd V .
	\label{eq:fluid_transfer_work_recovery_residual}
\end{align}
For a non-reference solid phase attached to the skeleton, the corresponding
algebraic residual is
\begin{align}
	R_{L_s^N}
	=&
	\int_{\refdomain}
	\test\ell_s
	\left[
		L_s^N
		-\left.\frac{\partial\tau}{\partial t}\right|_{\mbf X}
		+\frac{1}{2}\mbf v_s\cdot\mbf v_s
		-\frac{\pi_s^N}
		{\phi_s\bar\rho_s\eta_s^N}
	\right]\dd V .
	\label{eq:solid_transfer_work_recovery_residual}
\end{align}
The traditional mechanism affinity
\eqref{eq:implementation_affinity_projection} is then an algebraic material
property or constraint residual.  If chemical or electrochemical equilibrium
is enforced directly, the residual is
\begin{equation}
	R_{\mc A_{(m)}}=
	\int_{\refdomain}
	\test a_{(m)}
	\left[
		\mc A_{(m)}
		+\sum_\xi\sum_\alpha
		\mu_\xi^\alpha
		\nu_{\xi (m)}^\alpha
	\right]\dd V ,
	\label{eq:affinity_projection_residual}
\end{equation}
with \(\test a_{(m)}\) the test function for an affinity multiplier or local
constraint variable.  A kinetic phase-change or pore-filling mechanism is
instead driven by the full conversion coefficient in
\eqref{eq:implementation_generalized_conversion_affinity_relation}; replacing
that coefficient by \(\mc A_{(m)}\) is valid only when the phasewise
stoichiometric mass sums vanish or the transfer-work offsets have equilibrated.

\subsection{Overall momentum balance}
\label{sec:momentum_weak_form}

Let \(\mbf P^{\prime\prime}\) be the skeleton effective first
Piola--Kirchhoff stress and \(B\bar p_E J\mbf F^{-T}\) the equivalent pore
pressure contribution.  The quasi-static implementation target is
\begin{align}
	R_u
	=&
	\int_{\refdomain}
	\Grad \test{\mbf u} :
	\left(
		\mbf P^{\prime\prime}
		-
		B\bar p_E J\mbf F^{-T}
	\right)\dd V
	-
	\int_{\refdomain}
	\test{\mbf u}\cdot J\rho\mbf g\dd V
	\nonumber \\
	&-
	\int_{\refdomain}
	\test{\mbf u}\cdot J
	\left[
			\sum_f\sum_\alpha
			\dot c_f^\alpha
			\left(
				\grad\tau-\mbf v_f
			\right)
			+
			\sum_s
			\dot c_s^s
			\left(
				\grad\tau-\mbf v_s
			\right)
	\right]\dd V
	\nonumber \\
	&-
	\int_{\Gamma_{0t}}
	\test{\mbf u}\cdot \bar{\mbf t}_0\dd A .
	\label{eq:momentum_residual}
\end{align}
For low-speed reservoir applications the inertial term from the full momentum
balance is omitted.  If retained, the term
\(\int_{\refdomain}\test{\mbf u}\cdot J\sum_\xi\rho_\xi\mbf a_\xi\dd V\)
is added to \eqref{eq:momentum_residual} with the same sign convention as the
chosen time integrator.  In a fully reference-form implementation, the
transfer-potential gradient in the conversion-momentum terms is evaluated as
\(\grad\tau=\mbf F^{-T}\Grad\tau\) for \(\tau\) stored on the solid mesh, while
\(\mbf v_f\) is reconstructed from
\eqref{eq:implementation_fluid_velocity_reconstruction}.

\subsection{Constraint residuals}
\label{sec:constraint_residuals}

Volume fractions and phase/component mass fractions satisfy
\[
	\sum_\xi \phi_\xi = 1, \qquad
	\sum_\alpha \eta_\xi^\alpha = 1 .
\]
The implementation may enforce these constraints by eliminating dependent
variables, by solving pressure or potential multipliers, or by adding explicit
constraint residuals.  The selected strategy must preserve the thermodynamic
conjugacy used to define \(\bar p_f\), \(\bar p_E\), \(B\), and
\(\mu_\xi^\alpha\).  This is a design constraint, not a numerical preference.

\section{MOOSE implementation strategy}
\label{sec:moose_implementation_strategy}

The MOOSE implementation should be built from residual objects that map
directly to the weak forms in Section~\ref{sec:reference_solid_weak_forms}.  The
first implementation milestone is not broad feature coverage; it is one
validated residual path on the solid reference configuration.

\subsection{Proposed object responsibilities}
\label{sec:object_responsibilities}

\begin{table}[h]
	\centering
	\caption{Initial implementation interfaces.  Names are descriptive
	placeholders, not required C++ class names.}
	\label{tab:implementation_interfaces}
	\small
	\begin{tabular}{p{0.23\linewidth}p{0.34\linewidth}p{0.26\linewidth}}
		\toprule
		Object role & Computes & Weak-form target \\
		\midrule
		Solid kinematics material & \(\mbf F\), \(J\), \(\mbf F^{-1}\), \(\mbf v_s\) & all solid-reference residuals \\
		Phase advection material & \(\mbf v_f\) and \(D_f(\cdot)/Dt\) from \(\mbf W_f\) & Eq.~\eqref{eq:fluid_material_derivative_on_solid_frame} \\
		Thermodynamics material & \(\bar p_f\), \(\mu_\xi^\alpha\), \(\bar p_E\), \(B\), \(\pi_\xi^\alpha\), \(\Pi_\xi\), \(\lambda\) & flux, momentum, algebraic constraints \\
		Reaction material & \(\dot r_{(m)}\), \(\dot c_\xi^\alpha\), affinities & source and conversion terms \\
		Reference flux material & \(\mbf W_f\), \(\widetilde{\mbf K}_f\) & Eq.~\eqref{eq:fluid_component_residual} \\
		Stress material & \(\mbf P^{\prime\prime}\), pressure stress split & Eq.~\eqref{eq:momentum_residual} \\
		Component kernels & storage, flux, source residuals & Eq.~\eqref{eq:fluid_component_residual} \\
		Solid kernels & solid storage and conversion residuals & Eq.~\eqref{eq:solid_component_residual} \\
		Conversion-potential kernels & solid-frame evolution of \(\mu_\xi^\alpha\) & Eqs.~\eqref{eq:fluid_mu_residual}--\eqref{eq:solid_mu_residual} \\
		Constraint kernels/material solves & partition, pressure, and composition relations & Eqs.~\eqref{eq:implementation_partition_relations}--\eqref{eq:implementation_pi_difference_system} \\
		Momentum kernels & skeleton force balance & Eq.~\eqref{eq:momentum_residual} \\
		\bottomrule
	\end{tabular}
\end{table}

All objects should use automatic differentiation from the start.  Non-AD helper
properties may be acceptable only for constants or diagnostics that do not enter
Newton derivatives.

\subsection{Automatic-differentiation boundary}
\label{sec:ad_boundary}

The implementation should separate scalar thermodynamic differentiation from
tensorial kinematics and residual assembly.  The current MOOSE finding is that
\texttt{ADDerivativeParsedMaterial} can represent scalar Helmholtz-free-energy
closures and expose derivatives such as pressures, chemical potentials, and
selected second derivatives when the independent scalar arguments are chosen
explicitly.  That capability is useful for local thermodynamic derivatives, but
it does not remove the need for explicit AD C++ code for solid and fluid
deformation-gradient kinematics, Piola transformations, pulled-back modified
permeability, source closures, constraint residuals, conversion-potential
evolution, and the weak forms in
Section~\ref{sec:reference_solid_weak_forms}.

Consequently, parsed derivative materials should be treated as thermodynamic
building blocks rather than as the implementation architecture.  Kernels and
materials that compute \(\mbf F\), \(J\), \(\mbf F^{-1}\), \(\mbf W_f\),
\(\mbf v_f\), \(\mbf P^{\prime\prime}\),
\(B\bar p_EJ\mbf F^{-T}\), and the storage/source terms must retain full AD
dependence on the finite-element variables and their gradients.

This boundary gives a practical rule for the first code review: a parsed free
energy may define a scalar potential and its partial derivatives, but a MOOSE
kernel must still show exactly which weak-form term it contributes to.  If a
derivative is produced by a parsed material, the input file must document the
independent variables and the maximum derivative order requested.  If a
derivative is produced by compiled C++, the material header or source should
state the thermodynamic identity it implements.

\subsection{Primitive variables}
\label{sec:primitive_variables}

The scaffold leaves the primitive variables open.  The choice should be made by
checking which variables give a nonsingular thermodynamic Jacobian and a clean
mapping to standard reservoir limits.  Candidate sets include:
\begin{itemize}[leftmargin=*]
	\item displacement plus independent Lagrangian component masses;
	\item displacement, phase pressures, saturations, and component fractions with algebraic closures;
	\item displacement plus thermodynamic potentials and constraint multipliers.
\end{itemize}

Whichever set is chosen, materials must provide the physical quantities used in
the residuals: \(M_f^\alpha\), \(M_s\), \(\mbf W_f\),
\(\mbf P^{\prime\prime}\), \(\bar p_E\), \(B\), \(\mu_\xi^\alpha\),
\(\pi_\xi^\alpha\), \(\Pi_\xi\), \(\lambda\), and \(\dot r_{(m)}\).

The first reviewable implementation plan is closed only for the restricted
problem stated in Section~\ref{sec:closure_requirement}.  The broader
primitive-variable contract remains open; once chosen, it must be recorded as
the map from finite-element variables to the physical quantities appearing in
the residuals.

\subsection{Discretization defaults}
\label{sec:discretization_defaults}

The first finite-element discretization should use continuous Lagrange elements
with a mixed order sufficient to avoid locking or pressure oscillations in the
poroelastic limit.  A conservative starting point is quadratic displacement and
linear scalar variables, matching the Taylor--Hood pattern already used in the
three-phase implementation.  Stabilization should be introduced only after a
specific failure mode is identified and tied to a documented weak-form
approximation.

\subsection{Traceability rules}
\label{sec:traceability_rules}

Each eventual MOOSE object should cite the controlling equation in this paper
and the corresponding equation or section in the theory manuscript.  Each input
deck should state the reduction it exercises: compositional flow, black-oil
limit, poromechanics, reaction-only conversion, phase equilibrium, or a coupled
benchmark.  Generated exodus output is not validation by itself; every
verification case needs an observable, a reference result, and a tolerance.

\section{Validation matrix}
\label{sec:validation_matrix}

The validation suite should grow from reductions with analytic or trusted
reference solutions toward coupled reservoir-scale benchmarks.  Each entry must
exercise a defined subset of the closed interface in
Section~\ref{sec:closure_requirement}; otherwise a passing run would not
identify which part of the theory has been verified.  Table
\ref{tab:validation_matrix} is the initial planning matrix.

\begin{table}[ht]
	\centering
	\caption{Initial validation and verification plan.}
	\label{tab:validation_matrix}
	\small
	\begin{tabular}{p{0.18\linewidth}p{0.22\linewidth}p{0.24\linewidth}p{0.20\linewidth}}
		\toprule
		Case & Active residuals/interfaces & Reference result & Pass criterion \\
		\midrule
		Reference kinematics patch & solid kinematics material only & analytic \(\mbf F\), \(J\), and Piola tensor maps & pointwise tensor error below tolerance \\
		Fluid advection patch & Eq.~\eqref{eq:fluid_material_derivative_on_solid_frame} & prescribed \(\mbf W_f/(J\rho_f)\) and manufactured scalar field & advective derivative error below tolerance \\
		Thermodynamic derivative check & parsed or compiled free-energy material & analytic or finite-difference derivatives of \(\bar p_f\), \(\mu_\xi^\alpha\), \(\bar p_E\), and \(B\) & derivative error and AD Jacobian error below tolerance \\
		Algebraic partition check & Eqs.~\eqref{eq:implementation_partition_relations}--\eqref{eq:implementation_pi_difference_system} & manufactured phase split and pressures & constraint residuals and multiplier relations below tolerance \\
		Single-phase diffusion through fixed skeleton & Eq.~\eqref{eq:fluid_component_residual}, Eq.~\eqref{eq:reference_flux_closure} with no conversion & manufactured diffusion solution & mass error and convergence rate within tolerance \\
		Reaction-only closed box & Eq.~\eqref{eq:fluid_component_residual}, Eq.~\eqref{eq:solid_component_residual}, Eq.~\eqref{eq:implementation_conversion_rate_constraint}, no flux & analytic ordinary differential equation for Lagrangian masses & storage and volume-constraint errors below tolerance \\
		Retained transfer-potential block & Eqs.~\eqref{eq:tau_evolution_residual}--\eqref{eq:affinity_projection_residual} & manufactured \(\tau(t,\mbf X)\), \(L_\xi^N\), traditional affinity, and generalized conversion coefficient & evolution, algebraic-recovery, affinity, and transfer-work-correction errors below tolerance \\
		Drained reaction cell & component and solid residuals with pressure boundary & constant-pressure analytic reduction & pressure boundary maintained and storage history within tolerance \\
		Undrained reaction cell & component, solid, and momentum residuals with no flux & analytic or semi-analytic \(J(t)\) and pressure response & displacement, pressure, and mass errors below tolerance \\
		Terzaghi or Mandel limit & Eq.~\eqref{eq:momentum_residual} plus single-pressure storage reduction & classical poroelastic solution & pressure dissipation and displacement within tolerance \\
		Classical compositional-flow limit & rigid skeleton, component residuals, Darcy flux, thermodynamic material & trusted compositional simulator or manufactured solution & component inventories and pressure within tolerance \\
		Phase-equilibrium flash limit & thermodynamic material or user object only & independent flash calculation & phase split and chemical potentials within tolerance \\
		Coupled reaction--mechanics front & first closed problem plus spatial reaction source & benchmark or refined-solution comparison & front location, stress, and mass balance within tolerance \\
		\bottomrule
	\end{tabular}
\end{table}

\FloatBarrier

For each case, the repository should eventually contain the input deck,
postprocessing script, reference data or analytic calculation, and a short note
identifying which residual terms are active.  Expensive examples should have a
smaller regression analogue that protects the same implementation path.

\subsection{Current verification results}
\label{sec:current_verification_results}

The Mandel calculation exercises the implicit nonlinear Biot coefficient,
finite-deformation skeleton momentum balance, pressure storage, and Darcy flux
in the single-solid, single-fluid reduction.  Figure
\ref{fig:mandel_pressure_profiles} compares the finite-element water-pressure
profiles with the analytical Mandel series at five times.  The open circles are
the Q2/Q1 finite-element solution, and the solid curves are the analytical
profiles.

\begin{figure}[ht]
	\centering
	\resizebox{\textwidth}{!}{%
		%% Creator: Matplotlib, PGF backend
%%
%% To include the figure in your LaTeX document, write
%%   \input{<filename>.pgf}
%%
%% Make sure the required packages are loaded in your preamble
%%   \usepackage{pgf}
%%
%% Also ensure that all the required font packages are loaded; for instance,
%% the lmodern package is sometimes necessary when using math font.
%%   \usepackage{lmodern}
%%
%% Figures using additional raster images can only be included by \input if
%% they are in the same directory as the main LaTeX file. For loading figures
%% from other directories you can use the `import` package
%%   \usepackage{import}
%%
%% and then include the figures with
%%   \import{<path to file>}{<filename>.pgf}
%%
%% Matplotlib used the following preamble
%%   \def\mathdefault#1{#1}
%%   \everymath=\expandafter{\the\everymath\displaystyle}
%%   \IfFileExists{scrextend.sty}{
%%     \usepackage[fontsize=10.000000pt]{scrextend}
%%   }{
%%     \renewcommand{\normalsize}{\fontsize{10.000000}{12.000000}\selectfont}
%%     \normalsize
%%   }
%%
%%   \ifdefined\pdftexversion\else  % non-pdftex case.
%%     \usepackage{fontspec}
%%     \setmainfont{DejaVuSerif.ttf}[Path=\detokenize{/home/john/miniconda3/envs/moose/lib/python3.14/site-packages/matplotlib/mpl-data/fonts/ttf/}]
%%     \setsansfont{DejaVuSans.ttf}[Path=\detokenize{/home/john/miniconda3/envs/moose/lib/python3.14/site-packages/matplotlib/mpl-data/fonts/ttf/}]
%%     \setmonofont{DejaVuSansMono.ttf}[Path=\detokenize{/home/john/miniconda3/envs/moose/lib/python3.14/site-packages/matplotlib/mpl-data/fonts/ttf/}]
%%   \fi
%%   \makeatletter\@ifpackageloaded{underscore}{}{\usepackage[strings]{underscore}}\makeatother
%%
\begingroup%
\makeatletter%
\begin{pgfpicture}%
\pgfpathrectangle{\pgfpointorigin}{\pgfqpoint{7.400000in}{4.800000in}}%
\pgfusepath{use as bounding box, clip}%
\begin{pgfscope}%
\pgfsetbuttcap%
\pgfsetmiterjoin%
\definecolor{currentfill}{rgb}{1.000000,1.000000,1.000000}%
\pgfsetfillcolor{currentfill}%
\pgfsetlinewidth{0.000000pt}%
\definecolor{currentstroke}{rgb}{1.000000,1.000000,1.000000}%
\pgfsetstrokecolor{currentstroke}%
\pgfsetdash{}{0pt}%
\pgfpathmoveto{\pgfqpoint{0.000000in}{0.000000in}}%
\pgfpathlineto{\pgfqpoint{7.400000in}{0.000000in}}%
\pgfpathlineto{\pgfqpoint{7.400000in}{4.800000in}}%
\pgfpathlineto{\pgfqpoint{0.000000in}{4.800000in}}%
\pgfpathlineto{\pgfqpoint{0.000000in}{0.000000in}}%
\pgfpathclose%
\pgfusepath{fill}%
\end{pgfscope}%
\begin{pgfscope}%
\pgfsetbuttcap%
\pgfsetmiterjoin%
\definecolor{currentfill}{rgb}{1.000000,1.000000,1.000000}%
\pgfsetfillcolor{currentfill}%
\pgfsetlinewidth{0.000000pt}%
\definecolor{currentstroke}{rgb}{0.000000,0.000000,0.000000}%
\pgfsetstrokecolor{currentstroke}%
\pgfsetstrokeopacity{0.000000}%
\pgfsetdash{}{0pt}%
\pgfpathmoveto{\pgfqpoint{0.510067in}{0.472226in}}%
\pgfpathlineto{\pgfqpoint{7.247890in}{0.472226in}}%
\pgfpathlineto{\pgfqpoint{7.247890in}{4.758330in}}%
\pgfpathlineto{\pgfqpoint{0.510067in}{4.758330in}}%
\pgfpathlineto{\pgfqpoint{0.510067in}{0.472226in}}%
\pgfpathclose%
\pgfusepath{fill}%
\end{pgfscope}%
\begin{pgfscope}%
\pgfpathrectangle{\pgfqpoint{0.510067in}{0.472226in}}{\pgfqpoint{6.737823in}{4.286104in}}%
\pgfusepath{clip}%
\pgfsetrectcap%
\pgfsetroundjoin%
\pgfsetlinewidth{0.803000pt}%
\definecolor{currentstroke}{rgb}{0.690196,0.690196,0.690196}%
\pgfsetstrokecolor{currentstroke}%
\pgfsetstrokeopacity{0.220000}%
\pgfsetdash{}{0pt}%
\pgfpathmoveto{\pgfqpoint{0.510067in}{0.472226in}}%
\pgfpathlineto{\pgfqpoint{0.510067in}{4.758330in}}%
\pgfusepath{stroke}%
\end{pgfscope}%
\begin{pgfscope}%
\pgfsetbuttcap%
\pgfsetroundjoin%
\definecolor{currentfill}{rgb}{0.000000,0.000000,0.000000}%
\pgfsetfillcolor{currentfill}%
\pgfsetlinewidth{0.803000pt}%
\definecolor{currentstroke}{rgb}{0.000000,0.000000,0.000000}%
\pgfsetstrokecolor{currentstroke}%
\pgfsetdash{}{0pt}%
\pgfsys@defobject{currentmarker}{\pgfqpoint{0.000000in}{-0.048611in}}{\pgfqpoint{0.000000in}{0.000000in}}{%
\pgfpathmoveto{\pgfqpoint{0.000000in}{0.000000in}}%
\pgfpathlineto{\pgfqpoint{0.000000in}{-0.048611in}}%
\pgfusepath{stroke,fill}%
}%
\begin{pgfscope}%
\pgfsys@transformshift{0.510067in}{0.472226in}%
\pgfsys@useobject{currentmarker}{}%
\end{pgfscope}%
\end{pgfscope}%
\begin{pgfscope}%
\definecolor{textcolor}{rgb}{0.000000,0.000000,0.000000}%
\pgfsetstrokecolor{textcolor}%
\pgfsetfillcolor{textcolor}%
\pgftext[x=0.510067in,y=0.375003in,,top]{\color{textcolor}{\sffamily\fontsize{10.000000}{12.000000}\selectfont\catcode`\^=\active\def^{\ifmmode\sp\else\^{}\fi}\catcode`\%=\active\def%{\%}0.0}}%
\end{pgfscope}%
\begin{pgfscope}%
\pgfpathrectangle{\pgfqpoint{0.510067in}{0.472226in}}{\pgfqpoint{6.737823in}{4.286104in}}%
\pgfusepath{clip}%
\pgfsetrectcap%
\pgfsetroundjoin%
\pgfsetlinewidth{0.803000pt}%
\definecolor{currentstroke}{rgb}{0.690196,0.690196,0.690196}%
\pgfsetstrokecolor{currentstroke}%
\pgfsetstrokeopacity{0.220000}%
\pgfsetdash{}{0pt}%
\pgfpathmoveto{\pgfqpoint{1.857632in}{0.472226in}}%
\pgfpathlineto{\pgfqpoint{1.857632in}{4.758330in}}%
\pgfusepath{stroke}%
\end{pgfscope}%
\begin{pgfscope}%
\pgfsetbuttcap%
\pgfsetroundjoin%
\definecolor{currentfill}{rgb}{0.000000,0.000000,0.000000}%
\pgfsetfillcolor{currentfill}%
\pgfsetlinewidth{0.803000pt}%
\definecolor{currentstroke}{rgb}{0.000000,0.000000,0.000000}%
\pgfsetstrokecolor{currentstroke}%
\pgfsetdash{}{0pt}%
\pgfsys@defobject{currentmarker}{\pgfqpoint{0.000000in}{-0.048611in}}{\pgfqpoint{0.000000in}{0.000000in}}{%
\pgfpathmoveto{\pgfqpoint{0.000000in}{0.000000in}}%
\pgfpathlineto{\pgfqpoint{0.000000in}{-0.048611in}}%
\pgfusepath{stroke,fill}%
}%
\begin{pgfscope}%
\pgfsys@transformshift{1.857632in}{0.472226in}%
\pgfsys@useobject{currentmarker}{}%
\end{pgfscope}%
\end{pgfscope}%
\begin{pgfscope}%
\definecolor{textcolor}{rgb}{0.000000,0.000000,0.000000}%
\pgfsetstrokecolor{textcolor}%
\pgfsetfillcolor{textcolor}%
\pgftext[x=1.857632in,y=0.375003in,,top]{\color{textcolor}{\sffamily\fontsize{10.000000}{12.000000}\selectfont\catcode`\^=\active\def^{\ifmmode\sp\else\^{}\fi}\catcode`\%=\active\def%{\%}0.2}}%
\end{pgfscope}%
\begin{pgfscope}%
\pgfpathrectangle{\pgfqpoint{0.510067in}{0.472226in}}{\pgfqpoint{6.737823in}{4.286104in}}%
\pgfusepath{clip}%
\pgfsetrectcap%
\pgfsetroundjoin%
\pgfsetlinewidth{0.803000pt}%
\definecolor{currentstroke}{rgb}{0.690196,0.690196,0.690196}%
\pgfsetstrokecolor{currentstroke}%
\pgfsetstrokeopacity{0.220000}%
\pgfsetdash{}{0pt}%
\pgfpathmoveto{\pgfqpoint{3.205197in}{0.472226in}}%
\pgfpathlineto{\pgfqpoint{3.205197in}{4.758330in}}%
\pgfusepath{stroke}%
\end{pgfscope}%
\begin{pgfscope}%
\pgfsetbuttcap%
\pgfsetroundjoin%
\definecolor{currentfill}{rgb}{0.000000,0.000000,0.000000}%
\pgfsetfillcolor{currentfill}%
\pgfsetlinewidth{0.803000pt}%
\definecolor{currentstroke}{rgb}{0.000000,0.000000,0.000000}%
\pgfsetstrokecolor{currentstroke}%
\pgfsetdash{}{0pt}%
\pgfsys@defobject{currentmarker}{\pgfqpoint{0.000000in}{-0.048611in}}{\pgfqpoint{0.000000in}{0.000000in}}{%
\pgfpathmoveto{\pgfqpoint{0.000000in}{0.000000in}}%
\pgfpathlineto{\pgfqpoint{0.000000in}{-0.048611in}}%
\pgfusepath{stroke,fill}%
}%
\begin{pgfscope}%
\pgfsys@transformshift{3.205197in}{0.472226in}%
\pgfsys@useobject{currentmarker}{}%
\end{pgfscope}%
\end{pgfscope}%
\begin{pgfscope}%
\definecolor{textcolor}{rgb}{0.000000,0.000000,0.000000}%
\pgfsetstrokecolor{textcolor}%
\pgfsetfillcolor{textcolor}%
\pgftext[x=3.205197in,y=0.375003in,,top]{\color{textcolor}{\sffamily\fontsize{10.000000}{12.000000}\selectfont\catcode`\^=\active\def^{\ifmmode\sp\else\^{}\fi}\catcode`\%=\active\def%{\%}0.4}}%
\end{pgfscope}%
\begin{pgfscope}%
\pgfpathrectangle{\pgfqpoint{0.510067in}{0.472226in}}{\pgfqpoint{6.737823in}{4.286104in}}%
\pgfusepath{clip}%
\pgfsetrectcap%
\pgfsetroundjoin%
\pgfsetlinewidth{0.803000pt}%
\definecolor{currentstroke}{rgb}{0.690196,0.690196,0.690196}%
\pgfsetstrokecolor{currentstroke}%
\pgfsetstrokeopacity{0.220000}%
\pgfsetdash{}{0pt}%
\pgfpathmoveto{\pgfqpoint{4.552761in}{0.472226in}}%
\pgfpathlineto{\pgfqpoint{4.552761in}{4.758330in}}%
\pgfusepath{stroke}%
\end{pgfscope}%
\begin{pgfscope}%
\pgfsetbuttcap%
\pgfsetroundjoin%
\definecolor{currentfill}{rgb}{0.000000,0.000000,0.000000}%
\pgfsetfillcolor{currentfill}%
\pgfsetlinewidth{0.803000pt}%
\definecolor{currentstroke}{rgb}{0.000000,0.000000,0.000000}%
\pgfsetstrokecolor{currentstroke}%
\pgfsetdash{}{0pt}%
\pgfsys@defobject{currentmarker}{\pgfqpoint{0.000000in}{-0.048611in}}{\pgfqpoint{0.000000in}{0.000000in}}{%
\pgfpathmoveto{\pgfqpoint{0.000000in}{0.000000in}}%
\pgfpathlineto{\pgfqpoint{0.000000in}{-0.048611in}}%
\pgfusepath{stroke,fill}%
}%
\begin{pgfscope}%
\pgfsys@transformshift{4.552761in}{0.472226in}%
\pgfsys@useobject{currentmarker}{}%
\end{pgfscope}%
\end{pgfscope}%
\begin{pgfscope}%
\definecolor{textcolor}{rgb}{0.000000,0.000000,0.000000}%
\pgfsetstrokecolor{textcolor}%
\pgfsetfillcolor{textcolor}%
\pgftext[x=4.552761in,y=0.375003in,,top]{\color{textcolor}{\sffamily\fontsize{10.000000}{12.000000}\selectfont\catcode`\^=\active\def^{\ifmmode\sp\else\^{}\fi}\catcode`\%=\active\def%{\%}0.6}}%
\end{pgfscope}%
\begin{pgfscope}%
\pgfpathrectangle{\pgfqpoint{0.510067in}{0.472226in}}{\pgfqpoint{6.737823in}{4.286104in}}%
\pgfusepath{clip}%
\pgfsetrectcap%
\pgfsetroundjoin%
\pgfsetlinewidth{0.803000pt}%
\definecolor{currentstroke}{rgb}{0.690196,0.690196,0.690196}%
\pgfsetstrokecolor{currentstroke}%
\pgfsetstrokeopacity{0.220000}%
\pgfsetdash{}{0pt}%
\pgfpathmoveto{\pgfqpoint{5.900326in}{0.472226in}}%
\pgfpathlineto{\pgfqpoint{5.900326in}{4.758330in}}%
\pgfusepath{stroke}%
\end{pgfscope}%
\begin{pgfscope}%
\pgfsetbuttcap%
\pgfsetroundjoin%
\definecolor{currentfill}{rgb}{0.000000,0.000000,0.000000}%
\pgfsetfillcolor{currentfill}%
\pgfsetlinewidth{0.803000pt}%
\definecolor{currentstroke}{rgb}{0.000000,0.000000,0.000000}%
\pgfsetstrokecolor{currentstroke}%
\pgfsetdash{}{0pt}%
\pgfsys@defobject{currentmarker}{\pgfqpoint{0.000000in}{-0.048611in}}{\pgfqpoint{0.000000in}{0.000000in}}{%
\pgfpathmoveto{\pgfqpoint{0.000000in}{0.000000in}}%
\pgfpathlineto{\pgfqpoint{0.000000in}{-0.048611in}}%
\pgfusepath{stroke,fill}%
}%
\begin{pgfscope}%
\pgfsys@transformshift{5.900326in}{0.472226in}%
\pgfsys@useobject{currentmarker}{}%
\end{pgfscope}%
\end{pgfscope}%
\begin{pgfscope}%
\definecolor{textcolor}{rgb}{0.000000,0.000000,0.000000}%
\pgfsetstrokecolor{textcolor}%
\pgfsetfillcolor{textcolor}%
\pgftext[x=5.900326in,y=0.375003in,,top]{\color{textcolor}{\sffamily\fontsize{10.000000}{12.000000}\selectfont\catcode`\^=\active\def^{\ifmmode\sp\else\^{}\fi}\catcode`\%=\active\def%{\%}0.8}}%
\end{pgfscope}%
\begin{pgfscope}%
\pgfpathrectangle{\pgfqpoint{0.510067in}{0.472226in}}{\pgfqpoint{6.737823in}{4.286104in}}%
\pgfusepath{clip}%
\pgfsetrectcap%
\pgfsetroundjoin%
\pgfsetlinewidth{0.803000pt}%
\definecolor{currentstroke}{rgb}{0.690196,0.690196,0.690196}%
\pgfsetstrokecolor{currentstroke}%
\pgfsetstrokeopacity{0.220000}%
\pgfsetdash{}{0pt}%
\pgfpathmoveto{\pgfqpoint{7.247890in}{0.472226in}}%
\pgfpathlineto{\pgfqpoint{7.247890in}{4.758330in}}%
\pgfusepath{stroke}%
\end{pgfscope}%
\begin{pgfscope}%
\pgfsetbuttcap%
\pgfsetroundjoin%
\definecolor{currentfill}{rgb}{0.000000,0.000000,0.000000}%
\pgfsetfillcolor{currentfill}%
\pgfsetlinewidth{0.803000pt}%
\definecolor{currentstroke}{rgb}{0.000000,0.000000,0.000000}%
\pgfsetstrokecolor{currentstroke}%
\pgfsetdash{}{0pt}%
\pgfsys@defobject{currentmarker}{\pgfqpoint{0.000000in}{-0.048611in}}{\pgfqpoint{0.000000in}{0.000000in}}{%
\pgfpathmoveto{\pgfqpoint{0.000000in}{0.000000in}}%
\pgfpathlineto{\pgfqpoint{0.000000in}{-0.048611in}}%
\pgfusepath{stroke,fill}%
}%
\begin{pgfscope}%
\pgfsys@transformshift{7.247890in}{0.472226in}%
\pgfsys@useobject{currentmarker}{}%
\end{pgfscope}%
\end{pgfscope}%
\begin{pgfscope}%
\definecolor{textcolor}{rgb}{0.000000,0.000000,0.000000}%
\pgfsetstrokecolor{textcolor}%
\pgfsetfillcolor{textcolor}%
\pgftext[x=7.247890in,y=0.375003in,,top]{\color{textcolor}{\sffamily\fontsize{10.000000}{12.000000}\selectfont\catcode`\^=\active\def^{\ifmmode\sp\else\^{}\fi}\catcode`\%=\active\def%{\%}1.0}}%
\end{pgfscope}%
\begin{pgfscope}%
\definecolor{textcolor}{rgb}{0.000000,0.000000,0.000000}%
\pgfsetstrokecolor{textcolor}%
\pgfsetfillcolor{textcolor}%
\pgftext[x=3.878979in,y=0.180559in,,top]{\color{textcolor}{\sffamily\fontsize{10.000000}{12.000000}\selectfont\catcode`\^=\active\def^{\ifmmode\sp\else\^{}\fi}\catcode`\%=\active\def%{\%}distance from the undrained center, $x$ [m]}}%
\end{pgfscope}%
\begin{pgfscope}%
\pgfpathrectangle{\pgfqpoint{0.510067in}{0.472226in}}{\pgfqpoint{6.737823in}{4.286104in}}%
\pgfusepath{clip}%
\pgfsetrectcap%
\pgfsetroundjoin%
\pgfsetlinewidth{0.803000pt}%
\definecolor{currentstroke}{rgb}{0.690196,0.690196,0.690196}%
\pgfsetstrokecolor{currentstroke}%
\pgfsetstrokeopacity{0.220000}%
\pgfsetdash{}{0pt}%
\pgfpathmoveto{\pgfqpoint{0.510067in}{0.472226in}}%
\pgfpathlineto{\pgfqpoint{7.247890in}{0.472226in}}%
\pgfusepath{stroke}%
\end{pgfscope}%
\begin{pgfscope}%
\pgfsetbuttcap%
\pgfsetroundjoin%
\definecolor{currentfill}{rgb}{0.000000,0.000000,0.000000}%
\pgfsetfillcolor{currentfill}%
\pgfsetlinewidth{0.803000pt}%
\definecolor{currentstroke}{rgb}{0.000000,0.000000,0.000000}%
\pgfsetstrokecolor{currentstroke}%
\pgfsetdash{}{0pt}%
\pgfsys@defobject{currentmarker}{\pgfqpoint{-0.048611in}{0.000000in}}{\pgfqpoint{-0.000000in}{0.000000in}}{%
\pgfpathmoveto{\pgfqpoint{-0.000000in}{0.000000in}}%
\pgfpathlineto{\pgfqpoint{-0.048611in}{0.000000in}}%
\pgfusepath{stroke,fill}%
}%
\begin{pgfscope}%
\pgfsys@transformshift{0.510067in}{0.472226in}%
\pgfsys@useobject{currentmarker}{}%
\end{pgfscope}%
\end{pgfscope}%
\begin{pgfscope}%
\definecolor{textcolor}{rgb}{0.000000,0.000000,0.000000}%
\pgfsetstrokecolor{textcolor}%
\pgfsetfillcolor{textcolor}%
\pgftext[x=0.324480in, y=0.419464in, left, base]{\color{textcolor}{\sffamily\fontsize{10.000000}{12.000000}\selectfont\catcode`\^=\active\def^{\ifmmode\sp\else\^{}\fi}\catcode`\%=\active\def%{\%}0}}%
\end{pgfscope}%
\begin{pgfscope}%
\pgfpathrectangle{\pgfqpoint{0.510067in}{0.472226in}}{\pgfqpoint{6.737823in}{4.286104in}}%
\pgfusepath{clip}%
\pgfsetrectcap%
\pgfsetroundjoin%
\pgfsetlinewidth{0.803000pt}%
\definecolor{currentstroke}{rgb}{0.690196,0.690196,0.690196}%
\pgfsetstrokecolor{currentstroke}%
\pgfsetstrokeopacity{0.220000}%
\pgfsetdash{}{0pt}%
\pgfpathmoveto{\pgfqpoint{0.510067in}{1.279614in}}%
\pgfpathlineto{\pgfqpoint{7.247890in}{1.279614in}}%
\pgfusepath{stroke}%
\end{pgfscope}%
\begin{pgfscope}%
\pgfsetbuttcap%
\pgfsetroundjoin%
\definecolor{currentfill}{rgb}{0.000000,0.000000,0.000000}%
\pgfsetfillcolor{currentfill}%
\pgfsetlinewidth{0.803000pt}%
\definecolor{currentstroke}{rgb}{0.000000,0.000000,0.000000}%
\pgfsetstrokecolor{currentstroke}%
\pgfsetdash{}{0pt}%
\pgfsys@defobject{currentmarker}{\pgfqpoint{-0.048611in}{0.000000in}}{\pgfqpoint{-0.000000in}{0.000000in}}{%
\pgfpathmoveto{\pgfqpoint{-0.000000in}{0.000000in}}%
\pgfpathlineto{\pgfqpoint{-0.048611in}{0.000000in}}%
\pgfusepath{stroke,fill}%
}%
\begin{pgfscope}%
\pgfsys@transformshift{0.510067in}{1.279614in}%
\pgfsys@useobject{currentmarker}{}%
\end{pgfscope}%
\end{pgfscope}%
\begin{pgfscope}%
\definecolor{textcolor}{rgb}{0.000000,0.000000,0.000000}%
\pgfsetstrokecolor{textcolor}%
\pgfsetfillcolor{textcolor}%
\pgftext[x=0.236114in, y=1.226853in, left, base]{\color{textcolor}{\sffamily\fontsize{10.000000}{12.000000}\selectfont\catcode`\^=\active\def^{\ifmmode\sp\else\^{}\fi}\catcode`\%=\active\def%{\%}10}}%
\end{pgfscope}%
\begin{pgfscope}%
\pgfpathrectangle{\pgfqpoint{0.510067in}{0.472226in}}{\pgfqpoint{6.737823in}{4.286104in}}%
\pgfusepath{clip}%
\pgfsetrectcap%
\pgfsetroundjoin%
\pgfsetlinewidth{0.803000pt}%
\definecolor{currentstroke}{rgb}{0.690196,0.690196,0.690196}%
\pgfsetstrokecolor{currentstroke}%
\pgfsetstrokeopacity{0.220000}%
\pgfsetdash{}{0pt}%
\pgfpathmoveto{\pgfqpoint{0.510067in}{2.087003in}}%
\pgfpathlineto{\pgfqpoint{7.247890in}{2.087003in}}%
\pgfusepath{stroke}%
\end{pgfscope}%
\begin{pgfscope}%
\pgfsetbuttcap%
\pgfsetroundjoin%
\definecolor{currentfill}{rgb}{0.000000,0.000000,0.000000}%
\pgfsetfillcolor{currentfill}%
\pgfsetlinewidth{0.803000pt}%
\definecolor{currentstroke}{rgb}{0.000000,0.000000,0.000000}%
\pgfsetstrokecolor{currentstroke}%
\pgfsetdash{}{0pt}%
\pgfsys@defobject{currentmarker}{\pgfqpoint{-0.048611in}{0.000000in}}{\pgfqpoint{-0.000000in}{0.000000in}}{%
\pgfpathmoveto{\pgfqpoint{-0.000000in}{0.000000in}}%
\pgfpathlineto{\pgfqpoint{-0.048611in}{0.000000in}}%
\pgfusepath{stroke,fill}%
}%
\begin{pgfscope}%
\pgfsys@transformshift{0.510067in}{2.087003in}%
\pgfsys@useobject{currentmarker}{}%
\end{pgfscope}%
\end{pgfscope}%
\begin{pgfscope}%
\definecolor{textcolor}{rgb}{0.000000,0.000000,0.000000}%
\pgfsetstrokecolor{textcolor}%
\pgfsetfillcolor{textcolor}%
\pgftext[x=0.236114in, y=2.034241in, left, base]{\color{textcolor}{\sffamily\fontsize{10.000000}{12.000000}\selectfont\catcode`\^=\active\def^{\ifmmode\sp\else\^{}\fi}\catcode`\%=\active\def%{\%}20}}%
\end{pgfscope}%
\begin{pgfscope}%
\pgfpathrectangle{\pgfqpoint{0.510067in}{0.472226in}}{\pgfqpoint{6.737823in}{4.286104in}}%
\pgfusepath{clip}%
\pgfsetrectcap%
\pgfsetroundjoin%
\pgfsetlinewidth{0.803000pt}%
\definecolor{currentstroke}{rgb}{0.690196,0.690196,0.690196}%
\pgfsetstrokecolor{currentstroke}%
\pgfsetstrokeopacity{0.220000}%
\pgfsetdash{}{0pt}%
\pgfpathmoveto{\pgfqpoint{0.510067in}{2.894392in}}%
\pgfpathlineto{\pgfqpoint{7.247890in}{2.894392in}}%
\pgfusepath{stroke}%
\end{pgfscope}%
\begin{pgfscope}%
\pgfsetbuttcap%
\pgfsetroundjoin%
\definecolor{currentfill}{rgb}{0.000000,0.000000,0.000000}%
\pgfsetfillcolor{currentfill}%
\pgfsetlinewidth{0.803000pt}%
\definecolor{currentstroke}{rgb}{0.000000,0.000000,0.000000}%
\pgfsetstrokecolor{currentstroke}%
\pgfsetdash{}{0pt}%
\pgfsys@defobject{currentmarker}{\pgfqpoint{-0.048611in}{0.000000in}}{\pgfqpoint{-0.000000in}{0.000000in}}{%
\pgfpathmoveto{\pgfqpoint{-0.000000in}{0.000000in}}%
\pgfpathlineto{\pgfqpoint{-0.048611in}{0.000000in}}%
\pgfusepath{stroke,fill}%
}%
\begin{pgfscope}%
\pgfsys@transformshift{0.510067in}{2.894392in}%
\pgfsys@useobject{currentmarker}{}%
\end{pgfscope}%
\end{pgfscope}%
\begin{pgfscope}%
\definecolor{textcolor}{rgb}{0.000000,0.000000,0.000000}%
\pgfsetstrokecolor{textcolor}%
\pgfsetfillcolor{textcolor}%
\pgftext[x=0.236114in, y=2.841630in, left, base]{\color{textcolor}{\sffamily\fontsize{10.000000}{12.000000}\selectfont\catcode`\^=\active\def^{\ifmmode\sp\else\^{}\fi}\catcode`\%=\active\def%{\%}30}}%
\end{pgfscope}%
\begin{pgfscope}%
\pgfpathrectangle{\pgfqpoint{0.510067in}{0.472226in}}{\pgfqpoint{6.737823in}{4.286104in}}%
\pgfusepath{clip}%
\pgfsetrectcap%
\pgfsetroundjoin%
\pgfsetlinewidth{0.803000pt}%
\definecolor{currentstroke}{rgb}{0.690196,0.690196,0.690196}%
\pgfsetstrokecolor{currentstroke}%
\pgfsetstrokeopacity{0.220000}%
\pgfsetdash{}{0pt}%
\pgfpathmoveto{\pgfqpoint{0.510067in}{3.701780in}}%
\pgfpathlineto{\pgfqpoint{7.247890in}{3.701780in}}%
\pgfusepath{stroke}%
\end{pgfscope}%
\begin{pgfscope}%
\pgfsetbuttcap%
\pgfsetroundjoin%
\definecolor{currentfill}{rgb}{0.000000,0.000000,0.000000}%
\pgfsetfillcolor{currentfill}%
\pgfsetlinewidth{0.803000pt}%
\definecolor{currentstroke}{rgb}{0.000000,0.000000,0.000000}%
\pgfsetstrokecolor{currentstroke}%
\pgfsetdash{}{0pt}%
\pgfsys@defobject{currentmarker}{\pgfqpoint{-0.048611in}{0.000000in}}{\pgfqpoint{-0.000000in}{0.000000in}}{%
\pgfpathmoveto{\pgfqpoint{-0.000000in}{0.000000in}}%
\pgfpathlineto{\pgfqpoint{-0.048611in}{0.000000in}}%
\pgfusepath{stroke,fill}%
}%
\begin{pgfscope}%
\pgfsys@transformshift{0.510067in}{3.701780in}%
\pgfsys@useobject{currentmarker}{}%
\end{pgfscope}%
\end{pgfscope}%
\begin{pgfscope}%
\definecolor{textcolor}{rgb}{0.000000,0.000000,0.000000}%
\pgfsetstrokecolor{textcolor}%
\pgfsetfillcolor{textcolor}%
\pgftext[x=0.236114in, y=3.649019in, left, base]{\color{textcolor}{\sffamily\fontsize{10.000000}{12.000000}\selectfont\catcode`\^=\active\def^{\ifmmode\sp\else\^{}\fi}\catcode`\%=\active\def%{\%}40}}%
\end{pgfscope}%
\begin{pgfscope}%
\pgfpathrectangle{\pgfqpoint{0.510067in}{0.472226in}}{\pgfqpoint{6.737823in}{4.286104in}}%
\pgfusepath{clip}%
\pgfsetrectcap%
\pgfsetroundjoin%
\pgfsetlinewidth{0.803000pt}%
\definecolor{currentstroke}{rgb}{0.690196,0.690196,0.690196}%
\pgfsetstrokecolor{currentstroke}%
\pgfsetstrokeopacity{0.220000}%
\pgfsetdash{}{0pt}%
\pgfpathmoveto{\pgfqpoint{0.510067in}{4.509169in}}%
\pgfpathlineto{\pgfqpoint{7.247890in}{4.509169in}}%
\pgfusepath{stroke}%
\end{pgfscope}%
\begin{pgfscope}%
\pgfsetbuttcap%
\pgfsetroundjoin%
\definecolor{currentfill}{rgb}{0.000000,0.000000,0.000000}%
\pgfsetfillcolor{currentfill}%
\pgfsetlinewidth{0.803000pt}%
\definecolor{currentstroke}{rgb}{0.000000,0.000000,0.000000}%
\pgfsetstrokecolor{currentstroke}%
\pgfsetdash{}{0pt}%
\pgfsys@defobject{currentmarker}{\pgfqpoint{-0.048611in}{0.000000in}}{\pgfqpoint{-0.000000in}{0.000000in}}{%
\pgfpathmoveto{\pgfqpoint{-0.000000in}{0.000000in}}%
\pgfpathlineto{\pgfqpoint{-0.048611in}{0.000000in}}%
\pgfusepath{stroke,fill}%
}%
\begin{pgfscope}%
\pgfsys@transformshift{0.510067in}{4.509169in}%
\pgfsys@useobject{currentmarker}{}%
\end{pgfscope}%
\end{pgfscope}%
\begin{pgfscope}%
\definecolor{textcolor}{rgb}{0.000000,0.000000,0.000000}%
\pgfsetstrokecolor{textcolor}%
\pgfsetfillcolor{textcolor}%
\pgftext[x=0.236114in, y=4.456407in, left, base]{\color{textcolor}{\sffamily\fontsize{10.000000}{12.000000}\selectfont\catcode`\^=\active\def^{\ifmmode\sp\else\^{}\fi}\catcode`\%=\active\def%{\%}50}}%
\end{pgfscope}%
\begin{pgfscope}%
\definecolor{textcolor}{rgb}{0.000000,0.000000,0.000000}%
\pgfsetstrokecolor{textcolor}%
\pgfsetfillcolor{textcolor}%
\pgftext[x=0.180559in,y=2.615278in,,bottom,rotate=90.000000]{\color{textcolor}{\sffamily\fontsize{10.000000}{12.000000}\selectfont\catcode`\^=\active\def^{\ifmmode\sp\else\^{}\fi}\catcode`\%=\active\def%{\%}water pressure, $p$ [kPa]}}%
\end{pgfscope}%
\begin{pgfscope}%
\pgfpathrectangle{\pgfqpoint{0.510067in}{0.472226in}}{\pgfqpoint{6.737823in}{4.286104in}}%
\pgfusepath{clip}%
\pgfsetrectcap%
\pgfsetroundjoin%
\pgfsetlinewidth{1.806750pt}%
\definecolor{currentstroke}{rgb}{0.267004,0.004874,0.329415}%
\pgfsetstrokecolor{currentstroke}%
\pgfsetdash{}{0pt}%
\pgfpathmoveto{\pgfqpoint{0.510067in}{4.554230in}}%
\pgfpathlineto{\pgfqpoint{0.745891in}{4.553189in}}%
\pgfpathlineto{\pgfqpoint{0.981715in}{4.549991in}}%
\pgfpathlineto{\pgfqpoint{1.200694in}{4.544892in}}%
\pgfpathlineto{\pgfqpoint{1.402829in}{4.538100in}}%
\pgfpathlineto{\pgfqpoint{1.588119in}{4.529828in}}%
\pgfpathlineto{\pgfqpoint{1.773409in}{4.519271in}}%
\pgfpathlineto{\pgfqpoint{1.941855in}{4.507369in}}%
\pgfpathlineto{\pgfqpoint{2.093456in}{4.494495in}}%
\pgfpathlineto{\pgfqpoint{2.245057in}{4.479285in}}%
\pgfpathlineto{\pgfqpoint{2.379813in}{4.463557in}}%
\pgfpathlineto{\pgfqpoint{2.514570in}{4.445507in}}%
\pgfpathlineto{\pgfqpoint{2.632482in}{4.427611in}}%
\pgfpathlineto{\pgfqpoint{2.750394in}{4.407564in}}%
\pgfpathlineto{\pgfqpoint{2.868305in}{4.385178in}}%
\pgfpathlineto{\pgfqpoint{2.986217in}{4.360255in}}%
\pgfpathlineto{\pgfqpoint{3.087285in}{4.336717in}}%
\pgfpathlineto{\pgfqpoint{3.188352in}{4.311032in}}%
\pgfpathlineto{\pgfqpoint{3.289419in}{4.283063in}}%
\pgfpathlineto{\pgfqpoint{3.390487in}{4.252671in}}%
\pgfpathlineto{\pgfqpoint{3.491554in}{4.219717in}}%
\pgfpathlineto{\pgfqpoint{3.592621in}{4.184059in}}%
\pgfpathlineto{\pgfqpoint{3.676844in}{4.152178in}}%
\pgfpathlineto{\pgfqpoint{3.761067in}{4.118240in}}%
\pgfpathlineto{\pgfqpoint{3.845290in}{4.082165in}}%
\pgfpathlineto{\pgfqpoint{3.929513in}{4.043875in}}%
\pgfpathlineto{\pgfqpoint{4.013735in}{4.003291in}}%
\pgfpathlineto{\pgfqpoint{4.097958in}{3.960340in}}%
\pgfpathlineto{\pgfqpoint{4.182181in}{3.914947in}}%
\pgfpathlineto{\pgfqpoint{4.266404in}{3.867044in}}%
\pgfpathlineto{\pgfqpoint{4.350626in}{3.816564in}}%
\pgfpathlineto{\pgfqpoint{4.434849in}{3.763446in}}%
\pgfpathlineto{\pgfqpoint{4.519072in}{3.707631in}}%
\pgfpathlineto{\pgfqpoint{4.603295in}{3.649067in}}%
\pgfpathlineto{\pgfqpoint{4.687518in}{3.587708in}}%
\pgfpathlineto{\pgfqpoint{4.771740in}{3.523510in}}%
\pgfpathlineto{\pgfqpoint{4.855963in}{3.456440in}}%
\pgfpathlineto{\pgfqpoint{4.940186in}{3.386470in}}%
\pgfpathlineto{\pgfqpoint{5.024409in}{3.313577in}}%
\pgfpathlineto{\pgfqpoint{5.108632in}{3.237748in}}%
\pgfpathlineto{\pgfqpoint{5.192854in}{3.158979in}}%
\pgfpathlineto{\pgfqpoint{5.277077in}{3.077270in}}%
\pgfpathlineto{\pgfqpoint{5.361300in}{2.992634in}}%
\pgfpathlineto{\pgfqpoint{5.445523in}{2.905089in}}%
\pgfpathlineto{\pgfqpoint{5.529745in}{2.814665in}}%
\pgfpathlineto{\pgfqpoint{5.613968in}{2.721399in}}%
\pgfpathlineto{\pgfqpoint{5.715036in}{2.605795in}}%
\pgfpathlineto{\pgfqpoint{5.816103in}{2.486265in}}%
\pgfpathlineto{\pgfqpoint{5.917170in}{2.362927in}}%
\pgfpathlineto{\pgfqpoint{6.018238in}{2.235917in}}%
\pgfpathlineto{\pgfqpoint{6.119305in}{2.105390in}}%
\pgfpathlineto{\pgfqpoint{6.220372in}{1.971519in}}%
\pgfpathlineto{\pgfqpoint{6.338284in}{1.811365in}}%
\pgfpathlineto{\pgfqpoint{6.456196in}{1.647250in}}%
\pgfpathlineto{\pgfqpoint{6.574108in}{1.479533in}}%
\pgfpathlineto{\pgfqpoint{6.708864in}{1.283941in}}%
\pgfpathlineto{\pgfqpoint{6.860465in}{1.059640in}}%
\pgfpathlineto{\pgfqpoint{7.028911in}{0.806207in}}%
\pgfpathlineto{\pgfqpoint{7.247890in}{0.472226in}}%
\pgfpathlineto{\pgfqpoint{7.247890in}{0.472226in}}%
\pgfusepath{stroke}%
\end{pgfscope}%
\begin{pgfscope}%
\pgfpathrectangle{\pgfqpoint{0.510067in}{0.472226in}}{\pgfqpoint{6.737823in}{4.286104in}}%
\pgfusepath{clip}%
\pgfsetbuttcap%
\pgfsetroundjoin%
\definecolor{currentfill}{rgb}{1.000000,1.000000,1.000000}%
\pgfsetfillcolor{currentfill}%
\pgfsetlinewidth{0.903375pt}%
\definecolor{currentstroke}{rgb}{0.267004,0.004874,0.329415}%
\pgfsetstrokecolor{currentstroke}%
\pgfsetdash{}{0pt}%
\pgfsys@defobject{currentmarker}{\pgfqpoint{-0.026389in}{-0.026389in}}{\pgfqpoint{0.026389in}{0.026389in}}{%
\pgfpathmoveto{\pgfqpoint{0.000000in}{-0.026389in}}%
\pgfpathcurveto{\pgfqpoint{0.006998in}{-0.026389in}}{\pgfqpoint{0.013711in}{-0.023608in}}{\pgfqpoint{0.018660in}{-0.018660in}}%
\pgfpathcurveto{\pgfqpoint{0.023608in}{-0.013711in}}{\pgfqpoint{0.026389in}{-0.006998in}}{\pgfqpoint{0.026389in}{0.000000in}}%
\pgfpathcurveto{\pgfqpoint{0.026389in}{0.006998in}}{\pgfqpoint{0.023608in}{0.013711in}}{\pgfqpoint{0.018660in}{0.018660in}}%
\pgfpathcurveto{\pgfqpoint{0.013711in}{0.023608in}}{\pgfqpoint{0.006998in}{0.026389in}}{\pgfqpoint{0.000000in}{0.026389in}}%
\pgfpathcurveto{\pgfqpoint{-0.006998in}{0.026389in}}{\pgfqpoint{-0.013711in}{0.023608in}}{\pgfqpoint{-0.018660in}{0.018660in}}%
\pgfpathcurveto{\pgfqpoint{-0.023608in}{0.013711in}}{\pgfqpoint{-0.026389in}{0.006998in}}{\pgfqpoint{-0.026389in}{0.000000in}}%
\pgfpathcurveto{\pgfqpoint{-0.026389in}{-0.006998in}}{\pgfqpoint{-0.023608in}{-0.013711in}}{\pgfqpoint{-0.018660in}{-0.018660in}}%
\pgfpathcurveto{\pgfqpoint{-0.013711in}{-0.023608in}}{\pgfqpoint{-0.006998in}{-0.026389in}}{\pgfqpoint{0.000000in}{-0.026389in}}%
\pgfpathlineto{\pgfqpoint{0.000000in}{-0.026389in}}%
\pgfpathclose%
\pgfusepath{stroke,fill}%
}%
\begin{pgfscope}%
\pgfsys@transformshift{0.510067in}{4.536176in}%
\pgfsys@useobject{currentmarker}{}%
\end{pgfscope}%
\begin{pgfscope}%
\pgfsys@transformshift{0.846959in}{4.533797in}%
\pgfsys@useobject{currentmarker}{}%
\end{pgfscope}%
\begin{pgfscope}%
\pgfsys@transformshift{1.183850in}{4.526424in}%
\pgfsys@useobject{currentmarker}{}%
\end{pgfscope}%
\begin{pgfscope}%
\pgfsys@transformshift{1.520741in}{4.513329in}%
\pgfsys@useobject{currentmarker}{}%
\end{pgfscope}%
\begin{pgfscope}%
\pgfsys@transformshift{1.857632in}{4.493267in}%
\pgfsys@useobject{currentmarker}{}%
\end{pgfscope}%
\begin{pgfscope}%
\pgfsys@transformshift{2.194523in}{4.464422in}%
\pgfsys@useobject{currentmarker}{}%
\end{pgfscope}%
\begin{pgfscope}%
\pgfsys@transformshift{2.531414in}{4.424350in}%
\pgfsys@useobject{currentmarker}{}%
\end{pgfscope}%
\begin{pgfscope}%
\pgfsys@transformshift{2.868305in}{4.369923in}%
\pgfsys@useobject{currentmarker}{}%
\end{pgfscope}%
\begin{pgfscope}%
\pgfsys@transformshift{3.205197in}{4.297294in}%
\pgfsys@useobject{currentmarker}{}%
\end{pgfscope}%
\begin{pgfscope}%
\pgfsys@transformshift{3.542088in}{4.201906in}%
\pgfsys@useobject{currentmarker}{}%
\end{pgfscope}%
\begin{pgfscope}%
\pgfsys@transformshift{3.878979in}{4.078575in}%
\pgfsys@useobject{currentmarker}{}%
\end{pgfscope}%
\begin{pgfscope}%
\pgfsys@transformshift{4.215870in}{3.921674in}%
\pgfsys@useobject{currentmarker}{}%
\end{pgfscope}%
\begin{pgfscope}%
\pgfsys@transformshift{4.552761in}{3.725458in}%
\pgfsys@useobject{currentmarker}{}%
\end{pgfscope}%
\begin{pgfscope}%
\pgfsys@transformshift{4.889652in}{3.484536in}%
\pgfsys@useobject{currentmarker}{}%
\end{pgfscope}%
\begin{pgfscope}%
\pgfsys@transformshift{5.226543in}{3.194500in}%
\pgfsys@useobject{currentmarker}{}%
\end{pgfscope}%
\begin{pgfscope}%
\pgfsys@transformshift{5.563435in}{2.852661in}%
\pgfsys@useobject{currentmarker}{}%
\end{pgfscope}%
\begin{pgfscope}%
\pgfsys@transformshift{5.900326in}{2.458808in}%
\pgfsys@useobject{currentmarker}{}%
\end{pgfscope}%
\begin{pgfscope}%
\pgfsys@transformshift{6.237217in}{2.015848in}%
\pgfsys@useobject{currentmarker}{}%
\end{pgfscope}%
\begin{pgfscope}%
\pgfsys@transformshift{6.574108in}{1.530157in}%
\pgfsys@useobject{currentmarker}{}%
\end{pgfscope}%
\begin{pgfscope}%
\pgfsys@transformshift{6.910999in}{1.011474in}%
\pgfsys@useobject{currentmarker}{}%
\end{pgfscope}%
\begin{pgfscope}%
\pgfsys@transformshift{7.247890in}{0.472226in}%
\pgfsys@useobject{currentmarker}{}%
\end{pgfscope}%
\end{pgfscope}%
\begin{pgfscope}%
\pgfpathrectangle{\pgfqpoint{0.510067in}{0.472226in}}{\pgfqpoint{6.737823in}{4.286104in}}%
\pgfusepath{clip}%
\pgfsetrectcap%
\pgfsetroundjoin%
\pgfsetlinewidth{1.806750pt}%
\definecolor{currentstroke}{rgb}{0.229739,0.322361,0.545706}%
\pgfsetstrokecolor{currentstroke}%
\pgfsetdash{}{0pt}%
\pgfpathmoveto{\pgfqpoint{0.510067in}{3.925301in}}%
\pgfpathlineto{\pgfqpoint{0.627979in}{3.924123in}}%
\pgfpathlineto{\pgfqpoint{0.745891in}{3.920591in}}%
\pgfpathlineto{\pgfqpoint{0.863803in}{3.914704in}}%
\pgfpathlineto{\pgfqpoint{0.981715in}{3.906464in}}%
\pgfpathlineto{\pgfqpoint{1.099627in}{3.895872in}}%
\pgfpathlineto{\pgfqpoint{1.217539in}{3.882930in}}%
\pgfpathlineto{\pgfqpoint{1.335451in}{3.867641in}}%
\pgfpathlineto{\pgfqpoint{1.453363in}{3.850007in}}%
\pgfpathlineto{\pgfqpoint{1.571275in}{3.830032in}}%
\pgfpathlineto{\pgfqpoint{1.689186in}{3.807720in}}%
\pgfpathlineto{\pgfqpoint{1.807098in}{3.783076in}}%
\pgfpathlineto{\pgfqpoint{1.925010in}{3.756106in}}%
\pgfpathlineto{\pgfqpoint{2.042922in}{3.726814in}}%
\pgfpathlineto{\pgfqpoint{2.160834in}{3.695210in}}%
\pgfpathlineto{\pgfqpoint{2.278746in}{3.661300in}}%
\pgfpathlineto{\pgfqpoint{2.396658in}{3.625094in}}%
\pgfpathlineto{\pgfqpoint{2.514570in}{3.586602in}}%
\pgfpathlineto{\pgfqpoint{2.632482in}{3.545834in}}%
\pgfpathlineto{\pgfqpoint{2.750394in}{3.502804in}}%
\pgfpathlineto{\pgfqpoint{2.868305in}{3.457526in}}%
\pgfpathlineto{\pgfqpoint{2.986217in}{3.410013in}}%
\pgfpathlineto{\pgfqpoint{3.104129in}{3.360284in}}%
\pgfpathlineto{\pgfqpoint{3.222041in}{3.308355in}}%
\pgfpathlineto{\pgfqpoint{3.339953in}{3.254246in}}%
\pgfpathlineto{\pgfqpoint{3.457865in}{3.197979in}}%
\pgfpathlineto{\pgfqpoint{3.592621in}{3.131061in}}%
\pgfpathlineto{\pgfqpoint{3.727378in}{3.061392in}}%
\pgfpathlineto{\pgfqpoint{3.862134in}{2.989011in}}%
\pgfpathlineto{\pgfqpoint{3.996891in}{2.913962in}}%
\pgfpathlineto{\pgfqpoint{4.131647in}{2.836293in}}%
\pgfpathlineto{\pgfqpoint{4.266404in}{2.756054in}}%
\pgfpathlineto{\pgfqpoint{4.401160in}{2.673299in}}%
\pgfpathlineto{\pgfqpoint{4.535917in}{2.588086in}}%
\pgfpathlineto{\pgfqpoint{4.670673in}{2.500476in}}%
\pgfpathlineto{\pgfqpoint{4.822274in}{2.399130in}}%
\pgfpathlineto{\pgfqpoint{4.973875in}{2.294930in}}%
\pgfpathlineto{\pgfqpoint{5.125476in}{2.187979in}}%
\pgfpathlineto{\pgfqpoint{5.277077in}{2.078387in}}%
\pgfpathlineto{\pgfqpoint{5.445523in}{1.953661in}}%
\pgfpathlineto{\pgfqpoint{5.613968in}{1.825980in}}%
\pgfpathlineto{\pgfqpoint{5.782414in}{1.695519in}}%
\pgfpathlineto{\pgfqpoint{5.967704in}{1.549016in}}%
\pgfpathlineto{\pgfqpoint{6.152994in}{1.399618in}}%
\pgfpathlineto{\pgfqpoint{6.355129in}{1.233642in}}%
\pgfpathlineto{\pgfqpoint{6.574108in}{1.050702in}}%
\pgfpathlineto{\pgfqpoint{6.809932in}{0.850558in}}%
\pgfpathlineto{\pgfqpoint{7.062600in}{0.633168in}}%
\pgfpathlineto{\pgfqpoint{7.247890in}{0.472226in}}%
\pgfpathlineto{\pgfqpoint{7.247890in}{0.472226in}}%
\pgfusepath{stroke}%
\end{pgfscope}%
\begin{pgfscope}%
\pgfpathrectangle{\pgfqpoint{0.510067in}{0.472226in}}{\pgfqpoint{6.737823in}{4.286104in}}%
\pgfusepath{clip}%
\pgfsetbuttcap%
\pgfsetroundjoin%
\definecolor{currentfill}{rgb}{1.000000,1.000000,1.000000}%
\pgfsetfillcolor{currentfill}%
\pgfsetlinewidth{0.903375pt}%
\definecolor{currentstroke}{rgb}{0.229739,0.322361,0.545706}%
\pgfsetstrokecolor{currentstroke}%
\pgfsetdash{}{0pt}%
\pgfsys@defobject{currentmarker}{\pgfqpoint{-0.026389in}{-0.026389in}}{\pgfqpoint{0.026389in}{0.026389in}}{%
\pgfpathmoveto{\pgfqpoint{0.000000in}{-0.026389in}}%
\pgfpathcurveto{\pgfqpoint{0.006998in}{-0.026389in}}{\pgfqpoint{0.013711in}{-0.023608in}}{\pgfqpoint{0.018660in}{-0.018660in}}%
\pgfpathcurveto{\pgfqpoint{0.023608in}{-0.013711in}}{\pgfqpoint{0.026389in}{-0.006998in}}{\pgfqpoint{0.026389in}{0.000000in}}%
\pgfpathcurveto{\pgfqpoint{0.026389in}{0.006998in}}{\pgfqpoint{0.023608in}{0.013711in}}{\pgfqpoint{0.018660in}{0.018660in}}%
\pgfpathcurveto{\pgfqpoint{0.013711in}{0.023608in}}{\pgfqpoint{0.006998in}{0.026389in}}{\pgfqpoint{0.000000in}{0.026389in}}%
\pgfpathcurveto{\pgfqpoint{-0.006998in}{0.026389in}}{\pgfqpoint{-0.013711in}{0.023608in}}{\pgfqpoint{-0.018660in}{0.018660in}}%
\pgfpathcurveto{\pgfqpoint{-0.023608in}{0.013711in}}{\pgfqpoint{-0.026389in}{0.006998in}}{\pgfqpoint{-0.026389in}{0.000000in}}%
\pgfpathcurveto{\pgfqpoint{-0.026389in}{-0.006998in}}{\pgfqpoint{-0.023608in}{-0.013711in}}{\pgfqpoint{-0.018660in}{-0.018660in}}%
\pgfpathcurveto{\pgfqpoint{-0.013711in}{-0.023608in}}{\pgfqpoint{-0.006998in}{-0.026389in}}{\pgfqpoint{0.000000in}{-0.026389in}}%
\pgfpathlineto{\pgfqpoint{0.000000in}{-0.026389in}}%
\pgfpathclose%
\pgfusepath{stroke,fill}%
}%
\begin{pgfscope}%
\pgfsys@transformshift{0.510067in}{3.927107in}%
\pgfsys@useobject{currentmarker}{}%
\end{pgfscope}%
\begin{pgfscope}%
\pgfsys@transformshift{0.846959in}{3.917714in}%
\pgfsys@useobject{currentmarker}{}%
\end{pgfscope}%
\begin{pgfscope}%
\pgfsys@transformshift{1.183850in}{3.889535in}%
\pgfsys@useobject{currentmarker}{}%
\end{pgfscope}%
\begin{pgfscope}%
\pgfsys@transformshift{1.520741in}{3.842574in}%
\pgfsys@useobject{currentmarker}{}%
\end{pgfscope}%
\begin{pgfscope}%
\pgfsys@transformshift{1.857632in}{3.776847in}%
\pgfsys@useobject{currentmarker}{}%
\end{pgfscope}%
\begin{pgfscope}%
\pgfsys@transformshift{2.194523in}{3.692392in}%
\pgfsys@useobject{currentmarker}{}%
\end{pgfscope}%
\begin{pgfscope}%
\pgfsys@transformshift{2.531414in}{3.589285in}%
\pgfsys@useobject{currentmarker}{}%
\end{pgfscope}%
\begin{pgfscope}%
\pgfsys@transformshift{2.868305in}{3.467658in}%
\pgfsys@useobject{currentmarker}{}%
\end{pgfscope}%
\begin{pgfscope}%
\pgfsys@transformshift{3.205197in}{3.327718in}%
\pgfsys@useobject{currentmarker}{}%
\end{pgfscope}%
\begin{pgfscope}%
\pgfsys@transformshift{3.542088in}{3.169771in}%
\pgfsys@useobject{currentmarker}{}%
\end{pgfscope}%
\begin{pgfscope}%
\pgfsys@transformshift{3.878979in}{2.994236in}%
\pgfsys@useobject{currentmarker}{}%
\end{pgfscope}%
\begin{pgfscope}%
\pgfsys@transformshift{4.215870in}{2.801670in}%
\pgfsys@useobject{currentmarker}{}%
\end{pgfscope}%
\begin{pgfscope}%
\pgfsys@transformshift{4.552761in}{2.592778in}%
\pgfsys@useobject{currentmarker}{}%
\end{pgfscope}%
\begin{pgfscope}%
\pgfsys@transformshift{4.889652in}{2.368426in}%
\pgfsys@useobject{currentmarker}{}%
\end{pgfscope}%
\begin{pgfscope}%
\pgfsys@transformshift{5.226543in}{2.129648in}%
\pgfsys@useobject{currentmarker}{}%
\end{pgfscope}%
\begin{pgfscope}%
\pgfsys@transformshift{5.563435in}{1.877646in}%
\pgfsys@useobject{currentmarker}{}%
\end{pgfscope}%
\begin{pgfscope}%
\pgfsys@transformshift{5.900326in}{1.613787in}%
\pgfsys@useobject{currentmarker}{}%
\end{pgfscope}%
\begin{pgfscope}%
\pgfsys@transformshift{6.237217in}{1.339592in}%
\pgfsys@useobject{currentmarker}{}%
\end{pgfscope}%
\begin{pgfscope}%
\pgfsys@transformshift{6.574108in}{1.056727in}%
\pgfsys@useobject{currentmarker}{}%
\end{pgfscope}%
\begin{pgfscope}%
\pgfsys@transformshift{6.910999in}{0.766977in}%
\pgfsys@useobject{currentmarker}{}%
\end{pgfscope}%
\begin{pgfscope}%
\pgfsys@transformshift{7.247890in}{0.472226in}%
\pgfsys@useobject{currentmarker}{}%
\end{pgfscope}%
\end{pgfscope}%
\begin{pgfscope}%
\pgfpathrectangle{\pgfqpoint{0.510067in}{0.472226in}}{\pgfqpoint{6.737823in}{4.286104in}}%
\pgfusepath{clip}%
\pgfsetrectcap%
\pgfsetroundjoin%
\pgfsetlinewidth{1.806750pt}%
\definecolor{currentstroke}{rgb}{0.127568,0.566949,0.550556}%
\pgfsetstrokecolor{currentstroke}%
\pgfsetdash{}{0pt}%
\pgfpathmoveto{\pgfqpoint{0.510067in}{2.839340in}}%
\pgfpathlineto{\pgfqpoint{0.644824in}{2.838219in}}%
\pgfpathlineto{\pgfqpoint{0.779580in}{2.834854in}}%
\pgfpathlineto{\pgfqpoint{0.914337in}{2.829250in}}%
\pgfpathlineto{\pgfqpoint{1.049093in}{2.821411in}}%
\pgfpathlineto{\pgfqpoint{1.183850in}{2.811342in}}%
\pgfpathlineto{\pgfqpoint{1.318606in}{2.799052in}}%
\pgfpathlineto{\pgfqpoint{1.453363in}{2.784551in}}%
\pgfpathlineto{\pgfqpoint{1.588119in}{2.767849in}}%
\pgfpathlineto{\pgfqpoint{1.722876in}{2.748961in}}%
\pgfpathlineto{\pgfqpoint{1.857632in}{2.727902in}}%
\pgfpathlineto{\pgfqpoint{1.992388in}{2.704687in}}%
\pgfpathlineto{\pgfqpoint{2.127145in}{2.679336in}}%
\pgfpathlineto{\pgfqpoint{2.261901in}{2.651869in}}%
\pgfpathlineto{\pgfqpoint{2.396658in}{2.622308in}}%
\pgfpathlineto{\pgfqpoint{2.531414in}{2.590676in}}%
\pgfpathlineto{\pgfqpoint{2.666171in}{2.556998in}}%
\pgfpathlineto{\pgfqpoint{2.817772in}{2.516699in}}%
\pgfpathlineto{\pgfqpoint{2.969373in}{2.473886in}}%
\pgfpathlineto{\pgfqpoint{3.120974in}{2.428602in}}%
\pgfpathlineto{\pgfqpoint{3.272575in}{2.380893in}}%
\pgfpathlineto{\pgfqpoint{3.424176in}{2.330807in}}%
\pgfpathlineto{\pgfqpoint{3.575777in}{2.278395in}}%
\pgfpathlineto{\pgfqpoint{3.727378in}{2.223710in}}%
\pgfpathlineto{\pgfqpoint{3.878979in}{2.166807in}}%
\pgfpathlineto{\pgfqpoint{4.047424in}{2.101052in}}%
\pgfpathlineto{\pgfqpoint{4.215870in}{2.032713in}}%
\pgfpathlineto{\pgfqpoint{4.384316in}{1.961875in}}%
\pgfpathlineto{\pgfqpoint{4.552761in}{1.888629in}}%
\pgfpathlineto{\pgfqpoint{4.721207in}{1.813066in}}%
\pgfpathlineto{\pgfqpoint{4.906497in}{1.727383in}}%
\pgfpathlineto{\pgfqpoint{5.091787in}{1.639143in}}%
\pgfpathlineto{\pgfqpoint{5.277077in}{1.548479in}}%
\pgfpathlineto{\pgfqpoint{5.479212in}{1.446972in}}%
\pgfpathlineto{\pgfqpoint{5.681346in}{1.342929in}}%
\pgfpathlineto{\pgfqpoint{5.900326in}{1.227575in}}%
\pgfpathlineto{\pgfqpoint{6.119305in}{1.109715in}}%
\pgfpathlineto{\pgfqpoint{6.355129in}{0.980273in}}%
\pgfpathlineto{\pgfqpoint{6.607797in}{0.839048in}}%
\pgfpathlineto{\pgfqpoint{6.894155in}{0.676317in}}%
\pgfpathlineto{\pgfqpoint{7.214201in}{0.491781in}}%
\pgfpathlineto{\pgfqpoint{7.247890in}{0.472226in}}%
\pgfpathlineto{\pgfqpoint{7.247890in}{0.472226in}}%
\pgfusepath{stroke}%
\end{pgfscope}%
\begin{pgfscope}%
\pgfpathrectangle{\pgfqpoint{0.510067in}{0.472226in}}{\pgfqpoint{6.737823in}{4.286104in}}%
\pgfusepath{clip}%
\pgfsetbuttcap%
\pgfsetroundjoin%
\definecolor{currentfill}{rgb}{1.000000,1.000000,1.000000}%
\pgfsetfillcolor{currentfill}%
\pgfsetlinewidth{0.903375pt}%
\definecolor{currentstroke}{rgb}{0.127568,0.566949,0.550556}%
\pgfsetstrokecolor{currentstroke}%
\pgfsetdash{}{0pt}%
\pgfsys@defobject{currentmarker}{\pgfqpoint{-0.026389in}{-0.026389in}}{\pgfqpoint{0.026389in}{0.026389in}}{%
\pgfpathmoveto{\pgfqpoint{0.000000in}{-0.026389in}}%
\pgfpathcurveto{\pgfqpoint{0.006998in}{-0.026389in}}{\pgfqpoint{0.013711in}{-0.023608in}}{\pgfqpoint{0.018660in}{-0.018660in}}%
\pgfpathcurveto{\pgfqpoint{0.023608in}{-0.013711in}}{\pgfqpoint{0.026389in}{-0.006998in}}{\pgfqpoint{0.026389in}{0.000000in}}%
\pgfpathcurveto{\pgfqpoint{0.026389in}{0.006998in}}{\pgfqpoint{0.023608in}{0.013711in}}{\pgfqpoint{0.018660in}{0.018660in}}%
\pgfpathcurveto{\pgfqpoint{0.013711in}{0.023608in}}{\pgfqpoint{0.006998in}{0.026389in}}{\pgfqpoint{0.000000in}{0.026389in}}%
\pgfpathcurveto{\pgfqpoint{-0.006998in}{0.026389in}}{\pgfqpoint{-0.013711in}{0.023608in}}{\pgfqpoint{-0.018660in}{0.018660in}}%
\pgfpathcurveto{\pgfqpoint{-0.023608in}{0.013711in}}{\pgfqpoint{-0.026389in}{0.006998in}}{\pgfqpoint{-0.026389in}{0.000000in}}%
\pgfpathcurveto{\pgfqpoint{-0.026389in}{-0.006998in}}{\pgfqpoint{-0.023608in}{-0.013711in}}{\pgfqpoint{-0.018660in}{-0.018660in}}%
\pgfpathcurveto{\pgfqpoint{-0.013711in}{-0.023608in}}{\pgfqpoint{-0.006998in}{-0.026389in}}{\pgfqpoint{0.000000in}{-0.026389in}}%
\pgfpathlineto{\pgfqpoint{0.000000in}{-0.026389in}}%
\pgfpathclose%
\pgfusepath{stroke,fill}%
}%
\begin{pgfscope}%
\pgfsys@transformshift{0.510067in}{2.854195in}%
\pgfsys@useobject{currentmarker}{}%
\end{pgfscope}%
\begin{pgfscope}%
\pgfsys@transformshift{0.846959in}{2.847151in}%
\pgfsys@useobject{currentmarker}{}%
\end{pgfscope}%
\begin{pgfscope}%
\pgfsys@transformshift{1.183850in}{2.826053in}%
\pgfsys@useobject{currentmarker}{}%
\end{pgfscope}%
\begin{pgfscope}%
\pgfsys@transformshift{1.520741in}{2.791006in}%
\pgfsys@useobject{currentmarker}{}%
\end{pgfscope}%
\begin{pgfscope}%
\pgfsys@transformshift{1.857632in}{2.742179in}%
\pgfsys@useobject{currentmarker}{}%
\end{pgfscope}%
\begin{pgfscope}%
\pgfsys@transformshift{2.194523in}{2.679814in}%
\pgfsys@useobject{currentmarker}{}%
\end{pgfscope}%
\begin{pgfscope}%
\pgfsys@transformshift{2.531414in}{2.604216in}%
\pgfsys@useobject{currentmarker}{}%
\end{pgfscope}%
\begin{pgfscope}%
\pgfsys@transformshift{2.868305in}{2.515758in}%
\pgfsys@useobject{currentmarker}{}%
\end{pgfscope}%
\begin{pgfscope}%
\pgfsys@transformshift{3.205197in}{2.414877in}%
\pgfsys@useobject{currentmarker}{}%
\end{pgfscope}%
\begin{pgfscope}%
\pgfsys@transformshift{3.542088in}{2.302071in}%
\pgfsys@useobject{currentmarker}{}%
\end{pgfscope}%
\begin{pgfscope}%
\pgfsys@transformshift{3.878979in}{2.177901in}%
\pgfsys@useobject{currentmarker}{}%
\end{pgfscope}%
\begin{pgfscope}%
\pgfsys@transformshift{4.215870in}{2.042986in}%
\pgfsys@useobject{currentmarker}{}%
\end{pgfscope}%
\begin{pgfscope}%
\pgfsys@transformshift{4.552761in}{1.897998in}%
\pgfsys@useobject{currentmarker}{}%
\end{pgfscope}%
\begin{pgfscope}%
\pgfsys@transformshift{4.889652in}{1.743666in}%
\pgfsys@useobject{currentmarker}{}%
\end{pgfscope}%
\begin{pgfscope}%
\pgfsys@transformshift{5.226543in}{1.580764in}%
\pgfsys@useobject{currentmarker}{}%
\end{pgfscope}%
\begin{pgfscope}%
\pgfsys@transformshift{5.563435in}{1.410115in}%
\pgfsys@useobject{currentmarker}{}%
\end{pgfscope}%
\begin{pgfscope}%
\pgfsys@transformshift{5.900326in}{1.232586in}%
\pgfsys@useobject{currentmarker}{}%
\end{pgfscope}%
\begin{pgfscope}%
\pgfsys@transformshift{6.237217in}{1.049080in}%
\pgfsys@useobject{currentmarker}{}%
\end{pgfscope}%
\begin{pgfscope}%
\pgfsys@transformshift{6.574108in}{0.860536in}%
\pgfsys@useobject{currentmarker}{}%
\end{pgfscope}%
\begin{pgfscope}%
\pgfsys@transformshift{6.910999in}{0.667921in}%
\pgfsys@useobject{currentmarker}{}%
\end{pgfscope}%
\begin{pgfscope}%
\pgfsys@transformshift{7.247890in}{0.472226in}%
\pgfsys@useobject{currentmarker}{}%
\end{pgfscope}%
\end{pgfscope}%
\begin{pgfscope}%
\pgfpathrectangle{\pgfqpoint{0.510067in}{0.472226in}}{\pgfqpoint{6.737823in}{4.286104in}}%
\pgfusepath{clip}%
\pgfsetrectcap%
\pgfsetroundjoin%
\pgfsetlinewidth{1.806750pt}%
\definecolor{currentstroke}{rgb}{0.369214,0.788888,0.382914}%
\pgfsetstrokecolor{currentstroke}%
\pgfsetdash{}{0pt}%
\pgfpathmoveto{\pgfqpoint{0.510067in}{1.567786in}}%
\pgfpathlineto{\pgfqpoint{0.712202in}{1.566616in}}%
\pgfpathlineto{\pgfqpoint{0.914337in}{1.563110in}}%
\pgfpathlineto{\pgfqpoint{1.116471in}{1.557272in}}%
\pgfpathlineto{\pgfqpoint{1.318606in}{1.549114in}}%
\pgfpathlineto{\pgfqpoint{1.520741in}{1.538651in}}%
\pgfpathlineto{\pgfqpoint{1.722876in}{1.525901in}}%
\pgfpathlineto{\pgfqpoint{1.925010in}{1.510888in}}%
\pgfpathlineto{\pgfqpoint{2.127145in}{1.493639in}}%
\pgfpathlineto{\pgfqpoint{2.329280in}{1.474185in}}%
\pgfpathlineto{\pgfqpoint{2.531414in}{1.452562in}}%
\pgfpathlineto{\pgfqpoint{2.733549in}{1.428809in}}%
\pgfpathlineto{\pgfqpoint{2.935684in}{1.402968in}}%
\pgfpathlineto{\pgfqpoint{3.154663in}{1.372673in}}%
\pgfpathlineto{\pgfqpoint{3.373642in}{1.340049in}}%
\pgfpathlineto{\pgfqpoint{3.592621in}{1.305164in}}%
\pgfpathlineto{\pgfqpoint{3.811601in}{1.268092in}}%
\pgfpathlineto{\pgfqpoint{4.047424in}{1.225815in}}%
\pgfpathlineto{\pgfqpoint{4.283248in}{1.181199in}}%
\pgfpathlineto{\pgfqpoint{4.519072in}{1.134353in}}%
\pgfpathlineto{\pgfqpoint{4.771740in}{1.081820in}}%
\pgfpathlineto{\pgfqpoint{5.024409in}{1.027009in}}%
\pgfpathlineto{\pgfqpoint{5.293922in}{0.966209in}}%
\pgfpathlineto{\pgfqpoint{5.580279in}{0.899183in}}%
\pgfpathlineto{\pgfqpoint{5.866637in}{0.829898in}}%
\pgfpathlineto{\pgfqpoint{6.186683in}{0.750098in}}%
\pgfpathlineto{\pgfqpoint{6.523574in}{0.663788in}}%
\pgfpathlineto{\pgfqpoint{6.910999in}{0.562159in}}%
\pgfpathlineto{\pgfqpoint{7.247890in}{0.472226in}}%
\pgfpathlineto{\pgfqpoint{7.247890in}{0.472226in}}%
\pgfusepath{stroke}%
\end{pgfscope}%
\begin{pgfscope}%
\pgfpathrectangle{\pgfqpoint{0.510067in}{0.472226in}}{\pgfqpoint{6.737823in}{4.286104in}}%
\pgfusepath{clip}%
\pgfsetbuttcap%
\pgfsetroundjoin%
\definecolor{currentfill}{rgb}{1.000000,1.000000,1.000000}%
\pgfsetfillcolor{currentfill}%
\pgfsetlinewidth{0.903375pt}%
\definecolor{currentstroke}{rgb}{0.369214,0.788888,0.382914}%
\pgfsetstrokecolor{currentstroke}%
\pgfsetdash{}{0pt}%
\pgfsys@defobject{currentmarker}{\pgfqpoint{-0.026389in}{-0.026389in}}{\pgfqpoint{0.026389in}{0.026389in}}{%
\pgfpathmoveto{\pgfqpoint{0.000000in}{-0.026389in}}%
\pgfpathcurveto{\pgfqpoint{0.006998in}{-0.026389in}}{\pgfqpoint{0.013711in}{-0.023608in}}{\pgfqpoint{0.018660in}{-0.018660in}}%
\pgfpathcurveto{\pgfqpoint{0.023608in}{-0.013711in}}{\pgfqpoint{0.026389in}{-0.006998in}}{\pgfqpoint{0.026389in}{0.000000in}}%
\pgfpathcurveto{\pgfqpoint{0.026389in}{0.006998in}}{\pgfqpoint{0.023608in}{0.013711in}}{\pgfqpoint{0.018660in}{0.018660in}}%
\pgfpathcurveto{\pgfqpoint{0.013711in}{0.023608in}}{\pgfqpoint{0.006998in}{0.026389in}}{\pgfqpoint{0.000000in}{0.026389in}}%
\pgfpathcurveto{\pgfqpoint{-0.006998in}{0.026389in}}{\pgfqpoint{-0.013711in}{0.023608in}}{\pgfqpoint{-0.018660in}{0.018660in}}%
\pgfpathcurveto{\pgfqpoint{-0.023608in}{0.013711in}}{\pgfqpoint{-0.026389in}{0.006998in}}{\pgfqpoint{-0.026389in}{0.000000in}}%
\pgfpathcurveto{\pgfqpoint{-0.026389in}{-0.006998in}}{\pgfqpoint{-0.023608in}{-0.013711in}}{\pgfqpoint{-0.018660in}{-0.018660in}}%
\pgfpathcurveto{\pgfqpoint{-0.013711in}{-0.023608in}}{\pgfqpoint{-0.006998in}{-0.026389in}}{\pgfqpoint{0.000000in}{-0.026389in}}%
\pgfpathlineto{\pgfqpoint{0.000000in}{-0.026389in}}%
\pgfpathclose%
\pgfusepath{stroke,fill}%
}%
\begin{pgfscope}%
\pgfsys@transformshift{0.510067in}{1.581122in}%
\pgfsys@useobject{currentmarker}{}%
\end{pgfscope}%
\begin{pgfscope}%
\pgfsys@transformshift{0.846959in}{1.577833in}%
\pgfsys@useobject{currentmarker}{}%
\end{pgfscope}%
\begin{pgfscope}%
\pgfsys@transformshift{1.183850in}{1.567985in}%
\pgfsys@useobject{currentmarker}{}%
\end{pgfscope}%
\begin{pgfscope}%
\pgfsys@transformshift{1.520741in}{1.551626in}%
\pgfsys@useobject{currentmarker}{}%
\end{pgfscope}%
\begin{pgfscope}%
\pgfsys@transformshift{1.857632in}{1.528839in}%
\pgfsys@useobject{currentmarker}{}%
\end{pgfscope}%
\begin{pgfscope}%
\pgfsys@transformshift{2.194523in}{1.499739in}%
\pgfsys@useobject{currentmarker}{}%
\end{pgfscope}%
\begin{pgfscope}%
\pgfsys@transformshift{2.531414in}{1.464473in}%
\pgfsys@useobject{currentmarker}{}%
\end{pgfscope}%
\begin{pgfscope}%
\pgfsys@transformshift{2.868305in}{1.423218in}%
\pgfsys@useobject{currentmarker}{}%
\end{pgfscope}%
\begin{pgfscope}%
\pgfsys@transformshift{3.205197in}{1.376182in}%
\pgfsys@useobject{currentmarker}{}%
\end{pgfscope}%
\begin{pgfscope}%
\pgfsys@transformshift{3.542088in}{1.323603in}%
\pgfsys@useobject{currentmarker}{}%
\end{pgfscope}%
\begin{pgfscope}%
\pgfsys@transformshift{3.878979in}{1.265746in}%
\pgfsys@useobject{currentmarker}{}%
\end{pgfscope}%
\begin{pgfscope}%
\pgfsys@transformshift{4.215870in}{1.202902in}%
\pgfsys@useobject{currentmarker}{}%
\end{pgfscope}%
\begin{pgfscope}%
\pgfsys@transformshift{4.552761in}{1.135390in}%
\pgfsys@useobject{currentmarker}{}%
\end{pgfscope}%
\begin{pgfscope}%
\pgfsys@transformshift{4.889652in}{1.063548in}%
\pgfsys@useobject{currentmarker}{}%
\end{pgfscope}%
\begin{pgfscope}%
\pgfsys@transformshift{5.226543in}{0.987742in}%
\pgfsys@useobject{currentmarker}{}%
\end{pgfscope}%
\begin{pgfscope}%
\pgfsys@transformshift{5.563435in}{0.908354in}%
\pgfsys@useobject{currentmarker}{}%
\end{pgfscope}%
\begin{pgfscope}%
\pgfsys@transformshift{5.900326in}{0.825785in}%
\pgfsys@useobject{currentmarker}{}%
\end{pgfscope}%
\begin{pgfscope}%
\pgfsys@transformshift{6.237217in}{0.740452in}%
\pgfsys@useobject{currentmarker}{}%
\end{pgfscope}%
\begin{pgfscope}%
\pgfsys@transformshift{6.574108in}{0.652786in}%
\pgfsys@useobject{currentmarker}{}%
\end{pgfscope}%
\begin{pgfscope}%
\pgfsys@transformshift{6.910999in}{0.563228in}%
\pgfsys@useobject{currentmarker}{}%
\end{pgfscope}%
\begin{pgfscope}%
\pgfsys@transformshift{7.247890in}{0.472226in}%
\pgfsys@useobject{currentmarker}{}%
\end{pgfscope}%
\end{pgfscope}%
\begin{pgfscope}%
\pgfpathrectangle{\pgfqpoint{0.510067in}{0.472226in}}{\pgfqpoint{6.737823in}{4.286104in}}%
\pgfusepath{clip}%
\pgfsetrectcap%
\pgfsetroundjoin%
\pgfsetlinewidth{1.806750pt}%
\definecolor{currentstroke}{rgb}{0.993248,0.906157,0.143936}%
\pgfsetstrokecolor{currentstroke}%
\pgfsetdash{}{0pt}%
\pgfpathmoveto{\pgfqpoint{0.510067in}{0.706818in}}%
\pgfpathlineto{\pgfqpoint{0.931181in}{0.705732in}}%
\pgfpathlineto{\pgfqpoint{1.352295in}{0.702481in}}%
\pgfpathlineto{\pgfqpoint{1.790254in}{0.696831in}}%
\pgfpathlineto{\pgfqpoint{2.228212in}{0.688917in}}%
\pgfpathlineto{\pgfqpoint{2.666171in}{0.678804in}}%
\pgfpathlineto{\pgfqpoint{3.120974in}{0.666069in}}%
\pgfpathlineto{\pgfqpoint{3.592621in}{0.650583in}}%
\pgfpathlineto{\pgfqpoint{4.081114in}{0.632257in}}%
\pgfpathlineto{\pgfqpoint{4.586450in}{0.611057in}}%
\pgfpathlineto{\pgfqpoint{5.125476in}{0.586198in}}%
\pgfpathlineto{\pgfqpoint{5.715036in}{0.556725in}}%
\pgfpathlineto{\pgfqpoint{6.371973in}{0.521619in}}%
\pgfpathlineto{\pgfqpoint{7.180512in}{0.476097in}}%
\pgfpathlineto{\pgfqpoint{7.247890in}{0.472226in}}%
\pgfpathlineto{\pgfqpoint{7.247890in}{0.472226in}}%
\pgfusepath{stroke}%
\end{pgfscope}%
\begin{pgfscope}%
\pgfpathrectangle{\pgfqpoint{0.510067in}{0.472226in}}{\pgfqpoint{6.737823in}{4.286104in}}%
\pgfusepath{clip}%
\pgfsetbuttcap%
\pgfsetroundjoin%
\definecolor{currentfill}{rgb}{1.000000,1.000000,1.000000}%
\pgfsetfillcolor{currentfill}%
\pgfsetlinewidth{0.903375pt}%
\definecolor{currentstroke}{rgb}{0.993248,0.906157,0.143936}%
\pgfsetstrokecolor{currentstroke}%
\pgfsetdash{}{0pt}%
\pgfsys@defobject{currentmarker}{\pgfqpoint{-0.026389in}{-0.026389in}}{\pgfqpoint{0.026389in}{0.026389in}}{%
\pgfpathmoveto{\pgfqpoint{0.000000in}{-0.026389in}}%
\pgfpathcurveto{\pgfqpoint{0.006998in}{-0.026389in}}{\pgfqpoint{0.013711in}{-0.023608in}}{\pgfqpoint{0.018660in}{-0.018660in}}%
\pgfpathcurveto{\pgfqpoint{0.023608in}{-0.013711in}}{\pgfqpoint{0.026389in}{-0.006998in}}{\pgfqpoint{0.026389in}{0.000000in}}%
\pgfpathcurveto{\pgfqpoint{0.026389in}{0.006998in}}{\pgfqpoint{0.023608in}{0.013711in}}{\pgfqpoint{0.018660in}{0.018660in}}%
\pgfpathcurveto{\pgfqpoint{0.013711in}{0.023608in}}{\pgfqpoint{0.006998in}{0.026389in}}{\pgfqpoint{0.000000in}{0.026389in}}%
\pgfpathcurveto{\pgfqpoint{-0.006998in}{0.026389in}}{\pgfqpoint{-0.013711in}{0.023608in}}{\pgfqpoint{-0.018660in}{0.018660in}}%
\pgfpathcurveto{\pgfqpoint{-0.023608in}{0.013711in}}{\pgfqpoint{-0.026389in}{0.006998in}}{\pgfqpoint{-0.026389in}{0.000000in}}%
\pgfpathcurveto{\pgfqpoint{-0.026389in}{-0.006998in}}{\pgfqpoint{-0.023608in}{-0.013711in}}{\pgfqpoint{-0.018660in}{-0.018660in}}%
\pgfpathcurveto{\pgfqpoint{-0.013711in}{-0.023608in}}{\pgfqpoint{-0.006998in}{-0.026389in}}{\pgfqpoint{0.000000in}{-0.026389in}}%
\pgfpathlineto{\pgfqpoint{0.000000in}{-0.026389in}}%
\pgfpathclose%
\pgfusepath{stroke,fill}%
}%
\begin{pgfscope}%
\pgfsys@transformshift{0.510067in}{0.711830in}%
\pgfsys@useobject{currentmarker}{}%
\end{pgfscope}%
\begin{pgfscope}%
\pgfsys@transformshift{0.846959in}{0.711119in}%
\pgfsys@useobject{currentmarker}{}%
\end{pgfscope}%
\begin{pgfscope}%
\pgfsys@transformshift{1.183850in}{0.708991in}%
\pgfsys@useobject{currentmarker}{}%
\end{pgfscope}%
\begin{pgfscope}%
\pgfsys@transformshift{1.520741in}{0.705455in}%
\pgfsys@useobject{currentmarker}{}%
\end{pgfscope}%
\begin{pgfscope}%
\pgfsys@transformshift{1.857632in}{0.700531in}%
\pgfsys@useobject{currentmarker}{}%
\end{pgfscope}%
\begin{pgfscope}%
\pgfsys@transformshift{2.194523in}{0.694242in}%
\pgfsys@useobject{currentmarker}{}%
\end{pgfscope}%
\begin{pgfscope}%
\pgfsys@transformshift{2.531414in}{0.686621in}%
\pgfsys@useobject{currentmarker}{}%
\end{pgfscope}%
\begin{pgfscope}%
\pgfsys@transformshift{2.868305in}{0.677706in}%
\pgfsys@useobject{currentmarker}{}%
\end{pgfscope}%
\begin{pgfscope}%
\pgfsys@transformshift{3.205197in}{0.667542in}%
\pgfsys@useobject{currentmarker}{}%
\end{pgfscope}%
\begin{pgfscope}%
\pgfsys@transformshift{3.542088in}{0.656181in}%
\pgfsys@useobject{currentmarker}{}%
\end{pgfscope}%
\begin{pgfscope}%
\pgfsys@transformshift{3.878979in}{0.643679in}%
\pgfsys@useobject{currentmarker}{}%
\end{pgfscope}%
\begin{pgfscope}%
\pgfsys@transformshift{4.215870in}{0.630101in}%
\pgfsys@useobject{currentmarker}{}%
\end{pgfscope}%
\begin{pgfscope}%
\pgfsys@transformshift{4.552761in}{0.615513in}%
\pgfsys@useobject{currentmarker}{}%
\end{pgfscope}%
\begin{pgfscope}%
\pgfsys@transformshift{4.889652in}{0.599991in}%
\pgfsys@useobject{currentmarker}{}%
\end{pgfscope}%
\begin{pgfscope}%
\pgfsys@transformshift{5.226543in}{0.583611in}%
\pgfsys@useobject{currentmarker}{}%
\end{pgfscope}%
\begin{pgfscope}%
\pgfsys@transformshift{5.563435in}{0.566458in}%
\pgfsys@useobject{currentmarker}{}%
\end{pgfscope}%
\begin{pgfscope}%
\pgfsys@transformshift{5.900326in}{0.548616in}%
\pgfsys@useobject{currentmarker}{}%
\end{pgfscope}%
\begin{pgfscope}%
\pgfsys@transformshift{6.237217in}{0.530177in}%
\pgfsys@useobject{currentmarker}{}%
\end{pgfscope}%
\begin{pgfscope}%
\pgfsys@transformshift{6.574108in}{0.511233in}%
\pgfsys@useobject{currentmarker}{}%
\end{pgfscope}%
\begin{pgfscope}%
\pgfsys@transformshift{6.910999in}{0.491882in}%
\pgfsys@useobject{currentmarker}{}%
\end{pgfscope}%
\begin{pgfscope}%
\pgfsys@transformshift{7.247890in}{0.472226in}%
\pgfsys@useobject{currentmarker}{}%
\end{pgfscope}%
\end{pgfscope}%
\begin{pgfscope}%
\pgfsetrectcap%
\pgfsetmiterjoin%
\pgfsetlinewidth{0.803000pt}%
\definecolor{currentstroke}{rgb}{0.000000,0.000000,0.000000}%
\pgfsetstrokecolor{currentstroke}%
\pgfsetdash{}{0pt}%
\pgfpathmoveto{\pgfqpoint{0.510067in}{0.472226in}}%
\pgfpathlineto{\pgfqpoint{0.510067in}{4.758330in}}%
\pgfusepath{stroke}%
\end{pgfscope}%
\begin{pgfscope}%
\pgfsetrectcap%
\pgfsetmiterjoin%
\pgfsetlinewidth{0.803000pt}%
\definecolor{currentstroke}{rgb}{0.000000,0.000000,0.000000}%
\pgfsetstrokecolor{currentstroke}%
\pgfsetdash{}{0pt}%
\pgfpathmoveto{\pgfqpoint{7.247890in}{0.472226in}}%
\pgfpathlineto{\pgfqpoint{7.247890in}{4.758330in}}%
\pgfusepath{stroke}%
\end{pgfscope}%
\begin{pgfscope}%
\pgfsetrectcap%
\pgfsetmiterjoin%
\pgfsetlinewidth{0.803000pt}%
\definecolor{currentstroke}{rgb}{0.000000,0.000000,0.000000}%
\pgfsetstrokecolor{currentstroke}%
\pgfsetdash{}{0pt}%
\pgfpathmoveto{\pgfqpoint{0.510067in}{0.472226in}}%
\pgfpathlineto{\pgfqpoint{7.247890in}{0.472226in}}%
\pgfusepath{stroke}%
\end{pgfscope}%
\begin{pgfscope}%
\pgfsetrectcap%
\pgfsetmiterjoin%
\pgfsetlinewidth{0.803000pt}%
\definecolor{currentstroke}{rgb}{0.000000,0.000000,0.000000}%
\pgfsetstrokecolor{currentstroke}%
\pgfsetdash{}{0pt}%
\pgfpathmoveto{\pgfqpoint{0.510067in}{4.758330in}}%
\pgfpathlineto{\pgfqpoint{7.247890in}{4.758330in}}%
\pgfusepath{stroke}%
\end{pgfscope}%
\begin{pgfscope}%
\pgfsetbuttcap%
\pgfsetmiterjoin%
\definecolor{currentfill}{rgb}{1.000000,1.000000,1.000000}%
\pgfsetfillcolor{currentfill}%
\pgfsetfillopacity{0.800000}%
\pgfsetlinewidth{1.003750pt}%
\definecolor{currentstroke}{rgb}{0.800000,0.800000,0.800000}%
\pgfsetstrokecolor{currentstroke}%
\pgfsetstrokeopacity{0.800000}%
\pgfsetdash{}{0pt}%
\pgfpathmoveto{\pgfqpoint{0.587845in}{0.527781in}}%
\pgfpathlineto{\pgfqpoint{2.482991in}{0.527781in}}%
\pgfpathquadraticcurveto{\pgfqpoint{2.505213in}{0.527781in}}{\pgfqpoint{2.505213in}{0.550003in}}%
\pgfpathlineto{\pgfqpoint{2.505213in}{1.219448in}}%
\pgfpathquadraticcurveto{\pgfqpoint{2.505213in}{1.241670in}}{\pgfqpoint{2.482991in}{1.241670in}}%
\pgfpathlineto{\pgfqpoint{0.587845in}{1.241670in}}%
\pgfpathquadraticcurveto{\pgfqpoint{0.565623in}{1.241670in}}{\pgfqpoint{0.565623in}{1.219448in}}%
\pgfpathlineto{\pgfqpoint{0.565623in}{0.550003in}}%
\pgfpathquadraticcurveto{\pgfqpoint{0.565623in}{0.527781in}}{\pgfqpoint{0.587845in}{0.527781in}}%
\pgfpathlineto{\pgfqpoint{0.587845in}{0.527781in}}%
\pgfpathclose%
\pgfusepath{stroke,fill}%
\end{pgfscope}%
\begin{pgfscope}%
\definecolor{textcolor}{rgb}{0.000000,0.000000,0.000000}%
\pgfsetstrokecolor{textcolor}%
\pgfsetfillcolor{textcolor}%
\pgftext[x=1.180010in,y=1.102255in,left,base]{\color{textcolor}{\sffamily\fontsize{9.000000}{10.800000}\selectfont\catcode`\^=\active\def^{\ifmmode\sp\else\^{}\fi}\catcode`\%=\active\def%{\%}Profile time}}%
\end{pgfscope}%
\begin{pgfscope}%
\pgfsetrectcap%
\pgfsetroundjoin%
\pgfsetlinewidth{1.806750pt}%
\definecolor{currentstroke}{rgb}{0.267004,0.004874,0.329415}%
\pgfsetstrokecolor{currentstroke}%
\pgfsetdash{}{0pt}%
\pgfpathmoveto{\pgfqpoint{0.610067in}{0.971141in}}%
\pgfpathlineto{\pgfqpoint{0.721179in}{0.971141in}}%
\pgfpathlineto{\pgfqpoint{0.832290in}{0.971141in}}%
\pgfusepath{stroke}%
\end{pgfscope}%
\begin{pgfscope}%
\definecolor{textcolor}{rgb}{0.000000,0.000000,0.000000}%
\pgfsetstrokecolor{textcolor}%
\pgfsetfillcolor{textcolor}%
\pgftext[x=0.921179in,y=0.932252in,left,base]{\color{textcolor}{\sffamily\fontsize{8.000000}{9.600000}\selectfont\catcode`\^=\active\def^{\ifmmode\sp\else\^{}\fi}\catcode`\%=\active\def%{\%}$t=0.014$ s}}%
\end{pgfscope}%
\begin{pgfscope}%
\pgfsetrectcap%
\pgfsetroundjoin%
\pgfsetlinewidth{1.806750pt}%
\definecolor{currentstroke}{rgb}{0.229739,0.322361,0.545706}%
\pgfsetstrokecolor{currentstroke}%
\pgfsetdash{}{0pt}%
\pgfpathmoveto{\pgfqpoint{0.610067in}{0.804474in}}%
\pgfpathlineto{\pgfqpoint{0.721179in}{0.804474in}}%
\pgfpathlineto{\pgfqpoint{0.832290in}{0.804474in}}%
\pgfusepath{stroke}%
\end{pgfscope}%
\begin{pgfscope}%
\definecolor{textcolor}{rgb}{0.000000,0.000000,0.000000}%
\pgfsetstrokecolor{textcolor}%
\pgfsetfillcolor{textcolor}%
\pgftext[x=0.921179in,y=0.765585in,left,base]{\color{textcolor}{\sffamily\fontsize{8.000000}{9.600000}\selectfont\catcode`\^=\active\def^{\ifmmode\sp\else\^{}\fi}\catcode`\%=\active\def%{\%}$t=0.05$ s}}%
\end{pgfscope}%
\begin{pgfscope}%
\pgfsetrectcap%
\pgfsetroundjoin%
\pgfsetlinewidth{1.806750pt}%
\definecolor{currentstroke}{rgb}{0.127568,0.566949,0.550556}%
\pgfsetstrokecolor{currentstroke}%
\pgfsetdash{}{0pt}%
\pgfpathmoveto{\pgfqpoint{0.610067in}{0.637807in}}%
\pgfpathlineto{\pgfqpoint{0.721179in}{0.637807in}}%
\pgfpathlineto{\pgfqpoint{0.832290in}{0.637807in}}%
\pgfusepath{stroke}%
\end{pgfscope}%
\begin{pgfscope}%
\definecolor{textcolor}{rgb}{0.000000,0.000000,0.000000}%
\pgfsetstrokecolor{textcolor}%
\pgfsetfillcolor{textcolor}%
\pgftext[x=0.921179in,y=0.598918in,left,base]{\color{textcolor}{\sffamily\fontsize{8.000000}{9.600000}\selectfont\catcode`\^=\active\def^{\ifmmode\sp\else\^{}\fi}\catcode`\%=\active\def%{\%}$t=0.1$ s}}%
\end{pgfscope}%
\begin{pgfscope}%
\pgfsetrectcap%
\pgfsetroundjoin%
\pgfsetlinewidth{1.806750pt}%
\definecolor{currentstroke}{rgb}{0.369214,0.788888,0.382914}%
\pgfsetstrokecolor{currentstroke}%
\pgfsetdash{}{0pt}%
\pgfpathmoveto{\pgfqpoint{1.705558in}{0.971141in}}%
\pgfpathlineto{\pgfqpoint{1.816669in}{0.971141in}}%
\pgfpathlineto{\pgfqpoint{1.927780in}{0.971141in}}%
\pgfusepath{stroke}%
\end{pgfscope}%
\begin{pgfscope}%
\definecolor{textcolor}{rgb}{0.000000,0.000000,0.000000}%
\pgfsetstrokecolor{textcolor}%
\pgfsetfillcolor{textcolor}%
\pgftext[x=2.016669in,y=0.932252in,left,base]{\color{textcolor}{\sffamily\fontsize{8.000000}{9.600000}\selectfont\catcode`\^=\active\def^{\ifmmode\sp\else\^{}\fi}\catcode`\%=\active\def%{\%}$t=0.2$ s}}%
\end{pgfscope}%
\begin{pgfscope}%
\pgfsetrectcap%
\pgfsetroundjoin%
\pgfsetlinewidth{1.806750pt}%
\definecolor{currentstroke}{rgb}{0.993248,0.906157,0.143936}%
\pgfsetstrokecolor{currentstroke}%
\pgfsetdash{}{0pt}%
\pgfpathmoveto{\pgfqpoint{1.705558in}{0.804474in}}%
\pgfpathlineto{\pgfqpoint{1.816669in}{0.804474in}}%
\pgfpathlineto{\pgfqpoint{1.927780in}{0.804474in}}%
\pgfusepath{stroke}%
\end{pgfscope}%
\begin{pgfscope}%
\definecolor{textcolor}{rgb}{0.000000,0.000000,0.000000}%
\pgfsetstrokecolor{textcolor}%
\pgfsetfillcolor{textcolor}%
\pgftext[x=2.016669in,y=0.765585in,left,base]{\color{textcolor}{\sffamily\fontsize{8.000000}{9.600000}\selectfont\catcode`\^=\active\def^{\ifmmode\sp\else\^{}\fi}\catcode`\%=\active\def%{\%}$t=0.4$ s}}%
\end{pgfscope}%
\begin{pgfscope}%
\pgfsetbuttcap%
\pgfsetmiterjoin%
\definecolor{currentfill}{rgb}{1.000000,1.000000,1.000000}%
\pgfsetfillcolor{currentfill}%
\pgfsetfillopacity{0.800000}%
\pgfsetlinewidth{1.003750pt}%
\definecolor{currentstroke}{rgb}{0.800000,0.800000,0.800000}%
\pgfsetstrokecolor{currentstroke}%
\pgfsetstrokeopacity{0.800000}%
\pgfsetdash{}{0pt}%
\pgfpathmoveto{\pgfqpoint{5.648812in}{4.336108in}}%
\pgfpathlineto{\pgfqpoint{7.170113in}{4.336108in}}%
\pgfpathquadraticcurveto{\pgfqpoint{7.192335in}{4.336108in}}{\pgfqpoint{7.192335in}{4.358330in}}%
\pgfpathlineto{\pgfqpoint{7.192335in}{4.680552in}}%
\pgfpathquadraticcurveto{\pgfqpoint{7.192335in}{4.702774in}}{\pgfqpoint{7.170113in}{4.702774in}}%
\pgfpathlineto{\pgfqpoint{5.648812in}{4.702774in}}%
\pgfpathquadraticcurveto{\pgfqpoint{5.626590in}{4.702774in}}{\pgfqpoint{5.626590in}{4.680552in}}%
\pgfpathlineto{\pgfqpoint{5.626590in}{4.358330in}}%
\pgfpathquadraticcurveto{\pgfqpoint{5.626590in}{4.336108in}}{\pgfqpoint{5.648812in}{4.336108in}}%
\pgfpathlineto{\pgfqpoint{5.648812in}{4.336108in}}%
\pgfpathclose%
\pgfusepath{stroke,fill}%
\end{pgfscope}%
\begin{pgfscope}%
\pgfsetrectcap%
\pgfsetroundjoin%
\pgfsetlinewidth{1.806750pt}%
\definecolor{currentstroke}{rgb}{0.200000,0.200000,0.200000}%
\pgfsetstrokecolor{currentstroke}%
\pgfsetdash{}{0pt}%
\pgfpathmoveto{\pgfqpoint{5.671035in}{4.612800in}}%
\pgfpathlineto{\pgfqpoint{5.782146in}{4.612800in}}%
\pgfpathlineto{\pgfqpoint{5.893257in}{4.612800in}}%
\pgfusepath{stroke}%
\end{pgfscope}%
\begin{pgfscope}%
\definecolor{textcolor}{rgb}{0.000000,0.000000,0.000000}%
\pgfsetstrokecolor{textcolor}%
\pgfsetfillcolor{textcolor}%
\pgftext[x=5.982146in,y=4.573912in,left,base]{\color{textcolor}{\sffamily\fontsize{8.000000}{9.600000}\selectfont\catcode`\^=\active\def^{\ifmmode\sp\else\^{}\fi}\catcode`\%=\active\def%{\%}analytical series}}%
\end{pgfscope}%
\begin{pgfscope}%
\pgfsetbuttcap%
\pgfsetroundjoin%
\definecolor{currentfill}{rgb}{1.000000,1.000000,1.000000}%
\pgfsetfillcolor{currentfill}%
\pgfsetlinewidth{1.003750pt}%
\definecolor{currentstroke}{rgb}{0.200000,0.200000,0.200000}%
\pgfsetstrokecolor{currentstroke}%
\pgfsetdash{}{0pt}%
\pgfsys@defobject{currentmarker}{\pgfqpoint{-0.027778in}{-0.027778in}}{\pgfqpoint{0.027778in}{0.027778in}}{%
\pgfpathmoveto{\pgfqpoint{0.000000in}{-0.027778in}}%
\pgfpathcurveto{\pgfqpoint{0.007367in}{-0.027778in}}{\pgfqpoint{0.014433in}{-0.024851in}}{\pgfqpoint{0.019642in}{-0.019642in}}%
\pgfpathcurveto{\pgfqpoint{0.024851in}{-0.014433in}}{\pgfqpoint{0.027778in}{-0.007367in}}{\pgfqpoint{0.027778in}{0.000000in}}%
\pgfpathcurveto{\pgfqpoint{0.027778in}{0.007367in}}{\pgfqpoint{0.024851in}{0.014433in}}{\pgfqpoint{0.019642in}{0.019642in}}%
\pgfpathcurveto{\pgfqpoint{0.014433in}{0.024851in}}{\pgfqpoint{0.007367in}{0.027778in}}{\pgfqpoint{0.000000in}{0.027778in}}%
\pgfpathcurveto{\pgfqpoint{-0.007367in}{0.027778in}}{\pgfqpoint{-0.014433in}{0.024851in}}{\pgfqpoint{-0.019642in}{0.019642in}}%
\pgfpathcurveto{\pgfqpoint{-0.024851in}{0.014433in}}{\pgfqpoint{-0.027778in}{0.007367in}}{\pgfqpoint{-0.027778in}{0.000000in}}%
\pgfpathcurveto{\pgfqpoint{-0.027778in}{-0.007367in}}{\pgfqpoint{-0.024851in}{-0.014433in}}{\pgfqpoint{-0.019642in}{-0.019642in}}%
\pgfpathcurveto{\pgfqpoint{-0.014433in}{-0.024851in}}{\pgfqpoint{-0.007367in}{-0.027778in}}{\pgfqpoint{0.000000in}{-0.027778in}}%
\pgfpathlineto{\pgfqpoint{0.000000in}{-0.027778in}}%
\pgfpathclose%
\pgfusepath{stroke,fill}%
}%
\begin{pgfscope}%
\pgfsys@transformshift{5.782146in}{4.446134in}%
\pgfsys@useobject{currentmarker}{}%
\end{pgfscope}%
\end{pgfscope}%
\begin{pgfscope}%
\definecolor{textcolor}{rgb}{0.000000,0.000000,0.000000}%
\pgfsetstrokecolor{textcolor}%
\pgfsetfillcolor{textcolor}%
\pgftext[x=5.982146in,y=4.407245in,left,base]{\color{textcolor}{\sffamily\fontsize{8.000000}{9.600000}\selectfont\catcode`\^=\active\def^{\ifmmode\sp\else\^{}\fi}\catcode`\%=\active\def%{\%}Q2/Q1 finite element}}%
\end{pgfscope}%
\end{pgfpicture}%
\makeatother%
\endgroup%
%
	}
	\caption{Water-pressure profiles for the Mandel verification problem.  The
		solid curves are the analytical series, and the open circles are the
		finite-element solution sampled along the specimen centerline.}
	\label{fig:mandel_pressure_profiles}
\end{figure}

The one-day SPE1 calculation exercises the coupled Q2/CG--EG path through gas
appearance, component conservation, transfer-coordinate evolution, skeleton
momentum, and the two thermal-subsystem residuals.  Figure
\ref{fig:spe1_one_day_acceptance_history} reports the accepted scalar history
and the applicable numerical gates.  Figure
\ref{fig:spe1_one_day_opm_context} places its final well controls against the
reference schedule.  The latter is schedule context because the first
like-for-like OPM report occurs at day 31.

\begin{figure}[ht]
	\centering
	\resizebox{\textwidth}{!}{%
		%% Creator: Matplotlib, PGF backend
%%
%% To include the figure in your LaTeX document, write
%%   \input{<filename>.pgf}
%%
%% Make sure the required packages are loaded in your preamble
%%   \usepackage{pgf}
%%
%% Also ensure that all the required font packages are loaded; for instance,
%% the lmodern package is sometimes necessary when using math font.
%%   \usepackage{lmodern}
%%
%% Figures using additional raster images can only be included by \input if
%% they are in the same directory as the main LaTeX file. For loading figures
%% from other directories you can use the `import` package
%%   \usepackage{import}
%%
%% and then include the figures with
%%   \import{<path to file>}{<filename>.pgf}
%%
%% Matplotlib used the following preamble
%%   \def\mathdefault#1{#1}
%%   \everymath=\expandafter{\the\everymath\displaystyle}
%%   \IfFileExists{scrextend.sty}{
%%     \usepackage[fontsize=10.000000pt]{scrextend}
%%   }{
%%     \renewcommand{\normalsize}{\fontsize{10.000000}{12.000000}\selectfont}
%%     \normalsize
%%   }
%%
%%   \ifdefined\pdftexversion\else  % non-pdftex case.
%%     \usepackage{fontspec}
%%     \setmainfont{DejaVuSerif.ttf}[Path=\detokenize{/home/john/miniconda3/envs/moose/lib/python3.14/site-packages/matplotlib/mpl-data/fonts/ttf/}]
%%     \setsansfont{Arial.ttf}[Path=\detokenize{/usr/share/fonts/truetype/msttcorefonts/}]
%%     \setmonofont{DejaVuSansMono.ttf}[Path=\detokenize{/home/john/miniconda3/envs/moose/lib/python3.14/site-packages/matplotlib/mpl-data/fonts/ttf/}]
%%   \fi
%%   \makeatletter\@ifpackageloaded{underscore}{}{\usepackage[strings]{underscore}}\makeatother
%%
\begingroup%
\makeatletter%
\begin{pgfpicture}%
\pgfpathrectangle{\pgfpointorigin}{\pgfqpoint{9.691796in}{7.350855in}}%
\pgfusepath{use as bounding box, clip}%
\begin{pgfscope}%
\pgfsetbuttcap%
\pgfsetmiterjoin%
\definecolor{currentfill}{rgb}{1.000000,1.000000,1.000000}%
\pgfsetfillcolor{currentfill}%
\pgfsetlinewidth{0.000000pt}%
\definecolor{currentstroke}{rgb}{1.000000,1.000000,1.000000}%
\pgfsetstrokecolor{currentstroke}%
\pgfsetdash{}{0pt}%
\pgfpathmoveto{\pgfqpoint{0.000000in}{0.000000in}}%
\pgfpathlineto{\pgfqpoint{9.691796in}{0.000000in}}%
\pgfpathlineto{\pgfqpoint{9.691796in}{7.350855in}}%
\pgfpathlineto{\pgfqpoint{0.000000in}{7.350855in}}%
\pgfpathlineto{\pgfqpoint{0.000000in}{0.000000in}}%
\pgfpathclose%
\pgfusepath{fill}%
\end{pgfscope}%
\begin{pgfscope}%
\pgfsetbuttcap%
\pgfsetmiterjoin%
\definecolor{currentfill}{rgb}{1.000000,1.000000,1.000000}%
\pgfsetfillcolor{currentfill}%
\pgfsetlinewidth{0.000000pt}%
\definecolor{currentstroke}{rgb}{0.000000,0.000000,0.000000}%
\pgfsetstrokecolor{currentstroke}%
\pgfsetstrokeopacity{0.000000}%
\pgfsetdash{}{0pt}%
\pgfpathmoveto{\pgfqpoint{0.622513in}{3.740855in}}%
\pgfpathlineto{\pgfqpoint{4.673649in}{3.740855in}}%
\pgfpathlineto{\pgfqpoint{4.673649in}{6.470855in}}%
\pgfpathlineto{\pgfqpoint{0.622513in}{6.470855in}}%
\pgfpathlineto{\pgfqpoint{0.622513in}{3.740855in}}%
\pgfpathclose%
\pgfusepath{fill}%
\end{pgfscope}%
\begin{pgfscope}%
\pgfpathrectangle{\pgfqpoint{0.622513in}{3.740855in}}{\pgfqpoint{4.051136in}{2.730000in}}%
\pgfusepath{clip}%
\pgfsetroundcap%
\pgfsetroundjoin%
\pgfsetlinewidth{0.803000pt}%
\definecolor{currentstroke}{rgb}{0.800000,0.800000,0.800000}%
\pgfsetstrokecolor{currentstroke}%
\pgfsetdash{}{0pt}%
\pgfpathmoveto{\pgfqpoint{1.379544in}{3.740855in}}%
\pgfpathlineto{\pgfqpoint{1.379544in}{6.470855in}}%
\pgfusepath{stroke}%
\end{pgfscope}%
\begin{pgfscope}%
\definecolor{textcolor}{rgb}{0.150000,0.150000,0.150000}%
\pgfsetstrokecolor{textcolor}%
\pgfsetfillcolor{textcolor}%
\pgftext[x=1.379544in,y=3.692243in,,top]{\color{textcolor}{\sffamily\fontsize{10.000000}{12.000000}\selectfont\catcode`\^=\active\def^{\ifmmode\sp\else\^{}\fi}\catcode`\%=\active\def%{\%}5}}%
\end{pgfscope}%
\begin{pgfscope}%
\pgfpathrectangle{\pgfqpoint{0.622513in}{3.740855in}}{\pgfqpoint{4.051136in}{2.730000in}}%
\pgfusepath{clip}%
\pgfsetroundcap%
\pgfsetroundjoin%
\pgfsetlinewidth{0.803000pt}%
\definecolor{currentstroke}{rgb}{0.800000,0.800000,0.800000}%
\pgfsetstrokecolor{currentstroke}%
\pgfsetdash{}{0pt}%
\pgfpathmoveto{\pgfqpoint{2.197955in}{3.740855in}}%
\pgfpathlineto{\pgfqpoint{2.197955in}{6.470855in}}%
\pgfusepath{stroke}%
\end{pgfscope}%
\begin{pgfscope}%
\definecolor{textcolor}{rgb}{0.150000,0.150000,0.150000}%
\pgfsetstrokecolor{textcolor}%
\pgfsetfillcolor{textcolor}%
\pgftext[x=2.197955in,y=3.692243in,,top]{\color{textcolor}{\sffamily\fontsize{10.000000}{12.000000}\selectfont\catcode`\^=\active\def^{\ifmmode\sp\else\^{}\fi}\catcode`\%=\active\def%{\%}10}}%
\end{pgfscope}%
\begin{pgfscope}%
\pgfpathrectangle{\pgfqpoint{0.622513in}{3.740855in}}{\pgfqpoint{4.051136in}{2.730000in}}%
\pgfusepath{clip}%
\pgfsetroundcap%
\pgfsetroundjoin%
\pgfsetlinewidth{0.803000pt}%
\definecolor{currentstroke}{rgb}{0.800000,0.800000,0.800000}%
\pgfsetstrokecolor{currentstroke}%
\pgfsetdash{}{0pt}%
\pgfpathmoveto{\pgfqpoint{3.016366in}{3.740855in}}%
\pgfpathlineto{\pgfqpoint{3.016366in}{6.470855in}}%
\pgfusepath{stroke}%
\end{pgfscope}%
\begin{pgfscope}%
\definecolor{textcolor}{rgb}{0.150000,0.150000,0.150000}%
\pgfsetstrokecolor{textcolor}%
\pgfsetfillcolor{textcolor}%
\pgftext[x=3.016366in,y=3.692243in,,top]{\color{textcolor}{\sffamily\fontsize{10.000000}{12.000000}\selectfont\catcode`\^=\active\def^{\ifmmode\sp\else\^{}\fi}\catcode`\%=\active\def%{\%}15}}%
\end{pgfscope}%
\begin{pgfscope}%
\pgfpathrectangle{\pgfqpoint{0.622513in}{3.740855in}}{\pgfqpoint{4.051136in}{2.730000in}}%
\pgfusepath{clip}%
\pgfsetroundcap%
\pgfsetroundjoin%
\pgfsetlinewidth{0.803000pt}%
\definecolor{currentstroke}{rgb}{0.800000,0.800000,0.800000}%
\pgfsetstrokecolor{currentstroke}%
\pgfsetdash{}{0pt}%
\pgfpathmoveto{\pgfqpoint{3.834778in}{3.740855in}}%
\pgfpathlineto{\pgfqpoint{3.834778in}{6.470855in}}%
\pgfusepath{stroke}%
\end{pgfscope}%
\begin{pgfscope}%
\definecolor{textcolor}{rgb}{0.150000,0.150000,0.150000}%
\pgfsetstrokecolor{textcolor}%
\pgfsetfillcolor{textcolor}%
\pgftext[x=3.834778in,y=3.692243in,,top]{\color{textcolor}{\sffamily\fontsize{10.000000}{12.000000}\selectfont\catcode`\^=\active\def^{\ifmmode\sp\else\^{}\fi}\catcode`\%=\active\def%{\%}20}}%
\end{pgfscope}%
\begin{pgfscope}%
\pgfpathrectangle{\pgfqpoint{0.622513in}{3.740855in}}{\pgfqpoint{4.051136in}{2.730000in}}%
\pgfusepath{clip}%
\pgfsetroundcap%
\pgfsetroundjoin%
\pgfsetlinewidth{0.803000pt}%
\definecolor{currentstroke}{rgb}{0.800000,0.800000,0.800000}%
\pgfsetstrokecolor{currentstroke}%
\pgfsetdash{}{0pt}%
\pgfpathmoveto{\pgfqpoint{4.653189in}{3.740855in}}%
\pgfpathlineto{\pgfqpoint{4.653189in}{6.470855in}}%
\pgfusepath{stroke}%
\end{pgfscope}%
\begin{pgfscope}%
\definecolor{textcolor}{rgb}{0.150000,0.150000,0.150000}%
\pgfsetstrokecolor{textcolor}%
\pgfsetfillcolor{textcolor}%
\pgftext[x=4.653189in,y=3.692243in,,top]{\color{textcolor}{\sffamily\fontsize{10.000000}{12.000000}\selectfont\catcode`\^=\active\def^{\ifmmode\sp\else\^{}\fi}\catcode`\%=\active\def%{\%}25}}%
\end{pgfscope}%
\begin{pgfscope}%
\definecolor{textcolor}{rgb}{0.150000,0.150000,0.150000}%
\pgfsetstrokecolor{textcolor}%
\pgfsetfillcolor{textcolor}%
\pgftext[x=2.648081in,y=3.506344in,,top]{\color{textcolor}{\sffamily\fontsize{10.000000}{12.000000}\selectfont\catcode`\^=\active\def^{\ifmmode\sp\else\^{}\fi}\catcode`\%=\active\def%{\%}time [h]}}%
\end{pgfscope}%
\begin{pgfscope}%
\pgfpathrectangle{\pgfqpoint{0.622513in}{3.740855in}}{\pgfqpoint{4.051136in}{2.730000in}}%
\pgfusepath{clip}%
\pgfsetroundcap%
\pgfsetroundjoin%
\pgfsetlinewidth{0.803000pt}%
\definecolor{currentstroke}{rgb}{0.800000,0.800000,0.800000}%
\pgfsetstrokecolor{currentstroke}%
\pgfsetdash{}{0pt}%
\pgfpathmoveto{\pgfqpoint{0.622513in}{3.864699in}}%
\pgfpathlineto{\pgfqpoint{4.673649in}{3.864699in}}%
\pgfusepath{stroke}%
\end{pgfscope}%
\begin{pgfscope}%
\definecolor{textcolor}{rgb}{0.150000,0.150000,0.150000}%
\pgfsetstrokecolor{textcolor}%
\pgfsetfillcolor{textcolor}%
\pgftext[x=0.380827in, y=3.814142in, left, base]{\color{textcolor}{\sffamily\fontsize{10.000000}{12.000000}\selectfont\catcode`\^=\active\def^{\ifmmode\sp\else\^{}\fi}\catcode`\%=\active\def%{\%}0.0}}%
\end{pgfscope}%
\begin{pgfscope}%
\pgfpathrectangle{\pgfqpoint{0.622513in}{3.740855in}}{\pgfqpoint{4.051136in}{2.730000in}}%
\pgfusepath{clip}%
\pgfsetroundcap%
\pgfsetroundjoin%
\pgfsetlinewidth{0.803000pt}%
\definecolor{currentstroke}{rgb}{0.800000,0.800000,0.800000}%
\pgfsetstrokecolor{currentstroke}%
\pgfsetdash{}{0pt}%
\pgfpathmoveto{\pgfqpoint{0.622513in}{4.415441in}}%
\pgfpathlineto{\pgfqpoint{4.673649in}{4.415441in}}%
\pgfusepath{stroke}%
\end{pgfscope}%
\begin{pgfscope}%
\definecolor{textcolor}{rgb}{0.150000,0.150000,0.150000}%
\pgfsetstrokecolor{textcolor}%
\pgfsetfillcolor{textcolor}%
\pgftext[x=0.380827in, y=4.364883in, left, base]{\color{textcolor}{\sffamily\fontsize{10.000000}{12.000000}\selectfont\catcode`\^=\active\def^{\ifmmode\sp\else\^{}\fi}\catcode`\%=\active\def%{\%}0.1}}%
\end{pgfscope}%
\begin{pgfscope}%
\pgfpathrectangle{\pgfqpoint{0.622513in}{3.740855in}}{\pgfqpoint{4.051136in}{2.730000in}}%
\pgfusepath{clip}%
\pgfsetroundcap%
\pgfsetroundjoin%
\pgfsetlinewidth{0.803000pt}%
\definecolor{currentstroke}{rgb}{0.800000,0.800000,0.800000}%
\pgfsetstrokecolor{currentstroke}%
\pgfsetdash{}{0pt}%
\pgfpathmoveto{\pgfqpoint{0.622513in}{4.966182in}}%
\pgfpathlineto{\pgfqpoint{4.673649in}{4.966182in}}%
\pgfusepath{stroke}%
\end{pgfscope}%
\begin{pgfscope}%
\definecolor{textcolor}{rgb}{0.150000,0.150000,0.150000}%
\pgfsetstrokecolor{textcolor}%
\pgfsetfillcolor{textcolor}%
\pgftext[x=0.380827in, y=4.915625in, left, base]{\color{textcolor}{\sffamily\fontsize{10.000000}{12.000000}\selectfont\catcode`\^=\active\def^{\ifmmode\sp\else\^{}\fi}\catcode`\%=\active\def%{\%}0.2}}%
\end{pgfscope}%
\begin{pgfscope}%
\pgfpathrectangle{\pgfqpoint{0.622513in}{3.740855in}}{\pgfqpoint{4.051136in}{2.730000in}}%
\pgfusepath{clip}%
\pgfsetroundcap%
\pgfsetroundjoin%
\pgfsetlinewidth{0.803000pt}%
\definecolor{currentstroke}{rgb}{0.800000,0.800000,0.800000}%
\pgfsetstrokecolor{currentstroke}%
\pgfsetdash{}{0pt}%
\pgfpathmoveto{\pgfqpoint{0.622513in}{5.516924in}}%
\pgfpathlineto{\pgfqpoint{4.673649in}{5.516924in}}%
\pgfusepath{stroke}%
\end{pgfscope}%
\begin{pgfscope}%
\definecolor{textcolor}{rgb}{0.150000,0.150000,0.150000}%
\pgfsetstrokecolor{textcolor}%
\pgfsetfillcolor{textcolor}%
\pgftext[x=0.380827in, y=5.466367in, left, base]{\color{textcolor}{\sffamily\fontsize{10.000000}{12.000000}\selectfont\catcode`\^=\active\def^{\ifmmode\sp\else\^{}\fi}\catcode`\%=\active\def%{\%}0.3}}%
\end{pgfscope}%
\begin{pgfscope}%
\pgfpathrectangle{\pgfqpoint{0.622513in}{3.740855in}}{\pgfqpoint{4.051136in}{2.730000in}}%
\pgfusepath{clip}%
\pgfsetroundcap%
\pgfsetroundjoin%
\pgfsetlinewidth{0.803000pt}%
\definecolor{currentstroke}{rgb}{0.800000,0.800000,0.800000}%
\pgfsetstrokecolor{currentstroke}%
\pgfsetdash{}{0pt}%
\pgfpathmoveto{\pgfqpoint{0.622513in}{6.067666in}}%
\pgfpathlineto{\pgfqpoint{4.673649in}{6.067666in}}%
\pgfusepath{stroke}%
\end{pgfscope}%
\begin{pgfscope}%
\definecolor{textcolor}{rgb}{0.150000,0.150000,0.150000}%
\pgfsetstrokecolor{textcolor}%
\pgfsetfillcolor{textcolor}%
\pgftext[x=0.380827in, y=6.017108in, left, base]{\color{textcolor}{\sffamily\fontsize{10.000000}{12.000000}\selectfont\catcode`\^=\active\def^{\ifmmode\sp\else\^{}\fi}\catcode`\%=\active\def%{\%}0.4}}%
\end{pgfscope}%
\begin{pgfscope}%
\definecolor{textcolor}{rgb}{0.150000,0.150000,0.150000}%
\pgfsetstrokecolor{textcolor}%
\pgfsetfillcolor{textcolor}%
\pgftext[x=0.325272in,y=5.105855in,,bottom,rotate=90.000000]{\color{textcolor}{\sffamily\fontsize{10.000000}{12.000000}\selectfont\catcode`\^=\active\def^{\ifmmode\sp\else\^{}\fi}\catcode`\%=\active\def%{\%}gas saturation [-]}}%
\end{pgfscope}%
\begin{pgfscope}%
\pgfpathrectangle{\pgfqpoint{0.622513in}{3.740855in}}{\pgfqpoint{4.051136in}{2.730000in}}%
\pgfusepath{clip}%
\pgfsetroundcap%
\pgfsetroundjoin%
\pgfsetlinewidth{1.505625pt}%
\definecolor{currentstroke}{rgb}{0.121569,0.466667,0.705882}%
\pgfsetstrokecolor{currentstroke}%
\pgfsetdash{}{0pt}%
\pgfpathmoveto{\pgfqpoint{0.806656in}{3.864946in}}%
\pgfpathlineto{\pgfqpoint{1.052179in}{3.865128in}}%
\pgfpathlineto{\pgfqpoint{1.297703in}{3.865307in}}%
\pgfpathlineto{\pgfqpoint{1.543226in}{3.865485in}}%
\pgfpathlineto{\pgfqpoint{1.788749in}{3.865667in}}%
\pgfpathlineto{\pgfqpoint{2.034273in}{3.865851in}}%
\pgfpathlineto{\pgfqpoint{2.279796in}{3.866041in}}%
\pgfpathlineto{\pgfqpoint{2.525320in}{3.866240in}}%
\pgfpathlineto{\pgfqpoint{2.770843in}{3.866447in}}%
\pgfpathlineto{\pgfqpoint{3.016366in}{3.866665in}}%
\pgfpathlineto{\pgfqpoint{3.261890in}{3.866894in}}%
\pgfpathlineto{\pgfqpoint{3.507413in}{3.867139in}}%
\pgfpathlineto{\pgfqpoint{3.752937in}{3.867400in}}%
\pgfpathlineto{\pgfqpoint{3.998460in}{3.867683in}}%
\pgfpathlineto{\pgfqpoint{4.243984in}{3.867992in}}%
\pgfpathlineto{\pgfqpoint{4.489507in}{3.868340in}}%
\pgfusepath{stroke}%
\end{pgfscope}%
\begin{pgfscope}%
\pgfpathrectangle{\pgfqpoint{0.622513in}{3.740855in}}{\pgfqpoint{4.051136in}{2.730000in}}%
\pgfusepath{clip}%
\pgfsetbuttcap%
\pgfsetroundjoin%
\definecolor{currentfill}{rgb}{0.121569,0.466667,0.705882}%
\pgfsetfillcolor{currentfill}%
\pgfsetlinewidth{1.003750pt}%
\definecolor{currentstroke}{rgb}{0.121569,0.466667,0.705882}%
\pgfsetstrokecolor{currentstroke}%
\pgfsetdash{}{0pt}%
\pgfsys@defobject{currentmarker}{\pgfqpoint{-0.041667in}{-0.041667in}}{\pgfqpoint{0.041667in}{0.041667in}}{%
\pgfpathmoveto{\pgfqpoint{0.000000in}{-0.041667in}}%
\pgfpathcurveto{\pgfqpoint{0.011050in}{-0.041667in}}{\pgfqpoint{0.021649in}{-0.037276in}}{\pgfqpoint{0.029463in}{-0.029463in}}%
\pgfpathcurveto{\pgfqpoint{0.037276in}{-0.021649in}}{\pgfqpoint{0.041667in}{-0.011050in}}{\pgfqpoint{0.041667in}{0.000000in}}%
\pgfpathcurveto{\pgfqpoint{0.041667in}{0.011050in}}{\pgfqpoint{0.037276in}{0.021649in}}{\pgfqpoint{0.029463in}{0.029463in}}%
\pgfpathcurveto{\pgfqpoint{0.021649in}{0.037276in}}{\pgfqpoint{0.011050in}{0.041667in}}{\pgfqpoint{0.000000in}{0.041667in}}%
\pgfpathcurveto{\pgfqpoint{-0.011050in}{0.041667in}}{\pgfqpoint{-0.021649in}{0.037276in}}{\pgfqpoint{-0.029463in}{0.029463in}}%
\pgfpathcurveto{\pgfqpoint{-0.037276in}{0.021649in}}{\pgfqpoint{-0.041667in}{0.011050in}}{\pgfqpoint{-0.041667in}{0.000000in}}%
\pgfpathcurveto{\pgfqpoint{-0.041667in}{-0.011050in}}{\pgfqpoint{-0.037276in}{-0.021649in}}{\pgfqpoint{-0.029463in}{-0.029463in}}%
\pgfpathcurveto{\pgfqpoint{-0.021649in}{-0.037276in}}{\pgfqpoint{-0.011050in}{-0.041667in}}{\pgfqpoint{0.000000in}{-0.041667in}}%
\pgfpathlineto{\pgfqpoint{0.000000in}{-0.041667in}}%
\pgfpathclose%
\pgfusepath{stroke,fill}%
}%
\begin{pgfscope}%
\pgfsys@transformshift{0.806656in}{3.864946in}%
\pgfsys@useobject{currentmarker}{}%
\end{pgfscope}%
\begin{pgfscope}%
\pgfsys@transformshift{1.052179in}{3.865128in}%
\pgfsys@useobject{currentmarker}{}%
\end{pgfscope}%
\begin{pgfscope}%
\pgfsys@transformshift{1.297703in}{3.865307in}%
\pgfsys@useobject{currentmarker}{}%
\end{pgfscope}%
\begin{pgfscope}%
\pgfsys@transformshift{1.543226in}{3.865485in}%
\pgfsys@useobject{currentmarker}{}%
\end{pgfscope}%
\begin{pgfscope}%
\pgfsys@transformshift{1.788749in}{3.865667in}%
\pgfsys@useobject{currentmarker}{}%
\end{pgfscope}%
\begin{pgfscope}%
\pgfsys@transformshift{2.034273in}{3.865851in}%
\pgfsys@useobject{currentmarker}{}%
\end{pgfscope}%
\begin{pgfscope}%
\pgfsys@transformshift{2.279796in}{3.866041in}%
\pgfsys@useobject{currentmarker}{}%
\end{pgfscope}%
\begin{pgfscope}%
\pgfsys@transformshift{2.525320in}{3.866240in}%
\pgfsys@useobject{currentmarker}{}%
\end{pgfscope}%
\begin{pgfscope}%
\pgfsys@transformshift{2.770843in}{3.866447in}%
\pgfsys@useobject{currentmarker}{}%
\end{pgfscope}%
\begin{pgfscope}%
\pgfsys@transformshift{3.016366in}{3.866665in}%
\pgfsys@useobject{currentmarker}{}%
\end{pgfscope}%
\begin{pgfscope}%
\pgfsys@transformshift{3.261890in}{3.866894in}%
\pgfsys@useobject{currentmarker}{}%
\end{pgfscope}%
\begin{pgfscope}%
\pgfsys@transformshift{3.507413in}{3.867139in}%
\pgfsys@useobject{currentmarker}{}%
\end{pgfscope}%
\begin{pgfscope}%
\pgfsys@transformshift{3.752937in}{3.867400in}%
\pgfsys@useobject{currentmarker}{}%
\end{pgfscope}%
\begin{pgfscope}%
\pgfsys@transformshift{3.998460in}{3.867683in}%
\pgfsys@useobject{currentmarker}{}%
\end{pgfscope}%
\begin{pgfscope}%
\pgfsys@transformshift{4.243984in}{3.867992in}%
\pgfsys@useobject{currentmarker}{}%
\end{pgfscope}%
\begin{pgfscope}%
\pgfsys@transformshift{4.489507in}{3.868340in}%
\pgfsys@useobject{currentmarker}{}%
\end{pgfscope}%
\end{pgfscope}%
\begin{pgfscope}%
\pgfpathrectangle{\pgfqpoint{0.622513in}{3.740855in}}{\pgfqpoint{4.051136in}{2.730000in}}%
\pgfusepath{clip}%
\pgfsetroundcap%
\pgfsetroundjoin%
\pgfsetlinewidth{1.505625pt}%
\definecolor{currentstroke}{rgb}{1.000000,0.498039,0.054902}%
\pgfsetstrokecolor{currentstroke}%
\pgfsetdash{}{0pt}%
\pgfpathmoveto{\pgfqpoint{0.806656in}{3.917477in}}%
\pgfpathlineto{\pgfqpoint{1.052179in}{3.970575in}}%
\pgfpathlineto{\pgfqpoint{1.297703in}{4.020526in}}%
\pgfpathlineto{\pgfqpoint{1.543226in}{4.065222in}}%
\pgfpathlineto{\pgfqpoint{1.788749in}{4.110784in}}%
\pgfpathlineto{\pgfqpoint{2.034273in}{4.161984in}}%
\pgfpathlineto{\pgfqpoint{2.279796in}{4.243310in}}%
\pgfpathlineto{\pgfqpoint{2.525320in}{4.341532in}}%
\pgfpathlineto{\pgfqpoint{2.770843in}{4.455715in}}%
\pgfpathlineto{\pgfqpoint{3.016366in}{4.588093in}}%
\pgfpathlineto{\pgfqpoint{3.261890in}{4.744944in}}%
\pgfpathlineto{\pgfqpoint{3.507413in}{4.929622in}}%
\pgfpathlineto{\pgfqpoint{3.752937in}{5.144662in}}%
\pgfpathlineto{\pgfqpoint{3.998460in}{5.394077in}}%
\pgfpathlineto{\pgfqpoint{4.243984in}{5.798399in}}%
\pgfpathlineto{\pgfqpoint{4.489507in}{6.346764in}}%
\pgfusepath{stroke}%
\end{pgfscope}%
\begin{pgfscope}%
\pgfpathrectangle{\pgfqpoint{0.622513in}{3.740855in}}{\pgfqpoint{4.051136in}{2.730000in}}%
\pgfusepath{clip}%
\pgfsetbuttcap%
\pgfsetroundjoin%
\definecolor{currentfill}{rgb}{1.000000,0.498039,0.054902}%
\pgfsetfillcolor{currentfill}%
\pgfsetlinewidth{1.003750pt}%
\definecolor{currentstroke}{rgb}{1.000000,0.498039,0.054902}%
\pgfsetstrokecolor{currentstroke}%
\pgfsetdash{}{0pt}%
\pgfsys@defobject{currentmarker}{\pgfqpoint{-0.041667in}{-0.041667in}}{\pgfqpoint{0.041667in}{0.041667in}}{%
\pgfpathmoveto{\pgfqpoint{0.000000in}{-0.041667in}}%
\pgfpathcurveto{\pgfqpoint{0.011050in}{-0.041667in}}{\pgfqpoint{0.021649in}{-0.037276in}}{\pgfqpoint{0.029463in}{-0.029463in}}%
\pgfpathcurveto{\pgfqpoint{0.037276in}{-0.021649in}}{\pgfqpoint{0.041667in}{-0.011050in}}{\pgfqpoint{0.041667in}{0.000000in}}%
\pgfpathcurveto{\pgfqpoint{0.041667in}{0.011050in}}{\pgfqpoint{0.037276in}{0.021649in}}{\pgfqpoint{0.029463in}{0.029463in}}%
\pgfpathcurveto{\pgfqpoint{0.021649in}{0.037276in}}{\pgfqpoint{0.011050in}{0.041667in}}{\pgfqpoint{0.000000in}{0.041667in}}%
\pgfpathcurveto{\pgfqpoint{-0.011050in}{0.041667in}}{\pgfqpoint{-0.021649in}{0.037276in}}{\pgfqpoint{-0.029463in}{0.029463in}}%
\pgfpathcurveto{\pgfqpoint{-0.037276in}{0.021649in}}{\pgfqpoint{-0.041667in}{0.011050in}}{\pgfqpoint{-0.041667in}{0.000000in}}%
\pgfpathcurveto{\pgfqpoint{-0.041667in}{-0.011050in}}{\pgfqpoint{-0.037276in}{-0.021649in}}{\pgfqpoint{-0.029463in}{-0.029463in}}%
\pgfpathcurveto{\pgfqpoint{-0.021649in}{-0.037276in}}{\pgfqpoint{-0.011050in}{-0.041667in}}{\pgfqpoint{0.000000in}{-0.041667in}}%
\pgfpathlineto{\pgfqpoint{0.000000in}{-0.041667in}}%
\pgfpathclose%
\pgfusepath{stroke,fill}%
}%
\begin{pgfscope}%
\pgfsys@transformshift{0.806656in}{3.917477in}%
\pgfsys@useobject{currentmarker}{}%
\end{pgfscope}%
\begin{pgfscope}%
\pgfsys@transformshift{1.052179in}{3.970575in}%
\pgfsys@useobject{currentmarker}{}%
\end{pgfscope}%
\begin{pgfscope}%
\pgfsys@transformshift{1.297703in}{4.020526in}%
\pgfsys@useobject{currentmarker}{}%
\end{pgfscope}%
\begin{pgfscope}%
\pgfsys@transformshift{1.543226in}{4.065222in}%
\pgfsys@useobject{currentmarker}{}%
\end{pgfscope}%
\begin{pgfscope}%
\pgfsys@transformshift{1.788749in}{4.110784in}%
\pgfsys@useobject{currentmarker}{}%
\end{pgfscope}%
\begin{pgfscope}%
\pgfsys@transformshift{2.034273in}{4.161984in}%
\pgfsys@useobject{currentmarker}{}%
\end{pgfscope}%
\begin{pgfscope}%
\pgfsys@transformshift{2.279796in}{4.243310in}%
\pgfsys@useobject{currentmarker}{}%
\end{pgfscope}%
\begin{pgfscope}%
\pgfsys@transformshift{2.525320in}{4.341532in}%
\pgfsys@useobject{currentmarker}{}%
\end{pgfscope}%
\begin{pgfscope}%
\pgfsys@transformshift{2.770843in}{4.455715in}%
\pgfsys@useobject{currentmarker}{}%
\end{pgfscope}%
\begin{pgfscope}%
\pgfsys@transformshift{3.016366in}{4.588093in}%
\pgfsys@useobject{currentmarker}{}%
\end{pgfscope}%
\begin{pgfscope}%
\pgfsys@transformshift{3.261890in}{4.744944in}%
\pgfsys@useobject{currentmarker}{}%
\end{pgfscope}%
\begin{pgfscope}%
\pgfsys@transformshift{3.507413in}{4.929622in}%
\pgfsys@useobject{currentmarker}{}%
\end{pgfscope}%
\begin{pgfscope}%
\pgfsys@transformshift{3.752937in}{5.144662in}%
\pgfsys@useobject{currentmarker}{}%
\end{pgfscope}%
\begin{pgfscope}%
\pgfsys@transformshift{3.998460in}{5.394077in}%
\pgfsys@useobject{currentmarker}{}%
\end{pgfscope}%
\begin{pgfscope}%
\pgfsys@transformshift{4.243984in}{5.798399in}%
\pgfsys@useobject{currentmarker}{}%
\end{pgfscope}%
\begin{pgfscope}%
\pgfsys@transformshift{4.489507in}{6.346764in}%
\pgfsys@useobject{currentmarker}{}%
\end{pgfscope}%
\end{pgfscope}%
\begin{pgfscope}%
\pgfsetrectcap%
\pgfsetmiterjoin%
\pgfsetlinewidth{1.003750pt}%
\definecolor{currentstroke}{rgb}{0.800000,0.800000,0.800000}%
\pgfsetstrokecolor{currentstroke}%
\pgfsetdash{}{0pt}%
\pgfpathmoveto{\pgfqpoint{0.622513in}{3.740855in}}%
\pgfpathlineto{\pgfqpoint{0.622513in}{6.470855in}}%
\pgfusepath{stroke}%
\end{pgfscope}%
\begin{pgfscope}%
\pgfsetrectcap%
\pgfsetmiterjoin%
\pgfsetlinewidth{1.003750pt}%
\definecolor{currentstroke}{rgb}{0.800000,0.800000,0.800000}%
\pgfsetstrokecolor{currentstroke}%
\pgfsetdash{}{0pt}%
\pgfpathmoveto{\pgfqpoint{4.673649in}{3.740855in}}%
\pgfpathlineto{\pgfqpoint{4.673649in}{6.470855in}}%
\pgfusepath{stroke}%
\end{pgfscope}%
\begin{pgfscope}%
\pgfsetrectcap%
\pgfsetmiterjoin%
\pgfsetlinewidth{1.003750pt}%
\definecolor{currentstroke}{rgb}{0.800000,0.800000,0.800000}%
\pgfsetstrokecolor{currentstroke}%
\pgfsetdash{}{0pt}%
\pgfpathmoveto{\pgfqpoint{0.622513in}{3.740855in}}%
\pgfpathlineto{\pgfqpoint{4.673649in}{3.740855in}}%
\pgfusepath{stroke}%
\end{pgfscope}%
\begin{pgfscope}%
\pgfsetrectcap%
\pgfsetmiterjoin%
\pgfsetlinewidth{1.003750pt}%
\definecolor{currentstroke}{rgb}{0.800000,0.800000,0.800000}%
\pgfsetstrokecolor{currentstroke}%
\pgfsetdash{}{0pt}%
\pgfpathmoveto{\pgfqpoint{0.622513in}{6.470855in}}%
\pgfpathlineto{\pgfqpoint{4.673649in}{6.470855in}}%
\pgfusepath{stroke}%
\end{pgfscope}%
\begin{pgfscope}%
\definecolor{textcolor}{rgb}{0.150000,0.150000,0.150000}%
\pgfsetstrokecolor{textcolor}%
\pgfsetfillcolor{textcolor}%
\pgftext[x=2.648081in,y=6.554188in,,base]{\color{textcolor}{\sffamily\fontsize{12.000000}{14.400000}\selectfont\catcode`\^=\active\def^{\ifmmode\sp\else\^{}\fi}\catcode`\%=\active\def%{\%}Gas-phase appearance}}%
\end{pgfscope}%
\begin{pgfscope}%
\pgfsetroundcap%
\pgfsetroundjoin%
\pgfsetlinewidth{1.505625pt}%
\definecolor{currentstroke}{rgb}{0.121569,0.466667,0.705882}%
\pgfsetstrokecolor{currentstroke}%
\pgfsetdash{}{0pt}%
\pgfpathmoveto{\pgfqpoint{0.747513in}{6.293351in}}%
\pgfpathlineto{\pgfqpoint{0.886402in}{6.293351in}}%
\pgfpathlineto{\pgfqpoint{1.025291in}{6.293351in}}%
\pgfusepath{stroke}%
\end{pgfscope}%
\begin{pgfscope}%
\pgfsetbuttcap%
\pgfsetroundjoin%
\definecolor{currentfill}{rgb}{0.121569,0.466667,0.705882}%
\pgfsetfillcolor{currentfill}%
\pgfsetlinewidth{1.003750pt}%
\definecolor{currentstroke}{rgb}{0.121569,0.466667,0.705882}%
\pgfsetstrokecolor{currentstroke}%
\pgfsetdash{}{0pt}%
\pgfsys@defobject{currentmarker}{\pgfqpoint{-0.041667in}{-0.041667in}}{\pgfqpoint{0.041667in}{0.041667in}}{%
\pgfpathmoveto{\pgfqpoint{0.000000in}{-0.041667in}}%
\pgfpathcurveto{\pgfqpoint{0.011050in}{-0.041667in}}{\pgfqpoint{0.021649in}{-0.037276in}}{\pgfqpoint{0.029463in}{-0.029463in}}%
\pgfpathcurveto{\pgfqpoint{0.037276in}{-0.021649in}}{\pgfqpoint{0.041667in}{-0.011050in}}{\pgfqpoint{0.041667in}{0.000000in}}%
\pgfpathcurveto{\pgfqpoint{0.041667in}{0.011050in}}{\pgfqpoint{0.037276in}{0.021649in}}{\pgfqpoint{0.029463in}{0.029463in}}%
\pgfpathcurveto{\pgfqpoint{0.021649in}{0.037276in}}{\pgfqpoint{0.011050in}{0.041667in}}{\pgfqpoint{0.000000in}{0.041667in}}%
\pgfpathcurveto{\pgfqpoint{-0.011050in}{0.041667in}}{\pgfqpoint{-0.021649in}{0.037276in}}{\pgfqpoint{-0.029463in}{0.029463in}}%
\pgfpathcurveto{\pgfqpoint{-0.037276in}{0.021649in}}{\pgfqpoint{-0.041667in}{0.011050in}}{\pgfqpoint{-0.041667in}{0.000000in}}%
\pgfpathcurveto{\pgfqpoint{-0.041667in}{-0.011050in}}{\pgfqpoint{-0.037276in}{-0.021649in}}{\pgfqpoint{-0.029463in}{-0.029463in}}%
\pgfpathcurveto{\pgfqpoint{-0.021649in}{-0.037276in}}{\pgfqpoint{-0.011050in}{-0.041667in}}{\pgfqpoint{0.000000in}{-0.041667in}}%
\pgfpathlineto{\pgfqpoint{0.000000in}{-0.041667in}}%
\pgfpathclose%
\pgfusepath{stroke,fill}%
}%
\begin{pgfscope}%
\pgfsys@transformshift{0.886402in}{6.293351in}%
\pgfsys@useobject{currentmarker}{}%
\end{pgfscope}%
\end{pgfscope}%
\begin{pgfscope}%
\definecolor{textcolor}{rgb}{0.150000,0.150000,0.150000}%
\pgfsetstrokecolor{textcolor}%
\pgfsetfillcolor{textcolor}%
\pgftext[x=1.136402in,y=6.244740in,left,base]{\color{textcolor}{\sffamily\fontsize{10.000000}{12.000000}\selectfont\catcode`\^=\active\def^{\ifmmode\sp\else\^{}\fi}\catcode`\%=\active\def%{\%}domain average}}%
\end{pgfscope}%
\begin{pgfscope}%
\pgfsetroundcap%
\pgfsetroundjoin%
\pgfsetlinewidth{1.505625pt}%
\definecolor{currentstroke}{rgb}{1.000000,0.498039,0.054902}%
\pgfsetstrokecolor{currentstroke}%
\pgfsetdash{}{0pt}%
\pgfpathmoveto{\pgfqpoint{0.747513in}{6.093562in}}%
\pgfpathlineto{\pgfqpoint{0.886402in}{6.093562in}}%
\pgfpathlineto{\pgfqpoint{1.025291in}{6.093562in}}%
\pgfusepath{stroke}%
\end{pgfscope}%
\begin{pgfscope}%
\pgfsetbuttcap%
\pgfsetroundjoin%
\definecolor{currentfill}{rgb}{1.000000,0.498039,0.054902}%
\pgfsetfillcolor{currentfill}%
\pgfsetlinewidth{1.003750pt}%
\definecolor{currentstroke}{rgb}{1.000000,0.498039,0.054902}%
\pgfsetstrokecolor{currentstroke}%
\pgfsetdash{}{0pt}%
\pgfsys@defobject{currentmarker}{\pgfqpoint{-0.041667in}{-0.041667in}}{\pgfqpoint{0.041667in}{0.041667in}}{%
\pgfpathmoveto{\pgfqpoint{0.000000in}{-0.041667in}}%
\pgfpathcurveto{\pgfqpoint{0.011050in}{-0.041667in}}{\pgfqpoint{0.021649in}{-0.037276in}}{\pgfqpoint{0.029463in}{-0.029463in}}%
\pgfpathcurveto{\pgfqpoint{0.037276in}{-0.021649in}}{\pgfqpoint{0.041667in}{-0.011050in}}{\pgfqpoint{0.041667in}{0.000000in}}%
\pgfpathcurveto{\pgfqpoint{0.041667in}{0.011050in}}{\pgfqpoint{0.037276in}{0.021649in}}{\pgfqpoint{0.029463in}{0.029463in}}%
\pgfpathcurveto{\pgfqpoint{0.021649in}{0.037276in}}{\pgfqpoint{0.011050in}{0.041667in}}{\pgfqpoint{0.000000in}{0.041667in}}%
\pgfpathcurveto{\pgfqpoint{-0.011050in}{0.041667in}}{\pgfqpoint{-0.021649in}{0.037276in}}{\pgfqpoint{-0.029463in}{0.029463in}}%
\pgfpathcurveto{\pgfqpoint{-0.037276in}{0.021649in}}{\pgfqpoint{-0.041667in}{0.011050in}}{\pgfqpoint{-0.041667in}{0.000000in}}%
\pgfpathcurveto{\pgfqpoint{-0.041667in}{-0.011050in}}{\pgfqpoint{-0.037276in}{-0.021649in}}{\pgfqpoint{-0.029463in}{-0.029463in}}%
\pgfpathcurveto{\pgfqpoint{-0.021649in}{-0.037276in}}{\pgfqpoint{-0.011050in}{-0.041667in}}{\pgfqpoint{0.000000in}{-0.041667in}}%
\pgfpathlineto{\pgfqpoint{0.000000in}{-0.041667in}}%
\pgfpathclose%
\pgfusepath{stroke,fill}%
}%
\begin{pgfscope}%
\pgfsys@transformshift{0.886402in}{6.093562in}%
\pgfsys@useobject{currentmarker}{}%
\end{pgfscope}%
\end{pgfscope}%
\begin{pgfscope}%
\definecolor{textcolor}{rgb}{0.150000,0.150000,0.150000}%
\pgfsetstrokecolor{textcolor}%
\pgfsetfillcolor{textcolor}%
\pgftext[x=1.136402in,y=6.044951in,left,base]{\color{textcolor}{\sffamily\fontsize{10.000000}{12.000000}\selectfont\catcode`\^=\active\def^{\ifmmode\sp\else\^{}\fi}\catcode`\%=\active\def%{\%}domain maximum}}%
\end{pgfscope}%
\begin{pgfscope}%
\pgfsetbuttcap%
\pgfsetmiterjoin%
\definecolor{currentfill}{rgb}{1.000000,1.000000,1.000000}%
\pgfsetfillcolor{currentfill}%
\pgfsetlinewidth{0.000000pt}%
\definecolor{currentstroke}{rgb}{0.000000,0.000000,0.000000}%
\pgfsetstrokecolor{currentstroke}%
\pgfsetstrokeopacity{0.000000}%
\pgfsetdash{}{0pt}%
\pgfpathmoveto{\pgfqpoint{5.483877in}{3.740855in}}%
\pgfpathlineto{\pgfqpoint{9.535013in}{3.740855in}}%
\pgfpathlineto{\pgfqpoint{9.535013in}{6.470855in}}%
\pgfpathlineto{\pgfqpoint{5.483877in}{6.470855in}}%
\pgfpathlineto{\pgfqpoint{5.483877in}{3.740855in}}%
\pgfpathclose%
\pgfusepath{fill}%
\end{pgfscope}%
\begin{pgfscope}%
\pgfpathrectangle{\pgfqpoint{5.483877in}{3.740855in}}{\pgfqpoint{4.051136in}{2.730000in}}%
\pgfusepath{clip}%
\pgfsetroundcap%
\pgfsetroundjoin%
\pgfsetlinewidth{0.803000pt}%
\definecolor{currentstroke}{rgb}{0.800000,0.800000,0.800000}%
\pgfsetstrokecolor{currentstroke}%
\pgfsetdash{}{0pt}%
\pgfpathmoveto{\pgfqpoint{6.240907in}{3.740855in}}%
\pgfpathlineto{\pgfqpoint{6.240907in}{6.470855in}}%
\pgfusepath{stroke}%
\end{pgfscope}%
\begin{pgfscope}%
\definecolor{textcolor}{rgb}{0.150000,0.150000,0.150000}%
\pgfsetstrokecolor{textcolor}%
\pgfsetfillcolor{textcolor}%
\pgftext[x=6.240907in,y=3.692243in,,top]{\color{textcolor}{\sffamily\fontsize{10.000000}{12.000000}\selectfont\catcode`\^=\active\def^{\ifmmode\sp\else\^{}\fi}\catcode`\%=\active\def%{\%}5}}%
\end{pgfscope}%
\begin{pgfscope}%
\pgfpathrectangle{\pgfqpoint{5.483877in}{3.740855in}}{\pgfqpoint{4.051136in}{2.730000in}}%
\pgfusepath{clip}%
\pgfsetroundcap%
\pgfsetroundjoin%
\pgfsetlinewidth{0.803000pt}%
\definecolor{currentstroke}{rgb}{0.800000,0.800000,0.800000}%
\pgfsetstrokecolor{currentstroke}%
\pgfsetdash{}{0pt}%
\pgfpathmoveto{\pgfqpoint{7.059319in}{3.740855in}}%
\pgfpathlineto{\pgfqpoint{7.059319in}{6.470855in}}%
\pgfusepath{stroke}%
\end{pgfscope}%
\begin{pgfscope}%
\definecolor{textcolor}{rgb}{0.150000,0.150000,0.150000}%
\pgfsetstrokecolor{textcolor}%
\pgfsetfillcolor{textcolor}%
\pgftext[x=7.059319in,y=3.692243in,,top]{\color{textcolor}{\sffamily\fontsize{10.000000}{12.000000}\selectfont\catcode`\^=\active\def^{\ifmmode\sp\else\^{}\fi}\catcode`\%=\active\def%{\%}10}}%
\end{pgfscope}%
\begin{pgfscope}%
\pgfpathrectangle{\pgfqpoint{5.483877in}{3.740855in}}{\pgfqpoint{4.051136in}{2.730000in}}%
\pgfusepath{clip}%
\pgfsetroundcap%
\pgfsetroundjoin%
\pgfsetlinewidth{0.803000pt}%
\definecolor{currentstroke}{rgb}{0.800000,0.800000,0.800000}%
\pgfsetstrokecolor{currentstroke}%
\pgfsetdash{}{0pt}%
\pgfpathmoveto{\pgfqpoint{7.877730in}{3.740855in}}%
\pgfpathlineto{\pgfqpoint{7.877730in}{6.470855in}}%
\pgfusepath{stroke}%
\end{pgfscope}%
\begin{pgfscope}%
\definecolor{textcolor}{rgb}{0.150000,0.150000,0.150000}%
\pgfsetstrokecolor{textcolor}%
\pgfsetfillcolor{textcolor}%
\pgftext[x=7.877730in,y=3.692243in,,top]{\color{textcolor}{\sffamily\fontsize{10.000000}{12.000000}\selectfont\catcode`\^=\active\def^{\ifmmode\sp\else\^{}\fi}\catcode`\%=\active\def%{\%}15}}%
\end{pgfscope}%
\begin{pgfscope}%
\pgfpathrectangle{\pgfqpoint{5.483877in}{3.740855in}}{\pgfqpoint{4.051136in}{2.730000in}}%
\pgfusepath{clip}%
\pgfsetroundcap%
\pgfsetroundjoin%
\pgfsetlinewidth{0.803000pt}%
\definecolor{currentstroke}{rgb}{0.800000,0.800000,0.800000}%
\pgfsetstrokecolor{currentstroke}%
\pgfsetdash{}{0pt}%
\pgfpathmoveto{\pgfqpoint{8.696141in}{3.740855in}}%
\pgfpathlineto{\pgfqpoint{8.696141in}{6.470855in}}%
\pgfusepath{stroke}%
\end{pgfscope}%
\begin{pgfscope}%
\definecolor{textcolor}{rgb}{0.150000,0.150000,0.150000}%
\pgfsetstrokecolor{textcolor}%
\pgfsetfillcolor{textcolor}%
\pgftext[x=8.696141in,y=3.692243in,,top]{\color{textcolor}{\sffamily\fontsize{10.000000}{12.000000}\selectfont\catcode`\^=\active\def^{\ifmmode\sp\else\^{}\fi}\catcode`\%=\active\def%{\%}20}}%
\end{pgfscope}%
\begin{pgfscope}%
\pgfpathrectangle{\pgfqpoint{5.483877in}{3.740855in}}{\pgfqpoint{4.051136in}{2.730000in}}%
\pgfusepath{clip}%
\pgfsetroundcap%
\pgfsetroundjoin%
\pgfsetlinewidth{0.803000pt}%
\definecolor{currentstroke}{rgb}{0.800000,0.800000,0.800000}%
\pgfsetstrokecolor{currentstroke}%
\pgfsetdash{}{0pt}%
\pgfpathmoveto{\pgfqpoint{9.514553in}{3.740855in}}%
\pgfpathlineto{\pgfqpoint{9.514553in}{6.470855in}}%
\pgfusepath{stroke}%
\end{pgfscope}%
\begin{pgfscope}%
\definecolor{textcolor}{rgb}{0.150000,0.150000,0.150000}%
\pgfsetstrokecolor{textcolor}%
\pgfsetfillcolor{textcolor}%
\pgftext[x=9.514553in,y=3.692243in,,top]{\color{textcolor}{\sffamily\fontsize{10.000000}{12.000000}\selectfont\catcode`\^=\active\def^{\ifmmode\sp\else\^{}\fi}\catcode`\%=\active\def%{\%}25}}%
\end{pgfscope}%
\begin{pgfscope}%
\definecolor{textcolor}{rgb}{0.150000,0.150000,0.150000}%
\pgfsetstrokecolor{textcolor}%
\pgfsetfillcolor{textcolor}%
\pgftext[x=7.509445in,y=3.506344in,,top]{\color{textcolor}{\sffamily\fontsize{10.000000}{12.000000}\selectfont\catcode`\^=\active\def^{\ifmmode\sp\else\^{}\fi}\catcode`\%=\active\def%{\%}time [h]}}%
\end{pgfscope}%
\begin{pgfscope}%
\pgfpathrectangle{\pgfqpoint{5.483877in}{3.740855in}}{\pgfqpoint{4.051136in}{2.730000in}}%
\pgfusepath{clip}%
\pgfsetroundcap%
\pgfsetroundjoin%
\pgfsetlinewidth{0.803000pt}%
\definecolor{currentstroke}{rgb}{0.800000,0.800000,0.800000}%
\pgfsetstrokecolor{currentstroke}%
\pgfsetdash{}{0pt}%
\pgfpathmoveto{\pgfqpoint{5.483877in}{4.116359in}}%
\pgfpathlineto{\pgfqpoint{9.535013in}{4.116359in}}%
\pgfusepath{stroke}%
\end{pgfscope}%
\begin{pgfscope}%
\definecolor{textcolor}{rgb}{0.150000,0.150000,0.150000}%
\pgfsetstrokecolor{textcolor}%
\pgfsetfillcolor{textcolor}%
\pgftext[x=5.091900in, y=4.056359in, left, base]{\color{textcolor}{\sffamily\fontsize{10.000000}{12.000000}\selectfont\catcode`\^=\active\def^{\ifmmode\sp\else\^{}\fi}\catcode`\%=\active\def%{\%}$\mathdefault{10^{-10}}$}}%
\end{pgfscope}%
\begin{pgfscope}%
\pgfpathrectangle{\pgfqpoint{5.483877in}{3.740855in}}{\pgfqpoint{4.051136in}{2.730000in}}%
\pgfusepath{clip}%
\pgfsetroundcap%
\pgfsetroundjoin%
\pgfsetlinewidth{0.803000pt}%
\definecolor{currentstroke}{rgb}{0.800000,0.800000,0.800000}%
\pgfsetstrokecolor{currentstroke}%
\pgfsetdash{}{0pt}%
\pgfpathmoveto{\pgfqpoint{5.483877in}{4.669203in}}%
\pgfpathlineto{\pgfqpoint{9.535013in}{4.669203in}}%
\pgfusepath{stroke}%
\end{pgfscope}%
\begin{pgfscope}%
\definecolor{textcolor}{rgb}{0.150000,0.150000,0.150000}%
\pgfsetstrokecolor{textcolor}%
\pgfsetfillcolor{textcolor}%
\pgftext[x=5.147263in, y=4.609203in, left, base]{\color{textcolor}{\sffamily\fontsize{10.000000}{12.000000}\selectfont\catcode`\^=\active\def^{\ifmmode\sp\else\^{}\fi}\catcode`\%=\active\def%{\%}$\mathdefault{10^{-8}}$}}%
\end{pgfscope}%
\begin{pgfscope}%
\pgfpathrectangle{\pgfqpoint{5.483877in}{3.740855in}}{\pgfqpoint{4.051136in}{2.730000in}}%
\pgfusepath{clip}%
\pgfsetroundcap%
\pgfsetroundjoin%
\pgfsetlinewidth{0.803000pt}%
\definecolor{currentstroke}{rgb}{0.800000,0.800000,0.800000}%
\pgfsetstrokecolor{currentstroke}%
\pgfsetdash{}{0pt}%
\pgfpathmoveto{\pgfqpoint{5.483877in}{5.222047in}}%
\pgfpathlineto{\pgfqpoint{9.535013in}{5.222047in}}%
\pgfusepath{stroke}%
\end{pgfscope}%
\begin{pgfscope}%
\definecolor{textcolor}{rgb}{0.150000,0.150000,0.150000}%
\pgfsetstrokecolor{textcolor}%
\pgfsetfillcolor{textcolor}%
\pgftext[x=5.147263in, y=5.162047in, left, base]{\color{textcolor}{\sffamily\fontsize{10.000000}{12.000000}\selectfont\catcode`\^=\active\def^{\ifmmode\sp\else\^{}\fi}\catcode`\%=\active\def%{\%}$\mathdefault{10^{-6}}$}}%
\end{pgfscope}%
\begin{pgfscope}%
\pgfpathrectangle{\pgfqpoint{5.483877in}{3.740855in}}{\pgfqpoint{4.051136in}{2.730000in}}%
\pgfusepath{clip}%
\pgfsetroundcap%
\pgfsetroundjoin%
\pgfsetlinewidth{0.803000pt}%
\definecolor{currentstroke}{rgb}{0.800000,0.800000,0.800000}%
\pgfsetstrokecolor{currentstroke}%
\pgfsetdash{}{0pt}%
\pgfpathmoveto{\pgfqpoint{5.483877in}{5.774892in}}%
\pgfpathlineto{\pgfqpoint{9.535013in}{5.774892in}}%
\pgfusepath{stroke}%
\end{pgfscope}%
\begin{pgfscope}%
\definecolor{textcolor}{rgb}{0.150000,0.150000,0.150000}%
\pgfsetstrokecolor{textcolor}%
\pgfsetfillcolor{textcolor}%
\pgftext[x=5.147263in, y=5.714891in, left, base]{\color{textcolor}{\sffamily\fontsize{10.000000}{12.000000}\selectfont\catcode`\^=\active\def^{\ifmmode\sp\else\^{}\fi}\catcode`\%=\active\def%{\%}$\mathdefault{10^{-4}}$}}%
\end{pgfscope}%
\begin{pgfscope}%
\pgfpathrectangle{\pgfqpoint{5.483877in}{3.740855in}}{\pgfqpoint{4.051136in}{2.730000in}}%
\pgfusepath{clip}%
\pgfsetroundcap%
\pgfsetroundjoin%
\pgfsetlinewidth{0.803000pt}%
\definecolor{currentstroke}{rgb}{0.800000,0.800000,0.800000}%
\pgfsetstrokecolor{currentstroke}%
\pgfsetdash{}{0pt}%
\pgfpathmoveto{\pgfqpoint{5.483877in}{6.327736in}}%
\pgfpathlineto{\pgfqpoint{9.535013in}{6.327736in}}%
\pgfusepath{stroke}%
\end{pgfscope}%
\begin{pgfscope}%
\definecolor{textcolor}{rgb}{0.150000,0.150000,0.150000}%
\pgfsetstrokecolor{textcolor}%
\pgfsetfillcolor{textcolor}%
\pgftext[x=5.147263in, y=6.267735in, left, base]{\color{textcolor}{\sffamily\fontsize{10.000000}{12.000000}\selectfont\catcode`\^=\active\def^{\ifmmode\sp\else\^{}\fi}\catcode`\%=\active\def%{\%}$\mathdefault{10^{-2}}$}}%
\end{pgfscope}%
\begin{pgfscope}%
\definecolor{textcolor}{rgb}{0.150000,0.150000,0.150000}%
\pgfsetstrokecolor{textcolor}%
\pgfsetfillcolor{textcolor}%
\pgftext[x=5.036345in,y=5.105855in,,bottom,rotate=90.000000]{\color{textcolor}{\sffamily\fontsize{10.000000}{12.000000}\selectfont\catcode`\^=\active\def^{\ifmmode\sp\else\^{}\fi}\catcode`\%=\active\def%{\%}absolute defect [kg/s]}}%
\end{pgfscope}%
\begin{pgfscope}%
\pgfpathrectangle{\pgfqpoint{5.483877in}{3.740855in}}{\pgfqpoint{4.051136in}{2.730000in}}%
\pgfusepath{clip}%
\pgfsetroundcap%
\pgfsetroundjoin%
\pgfsetlinewidth{1.505625pt}%
\definecolor{currentstroke}{rgb}{0.121569,0.466667,0.705882}%
\pgfsetstrokecolor{currentstroke}%
\pgfsetdash{}{0pt}%
\pgfpathmoveto{\pgfqpoint{5.668019in}{5.924603in}}%
\pgfpathlineto{\pgfqpoint{5.913543in}{5.607999in}}%
\pgfpathlineto{\pgfqpoint{6.159066in}{5.143028in}}%
\pgfpathlineto{\pgfqpoint{6.404590in}{5.566459in}}%
\pgfpathlineto{\pgfqpoint{6.650113in}{5.783254in}}%
\pgfpathlineto{\pgfqpoint{6.895636in}{5.529329in}}%
\pgfpathlineto{\pgfqpoint{7.141160in}{5.061293in}}%
\pgfpathlineto{\pgfqpoint{7.386683in}{5.132143in}}%
\pgfpathlineto{\pgfqpoint{7.632207in}{4.593902in}}%
\pgfpathlineto{\pgfqpoint{7.877730in}{5.044131in}}%
\pgfpathlineto{\pgfqpoint{8.123253in}{5.469469in}}%
\pgfpathlineto{\pgfqpoint{8.368777in}{5.458486in}}%
\pgfpathlineto{\pgfqpoint{8.614300in}{4.308015in}}%
\pgfpathlineto{\pgfqpoint{8.859824in}{4.805148in}}%
\pgfpathlineto{\pgfqpoint{9.105347in}{5.472478in}}%
\pgfpathlineto{\pgfqpoint{9.350871in}{3.864946in}}%
\pgfusepath{stroke}%
\end{pgfscope}%
\begin{pgfscope}%
\pgfpathrectangle{\pgfqpoint{5.483877in}{3.740855in}}{\pgfqpoint{4.051136in}{2.730000in}}%
\pgfusepath{clip}%
\pgfsetbuttcap%
\pgfsetroundjoin%
\definecolor{currentfill}{rgb}{0.121569,0.466667,0.705882}%
\pgfsetfillcolor{currentfill}%
\pgfsetlinewidth{1.003750pt}%
\definecolor{currentstroke}{rgb}{0.121569,0.466667,0.705882}%
\pgfsetstrokecolor{currentstroke}%
\pgfsetdash{}{0pt}%
\pgfsys@defobject{currentmarker}{\pgfqpoint{-0.041667in}{-0.041667in}}{\pgfqpoint{0.041667in}{0.041667in}}{%
\pgfpathmoveto{\pgfqpoint{0.000000in}{-0.041667in}}%
\pgfpathcurveto{\pgfqpoint{0.011050in}{-0.041667in}}{\pgfqpoint{0.021649in}{-0.037276in}}{\pgfqpoint{0.029463in}{-0.029463in}}%
\pgfpathcurveto{\pgfqpoint{0.037276in}{-0.021649in}}{\pgfqpoint{0.041667in}{-0.011050in}}{\pgfqpoint{0.041667in}{0.000000in}}%
\pgfpathcurveto{\pgfqpoint{0.041667in}{0.011050in}}{\pgfqpoint{0.037276in}{0.021649in}}{\pgfqpoint{0.029463in}{0.029463in}}%
\pgfpathcurveto{\pgfqpoint{0.021649in}{0.037276in}}{\pgfqpoint{0.011050in}{0.041667in}}{\pgfqpoint{0.000000in}{0.041667in}}%
\pgfpathcurveto{\pgfqpoint{-0.011050in}{0.041667in}}{\pgfqpoint{-0.021649in}{0.037276in}}{\pgfqpoint{-0.029463in}{0.029463in}}%
\pgfpathcurveto{\pgfqpoint{-0.037276in}{0.021649in}}{\pgfqpoint{-0.041667in}{0.011050in}}{\pgfqpoint{-0.041667in}{0.000000in}}%
\pgfpathcurveto{\pgfqpoint{-0.041667in}{-0.011050in}}{\pgfqpoint{-0.037276in}{-0.021649in}}{\pgfqpoint{-0.029463in}{-0.029463in}}%
\pgfpathcurveto{\pgfqpoint{-0.021649in}{-0.037276in}}{\pgfqpoint{-0.011050in}{-0.041667in}}{\pgfqpoint{0.000000in}{-0.041667in}}%
\pgfpathlineto{\pgfqpoint{0.000000in}{-0.041667in}}%
\pgfpathclose%
\pgfusepath{stroke,fill}%
}%
\begin{pgfscope}%
\pgfsys@transformshift{5.668019in}{5.924603in}%
\pgfsys@useobject{currentmarker}{}%
\end{pgfscope}%
\begin{pgfscope}%
\pgfsys@transformshift{5.913543in}{5.607999in}%
\pgfsys@useobject{currentmarker}{}%
\end{pgfscope}%
\begin{pgfscope}%
\pgfsys@transformshift{6.159066in}{5.143028in}%
\pgfsys@useobject{currentmarker}{}%
\end{pgfscope}%
\begin{pgfscope}%
\pgfsys@transformshift{6.404590in}{5.566459in}%
\pgfsys@useobject{currentmarker}{}%
\end{pgfscope}%
\begin{pgfscope}%
\pgfsys@transformshift{6.650113in}{5.783254in}%
\pgfsys@useobject{currentmarker}{}%
\end{pgfscope}%
\begin{pgfscope}%
\pgfsys@transformshift{6.895636in}{5.529329in}%
\pgfsys@useobject{currentmarker}{}%
\end{pgfscope}%
\begin{pgfscope}%
\pgfsys@transformshift{7.141160in}{5.061293in}%
\pgfsys@useobject{currentmarker}{}%
\end{pgfscope}%
\begin{pgfscope}%
\pgfsys@transformshift{7.386683in}{5.132143in}%
\pgfsys@useobject{currentmarker}{}%
\end{pgfscope}%
\begin{pgfscope}%
\pgfsys@transformshift{7.632207in}{4.593902in}%
\pgfsys@useobject{currentmarker}{}%
\end{pgfscope}%
\begin{pgfscope}%
\pgfsys@transformshift{7.877730in}{5.044131in}%
\pgfsys@useobject{currentmarker}{}%
\end{pgfscope}%
\begin{pgfscope}%
\pgfsys@transformshift{8.123253in}{5.469469in}%
\pgfsys@useobject{currentmarker}{}%
\end{pgfscope}%
\begin{pgfscope}%
\pgfsys@transformshift{8.368777in}{5.458486in}%
\pgfsys@useobject{currentmarker}{}%
\end{pgfscope}%
\begin{pgfscope}%
\pgfsys@transformshift{8.614300in}{4.308015in}%
\pgfsys@useobject{currentmarker}{}%
\end{pgfscope}%
\begin{pgfscope}%
\pgfsys@transformshift{8.859824in}{4.805148in}%
\pgfsys@useobject{currentmarker}{}%
\end{pgfscope}%
\begin{pgfscope}%
\pgfsys@transformshift{9.105347in}{5.472478in}%
\pgfsys@useobject{currentmarker}{}%
\end{pgfscope}%
\begin{pgfscope}%
\pgfsys@transformshift{9.350871in}{3.864946in}%
\pgfsys@useobject{currentmarker}{}%
\end{pgfscope}%
\end{pgfscope}%
\begin{pgfscope}%
\pgfpathrectangle{\pgfqpoint{5.483877in}{3.740855in}}{\pgfqpoint{4.051136in}{2.730000in}}%
\pgfusepath{clip}%
\pgfsetroundcap%
\pgfsetroundjoin%
\pgfsetlinewidth{1.505625pt}%
\definecolor{currentstroke}{rgb}{1.000000,0.498039,0.054902}%
\pgfsetstrokecolor{currentstroke}%
\pgfsetdash{}{0pt}%
\pgfpathmoveto{\pgfqpoint{5.668019in}{5.906496in}}%
\pgfpathlineto{\pgfqpoint{5.913543in}{6.189660in}}%
\pgfpathlineto{\pgfqpoint{6.159066in}{5.530341in}}%
\pgfpathlineto{\pgfqpoint{6.404590in}{6.106527in}}%
\pgfpathlineto{\pgfqpoint{6.650113in}{5.718460in}}%
\pgfpathlineto{\pgfqpoint{6.895636in}{6.143031in}}%
\pgfpathlineto{\pgfqpoint{7.141160in}{5.419675in}}%
\pgfpathlineto{\pgfqpoint{7.386683in}{5.504352in}}%
\pgfpathlineto{\pgfqpoint{7.632207in}{5.057682in}}%
\pgfpathlineto{\pgfqpoint{7.877730in}{5.409377in}}%
\pgfpathlineto{\pgfqpoint{8.123253in}{5.820155in}}%
\pgfpathlineto{\pgfqpoint{8.368777in}{5.808985in}}%
\pgfpathlineto{\pgfqpoint{8.614300in}{4.964440in}}%
\pgfpathlineto{\pgfqpoint{8.859824in}{5.217523in}}%
\pgfpathlineto{\pgfqpoint{9.105347in}{5.396210in}}%
\pgfpathlineto{\pgfqpoint{9.350871in}{4.848854in}}%
\pgfusepath{stroke}%
\end{pgfscope}%
\begin{pgfscope}%
\pgfpathrectangle{\pgfqpoint{5.483877in}{3.740855in}}{\pgfqpoint{4.051136in}{2.730000in}}%
\pgfusepath{clip}%
\pgfsetbuttcap%
\pgfsetroundjoin%
\definecolor{currentfill}{rgb}{1.000000,0.498039,0.054902}%
\pgfsetfillcolor{currentfill}%
\pgfsetlinewidth{1.003750pt}%
\definecolor{currentstroke}{rgb}{1.000000,0.498039,0.054902}%
\pgfsetstrokecolor{currentstroke}%
\pgfsetdash{}{0pt}%
\pgfsys@defobject{currentmarker}{\pgfqpoint{-0.041667in}{-0.041667in}}{\pgfqpoint{0.041667in}{0.041667in}}{%
\pgfpathmoveto{\pgfqpoint{0.000000in}{-0.041667in}}%
\pgfpathcurveto{\pgfqpoint{0.011050in}{-0.041667in}}{\pgfqpoint{0.021649in}{-0.037276in}}{\pgfqpoint{0.029463in}{-0.029463in}}%
\pgfpathcurveto{\pgfqpoint{0.037276in}{-0.021649in}}{\pgfqpoint{0.041667in}{-0.011050in}}{\pgfqpoint{0.041667in}{0.000000in}}%
\pgfpathcurveto{\pgfqpoint{0.041667in}{0.011050in}}{\pgfqpoint{0.037276in}{0.021649in}}{\pgfqpoint{0.029463in}{0.029463in}}%
\pgfpathcurveto{\pgfqpoint{0.021649in}{0.037276in}}{\pgfqpoint{0.011050in}{0.041667in}}{\pgfqpoint{0.000000in}{0.041667in}}%
\pgfpathcurveto{\pgfqpoint{-0.011050in}{0.041667in}}{\pgfqpoint{-0.021649in}{0.037276in}}{\pgfqpoint{-0.029463in}{0.029463in}}%
\pgfpathcurveto{\pgfqpoint{-0.037276in}{0.021649in}}{\pgfqpoint{-0.041667in}{0.011050in}}{\pgfqpoint{-0.041667in}{0.000000in}}%
\pgfpathcurveto{\pgfqpoint{-0.041667in}{-0.011050in}}{\pgfqpoint{-0.037276in}{-0.021649in}}{\pgfqpoint{-0.029463in}{-0.029463in}}%
\pgfpathcurveto{\pgfqpoint{-0.021649in}{-0.037276in}}{\pgfqpoint{-0.011050in}{-0.041667in}}{\pgfqpoint{0.000000in}{-0.041667in}}%
\pgfpathlineto{\pgfqpoint{0.000000in}{-0.041667in}}%
\pgfpathclose%
\pgfusepath{stroke,fill}%
}%
\begin{pgfscope}%
\pgfsys@transformshift{5.668019in}{5.906496in}%
\pgfsys@useobject{currentmarker}{}%
\end{pgfscope}%
\begin{pgfscope}%
\pgfsys@transformshift{5.913543in}{6.189660in}%
\pgfsys@useobject{currentmarker}{}%
\end{pgfscope}%
\begin{pgfscope}%
\pgfsys@transformshift{6.159066in}{5.530341in}%
\pgfsys@useobject{currentmarker}{}%
\end{pgfscope}%
\begin{pgfscope}%
\pgfsys@transformshift{6.404590in}{6.106527in}%
\pgfsys@useobject{currentmarker}{}%
\end{pgfscope}%
\begin{pgfscope}%
\pgfsys@transformshift{6.650113in}{5.718460in}%
\pgfsys@useobject{currentmarker}{}%
\end{pgfscope}%
\begin{pgfscope}%
\pgfsys@transformshift{6.895636in}{6.143031in}%
\pgfsys@useobject{currentmarker}{}%
\end{pgfscope}%
\begin{pgfscope}%
\pgfsys@transformshift{7.141160in}{5.419675in}%
\pgfsys@useobject{currentmarker}{}%
\end{pgfscope}%
\begin{pgfscope}%
\pgfsys@transformshift{7.386683in}{5.504352in}%
\pgfsys@useobject{currentmarker}{}%
\end{pgfscope}%
\begin{pgfscope}%
\pgfsys@transformshift{7.632207in}{5.057682in}%
\pgfsys@useobject{currentmarker}{}%
\end{pgfscope}%
\begin{pgfscope}%
\pgfsys@transformshift{7.877730in}{5.409377in}%
\pgfsys@useobject{currentmarker}{}%
\end{pgfscope}%
\begin{pgfscope}%
\pgfsys@transformshift{8.123253in}{5.820155in}%
\pgfsys@useobject{currentmarker}{}%
\end{pgfscope}%
\begin{pgfscope}%
\pgfsys@transformshift{8.368777in}{5.808985in}%
\pgfsys@useobject{currentmarker}{}%
\end{pgfscope}%
\begin{pgfscope}%
\pgfsys@transformshift{8.614300in}{4.964440in}%
\pgfsys@useobject{currentmarker}{}%
\end{pgfscope}%
\begin{pgfscope}%
\pgfsys@transformshift{8.859824in}{5.217523in}%
\pgfsys@useobject{currentmarker}{}%
\end{pgfscope}%
\begin{pgfscope}%
\pgfsys@transformshift{9.105347in}{5.396210in}%
\pgfsys@useobject{currentmarker}{}%
\end{pgfscope}%
\begin{pgfscope}%
\pgfsys@transformshift{9.350871in}{4.848854in}%
\pgfsys@useobject{currentmarker}{}%
\end{pgfscope}%
\end{pgfscope}%
\begin{pgfscope}%
\pgfpathrectangle{\pgfqpoint{5.483877in}{3.740855in}}{\pgfqpoint{4.051136in}{2.730000in}}%
\pgfusepath{clip}%
\pgfsetroundcap%
\pgfsetroundjoin%
\pgfsetlinewidth{1.505625pt}%
\definecolor{currentstroke}{rgb}{0.172549,0.627451,0.172549}%
\pgfsetstrokecolor{currentstroke}%
\pgfsetdash{}{0pt}%
\pgfpathmoveto{\pgfqpoint{5.668019in}{5.971832in}}%
\pgfpathlineto{\pgfqpoint{5.913543in}{6.346764in}}%
\pgfpathlineto{\pgfqpoint{6.159066in}{5.713240in}}%
\pgfpathlineto{\pgfqpoint{6.404590in}{6.262737in}}%
\pgfpathlineto{\pgfqpoint{6.650113in}{5.582494in}}%
\pgfpathlineto{\pgfqpoint{6.895636in}{6.300690in}}%
\pgfpathlineto{\pgfqpoint{7.141160in}{6.027736in}}%
\pgfpathlineto{\pgfqpoint{7.386683in}{5.672320in}}%
\pgfpathlineto{\pgfqpoint{7.632207in}{5.237367in}}%
\pgfpathlineto{\pgfqpoint{7.877730in}{5.593646in}}%
\pgfpathlineto{\pgfqpoint{8.123253in}{5.941474in}}%
\pgfpathlineto{\pgfqpoint{8.368777in}{5.981847in}}%
\pgfpathlineto{\pgfqpoint{8.614300in}{5.158297in}}%
\pgfpathlineto{\pgfqpoint{8.859824in}{5.394072in}}%
\pgfpathlineto{\pgfqpoint{9.105347in}{5.093299in}}%
\pgfpathlineto{\pgfqpoint{9.350871in}{5.032806in}}%
\pgfusepath{stroke}%
\end{pgfscope}%
\begin{pgfscope}%
\pgfpathrectangle{\pgfqpoint{5.483877in}{3.740855in}}{\pgfqpoint{4.051136in}{2.730000in}}%
\pgfusepath{clip}%
\pgfsetbuttcap%
\pgfsetroundjoin%
\definecolor{currentfill}{rgb}{0.172549,0.627451,0.172549}%
\pgfsetfillcolor{currentfill}%
\pgfsetlinewidth{1.003750pt}%
\definecolor{currentstroke}{rgb}{0.172549,0.627451,0.172549}%
\pgfsetstrokecolor{currentstroke}%
\pgfsetdash{}{0pt}%
\pgfsys@defobject{currentmarker}{\pgfqpoint{-0.041667in}{-0.041667in}}{\pgfqpoint{0.041667in}{0.041667in}}{%
\pgfpathmoveto{\pgfqpoint{0.000000in}{-0.041667in}}%
\pgfpathcurveto{\pgfqpoint{0.011050in}{-0.041667in}}{\pgfqpoint{0.021649in}{-0.037276in}}{\pgfqpoint{0.029463in}{-0.029463in}}%
\pgfpathcurveto{\pgfqpoint{0.037276in}{-0.021649in}}{\pgfqpoint{0.041667in}{-0.011050in}}{\pgfqpoint{0.041667in}{0.000000in}}%
\pgfpathcurveto{\pgfqpoint{0.041667in}{0.011050in}}{\pgfqpoint{0.037276in}{0.021649in}}{\pgfqpoint{0.029463in}{0.029463in}}%
\pgfpathcurveto{\pgfqpoint{0.021649in}{0.037276in}}{\pgfqpoint{0.011050in}{0.041667in}}{\pgfqpoint{0.000000in}{0.041667in}}%
\pgfpathcurveto{\pgfqpoint{-0.011050in}{0.041667in}}{\pgfqpoint{-0.021649in}{0.037276in}}{\pgfqpoint{-0.029463in}{0.029463in}}%
\pgfpathcurveto{\pgfqpoint{-0.037276in}{0.021649in}}{\pgfqpoint{-0.041667in}{0.011050in}}{\pgfqpoint{-0.041667in}{0.000000in}}%
\pgfpathcurveto{\pgfqpoint{-0.041667in}{-0.011050in}}{\pgfqpoint{-0.037276in}{-0.021649in}}{\pgfqpoint{-0.029463in}{-0.029463in}}%
\pgfpathcurveto{\pgfqpoint{-0.021649in}{-0.037276in}}{\pgfqpoint{-0.011050in}{-0.041667in}}{\pgfqpoint{0.000000in}{-0.041667in}}%
\pgfpathlineto{\pgfqpoint{0.000000in}{-0.041667in}}%
\pgfpathclose%
\pgfusepath{stroke,fill}%
}%
\begin{pgfscope}%
\pgfsys@transformshift{5.668019in}{5.971832in}%
\pgfsys@useobject{currentmarker}{}%
\end{pgfscope}%
\begin{pgfscope}%
\pgfsys@transformshift{5.913543in}{6.346764in}%
\pgfsys@useobject{currentmarker}{}%
\end{pgfscope}%
\begin{pgfscope}%
\pgfsys@transformshift{6.159066in}{5.713240in}%
\pgfsys@useobject{currentmarker}{}%
\end{pgfscope}%
\begin{pgfscope}%
\pgfsys@transformshift{6.404590in}{6.262737in}%
\pgfsys@useobject{currentmarker}{}%
\end{pgfscope}%
\begin{pgfscope}%
\pgfsys@transformshift{6.650113in}{5.582494in}%
\pgfsys@useobject{currentmarker}{}%
\end{pgfscope}%
\begin{pgfscope}%
\pgfsys@transformshift{6.895636in}{6.300690in}%
\pgfsys@useobject{currentmarker}{}%
\end{pgfscope}%
\begin{pgfscope}%
\pgfsys@transformshift{7.141160in}{6.027736in}%
\pgfsys@useobject{currentmarker}{}%
\end{pgfscope}%
\begin{pgfscope}%
\pgfsys@transformshift{7.386683in}{5.672320in}%
\pgfsys@useobject{currentmarker}{}%
\end{pgfscope}%
\begin{pgfscope}%
\pgfsys@transformshift{7.632207in}{5.237367in}%
\pgfsys@useobject{currentmarker}{}%
\end{pgfscope}%
\begin{pgfscope}%
\pgfsys@transformshift{7.877730in}{5.593646in}%
\pgfsys@useobject{currentmarker}{}%
\end{pgfscope}%
\begin{pgfscope}%
\pgfsys@transformshift{8.123253in}{5.941474in}%
\pgfsys@useobject{currentmarker}{}%
\end{pgfscope}%
\begin{pgfscope}%
\pgfsys@transformshift{8.368777in}{5.981847in}%
\pgfsys@useobject{currentmarker}{}%
\end{pgfscope}%
\begin{pgfscope}%
\pgfsys@transformshift{8.614300in}{5.158297in}%
\pgfsys@useobject{currentmarker}{}%
\end{pgfscope}%
\begin{pgfscope}%
\pgfsys@transformshift{8.859824in}{5.394072in}%
\pgfsys@useobject{currentmarker}{}%
\end{pgfscope}%
\begin{pgfscope}%
\pgfsys@transformshift{9.105347in}{5.093299in}%
\pgfsys@useobject{currentmarker}{}%
\end{pgfscope}%
\begin{pgfscope}%
\pgfsys@transformshift{9.350871in}{5.032806in}%
\pgfsys@useobject{currentmarker}{}%
\end{pgfscope}%
\end{pgfscope}%
\begin{pgfscope}%
\pgfpathrectangle{\pgfqpoint{5.483877in}{3.740855in}}{\pgfqpoint{4.051136in}{2.730000in}}%
\pgfusepath{clip}%
\pgfsetbuttcap%
\pgfsetroundjoin%
\pgfsetlinewidth{1.505625pt}%
\definecolor{currentstroke}{rgb}{0.000000,0.000000,0.000000}%
\pgfsetstrokecolor{currentstroke}%
\pgfsetdash{{5.550000pt}{2.400000pt}}{0.000000pt}%
\pgfpathmoveto{\pgfqpoint{5.483877in}{5.222047in}}%
\pgfpathlineto{\pgfqpoint{9.535013in}{5.222047in}}%
\pgfusepath{stroke}%
\end{pgfscope}%
\begin{pgfscope}%
\pgfsetrectcap%
\pgfsetmiterjoin%
\pgfsetlinewidth{1.003750pt}%
\definecolor{currentstroke}{rgb}{0.800000,0.800000,0.800000}%
\pgfsetstrokecolor{currentstroke}%
\pgfsetdash{}{0pt}%
\pgfpathmoveto{\pgfqpoint{5.483877in}{3.740855in}}%
\pgfpathlineto{\pgfqpoint{5.483877in}{6.470855in}}%
\pgfusepath{stroke}%
\end{pgfscope}%
\begin{pgfscope}%
\pgfsetrectcap%
\pgfsetmiterjoin%
\pgfsetlinewidth{1.003750pt}%
\definecolor{currentstroke}{rgb}{0.800000,0.800000,0.800000}%
\pgfsetstrokecolor{currentstroke}%
\pgfsetdash{}{0pt}%
\pgfpathmoveto{\pgfqpoint{9.535013in}{3.740855in}}%
\pgfpathlineto{\pgfqpoint{9.535013in}{6.470855in}}%
\pgfusepath{stroke}%
\end{pgfscope}%
\begin{pgfscope}%
\pgfsetrectcap%
\pgfsetmiterjoin%
\pgfsetlinewidth{1.003750pt}%
\definecolor{currentstroke}{rgb}{0.800000,0.800000,0.800000}%
\pgfsetstrokecolor{currentstroke}%
\pgfsetdash{}{0pt}%
\pgfpathmoveto{\pgfqpoint{5.483877in}{3.740855in}}%
\pgfpathlineto{\pgfqpoint{9.535013in}{3.740855in}}%
\pgfusepath{stroke}%
\end{pgfscope}%
\begin{pgfscope}%
\pgfsetrectcap%
\pgfsetmiterjoin%
\pgfsetlinewidth{1.003750pt}%
\definecolor{currentstroke}{rgb}{0.800000,0.800000,0.800000}%
\pgfsetstrokecolor{currentstroke}%
\pgfsetdash{}{0pt}%
\pgfpathmoveto{\pgfqpoint{5.483877in}{6.470855in}}%
\pgfpathlineto{\pgfqpoint{9.535013in}{6.470855in}}%
\pgfusepath{stroke}%
\end{pgfscope}%
\begin{pgfscope}%
\definecolor{textcolor}{rgb}{0.150000,0.150000,0.150000}%
\pgfsetstrokecolor{textcolor}%
\pgfsetfillcolor{textcolor}%
\pgftext[x=7.509445in,y=6.554188in,,base]{\color{textcolor}{\sffamily\fontsize{12.000000}{14.400000}\selectfont\catcode`\^=\active\def^{\ifmmode\sp\else\^{}\fi}\catcode`\%=\active\def%{\%}Component-balance defects}}%
\end{pgfscope}%
\begin{pgfscope}%
\pgfsetroundcap%
\pgfsetroundjoin%
\pgfsetlinewidth{1.505625pt}%
\definecolor{currentstroke}{rgb}{0.121569,0.466667,0.705882}%
\pgfsetstrokecolor{currentstroke}%
\pgfsetdash{}{0pt}%
\pgfpathmoveto{\pgfqpoint{5.608877in}{4.543060in}}%
\pgfpathlineto{\pgfqpoint{5.747766in}{4.543060in}}%
\pgfpathlineto{\pgfqpoint{5.886655in}{4.543060in}}%
\pgfusepath{stroke}%
\end{pgfscope}%
\begin{pgfscope}%
\pgfsetbuttcap%
\pgfsetroundjoin%
\definecolor{currentfill}{rgb}{0.121569,0.466667,0.705882}%
\pgfsetfillcolor{currentfill}%
\pgfsetlinewidth{1.003750pt}%
\definecolor{currentstroke}{rgb}{0.121569,0.466667,0.705882}%
\pgfsetstrokecolor{currentstroke}%
\pgfsetdash{}{0pt}%
\pgfsys@defobject{currentmarker}{\pgfqpoint{-0.041667in}{-0.041667in}}{\pgfqpoint{0.041667in}{0.041667in}}{%
\pgfpathmoveto{\pgfqpoint{0.000000in}{-0.041667in}}%
\pgfpathcurveto{\pgfqpoint{0.011050in}{-0.041667in}}{\pgfqpoint{0.021649in}{-0.037276in}}{\pgfqpoint{0.029463in}{-0.029463in}}%
\pgfpathcurveto{\pgfqpoint{0.037276in}{-0.021649in}}{\pgfqpoint{0.041667in}{-0.011050in}}{\pgfqpoint{0.041667in}{0.000000in}}%
\pgfpathcurveto{\pgfqpoint{0.041667in}{0.011050in}}{\pgfqpoint{0.037276in}{0.021649in}}{\pgfqpoint{0.029463in}{0.029463in}}%
\pgfpathcurveto{\pgfqpoint{0.021649in}{0.037276in}}{\pgfqpoint{0.011050in}{0.041667in}}{\pgfqpoint{0.000000in}{0.041667in}}%
\pgfpathcurveto{\pgfqpoint{-0.011050in}{0.041667in}}{\pgfqpoint{-0.021649in}{0.037276in}}{\pgfqpoint{-0.029463in}{0.029463in}}%
\pgfpathcurveto{\pgfqpoint{-0.037276in}{0.021649in}}{\pgfqpoint{-0.041667in}{0.011050in}}{\pgfqpoint{-0.041667in}{0.000000in}}%
\pgfpathcurveto{\pgfqpoint{-0.041667in}{-0.011050in}}{\pgfqpoint{-0.037276in}{-0.021649in}}{\pgfqpoint{-0.029463in}{-0.029463in}}%
\pgfpathcurveto{\pgfqpoint{-0.021649in}{-0.037276in}}{\pgfqpoint{-0.011050in}{-0.041667in}}{\pgfqpoint{0.000000in}{-0.041667in}}%
\pgfpathlineto{\pgfqpoint{0.000000in}{-0.041667in}}%
\pgfpathclose%
\pgfusepath{stroke,fill}%
}%
\begin{pgfscope}%
\pgfsys@transformshift{5.747766in}{4.543060in}%
\pgfsys@useobject{currentmarker}{}%
\end{pgfscope}%
\end{pgfscope}%
\begin{pgfscope}%
\definecolor{textcolor}{rgb}{0.150000,0.150000,0.150000}%
\pgfsetstrokecolor{textcolor}%
\pgfsetfillcolor{textcolor}%
\pgftext[x=5.997766in,y=4.494449in,left,base]{\color{textcolor}{\sffamily\fontsize{10.000000}{12.000000}\selectfont\catcode`\^=\active\def^{\ifmmode\sp\else\^{}\fi}\catcode`\%=\active\def%{\%}water}}%
\end{pgfscope}%
\begin{pgfscope}%
\pgfsetroundcap%
\pgfsetroundjoin%
\pgfsetlinewidth{1.505625pt}%
\definecolor{currentstroke}{rgb}{1.000000,0.498039,0.054902}%
\pgfsetstrokecolor{currentstroke}%
\pgfsetdash{}{0pt}%
\pgfpathmoveto{\pgfqpoint{5.608877in}{4.343272in}}%
\pgfpathlineto{\pgfqpoint{5.747766in}{4.343272in}}%
\pgfpathlineto{\pgfqpoint{5.886655in}{4.343272in}}%
\pgfusepath{stroke}%
\end{pgfscope}%
\begin{pgfscope}%
\pgfsetbuttcap%
\pgfsetroundjoin%
\definecolor{currentfill}{rgb}{1.000000,0.498039,0.054902}%
\pgfsetfillcolor{currentfill}%
\pgfsetlinewidth{1.003750pt}%
\definecolor{currentstroke}{rgb}{1.000000,0.498039,0.054902}%
\pgfsetstrokecolor{currentstroke}%
\pgfsetdash{}{0pt}%
\pgfsys@defobject{currentmarker}{\pgfqpoint{-0.041667in}{-0.041667in}}{\pgfqpoint{0.041667in}{0.041667in}}{%
\pgfpathmoveto{\pgfqpoint{0.000000in}{-0.041667in}}%
\pgfpathcurveto{\pgfqpoint{0.011050in}{-0.041667in}}{\pgfqpoint{0.021649in}{-0.037276in}}{\pgfqpoint{0.029463in}{-0.029463in}}%
\pgfpathcurveto{\pgfqpoint{0.037276in}{-0.021649in}}{\pgfqpoint{0.041667in}{-0.011050in}}{\pgfqpoint{0.041667in}{0.000000in}}%
\pgfpathcurveto{\pgfqpoint{0.041667in}{0.011050in}}{\pgfqpoint{0.037276in}{0.021649in}}{\pgfqpoint{0.029463in}{0.029463in}}%
\pgfpathcurveto{\pgfqpoint{0.021649in}{0.037276in}}{\pgfqpoint{0.011050in}{0.041667in}}{\pgfqpoint{0.000000in}{0.041667in}}%
\pgfpathcurveto{\pgfqpoint{-0.011050in}{0.041667in}}{\pgfqpoint{-0.021649in}{0.037276in}}{\pgfqpoint{-0.029463in}{0.029463in}}%
\pgfpathcurveto{\pgfqpoint{-0.037276in}{0.021649in}}{\pgfqpoint{-0.041667in}{0.011050in}}{\pgfqpoint{-0.041667in}{0.000000in}}%
\pgfpathcurveto{\pgfqpoint{-0.041667in}{-0.011050in}}{\pgfqpoint{-0.037276in}{-0.021649in}}{\pgfqpoint{-0.029463in}{-0.029463in}}%
\pgfpathcurveto{\pgfqpoint{-0.021649in}{-0.037276in}}{\pgfqpoint{-0.011050in}{-0.041667in}}{\pgfqpoint{0.000000in}{-0.041667in}}%
\pgfpathlineto{\pgfqpoint{0.000000in}{-0.041667in}}%
\pgfpathclose%
\pgfusepath{stroke,fill}%
}%
\begin{pgfscope}%
\pgfsys@transformshift{5.747766in}{4.343272in}%
\pgfsys@useobject{currentmarker}{}%
\end{pgfscope}%
\end{pgfscope}%
\begin{pgfscope}%
\definecolor{textcolor}{rgb}{0.150000,0.150000,0.150000}%
\pgfsetstrokecolor{textcolor}%
\pgfsetfillcolor{textcolor}%
\pgftext[x=5.997766in,y=4.294660in,left,base]{\color{textcolor}{\sffamily\fontsize{10.000000}{12.000000}\selectfont\catcode`\^=\active\def^{\ifmmode\sp\else\^{}\fi}\catcode`\%=\active\def%{\%}oil}}%
\end{pgfscope}%
\begin{pgfscope}%
\pgfsetroundcap%
\pgfsetroundjoin%
\pgfsetlinewidth{1.505625pt}%
\definecolor{currentstroke}{rgb}{0.172549,0.627451,0.172549}%
\pgfsetstrokecolor{currentstroke}%
\pgfsetdash{}{0pt}%
\pgfpathmoveto{\pgfqpoint{5.608877in}{4.143483in}}%
\pgfpathlineto{\pgfqpoint{5.747766in}{4.143483in}}%
\pgfpathlineto{\pgfqpoint{5.886655in}{4.143483in}}%
\pgfusepath{stroke}%
\end{pgfscope}%
\begin{pgfscope}%
\pgfsetbuttcap%
\pgfsetroundjoin%
\definecolor{currentfill}{rgb}{0.172549,0.627451,0.172549}%
\pgfsetfillcolor{currentfill}%
\pgfsetlinewidth{1.003750pt}%
\definecolor{currentstroke}{rgb}{0.172549,0.627451,0.172549}%
\pgfsetstrokecolor{currentstroke}%
\pgfsetdash{}{0pt}%
\pgfsys@defobject{currentmarker}{\pgfqpoint{-0.041667in}{-0.041667in}}{\pgfqpoint{0.041667in}{0.041667in}}{%
\pgfpathmoveto{\pgfqpoint{0.000000in}{-0.041667in}}%
\pgfpathcurveto{\pgfqpoint{0.011050in}{-0.041667in}}{\pgfqpoint{0.021649in}{-0.037276in}}{\pgfqpoint{0.029463in}{-0.029463in}}%
\pgfpathcurveto{\pgfqpoint{0.037276in}{-0.021649in}}{\pgfqpoint{0.041667in}{-0.011050in}}{\pgfqpoint{0.041667in}{0.000000in}}%
\pgfpathcurveto{\pgfqpoint{0.041667in}{0.011050in}}{\pgfqpoint{0.037276in}{0.021649in}}{\pgfqpoint{0.029463in}{0.029463in}}%
\pgfpathcurveto{\pgfqpoint{0.021649in}{0.037276in}}{\pgfqpoint{0.011050in}{0.041667in}}{\pgfqpoint{0.000000in}{0.041667in}}%
\pgfpathcurveto{\pgfqpoint{-0.011050in}{0.041667in}}{\pgfqpoint{-0.021649in}{0.037276in}}{\pgfqpoint{-0.029463in}{0.029463in}}%
\pgfpathcurveto{\pgfqpoint{-0.037276in}{0.021649in}}{\pgfqpoint{-0.041667in}{0.011050in}}{\pgfqpoint{-0.041667in}{0.000000in}}%
\pgfpathcurveto{\pgfqpoint{-0.041667in}{-0.011050in}}{\pgfqpoint{-0.037276in}{-0.021649in}}{\pgfqpoint{-0.029463in}{-0.029463in}}%
\pgfpathcurveto{\pgfqpoint{-0.021649in}{-0.037276in}}{\pgfqpoint{-0.011050in}{-0.041667in}}{\pgfqpoint{0.000000in}{-0.041667in}}%
\pgfpathlineto{\pgfqpoint{0.000000in}{-0.041667in}}%
\pgfpathclose%
\pgfusepath{stroke,fill}%
}%
\begin{pgfscope}%
\pgfsys@transformshift{5.747766in}{4.143483in}%
\pgfsys@useobject{currentmarker}{}%
\end{pgfscope}%
\end{pgfscope}%
\begin{pgfscope}%
\definecolor{textcolor}{rgb}{0.150000,0.150000,0.150000}%
\pgfsetstrokecolor{textcolor}%
\pgfsetfillcolor{textcolor}%
\pgftext[x=5.997766in,y=4.094872in,left,base]{\color{textcolor}{\sffamily\fontsize{10.000000}{12.000000}\selectfont\catcode`\^=\active\def^{\ifmmode\sp\else\^{}\fi}\catcode`\%=\active\def%{\%}gas}}%
\end{pgfscope}%
\begin{pgfscope}%
\pgfsetbuttcap%
\pgfsetroundjoin%
\pgfsetlinewidth{1.505625pt}%
\definecolor{currentstroke}{rgb}{0.000000,0.000000,0.000000}%
\pgfsetstrokecolor{currentstroke}%
\pgfsetdash{{5.550000pt}{2.400000pt}}{0.000000pt}%
\pgfpathmoveto{\pgfqpoint{5.608877in}{3.943695in}}%
\pgfpathlineto{\pgfqpoint{5.747766in}{3.943695in}}%
\pgfpathlineto{\pgfqpoint{5.886655in}{3.943695in}}%
\pgfusepath{stroke}%
\end{pgfscope}%
\begin{pgfscope}%
\definecolor{textcolor}{rgb}{0.150000,0.150000,0.150000}%
\pgfsetstrokecolor{textcolor}%
\pgfsetfillcolor{textcolor}%
\pgftext[x=5.997766in,y=3.895084in,left,base]{\color{textcolor}{\sffamily\fontsize{10.000000}{12.000000}\selectfont\catcode`\^=\active\def^{\ifmmode\sp\else\^{}\fi}\catcode`\%=\active\def%{\%}balance gate}}%
\end{pgfscope}%
\begin{pgfscope}%
\pgfsetbuttcap%
\pgfsetmiterjoin%
\definecolor{currentfill}{rgb}{1.000000,1.000000,1.000000}%
\pgfsetfillcolor{currentfill}%
\pgfsetlinewidth{0.000000pt}%
\definecolor{currentstroke}{rgb}{0.000000,0.000000,0.000000}%
\pgfsetstrokecolor{currentstroke}%
\pgfsetstrokeopacity{0.000000}%
\pgfsetdash{}{0pt}%
\pgfpathmoveto{\pgfqpoint{0.622513in}{0.464855in}}%
\pgfpathlineto{\pgfqpoint{4.673649in}{0.464855in}}%
\pgfpathlineto{\pgfqpoint{4.673649in}{3.194855in}}%
\pgfpathlineto{\pgfqpoint{0.622513in}{3.194855in}}%
\pgfpathlineto{\pgfqpoint{0.622513in}{0.464855in}}%
\pgfpathclose%
\pgfusepath{fill}%
\end{pgfscope}%
\begin{pgfscope}%
\pgfpathrectangle{\pgfqpoint{0.622513in}{0.464855in}}{\pgfqpoint{4.051136in}{2.730000in}}%
\pgfusepath{clip}%
\pgfsetroundcap%
\pgfsetroundjoin%
\pgfsetlinewidth{0.803000pt}%
\definecolor{currentstroke}{rgb}{0.800000,0.800000,0.800000}%
\pgfsetstrokecolor{currentstroke}%
\pgfsetdash{}{0pt}%
\pgfpathmoveto{\pgfqpoint{1.379544in}{0.464855in}}%
\pgfpathlineto{\pgfqpoint{1.379544in}{3.194855in}}%
\pgfusepath{stroke}%
\end{pgfscope}%
\begin{pgfscope}%
\definecolor{textcolor}{rgb}{0.150000,0.150000,0.150000}%
\pgfsetstrokecolor{textcolor}%
\pgfsetfillcolor{textcolor}%
\pgftext[x=1.379544in,y=0.416243in,,top]{\color{textcolor}{\sffamily\fontsize{10.000000}{12.000000}\selectfont\catcode`\^=\active\def^{\ifmmode\sp\else\^{}\fi}\catcode`\%=\active\def%{\%}5}}%
\end{pgfscope}%
\begin{pgfscope}%
\pgfpathrectangle{\pgfqpoint{0.622513in}{0.464855in}}{\pgfqpoint{4.051136in}{2.730000in}}%
\pgfusepath{clip}%
\pgfsetroundcap%
\pgfsetroundjoin%
\pgfsetlinewidth{0.803000pt}%
\definecolor{currentstroke}{rgb}{0.800000,0.800000,0.800000}%
\pgfsetstrokecolor{currentstroke}%
\pgfsetdash{}{0pt}%
\pgfpathmoveto{\pgfqpoint{2.197955in}{0.464855in}}%
\pgfpathlineto{\pgfqpoint{2.197955in}{3.194855in}}%
\pgfusepath{stroke}%
\end{pgfscope}%
\begin{pgfscope}%
\definecolor{textcolor}{rgb}{0.150000,0.150000,0.150000}%
\pgfsetstrokecolor{textcolor}%
\pgfsetfillcolor{textcolor}%
\pgftext[x=2.197955in,y=0.416243in,,top]{\color{textcolor}{\sffamily\fontsize{10.000000}{12.000000}\selectfont\catcode`\^=\active\def^{\ifmmode\sp\else\^{}\fi}\catcode`\%=\active\def%{\%}10}}%
\end{pgfscope}%
\begin{pgfscope}%
\pgfpathrectangle{\pgfqpoint{0.622513in}{0.464855in}}{\pgfqpoint{4.051136in}{2.730000in}}%
\pgfusepath{clip}%
\pgfsetroundcap%
\pgfsetroundjoin%
\pgfsetlinewidth{0.803000pt}%
\definecolor{currentstroke}{rgb}{0.800000,0.800000,0.800000}%
\pgfsetstrokecolor{currentstroke}%
\pgfsetdash{}{0pt}%
\pgfpathmoveto{\pgfqpoint{3.016366in}{0.464855in}}%
\pgfpathlineto{\pgfqpoint{3.016366in}{3.194855in}}%
\pgfusepath{stroke}%
\end{pgfscope}%
\begin{pgfscope}%
\definecolor{textcolor}{rgb}{0.150000,0.150000,0.150000}%
\pgfsetstrokecolor{textcolor}%
\pgfsetfillcolor{textcolor}%
\pgftext[x=3.016366in,y=0.416243in,,top]{\color{textcolor}{\sffamily\fontsize{10.000000}{12.000000}\selectfont\catcode`\^=\active\def^{\ifmmode\sp\else\^{}\fi}\catcode`\%=\active\def%{\%}15}}%
\end{pgfscope}%
\begin{pgfscope}%
\pgfpathrectangle{\pgfqpoint{0.622513in}{0.464855in}}{\pgfqpoint{4.051136in}{2.730000in}}%
\pgfusepath{clip}%
\pgfsetroundcap%
\pgfsetroundjoin%
\pgfsetlinewidth{0.803000pt}%
\definecolor{currentstroke}{rgb}{0.800000,0.800000,0.800000}%
\pgfsetstrokecolor{currentstroke}%
\pgfsetdash{}{0pt}%
\pgfpathmoveto{\pgfqpoint{3.834778in}{0.464855in}}%
\pgfpathlineto{\pgfqpoint{3.834778in}{3.194855in}}%
\pgfusepath{stroke}%
\end{pgfscope}%
\begin{pgfscope}%
\definecolor{textcolor}{rgb}{0.150000,0.150000,0.150000}%
\pgfsetstrokecolor{textcolor}%
\pgfsetfillcolor{textcolor}%
\pgftext[x=3.834778in,y=0.416243in,,top]{\color{textcolor}{\sffamily\fontsize{10.000000}{12.000000}\selectfont\catcode`\^=\active\def^{\ifmmode\sp\else\^{}\fi}\catcode`\%=\active\def%{\%}20}}%
\end{pgfscope}%
\begin{pgfscope}%
\pgfpathrectangle{\pgfqpoint{0.622513in}{0.464855in}}{\pgfqpoint{4.051136in}{2.730000in}}%
\pgfusepath{clip}%
\pgfsetroundcap%
\pgfsetroundjoin%
\pgfsetlinewidth{0.803000pt}%
\definecolor{currentstroke}{rgb}{0.800000,0.800000,0.800000}%
\pgfsetstrokecolor{currentstroke}%
\pgfsetdash{}{0pt}%
\pgfpathmoveto{\pgfqpoint{4.653189in}{0.464855in}}%
\pgfpathlineto{\pgfqpoint{4.653189in}{3.194855in}}%
\pgfusepath{stroke}%
\end{pgfscope}%
\begin{pgfscope}%
\definecolor{textcolor}{rgb}{0.150000,0.150000,0.150000}%
\pgfsetstrokecolor{textcolor}%
\pgfsetfillcolor{textcolor}%
\pgftext[x=4.653189in,y=0.416243in,,top]{\color{textcolor}{\sffamily\fontsize{10.000000}{12.000000}\selectfont\catcode`\^=\active\def^{\ifmmode\sp\else\^{}\fi}\catcode`\%=\active\def%{\%}25}}%
\end{pgfscope}%
\begin{pgfscope}%
\definecolor{textcolor}{rgb}{0.150000,0.150000,0.150000}%
\pgfsetstrokecolor{textcolor}%
\pgfsetfillcolor{textcolor}%
\pgftext[x=2.648081in,y=0.230344in,,top]{\color{textcolor}{\sffamily\fontsize{10.000000}{12.000000}\selectfont\catcode`\^=\active\def^{\ifmmode\sp\else\^{}\fi}\catcode`\%=\active\def%{\%}time [h]}}%
\end{pgfscope}%
\begin{pgfscope}%
\pgfpathrectangle{\pgfqpoint{0.622513in}{0.464855in}}{\pgfqpoint{4.051136in}{2.730000in}}%
\pgfusepath{clip}%
\pgfsetroundcap%
\pgfsetroundjoin%
\pgfsetlinewidth{0.803000pt}%
\definecolor{currentstroke}{rgb}{0.800000,0.800000,0.800000}%
\pgfsetstrokecolor{currentstroke}%
\pgfsetdash{}{0pt}%
\pgfpathmoveto{\pgfqpoint{0.622513in}{0.747937in}}%
\pgfpathlineto{\pgfqpoint{4.673649in}{0.747937in}}%
\pgfusepath{stroke}%
\end{pgfscope}%
\begin{pgfscope}%
\definecolor{textcolor}{rgb}{0.150000,0.150000,0.150000}%
\pgfsetstrokecolor{textcolor}%
\pgfsetfillcolor{textcolor}%
\pgftext[x=0.285900in, y=0.687937in, left, base]{\color{textcolor}{\sffamily\fontsize{10.000000}{12.000000}\selectfont\catcode`\^=\active\def^{\ifmmode\sp\else\^{}\fi}\catcode`\%=\active\def%{\%}$\mathdefault{10^{-4}}$}}%
\end{pgfscope}%
\begin{pgfscope}%
\pgfpathrectangle{\pgfqpoint{0.622513in}{0.464855in}}{\pgfqpoint{4.051136in}{2.730000in}}%
\pgfusepath{clip}%
\pgfsetroundcap%
\pgfsetroundjoin%
\pgfsetlinewidth{0.803000pt}%
\definecolor{currentstroke}{rgb}{0.800000,0.800000,0.800000}%
\pgfsetstrokecolor{currentstroke}%
\pgfsetdash{}{0pt}%
\pgfpathmoveto{\pgfqpoint{0.622513in}{1.112742in}}%
\pgfpathlineto{\pgfqpoint{4.673649in}{1.112742in}}%
\pgfusepath{stroke}%
\end{pgfscope}%
\begin{pgfscope}%
\definecolor{textcolor}{rgb}{0.150000,0.150000,0.150000}%
\pgfsetstrokecolor{textcolor}%
\pgfsetfillcolor{textcolor}%
\pgftext[x=0.285900in, y=1.052742in, left, base]{\color{textcolor}{\sffamily\fontsize{10.000000}{12.000000}\selectfont\catcode`\^=\active\def^{\ifmmode\sp\else\^{}\fi}\catcode`\%=\active\def%{\%}$\mathdefault{10^{-3}}$}}%
\end{pgfscope}%
\begin{pgfscope}%
\pgfpathrectangle{\pgfqpoint{0.622513in}{0.464855in}}{\pgfqpoint{4.051136in}{2.730000in}}%
\pgfusepath{clip}%
\pgfsetroundcap%
\pgfsetroundjoin%
\pgfsetlinewidth{0.803000pt}%
\definecolor{currentstroke}{rgb}{0.800000,0.800000,0.800000}%
\pgfsetstrokecolor{currentstroke}%
\pgfsetdash{}{0pt}%
\pgfpathmoveto{\pgfqpoint{0.622513in}{1.477547in}}%
\pgfpathlineto{\pgfqpoint{4.673649in}{1.477547in}}%
\pgfusepath{stroke}%
\end{pgfscope}%
\begin{pgfscope}%
\definecolor{textcolor}{rgb}{0.150000,0.150000,0.150000}%
\pgfsetstrokecolor{textcolor}%
\pgfsetfillcolor{textcolor}%
\pgftext[x=0.285900in, y=1.417547in, left, base]{\color{textcolor}{\sffamily\fontsize{10.000000}{12.000000}\selectfont\catcode`\^=\active\def^{\ifmmode\sp\else\^{}\fi}\catcode`\%=\active\def%{\%}$\mathdefault{10^{-2}}$}}%
\end{pgfscope}%
\begin{pgfscope}%
\pgfpathrectangle{\pgfqpoint{0.622513in}{0.464855in}}{\pgfqpoint{4.051136in}{2.730000in}}%
\pgfusepath{clip}%
\pgfsetroundcap%
\pgfsetroundjoin%
\pgfsetlinewidth{0.803000pt}%
\definecolor{currentstroke}{rgb}{0.800000,0.800000,0.800000}%
\pgfsetstrokecolor{currentstroke}%
\pgfsetdash{}{0pt}%
\pgfpathmoveto{\pgfqpoint{0.622513in}{1.842352in}}%
\pgfpathlineto{\pgfqpoint{4.673649in}{1.842352in}}%
\pgfusepath{stroke}%
\end{pgfscope}%
\begin{pgfscope}%
\definecolor{textcolor}{rgb}{0.150000,0.150000,0.150000}%
\pgfsetstrokecolor{textcolor}%
\pgfsetfillcolor{textcolor}%
\pgftext[x=0.285900in, y=1.782352in, left, base]{\color{textcolor}{\sffamily\fontsize{10.000000}{12.000000}\selectfont\catcode`\^=\active\def^{\ifmmode\sp\else\^{}\fi}\catcode`\%=\active\def%{\%}$\mathdefault{10^{-1}}$}}%
\end{pgfscope}%
\begin{pgfscope}%
\pgfpathrectangle{\pgfqpoint{0.622513in}{0.464855in}}{\pgfqpoint{4.051136in}{2.730000in}}%
\pgfusepath{clip}%
\pgfsetroundcap%
\pgfsetroundjoin%
\pgfsetlinewidth{0.803000pt}%
\definecolor{currentstroke}{rgb}{0.800000,0.800000,0.800000}%
\pgfsetstrokecolor{currentstroke}%
\pgfsetdash{}{0pt}%
\pgfpathmoveto{\pgfqpoint{0.622513in}{2.207157in}}%
\pgfpathlineto{\pgfqpoint{4.673649in}{2.207157in}}%
\pgfusepath{stroke}%
\end{pgfscope}%
\begin{pgfscope}%
\definecolor{textcolor}{rgb}{0.150000,0.150000,0.150000}%
\pgfsetstrokecolor{textcolor}%
\pgfsetfillcolor{textcolor}%
\pgftext[x=0.372705in, y=2.147157in, left, base]{\color{textcolor}{\sffamily\fontsize{10.000000}{12.000000}\selectfont\catcode`\^=\active\def^{\ifmmode\sp\else\^{}\fi}\catcode`\%=\active\def%{\%}$\mathdefault{10^{0}}$}}%
\end{pgfscope}%
\begin{pgfscope}%
\pgfpathrectangle{\pgfqpoint{0.622513in}{0.464855in}}{\pgfqpoint{4.051136in}{2.730000in}}%
\pgfusepath{clip}%
\pgfsetroundcap%
\pgfsetroundjoin%
\pgfsetlinewidth{0.803000pt}%
\definecolor{currentstroke}{rgb}{0.800000,0.800000,0.800000}%
\pgfsetstrokecolor{currentstroke}%
\pgfsetdash{}{0pt}%
\pgfpathmoveto{\pgfqpoint{0.622513in}{2.571962in}}%
\pgfpathlineto{\pgfqpoint{4.673649in}{2.571962in}}%
\pgfusepath{stroke}%
\end{pgfscope}%
\begin{pgfscope}%
\definecolor{textcolor}{rgb}{0.150000,0.150000,0.150000}%
\pgfsetstrokecolor{textcolor}%
\pgfsetfillcolor{textcolor}%
\pgftext[x=0.372705in, y=2.511962in, left, base]{\color{textcolor}{\sffamily\fontsize{10.000000}{12.000000}\selectfont\catcode`\^=\active\def^{\ifmmode\sp\else\^{}\fi}\catcode`\%=\active\def%{\%}$\mathdefault{10^{1}}$}}%
\end{pgfscope}%
\begin{pgfscope}%
\pgfpathrectangle{\pgfqpoint{0.622513in}{0.464855in}}{\pgfqpoint{4.051136in}{2.730000in}}%
\pgfusepath{clip}%
\pgfsetroundcap%
\pgfsetroundjoin%
\pgfsetlinewidth{0.803000pt}%
\definecolor{currentstroke}{rgb}{0.800000,0.800000,0.800000}%
\pgfsetstrokecolor{currentstroke}%
\pgfsetdash{}{0pt}%
\pgfpathmoveto{\pgfqpoint{0.622513in}{2.936767in}}%
\pgfpathlineto{\pgfqpoint{4.673649in}{2.936767in}}%
\pgfusepath{stroke}%
\end{pgfscope}%
\begin{pgfscope}%
\definecolor{textcolor}{rgb}{0.150000,0.150000,0.150000}%
\pgfsetstrokecolor{textcolor}%
\pgfsetfillcolor{textcolor}%
\pgftext[x=0.372705in, y=2.876767in, left, base]{\color{textcolor}{\sffamily\fontsize{10.000000}{12.000000}\selectfont\catcode`\^=\active\def^{\ifmmode\sp\else\^{}\fi}\catcode`\%=\active\def%{\%}$\mathdefault{10^{2}}$}}%
\end{pgfscope}%
\begin{pgfscope}%
\definecolor{textcolor}{rgb}{0.150000,0.150000,0.150000}%
\pgfsetstrokecolor{textcolor}%
\pgfsetfillcolor{textcolor}%
\pgftext[x=0.230344in,y=1.829855in,,bottom,rotate=90.000000]{\color{textcolor}{\sffamily\fontsize{10.000000}{12.000000}\selectfont\catcode`\^=\active\def^{\ifmmode\sp\else\^{}\fi}\catcode`\%=\active\def%{\%}residual / limit}}%
\end{pgfscope}%
\begin{pgfscope}%
\pgfpathrectangle{\pgfqpoint{0.622513in}{0.464855in}}{\pgfqpoint{4.051136in}{2.730000in}}%
\pgfusepath{clip}%
\pgfsetroundcap%
\pgfsetroundjoin%
\pgfsetlinewidth{1.505625pt}%
\definecolor{currentstroke}{rgb}{0.121569,0.466667,0.705882}%
\pgfsetstrokecolor{currentstroke}%
\pgfsetdash{}{0pt}%
\pgfpathmoveto{\pgfqpoint{0.806656in}{1.290129in}}%
\pgfpathlineto{\pgfqpoint{1.052179in}{1.283904in}}%
\pgfpathlineto{\pgfqpoint{1.297703in}{1.282850in}}%
\pgfpathlineto{\pgfqpoint{1.543226in}{1.290694in}}%
\pgfpathlineto{\pgfqpoint{1.788749in}{1.287272in}}%
\pgfpathlineto{\pgfqpoint{2.034273in}{1.284865in}}%
\pgfpathlineto{\pgfqpoint{2.279796in}{1.278233in}}%
\pgfpathlineto{\pgfqpoint{2.525320in}{1.285318in}}%
\pgfpathlineto{\pgfqpoint{2.770843in}{1.284949in}}%
\pgfpathlineto{\pgfqpoint{3.016366in}{1.285449in}}%
\pgfpathlineto{\pgfqpoint{3.261890in}{1.280916in}}%
\pgfpathlineto{\pgfqpoint{3.507413in}{1.286111in}}%
\pgfpathlineto{\pgfqpoint{3.752937in}{1.299100in}}%
\pgfpathlineto{\pgfqpoint{3.998460in}{1.253151in}}%
\pgfpathlineto{\pgfqpoint{4.243984in}{1.290217in}}%
\pgfpathlineto{\pgfqpoint{4.489507in}{1.322486in}}%
\pgfusepath{stroke}%
\end{pgfscope}%
\begin{pgfscope}%
\pgfpathrectangle{\pgfqpoint{0.622513in}{0.464855in}}{\pgfqpoint{4.051136in}{2.730000in}}%
\pgfusepath{clip}%
\pgfsetbuttcap%
\pgfsetroundjoin%
\definecolor{currentfill}{rgb}{0.121569,0.466667,0.705882}%
\pgfsetfillcolor{currentfill}%
\pgfsetlinewidth{1.003750pt}%
\definecolor{currentstroke}{rgb}{0.121569,0.466667,0.705882}%
\pgfsetstrokecolor{currentstroke}%
\pgfsetdash{}{0pt}%
\pgfsys@defobject{currentmarker}{\pgfqpoint{-0.041667in}{-0.041667in}}{\pgfqpoint{0.041667in}{0.041667in}}{%
\pgfpathmoveto{\pgfqpoint{0.000000in}{-0.041667in}}%
\pgfpathcurveto{\pgfqpoint{0.011050in}{-0.041667in}}{\pgfqpoint{0.021649in}{-0.037276in}}{\pgfqpoint{0.029463in}{-0.029463in}}%
\pgfpathcurveto{\pgfqpoint{0.037276in}{-0.021649in}}{\pgfqpoint{0.041667in}{-0.011050in}}{\pgfqpoint{0.041667in}{0.000000in}}%
\pgfpathcurveto{\pgfqpoint{0.041667in}{0.011050in}}{\pgfqpoint{0.037276in}{0.021649in}}{\pgfqpoint{0.029463in}{0.029463in}}%
\pgfpathcurveto{\pgfqpoint{0.021649in}{0.037276in}}{\pgfqpoint{0.011050in}{0.041667in}}{\pgfqpoint{0.000000in}{0.041667in}}%
\pgfpathcurveto{\pgfqpoint{-0.011050in}{0.041667in}}{\pgfqpoint{-0.021649in}{0.037276in}}{\pgfqpoint{-0.029463in}{0.029463in}}%
\pgfpathcurveto{\pgfqpoint{-0.037276in}{0.021649in}}{\pgfqpoint{-0.041667in}{0.011050in}}{\pgfqpoint{-0.041667in}{0.000000in}}%
\pgfpathcurveto{\pgfqpoint{-0.041667in}{-0.011050in}}{\pgfqpoint{-0.037276in}{-0.021649in}}{\pgfqpoint{-0.029463in}{-0.029463in}}%
\pgfpathcurveto{\pgfqpoint{-0.021649in}{-0.037276in}}{\pgfqpoint{-0.011050in}{-0.041667in}}{\pgfqpoint{0.000000in}{-0.041667in}}%
\pgfpathlineto{\pgfqpoint{0.000000in}{-0.041667in}}%
\pgfpathclose%
\pgfusepath{stroke,fill}%
}%
\begin{pgfscope}%
\pgfsys@transformshift{0.806656in}{1.290129in}%
\pgfsys@useobject{currentmarker}{}%
\end{pgfscope}%
\begin{pgfscope}%
\pgfsys@transformshift{1.052179in}{1.283904in}%
\pgfsys@useobject{currentmarker}{}%
\end{pgfscope}%
\begin{pgfscope}%
\pgfsys@transformshift{1.297703in}{1.282850in}%
\pgfsys@useobject{currentmarker}{}%
\end{pgfscope}%
\begin{pgfscope}%
\pgfsys@transformshift{1.543226in}{1.290694in}%
\pgfsys@useobject{currentmarker}{}%
\end{pgfscope}%
\begin{pgfscope}%
\pgfsys@transformshift{1.788749in}{1.287272in}%
\pgfsys@useobject{currentmarker}{}%
\end{pgfscope}%
\begin{pgfscope}%
\pgfsys@transformshift{2.034273in}{1.284865in}%
\pgfsys@useobject{currentmarker}{}%
\end{pgfscope}%
\begin{pgfscope}%
\pgfsys@transformshift{2.279796in}{1.278233in}%
\pgfsys@useobject{currentmarker}{}%
\end{pgfscope}%
\begin{pgfscope}%
\pgfsys@transformshift{2.525320in}{1.285318in}%
\pgfsys@useobject{currentmarker}{}%
\end{pgfscope}%
\begin{pgfscope}%
\pgfsys@transformshift{2.770843in}{1.284949in}%
\pgfsys@useobject{currentmarker}{}%
\end{pgfscope}%
\begin{pgfscope}%
\pgfsys@transformshift{3.016366in}{1.285449in}%
\pgfsys@useobject{currentmarker}{}%
\end{pgfscope}%
\begin{pgfscope}%
\pgfsys@transformshift{3.261890in}{1.280916in}%
\pgfsys@useobject{currentmarker}{}%
\end{pgfscope}%
\begin{pgfscope}%
\pgfsys@transformshift{3.507413in}{1.286111in}%
\pgfsys@useobject{currentmarker}{}%
\end{pgfscope}%
\begin{pgfscope}%
\pgfsys@transformshift{3.752937in}{1.299100in}%
\pgfsys@useobject{currentmarker}{}%
\end{pgfscope}%
\begin{pgfscope}%
\pgfsys@transformshift{3.998460in}{1.253151in}%
\pgfsys@useobject{currentmarker}{}%
\end{pgfscope}%
\begin{pgfscope}%
\pgfsys@transformshift{4.243984in}{1.290217in}%
\pgfsys@useobject{currentmarker}{}%
\end{pgfscope}%
\begin{pgfscope}%
\pgfsys@transformshift{4.489507in}{1.322486in}%
\pgfsys@useobject{currentmarker}{}%
\end{pgfscope}%
\end{pgfscope}%
\begin{pgfscope}%
\pgfpathrectangle{\pgfqpoint{0.622513in}{0.464855in}}{\pgfqpoint{4.051136in}{2.730000in}}%
\pgfusepath{clip}%
\pgfsetroundcap%
\pgfsetroundjoin%
\pgfsetlinewidth{1.505625pt}%
\definecolor{currentstroke}{rgb}{1.000000,0.498039,0.054902}%
\pgfsetstrokecolor{currentstroke}%
\pgfsetdash{}{0pt}%
\pgfpathmoveto{\pgfqpoint{0.806656in}{1.350730in}}%
\pgfpathlineto{\pgfqpoint{1.052179in}{1.425463in}}%
\pgfpathlineto{\pgfqpoint{1.297703in}{1.303039in}}%
\pgfpathlineto{\pgfqpoint{1.543226in}{1.374100in}}%
\pgfpathlineto{\pgfqpoint{1.788749in}{1.266468in}}%
\pgfpathlineto{\pgfqpoint{2.034273in}{1.327817in}}%
\pgfpathlineto{\pgfqpoint{2.279796in}{1.235118in}}%
\pgfpathlineto{\pgfqpoint{2.525320in}{1.221076in}}%
\pgfpathlineto{\pgfqpoint{2.770843in}{1.208195in}}%
\pgfpathlineto{\pgfqpoint{3.016366in}{1.195375in}}%
\pgfpathlineto{\pgfqpoint{3.261890in}{1.255208in}}%
\pgfpathlineto{\pgfqpoint{3.507413in}{1.241713in}}%
\pgfpathlineto{\pgfqpoint{3.752937in}{1.161875in}}%
\pgfpathlineto{\pgfqpoint{3.998460in}{1.151502in}}%
\pgfpathlineto{\pgfqpoint{4.243984in}{1.140976in}}%
\pgfpathlineto{\pgfqpoint{4.489507in}{1.129677in}}%
\pgfusepath{stroke}%
\end{pgfscope}%
\begin{pgfscope}%
\pgfpathrectangle{\pgfqpoint{0.622513in}{0.464855in}}{\pgfqpoint{4.051136in}{2.730000in}}%
\pgfusepath{clip}%
\pgfsetbuttcap%
\pgfsetroundjoin%
\definecolor{currentfill}{rgb}{1.000000,0.498039,0.054902}%
\pgfsetfillcolor{currentfill}%
\pgfsetlinewidth{1.003750pt}%
\definecolor{currentstroke}{rgb}{1.000000,0.498039,0.054902}%
\pgfsetstrokecolor{currentstroke}%
\pgfsetdash{}{0pt}%
\pgfsys@defobject{currentmarker}{\pgfqpoint{-0.041667in}{-0.041667in}}{\pgfqpoint{0.041667in}{0.041667in}}{%
\pgfpathmoveto{\pgfqpoint{0.000000in}{-0.041667in}}%
\pgfpathcurveto{\pgfqpoint{0.011050in}{-0.041667in}}{\pgfqpoint{0.021649in}{-0.037276in}}{\pgfqpoint{0.029463in}{-0.029463in}}%
\pgfpathcurveto{\pgfqpoint{0.037276in}{-0.021649in}}{\pgfqpoint{0.041667in}{-0.011050in}}{\pgfqpoint{0.041667in}{0.000000in}}%
\pgfpathcurveto{\pgfqpoint{0.041667in}{0.011050in}}{\pgfqpoint{0.037276in}{0.021649in}}{\pgfqpoint{0.029463in}{0.029463in}}%
\pgfpathcurveto{\pgfqpoint{0.021649in}{0.037276in}}{\pgfqpoint{0.011050in}{0.041667in}}{\pgfqpoint{0.000000in}{0.041667in}}%
\pgfpathcurveto{\pgfqpoint{-0.011050in}{0.041667in}}{\pgfqpoint{-0.021649in}{0.037276in}}{\pgfqpoint{-0.029463in}{0.029463in}}%
\pgfpathcurveto{\pgfqpoint{-0.037276in}{0.021649in}}{\pgfqpoint{-0.041667in}{0.011050in}}{\pgfqpoint{-0.041667in}{0.000000in}}%
\pgfpathcurveto{\pgfqpoint{-0.041667in}{-0.011050in}}{\pgfqpoint{-0.037276in}{-0.021649in}}{\pgfqpoint{-0.029463in}{-0.029463in}}%
\pgfpathcurveto{\pgfqpoint{-0.021649in}{-0.037276in}}{\pgfqpoint{-0.011050in}{-0.041667in}}{\pgfqpoint{0.000000in}{-0.041667in}}%
\pgfpathlineto{\pgfqpoint{0.000000in}{-0.041667in}}%
\pgfpathclose%
\pgfusepath{stroke,fill}%
}%
\begin{pgfscope}%
\pgfsys@transformshift{0.806656in}{1.350730in}%
\pgfsys@useobject{currentmarker}{}%
\end{pgfscope}%
\begin{pgfscope}%
\pgfsys@transformshift{1.052179in}{1.425463in}%
\pgfsys@useobject{currentmarker}{}%
\end{pgfscope}%
\begin{pgfscope}%
\pgfsys@transformshift{1.297703in}{1.303039in}%
\pgfsys@useobject{currentmarker}{}%
\end{pgfscope}%
\begin{pgfscope}%
\pgfsys@transformshift{1.543226in}{1.374100in}%
\pgfsys@useobject{currentmarker}{}%
\end{pgfscope}%
\begin{pgfscope}%
\pgfsys@transformshift{1.788749in}{1.266468in}%
\pgfsys@useobject{currentmarker}{}%
\end{pgfscope}%
\begin{pgfscope}%
\pgfsys@transformshift{2.034273in}{1.327817in}%
\pgfsys@useobject{currentmarker}{}%
\end{pgfscope}%
\begin{pgfscope}%
\pgfsys@transformshift{2.279796in}{1.235118in}%
\pgfsys@useobject{currentmarker}{}%
\end{pgfscope}%
\begin{pgfscope}%
\pgfsys@transformshift{2.525320in}{1.221076in}%
\pgfsys@useobject{currentmarker}{}%
\end{pgfscope}%
\begin{pgfscope}%
\pgfsys@transformshift{2.770843in}{1.208195in}%
\pgfsys@useobject{currentmarker}{}%
\end{pgfscope}%
\begin{pgfscope}%
\pgfsys@transformshift{3.016366in}{1.195375in}%
\pgfsys@useobject{currentmarker}{}%
\end{pgfscope}%
\begin{pgfscope}%
\pgfsys@transformshift{3.261890in}{1.255208in}%
\pgfsys@useobject{currentmarker}{}%
\end{pgfscope}%
\begin{pgfscope}%
\pgfsys@transformshift{3.507413in}{1.241713in}%
\pgfsys@useobject{currentmarker}{}%
\end{pgfscope}%
\begin{pgfscope}%
\pgfsys@transformshift{3.752937in}{1.161875in}%
\pgfsys@useobject{currentmarker}{}%
\end{pgfscope}%
\begin{pgfscope}%
\pgfsys@transformshift{3.998460in}{1.151502in}%
\pgfsys@useobject{currentmarker}{}%
\end{pgfscope}%
\begin{pgfscope}%
\pgfsys@transformshift{4.243984in}{1.140976in}%
\pgfsys@useobject{currentmarker}{}%
\end{pgfscope}%
\begin{pgfscope}%
\pgfsys@transformshift{4.489507in}{1.129677in}%
\pgfsys@useobject{currentmarker}{}%
\end{pgfscope}%
\end{pgfscope}%
\begin{pgfscope}%
\pgfpathrectangle{\pgfqpoint{0.622513in}{0.464855in}}{\pgfqpoint{4.051136in}{2.730000in}}%
\pgfusepath{clip}%
\pgfsetroundcap%
\pgfsetroundjoin%
\pgfsetlinewidth{1.505625pt}%
\definecolor{currentstroke}{rgb}{0.172549,0.627451,0.172549}%
\pgfsetstrokecolor{currentstroke}%
\pgfsetdash{}{0pt}%
\pgfpathmoveto{\pgfqpoint{0.806656in}{2.659717in}}%
\pgfpathlineto{\pgfqpoint{1.052179in}{3.070764in}}%
\pgfpathlineto{\pgfqpoint{1.297703in}{2.447416in}}%
\pgfpathlineto{\pgfqpoint{1.543226in}{2.764870in}}%
\pgfpathlineto{\pgfqpoint{1.788749in}{2.331948in}}%
\pgfpathlineto{\pgfqpoint{2.034273in}{2.707491in}}%
\pgfpathlineto{\pgfqpoint{2.279796in}{2.166162in}}%
\pgfpathlineto{\pgfqpoint{2.525320in}{2.438753in}}%
\pgfpathlineto{\pgfqpoint{2.770843in}{2.188609in}}%
\pgfpathlineto{\pgfqpoint{3.016366in}{2.090595in}}%
\pgfpathlineto{\pgfqpoint{3.261890in}{2.630463in}}%
\pgfpathlineto{\pgfqpoint{3.507413in}{2.615443in}}%
\pgfpathlineto{\pgfqpoint{3.752937in}{1.902856in}}%
\pgfpathlineto{\pgfqpoint{3.998460in}{2.437355in}}%
\pgfpathlineto{\pgfqpoint{4.243984in}{1.601285in}}%
\pgfpathlineto{\pgfqpoint{4.489507in}{1.450490in}}%
\pgfusepath{stroke}%
\end{pgfscope}%
\begin{pgfscope}%
\pgfpathrectangle{\pgfqpoint{0.622513in}{0.464855in}}{\pgfqpoint{4.051136in}{2.730000in}}%
\pgfusepath{clip}%
\pgfsetbuttcap%
\pgfsetroundjoin%
\definecolor{currentfill}{rgb}{0.172549,0.627451,0.172549}%
\pgfsetfillcolor{currentfill}%
\pgfsetlinewidth{1.003750pt}%
\definecolor{currentstroke}{rgb}{0.172549,0.627451,0.172549}%
\pgfsetstrokecolor{currentstroke}%
\pgfsetdash{}{0pt}%
\pgfsys@defobject{currentmarker}{\pgfqpoint{-0.041667in}{-0.041667in}}{\pgfqpoint{0.041667in}{0.041667in}}{%
\pgfpathmoveto{\pgfqpoint{0.000000in}{-0.041667in}}%
\pgfpathcurveto{\pgfqpoint{0.011050in}{-0.041667in}}{\pgfqpoint{0.021649in}{-0.037276in}}{\pgfqpoint{0.029463in}{-0.029463in}}%
\pgfpathcurveto{\pgfqpoint{0.037276in}{-0.021649in}}{\pgfqpoint{0.041667in}{-0.011050in}}{\pgfqpoint{0.041667in}{0.000000in}}%
\pgfpathcurveto{\pgfqpoint{0.041667in}{0.011050in}}{\pgfqpoint{0.037276in}{0.021649in}}{\pgfqpoint{0.029463in}{0.029463in}}%
\pgfpathcurveto{\pgfqpoint{0.021649in}{0.037276in}}{\pgfqpoint{0.011050in}{0.041667in}}{\pgfqpoint{0.000000in}{0.041667in}}%
\pgfpathcurveto{\pgfqpoint{-0.011050in}{0.041667in}}{\pgfqpoint{-0.021649in}{0.037276in}}{\pgfqpoint{-0.029463in}{0.029463in}}%
\pgfpathcurveto{\pgfqpoint{-0.037276in}{0.021649in}}{\pgfqpoint{-0.041667in}{0.011050in}}{\pgfqpoint{-0.041667in}{0.000000in}}%
\pgfpathcurveto{\pgfqpoint{-0.041667in}{-0.011050in}}{\pgfqpoint{-0.037276in}{-0.021649in}}{\pgfqpoint{-0.029463in}{-0.029463in}}%
\pgfpathcurveto{\pgfqpoint{-0.021649in}{-0.037276in}}{\pgfqpoint{-0.011050in}{-0.041667in}}{\pgfqpoint{0.000000in}{-0.041667in}}%
\pgfpathlineto{\pgfqpoint{0.000000in}{-0.041667in}}%
\pgfpathclose%
\pgfusepath{stroke,fill}%
}%
\begin{pgfscope}%
\pgfsys@transformshift{0.806656in}{2.659717in}%
\pgfsys@useobject{currentmarker}{}%
\end{pgfscope}%
\begin{pgfscope}%
\pgfsys@transformshift{1.052179in}{3.070764in}%
\pgfsys@useobject{currentmarker}{}%
\end{pgfscope}%
\begin{pgfscope}%
\pgfsys@transformshift{1.297703in}{2.447416in}%
\pgfsys@useobject{currentmarker}{}%
\end{pgfscope}%
\begin{pgfscope}%
\pgfsys@transformshift{1.543226in}{2.764870in}%
\pgfsys@useobject{currentmarker}{}%
\end{pgfscope}%
\begin{pgfscope}%
\pgfsys@transformshift{1.788749in}{2.331948in}%
\pgfsys@useobject{currentmarker}{}%
\end{pgfscope}%
\begin{pgfscope}%
\pgfsys@transformshift{2.034273in}{2.707491in}%
\pgfsys@useobject{currentmarker}{}%
\end{pgfscope}%
\begin{pgfscope}%
\pgfsys@transformshift{2.279796in}{2.166162in}%
\pgfsys@useobject{currentmarker}{}%
\end{pgfscope}%
\begin{pgfscope}%
\pgfsys@transformshift{2.525320in}{2.438753in}%
\pgfsys@useobject{currentmarker}{}%
\end{pgfscope}%
\begin{pgfscope}%
\pgfsys@transformshift{2.770843in}{2.188609in}%
\pgfsys@useobject{currentmarker}{}%
\end{pgfscope}%
\begin{pgfscope}%
\pgfsys@transformshift{3.016366in}{2.090595in}%
\pgfsys@useobject{currentmarker}{}%
\end{pgfscope}%
\begin{pgfscope}%
\pgfsys@transformshift{3.261890in}{2.630463in}%
\pgfsys@useobject{currentmarker}{}%
\end{pgfscope}%
\begin{pgfscope}%
\pgfsys@transformshift{3.507413in}{2.615443in}%
\pgfsys@useobject{currentmarker}{}%
\end{pgfscope}%
\begin{pgfscope}%
\pgfsys@transformshift{3.752937in}{1.902856in}%
\pgfsys@useobject{currentmarker}{}%
\end{pgfscope}%
\begin{pgfscope}%
\pgfsys@transformshift{3.998460in}{2.437355in}%
\pgfsys@useobject{currentmarker}{}%
\end{pgfscope}%
\begin{pgfscope}%
\pgfsys@transformshift{4.243984in}{1.601285in}%
\pgfsys@useobject{currentmarker}{}%
\end{pgfscope}%
\begin{pgfscope}%
\pgfsys@transformshift{4.489507in}{1.450490in}%
\pgfsys@useobject{currentmarker}{}%
\end{pgfscope}%
\end{pgfscope}%
\begin{pgfscope}%
\pgfpathrectangle{\pgfqpoint{0.622513in}{0.464855in}}{\pgfqpoint{4.051136in}{2.730000in}}%
\pgfusepath{clip}%
\pgfsetroundcap%
\pgfsetroundjoin%
\pgfsetlinewidth{1.505625pt}%
\definecolor{currentstroke}{rgb}{0.839216,0.152941,0.156863}%
\pgfsetstrokecolor{currentstroke}%
\pgfsetdash{}{0pt}%
\pgfpathmoveto{\pgfqpoint{0.806656in}{2.431418in}}%
\pgfpathlineto{\pgfqpoint{1.052179in}{2.878870in}}%
\pgfpathlineto{\pgfqpoint{1.297703in}{2.154591in}}%
\pgfpathlineto{\pgfqpoint{1.543226in}{2.832124in}}%
\pgfpathlineto{\pgfqpoint{1.788749in}{2.142110in}}%
\pgfpathlineto{\pgfqpoint{2.034273in}{2.791223in}}%
\pgfpathlineto{\pgfqpoint{2.279796in}{2.115966in}}%
\pgfpathlineto{\pgfqpoint{2.525320in}{2.126279in}}%
\pgfpathlineto{\pgfqpoint{2.770843in}{1.562774in}}%
\pgfpathlineto{\pgfqpoint{3.016366in}{2.015520in}}%
\pgfpathlineto{\pgfqpoint{3.261890in}{2.679113in}}%
\pgfpathlineto{\pgfqpoint{3.507413in}{2.678329in}}%
\pgfpathlineto{\pgfqpoint{3.752937in}{1.524761in}}%
\pgfpathlineto{\pgfqpoint{3.998460in}{1.788428in}}%
\pgfpathlineto{\pgfqpoint{4.243984in}{1.545053in}}%
\pgfpathlineto{\pgfqpoint{4.489507in}{1.222850in}}%
\pgfusepath{stroke}%
\end{pgfscope}%
\begin{pgfscope}%
\pgfpathrectangle{\pgfqpoint{0.622513in}{0.464855in}}{\pgfqpoint{4.051136in}{2.730000in}}%
\pgfusepath{clip}%
\pgfsetbuttcap%
\pgfsetroundjoin%
\definecolor{currentfill}{rgb}{0.839216,0.152941,0.156863}%
\pgfsetfillcolor{currentfill}%
\pgfsetlinewidth{1.003750pt}%
\definecolor{currentstroke}{rgb}{0.839216,0.152941,0.156863}%
\pgfsetstrokecolor{currentstroke}%
\pgfsetdash{}{0pt}%
\pgfsys@defobject{currentmarker}{\pgfqpoint{-0.041667in}{-0.041667in}}{\pgfqpoint{0.041667in}{0.041667in}}{%
\pgfpathmoveto{\pgfqpoint{0.000000in}{-0.041667in}}%
\pgfpathcurveto{\pgfqpoint{0.011050in}{-0.041667in}}{\pgfqpoint{0.021649in}{-0.037276in}}{\pgfqpoint{0.029463in}{-0.029463in}}%
\pgfpathcurveto{\pgfqpoint{0.037276in}{-0.021649in}}{\pgfqpoint{0.041667in}{-0.011050in}}{\pgfqpoint{0.041667in}{0.000000in}}%
\pgfpathcurveto{\pgfqpoint{0.041667in}{0.011050in}}{\pgfqpoint{0.037276in}{0.021649in}}{\pgfqpoint{0.029463in}{0.029463in}}%
\pgfpathcurveto{\pgfqpoint{0.021649in}{0.037276in}}{\pgfqpoint{0.011050in}{0.041667in}}{\pgfqpoint{0.000000in}{0.041667in}}%
\pgfpathcurveto{\pgfqpoint{-0.011050in}{0.041667in}}{\pgfqpoint{-0.021649in}{0.037276in}}{\pgfqpoint{-0.029463in}{0.029463in}}%
\pgfpathcurveto{\pgfqpoint{-0.037276in}{0.021649in}}{\pgfqpoint{-0.041667in}{0.011050in}}{\pgfqpoint{-0.041667in}{0.000000in}}%
\pgfpathcurveto{\pgfqpoint{-0.041667in}{-0.011050in}}{\pgfqpoint{-0.037276in}{-0.021649in}}{\pgfqpoint{-0.029463in}{-0.029463in}}%
\pgfpathcurveto{\pgfqpoint{-0.021649in}{-0.037276in}}{\pgfqpoint{-0.011050in}{-0.041667in}}{\pgfqpoint{0.000000in}{-0.041667in}}%
\pgfpathlineto{\pgfqpoint{0.000000in}{-0.041667in}}%
\pgfpathclose%
\pgfusepath{stroke,fill}%
}%
\begin{pgfscope}%
\pgfsys@transformshift{0.806656in}{2.431418in}%
\pgfsys@useobject{currentmarker}{}%
\end{pgfscope}%
\begin{pgfscope}%
\pgfsys@transformshift{1.052179in}{2.878870in}%
\pgfsys@useobject{currentmarker}{}%
\end{pgfscope}%
\begin{pgfscope}%
\pgfsys@transformshift{1.297703in}{2.154591in}%
\pgfsys@useobject{currentmarker}{}%
\end{pgfscope}%
\begin{pgfscope}%
\pgfsys@transformshift{1.543226in}{2.832124in}%
\pgfsys@useobject{currentmarker}{}%
\end{pgfscope}%
\begin{pgfscope}%
\pgfsys@transformshift{1.788749in}{2.142110in}%
\pgfsys@useobject{currentmarker}{}%
\end{pgfscope}%
\begin{pgfscope}%
\pgfsys@transformshift{2.034273in}{2.791223in}%
\pgfsys@useobject{currentmarker}{}%
\end{pgfscope}%
\begin{pgfscope}%
\pgfsys@transformshift{2.279796in}{2.115966in}%
\pgfsys@useobject{currentmarker}{}%
\end{pgfscope}%
\begin{pgfscope}%
\pgfsys@transformshift{2.525320in}{2.126279in}%
\pgfsys@useobject{currentmarker}{}%
\end{pgfscope}%
\begin{pgfscope}%
\pgfsys@transformshift{2.770843in}{1.562774in}%
\pgfsys@useobject{currentmarker}{}%
\end{pgfscope}%
\begin{pgfscope}%
\pgfsys@transformshift{3.016366in}{2.015520in}%
\pgfsys@useobject{currentmarker}{}%
\end{pgfscope}%
\begin{pgfscope}%
\pgfsys@transformshift{3.261890in}{2.679113in}%
\pgfsys@useobject{currentmarker}{}%
\end{pgfscope}%
\begin{pgfscope}%
\pgfsys@transformshift{3.507413in}{2.678329in}%
\pgfsys@useobject{currentmarker}{}%
\end{pgfscope}%
\begin{pgfscope}%
\pgfsys@transformshift{3.752937in}{1.524761in}%
\pgfsys@useobject{currentmarker}{}%
\end{pgfscope}%
\begin{pgfscope}%
\pgfsys@transformshift{3.998460in}{1.788428in}%
\pgfsys@useobject{currentmarker}{}%
\end{pgfscope}%
\begin{pgfscope}%
\pgfsys@transformshift{4.243984in}{1.545053in}%
\pgfsys@useobject{currentmarker}{}%
\end{pgfscope}%
\begin{pgfscope}%
\pgfsys@transformshift{4.489507in}{1.222850in}%
\pgfsys@useobject{currentmarker}{}%
\end{pgfscope}%
\end{pgfscope}%
\begin{pgfscope}%
\pgfpathrectangle{\pgfqpoint{0.622513in}{0.464855in}}{\pgfqpoint{4.051136in}{2.730000in}}%
\pgfusepath{clip}%
\pgfsetroundcap%
\pgfsetroundjoin%
\pgfsetlinewidth{1.505625pt}%
\definecolor{currentstroke}{rgb}{0.580392,0.403922,0.741176}%
\pgfsetstrokecolor{currentstroke}%
\pgfsetdash{}{0pt}%
\pgfpathmoveto{\pgfqpoint{0.806656in}{0.588946in}}%
\pgfpathlineto{\pgfqpoint{1.052179in}{0.977623in}}%
\pgfpathlineto{\pgfqpoint{1.297703in}{0.655555in}}%
\pgfpathlineto{\pgfqpoint{1.543226in}{1.012057in}}%
\pgfpathlineto{\pgfqpoint{1.788749in}{0.623034in}}%
\pgfpathlineto{\pgfqpoint{2.034273in}{0.988429in}}%
\pgfpathlineto{\pgfqpoint{2.279796in}{0.598855in}}%
\pgfpathlineto{\pgfqpoint{2.525320in}{0.628790in}}%
\pgfpathlineto{\pgfqpoint{2.770843in}{0.613021in}}%
\pgfpathlineto{\pgfqpoint{3.016366in}{0.629963in}}%
\pgfpathlineto{\pgfqpoint{3.261890in}{0.913541in}}%
\pgfpathlineto{\pgfqpoint{3.507413in}{0.900649in}}%
\pgfpathlineto{\pgfqpoint{3.752937in}{0.619592in}}%
\pgfpathlineto{\pgfqpoint{3.998460in}{0.665168in}}%
\pgfpathlineto{\pgfqpoint{4.243984in}{0.590909in}}%
\pgfpathlineto{\pgfqpoint{4.489507in}{0.637643in}}%
\pgfusepath{stroke}%
\end{pgfscope}%
\begin{pgfscope}%
\pgfpathrectangle{\pgfqpoint{0.622513in}{0.464855in}}{\pgfqpoint{4.051136in}{2.730000in}}%
\pgfusepath{clip}%
\pgfsetbuttcap%
\pgfsetroundjoin%
\definecolor{currentfill}{rgb}{0.580392,0.403922,0.741176}%
\pgfsetfillcolor{currentfill}%
\pgfsetlinewidth{1.003750pt}%
\definecolor{currentstroke}{rgb}{0.580392,0.403922,0.741176}%
\pgfsetstrokecolor{currentstroke}%
\pgfsetdash{}{0pt}%
\pgfsys@defobject{currentmarker}{\pgfqpoint{-0.041667in}{-0.041667in}}{\pgfqpoint{0.041667in}{0.041667in}}{%
\pgfpathmoveto{\pgfqpoint{0.000000in}{-0.041667in}}%
\pgfpathcurveto{\pgfqpoint{0.011050in}{-0.041667in}}{\pgfqpoint{0.021649in}{-0.037276in}}{\pgfqpoint{0.029463in}{-0.029463in}}%
\pgfpathcurveto{\pgfqpoint{0.037276in}{-0.021649in}}{\pgfqpoint{0.041667in}{-0.011050in}}{\pgfqpoint{0.041667in}{0.000000in}}%
\pgfpathcurveto{\pgfqpoint{0.041667in}{0.011050in}}{\pgfqpoint{0.037276in}{0.021649in}}{\pgfqpoint{0.029463in}{0.029463in}}%
\pgfpathcurveto{\pgfqpoint{0.021649in}{0.037276in}}{\pgfqpoint{0.011050in}{0.041667in}}{\pgfqpoint{0.000000in}{0.041667in}}%
\pgfpathcurveto{\pgfqpoint{-0.011050in}{0.041667in}}{\pgfqpoint{-0.021649in}{0.037276in}}{\pgfqpoint{-0.029463in}{0.029463in}}%
\pgfpathcurveto{\pgfqpoint{-0.037276in}{0.021649in}}{\pgfqpoint{-0.041667in}{0.011050in}}{\pgfqpoint{-0.041667in}{0.000000in}}%
\pgfpathcurveto{\pgfqpoint{-0.041667in}{-0.011050in}}{\pgfqpoint{-0.037276in}{-0.021649in}}{\pgfqpoint{-0.029463in}{-0.029463in}}%
\pgfpathcurveto{\pgfqpoint{-0.021649in}{-0.037276in}}{\pgfqpoint{-0.011050in}{-0.041667in}}{\pgfqpoint{0.000000in}{-0.041667in}}%
\pgfpathlineto{\pgfqpoint{0.000000in}{-0.041667in}}%
\pgfpathclose%
\pgfusepath{stroke,fill}%
}%
\begin{pgfscope}%
\pgfsys@transformshift{0.806656in}{0.588946in}%
\pgfsys@useobject{currentmarker}{}%
\end{pgfscope}%
\begin{pgfscope}%
\pgfsys@transformshift{1.052179in}{0.977623in}%
\pgfsys@useobject{currentmarker}{}%
\end{pgfscope}%
\begin{pgfscope}%
\pgfsys@transformshift{1.297703in}{0.655555in}%
\pgfsys@useobject{currentmarker}{}%
\end{pgfscope}%
\begin{pgfscope}%
\pgfsys@transformshift{1.543226in}{1.012057in}%
\pgfsys@useobject{currentmarker}{}%
\end{pgfscope}%
\begin{pgfscope}%
\pgfsys@transformshift{1.788749in}{0.623034in}%
\pgfsys@useobject{currentmarker}{}%
\end{pgfscope}%
\begin{pgfscope}%
\pgfsys@transformshift{2.034273in}{0.988429in}%
\pgfsys@useobject{currentmarker}{}%
\end{pgfscope}%
\begin{pgfscope}%
\pgfsys@transformshift{2.279796in}{0.598855in}%
\pgfsys@useobject{currentmarker}{}%
\end{pgfscope}%
\begin{pgfscope}%
\pgfsys@transformshift{2.525320in}{0.628790in}%
\pgfsys@useobject{currentmarker}{}%
\end{pgfscope}%
\begin{pgfscope}%
\pgfsys@transformshift{2.770843in}{0.613021in}%
\pgfsys@useobject{currentmarker}{}%
\end{pgfscope}%
\begin{pgfscope}%
\pgfsys@transformshift{3.016366in}{0.629963in}%
\pgfsys@useobject{currentmarker}{}%
\end{pgfscope}%
\begin{pgfscope}%
\pgfsys@transformshift{3.261890in}{0.913541in}%
\pgfsys@useobject{currentmarker}{}%
\end{pgfscope}%
\begin{pgfscope}%
\pgfsys@transformshift{3.507413in}{0.900649in}%
\pgfsys@useobject{currentmarker}{}%
\end{pgfscope}%
\begin{pgfscope}%
\pgfsys@transformshift{3.752937in}{0.619592in}%
\pgfsys@useobject{currentmarker}{}%
\end{pgfscope}%
\begin{pgfscope}%
\pgfsys@transformshift{3.998460in}{0.665168in}%
\pgfsys@useobject{currentmarker}{}%
\end{pgfscope}%
\begin{pgfscope}%
\pgfsys@transformshift{4.243984in}{0.590909in}%
\pgfsys@useobject{currentmarker}{}%
\end{pgfscope}%
\begin{pgfscope}%
\pgfsys@transformshift{4.489507in}{0.637643in}%
\pgfsys@useobject{currentmarker}{}%
\end{pgfscope}%
\end{pgfscope}%
\begin{pgfscope}%
\pgfpathrectangle{\pgfqpoint{0.622513in}{0.464855in}}{\pgfqpoint{4.051136in}{2.730000in}}%
\pgfusepath{clip}%
\pgfsetbuttcap%
\pgfsetroundjoin%
\pgfsetlinewidth{1.505625pt}%
\definecolor{currentstroke}{rgb}{0.000000,0.000000,0.000000}%
\pgfsetstrokecolor{currentstroke}%
\pgfsetdash{{5.550000pt}{2.400000pt}}{0.000000pt}%
\pgfpathmoveto{\pgfqpoint{0.622513in}{2.207157in}}%
\pgfpathlineto{\pgfqpoint{4.673649in}{2.207157in}}%
\pgfusepath{stroke}%
\end{pgfscope}%
\begin{pgfscope}%
\pgfsetrectcap%
\pgfsetmiterjoin%
\pgfsetlinewidth{1.003750pt}%
\definecolor{currentstroke}{rgb}{0.800000,0.800000,0.800000}%
\pgfsetstrokecolor{currentstroke}%
\pgfsetdash{}{0pt}%
\pgfpathmoveto{\pgfqpoint{0.622513in}{0.464855in}}%
\pgfpathlineto{\pgfqpoint{0.622513in}{3.194855in}}%
\pgfusepath{stroke}%
\end{pgfscope}%
\begin{pgfscope}%
\pgfsetrectcap%
\pgfsetmiterjoin%
\pgfsetlinewidth{1.003750pt}%
\definecolor{currentstroke}{rgb}{0.800000,0.800000,0.800000}%
\pgfsetstrokecolor{currentstroke}%
\pgfsetdash{}{0pt}%
\pgfpathmoveto{\pgfqpoint{4.673649in}{0.464855in}}%
\pgfpathlineto{\pgfqpoint{4.673649in}{3.194855in}}%
\pgfusepath{stroke}%
\end{pgfscope}%
\begin{pgfscope}%
\pgfsetrectcap%
\pgfsetmiterjoin%
\pgfsetlinewidth{1.003750pt}%
\definecolor{currentstroke}{rgb}{0.800000,0.800000,0.800000}%
\pgfsetstrokecolor{currentstroke}%
\pgfsetdash{}{0pt}%
\pgfpathmoveto{\pgfqpoint{0.622513in}{0.464855in}}%
\pgfpathlineto{\pgfqpoint{4.673649in}{0.464855in}}%
\pgfusepath{stroke}%
\end{pgfscope}%
\begin{pgfscope}%
\pgfsetrectcap%
\pgfsetmiterjoin%
\pgfsetlinewidth{1.003750pt}%
\definecolor{currentstroke}{rgb}{0.800000,0.800000,0.800000}%
\pgfsetstrokecolor{currentstroke}%
\pgfsetdash{}{0pt}%
\pgfpathmoveto{\pgfqpoint{0.622513in}{3.194855in}}%
\pgfpathlineto{\pgfqpoint{4.673649in}{3.194855in}}%
\pgfusepath{stroke}%
\end{pgfscope}%
\begin{pgfscope}%
\definecolor{textcolor}{rgb}{0.150000,0.150000,0.150000}%
\pgfsetstrokecolor{textcolor}%
\pgfsetfillcolor{textcolor}%
\pgftext[x=2.648081in,y=3.278188in,,base]{\color{textcolor}{\sffamily\fontsize{12.000000}{14.400000}\selectfont\catcode`\^=\active\def^{\ifmmode\sp\else\^{}\fi}\catcode`\%=\active\def%{\%}Physical residuals normalized by their gates}}%
\end{pgfscope}%
\begin{pgfscope}%
\pgfsetroundcap%
\pgfsetroundjoin%
\pgfsetlinewidth{1.505625pt}%
\definecolor{currentstroke}{rgb}{0.121569,0.466667,0.705882}%
\pgfsetstrokecolor{currentstroke}%
\pgfsetdash{}{0pt}%
\pgfpathmoveto{\pgfqpoint{3.336106in}{3.052852in}}%
\pgfpathlineto{\pgfqpoint{3.447217in}{3.052852in}}%
\pgfpathlineto{\pgfqpoint{3.558328in}{3.052852in}}%
\pgfusepath{stroke}%
\end{pgfscope}%
\begin{pgfscope}%
\pgfsetbuttcap%
\pgfsetroundjoin%
\definecolor{currentfill}{rgb}{0.121569,0.466667,0.705882}%
\pgfsetfillcolor{currentfill}%
\pgfsetlinewidth{1.003750pt}%
\definecolor{currentstroke}{rgb}{0.121569,0.466667,0.705882}%
\pgfsetstrokecolor{currentstroke}%
\pgfsetdash{}{0pt}%
\pgfsys@defobject{currentmarker}{\pgfqpoint{-0.041667in}{-0.041667in}}{\pgfqpoint{0.041667in}{0.041667in}}{%
\pgfpathmoveto{\pgfqpoint{0.000000in}{-0.041667in}}%
\pgfpathcurveto{\pgfqpoint{0.011050in}{-0.041667in}}{\pgfqpoint{0.021649in}{-0.037276in}}{\pgfqpoint{0.029463in}{-0.029463in}}%
\pgfpathcurveto{\pgfqpoint{0.037276in}{-0.021649in}}{\pgfqpoint{0.041667in}{-0.011050in}}{\pgfqpoint{0.041667in}{0.000000in}}%
\pgfpathcurveto{\pgfqpoint{0.041667in}{0.011050in}}{\pgfqpoint{0.037276in}{0.021649in}}{\pgfqpoint{0.029463in}{0.029463in}}%
\pgfpathcurveto{\pgfqpoint{0.021649in}{0.037276in}}{\pgfqpoint{0.011050in}{0.041667in}}{\pgfqpoint{0.000000in}{0.041667in}}%
\pgfpathcurveto{\pgfqpoint{-0.011050in}{0.041667in}}{\pgfqpoint{-0.021649in}{0.037276in}}{\pgfqpoint{-0.029463in}{0.029463in}}%
\pgfpathcurveto{\pgfqpoint{-0.037276in}{0.021649in}}{\pgfqpoint{-0.041667in}{0.011050in}}{\pgfqpoint{-0.041667in}{0.000000in}}%
\pgfpathcurveto{\pgfqpoint{-0.041667in}{-0.011050in}}{\pgfqpoint{-0.037276in}{-0.021649in}}{\pgfqpoint{-0.029463in}{-0.029463in}}%
\pgfpathcurveto{\pgfqpoint{-0.021649in}{-0.037276in}}{\pgfqpoint{-0.011050in}{-0.041667in}}{\pgfqpoint{0.000000in}{-0.041667in}}%
\pgfpathlineto{\pgfqpoint{0.000000in}{-0.041667in}}%
\pgfpathclose%
\pgfusepath{stroke,fill}%
}%
\begin{pgfscope}%
\pgfsys@transformshift{3.447217in}{3.052852in}%
\pgfsys@useobject{currentmarker}{}%
\end{pgfscope}%
\end{pgfscope}%
\begin{pgfscope}%
\definecolor{textcolor}{rgb}{0.150000,0.150000,0.150000}%
\pgfsetstrokecolor{textcolor}%
\pgfsetfillcolor{textcolor}%
\pgftext[x=3.647217in,y=3.013963in,left,base]{\color{textcolor}{\sffamily\fontsize{8.000000}{9.600000}\selectfont\catcode`\^=\active\def^{\ifmmode\sp\else\^{}\fi}\catcode`\%=\active\def%{\%}gas appearance}}%
\end{pgfscope}%
\begin{pgfscope}%
\pgfsetroundcap%
\pgfsetroundjoin%
\pgfsetlinewidth{1.505625pt}%
\definecolor{currentstroke}{rgb}{1.000000,0.498039,0.054902}%
\pgfsetstrokecolor{currentstroke}%
\pgfsetdash{}{0pt}%
\pgfpathmoveto{\pgfqpoint{3.336106in}{2.893021in}}%
\pgfpathlineto{\pgfqpoint{3.447217in}{2.893021in}}%
\pgfpathlineto{\pgfqpoint{3.558328in}{2.893021in}}%
\pgfusepath{stroke}%
\end{pgfscope}%
\begin{pgfscope}%
\pgfsetbuttcap%
\pgfsetroundjoin%
\definecolor{currentfill}{rgb}{1.000000,0.498039,0.054902}%
\pgfsetfillcolor{currentfill}%
\pgfsetlinewidth{1.003750pt}%
\definecolor{currentstroke}{rgb}{1.000000,0.498039,0.054902}%
\pgfsetstrokecolor{currentstroke}%
\pgfsetdash{}{0pt}%
\pgfsys@defobject{currentmarker}{\pgfqpoint{-0.041667in}{-0.041667in}}{\pgfqpoint{0.041667in}{0.041667in}}{%
\pgfpathmoveto{\pgfqpoint{0.000000in}{-0.041667in}}%
\pgfpathcurveto{\pgfqpoint{0.011050in}{-0.041667in}}{\pgfqpoint{0.021649in}{-0.037276in}}{\pgfqpoint{0.029463in}{-0.029463in}}%
\pgfpathcurveto{\pgfqpoint{0.037276in}{-0.021649in}}{\pgfqpoint{0.041667in}{-0.011050in}}{\pgfqpoint{0.041667in}{0.000000in}}%
\pgfpathcurveto{\pgfqpoint{0.041667in}{0.011050in}}{\pgfqpoint{0.037276in}{0.021649in}}{\pgfqpoint{0.029463in}{0.029463in}}%
\pgfpathcurveto{\pgfqpoint{0.021649in}{0.037276in}}{\pgfqpoint{0.011050in}{0.041667in}}{\pgfqpoint{0.000000in}{0.041667in}}%
\pgfpathcurveto{\pgfqpoint{-0.011050in}{0.041667in}}{\pgfqpoint{-0.021649in}{0.037276in}}{\pgfqpoint{-0.029463in}{0.029463in}}%
\pgfpathcurveto{\pgfqpoint{-0.037276in}{0.021649in}}{\pgfqpoint{-0.041667in}{0.011050in}}{\pgfqpoint{-0.041667in}{0.000000in}}%
\pgfpathcurveto{\pgfqpoint{-0.041667in}{-0.011050in}}{\pgfqpoint{-0.037276in}{-0.021649in}}{\pgfqpoint{-0.029463in}{-0.029463in}}%
\pgfpathcurveto{\pgfqpoint{-0.021649in}{-0.037276in}}{\pgfqpoint{-0.011050in}{-0.041667in}}{\pgfqpoint{0.000000in}{-0.041667in}}%
\pgfpathlineto{\pgfqpoint{0.000000in}{-0.041667in}}%
\pgfpathclose%
\pgfusepath{stroke,fill}%
}%
\begin{pgfscope}%
\pgfsys@transformshift{3.447217in}{2.893021in}%
\pgfsys@useobject{currentmarker}{}%
\end{pgfscope}%
\end{pgfscope}%
\begin{pgfscope}%
\definecolor{textcolor}{rgb}{0.150000,0.150000,0.150000}%
\pgfsetstrokecolor{textcolor}%
\pgfsetfillcolor{textcolor}%
\pgftext[x=3.647217in,y=2.854132in,left,base]{\color{textcolor}{\sffamily\fontsize{8.000000}{9.600000}\selectfont\catcode`\^=\active\def^{\ifmmode\sp\else\^{}\fi}\catcode`\%=\active\def%{\%}transfer coordinate}}%
\end{pgfscope}%
\begin{pgfscope}%
\pgfsetroundcap%
\pgfsetroundjoin%
\pgfsetlinewidth{1.505625pt}%
\definecolor{currentstroke}{rgb}{0.172549,0.627451,0.172549}%
\pgfsetstrokecolor{currentstroke}%
\pgfsetdash{}{0pt}%
\pgfpathmoveto{\pgfqpoint{3.336106in}{2.733190in}}%
\pgfpathlineto{\pgfqpoint{3.447217in}{2.733190in}}%
\pgfpathlineto{\pgfqpoint{3.558328in}{2.733190in}}%
\pgfusepath{stroke}%
\end{pgfscope}%
\begin{pgfscope}%
\pgfsetbuttcap%
\pgfsetroundjoin%
\definecolor{currentfill}{rgb}{0.172549,0.627451,0.172549}%
\pgfsetfillcolor{currentfill}%
\pgfsetlinewidth{1.003750pt}%
\definecolor{currentstroke}{rgb}{0.172549,0.627451,0.172549}%
\pgfsetstrokecolor{currentstroke}%
\pgfsetdash{}{0pt}%
\pgfsys@defobject{currentmarker}{\pgfqpoint{-0.041667in}{-0.041667in}}{\pgfqpoint{0.041667in}{0.041667in}}{%
\pgfpathmoveto{\pgfqpoint{0.000000in}{-0.041667in}}%
\pgfpathcurveto{\pgfqpoint{0.011050in}{-0.041667in}}{\pgfqpoint{0.021649in}{-0.037276in}}{\pgfqpoint{0.029463in}{-0.029463in}}%
\pgfpathcurveto{\pgfqpoint{0.037276in}{-0.021649in}}{\pgfqpoint{0.041667in}{-0.011050in}}{\pgfqpoint{0.041667in}{0.000000in}}%
\pgfpathcurveto{\pgfqpoint{0.041667in}{0.011050in}}{\pgfqpoint{0.037276in}{0.021649in}}{\pgfqpoint{0.029463in}{0.029463in}}%
\pgfpathcurveto{\pgfqpoint{0.021649in}{0.037276in}}{\pgfqpoint{0.011050in}{0.041667in}}{\pgfqpoint{0.000000in}{0.041667in}}%
\pgfpathcurveto{\pgfqpoint{-0.011050in}{0.041667in}}{\pgfqpoint{-0.021649in}{0.037276in}}{\pgfqpoint{-0.029463in}{0.029463in}}%
\pgfpathcurveto{\pgfqpoint{-0.037276in}{0.021649in}}{\pgfqpoint{-0.041667in}{0.011050in}}{\pgfqpoint{-0.041667in}{0.000000in}}%
\pgfpathcurveto{\pgfqpoint{-0.041667in}{-0.011050in}}{\pgfqpoint{-0.037276in}{-0.021649in}}{\pgfqpoint{-0.029463in}{-0.029463in}}%
\pgfpathcurveto{\pgfqpoint{-0.021649in}{-0.037276in}}{\pgfqpoint{-0.011050in}{-0.041667in}}{\pgfqpoint{0.000000in}{-0.041667in}}%
\pgfpathlineto{\pgfqpoint{0.000000in}{-0.041667in}}%
\pgfpathclose%
\pgfusepath{stroke,fill}%
}%
\begin{pgfscope}%
\pgfsys@transformshift{3.447217in}{2.733190in}%
\pgfsys@useobject{currentmarker}{}%
\end{pgfscope}%
\end{pgfscope}%
\begin{pgfscope}%
\definecolor{textcolor}{rgb}{0.150000,0.150000,0.150000}%
\pgfsetstrokecolor{textcolor}%
\pgfsetfillcolor{textcolor}%
\pgftext[x=3.647217in,y=2.694301in,left,base]{\color{textcolor}{\sffamily\fontsize{8.000000}{9.600000}\selectfont\catcode`\^=\active\def^{\ifmmode\sp\else\^{}\fi}\catcode`\%=\active\def%{\%}momentum}}%
\end{pgfscope}%
\begin{pgfscope}%
\pgfsetroundcap%
\pgfsetroundjoin%
\pgfsetlinewidth{1.505625pt}%
\definecolor{currentstroke}{rgb}{0.839216,0.152941,0.156863}%
\pgfsetstrokecolor{currentstroke}%
\pgfsetdash{}{0pt}%
\pgfpathmoveto{\pgfqpoint{3.336106in}{2.573359in}}%
\pgfpathlineto{\pgfqpoint{3.447217in}{2.573359in}}%
\pgfpathlineto{\pgfqpoint{3.558328in}{2.573359in}}%
\pgfusepath{stroke}%
\end{pgfscope}%
\begin{pgfscope}%
\pgfsetbuttcap%
\pgfsetroundjoin%
\definecolor{currentfill}{rgb}{0.839216,0.152941,0.156863}%
\pgfsetfillcolor{currentfill}%
\pgfsetlinewidth{1.003750pt}%
\definecolor{currentstroke}{rgb}{0.839216,0.152941,0.156863}%
\pgfsetstrokecolor{currentstroke}%
\pgfsetdash{}{0pt}%
\pgfsys@defobject{currentmarker}{\pgfqpoint{-0.041667in}{-0.041667in}}{\pgfqpoint{0.041667in}{0.041667in}}{%
\pgfpathmoveto{\pgfqpoint{0.000000in}{-0.041667in}}%
\pgfpathcurveto{\pgfqpoint{0.011050in}{-0.041667in}}{\pgfqpoint{0.021649in}{-0.037276in}}{\pgfqpoint{0.029463in}{-0.029463in}}%
\pgfpathcurveto{\pgfqpoint{0.037276in}{-0.021649in}}{\pgfqpoint{0.041667in}{-0.011050in}}{\pgfqpoint{0.041667in}{0.000000in}}%
\pgfpathcurveto{\pgfqpoint{0.041667in}{0.011050in}}{\pgfqpoint{0.037276in}{0.021649in}}{\pgfqpoint{0.029463in}{0.029463in}}%
\pgfpathcurveto{\pgfqpoint{0.021649in}{0.037276in}}{\pgfqpoint{0.011050in}{0.041667in}}{\pgfqpoint{0.000000in}{0.041667in}}%
\pgfpathcurveto{\pgfqpoint{-0.011050in}{0.041667in}}{\pgfqpoint{-0.021649in}{0.037276in}}{\pgfqpoint{-0.029463in}{0.029463in}}%
\pgfpathcurveto{\pgfqpoint{-0.037276in}{0.021649in}}{\pgfqpoint{-0.041667in}{0.011050in}}{\pgfqpoint{-0.041667in}{0.000000in}}%
\pgfpathcurveto{\pgfqpoint{-0.041667in}{-0.011050in}}{\pgfqpoint{-0.037276in}{-0.021649in}}{\pgfqpoint{-0.029463in}{-0.029463in}}%
\pgfpathcurveto{\pgfqpoint{-0.021649in}{-0.037276in}}{\pgfqpoint{-0.011050in}{-0.041667in}}{\pgfqpoint{0.000000in}{-0.041667in}}%
\pgfpathlineto{\pgfqpoint{0.000000in}{-0.041667in}}%
\pgfpathclose%
\pgfusepath{stroke,fill}%
}%
\begin{pgfscope}%
\pgfsys@transformshift{3.447217in}{2.573359in}%
\pgfsys@useobject{currentmarker}{}%
\end{pgfscope}%
\end{pgfscope}%
\begin{pgfscope}%
\definecolor{textcolor}{rgb}{0.150000,0.150000,0.150000}%
\pgfsetstrokecolor{textcolor}%
\pgfsetfillcolor{textcolor}%
\pgftext[x=3.647217in,y=2.534470in,left,base]{\color{textcolor}{\sffamily\fontsize{8.000000}{9.600000}\selectfont\catcode`\^=\active\def^{\ifmmode\sp\else\^{}\fi}\catcode`\%=\active\def%{\%}fluid energy}}%
\end{pgfscope}%
\begin{pgfscope}%
\pgfsetroundcap%
\pgfsetroundjoin%
\pgfsetlinewidth{1.505625pt}%
\definecolor{currentstroke}{rgb}{0.580392,0.403922,0.741176}%
\pgfsetstrokecolor{currentstroke}%
\pgfsetdash{}{0pt}%
\pgfpathmoveto{\pgfqpoint{3.336106in}{2.413529in}}%
\pgfpathlineto{\pgfqpoint{3.447217in}{2.413529in}}%
\pgfpathlineto{\pgfqpoint{3.558328in}{2.413529in}}%
\pgfusepath{stroke}%
\end{pgfscope}%
\begin{pgfscope}%
\pgfsetbuttcap%
\pgfsetroundjoin%
\definecolor{currentfill}{rgb}{0.580392,0.403922,0.741176}%
\pgfsetfillcolor{currentfill}%
\pgfsetlinewidth{1.003750pt}%
\definecolor{currentstroke}{rgb}{0.580392,0.403922,0.741176}%
\pgfsetstrokecolor{currentstroke}%
\pgfsetdash{}{0pt}%
\pgfsys@defobject{currentmarker}{\pgfqpoint{-0.041667in}{-0.041667in}}{\pgfqpoint{0.041667in}{0.041667in}}{%
\pgfpathmoveto{\pgfqpoint{0.000000in}{-0.041667in}}%
\pgfpathcurveto{\pgfqpoint{0.011050in}{-0.041667in}}{\pgfqpoint{0.021649in}{-0.037276in}}{\pgfqpoint{0.029463in}{-0.029463in}}%
\pgfpathcurveto{\pgfqpoint{0.037276in}{-0.021649in}}{\pgfqpoint{0.041667in}{-0.011050in}}{\pgfqpoint{0.041667in}{0.000000in}}%
\pgfpathcurveto{\pgfqpoint{0.041667in}{0.011050in}}{\pgfqpoint{0.037276in}{0.021649in}}{\pgfqpoint{0.029463in}{0.029463in}}%
\pgfpathcurveto{\pgfqpoint{0.021649in}{0.037276in}}{\pgfqpoint{0.011050in}{0.041667in}}{\pgfqpoint{0.000000in}{0.041667in}}%
\pgfpathcurveto{\pgfqpoint{-0.011050in}{0.041667in}}{\pgfqpoint{-0.021649in}{0.037276in}}{\pgfqpoint{-0.029463in}{0.029463in}}%
\pgfpathcurveto{\pgfqpoint{-0.037276in}{0.021649in}}{\pgfqpoint{-0.041667in}{0.011050in}}{\pgfqpoint{-0.041667in}{0.000000in}}%
\pgfpathcurveto{\pgfqpoint{-0.041667in}{-0.011050in}}{\pgfqpoint{-0.037276in}{-0.021649in}}{\pgfqpoint{-0.029463in}{-0.029463in}}%
\pgfpathcurveto{\pgfqpoint{-0.021649in}{-0.037276in}}{\pgfqpoint{-0.011050in}{-0.041667in}}{\pgfqpoint{0.000000in}{-0.041667in}}%
\pgfpathlineto{\pgfqpoint{0.000000in}{-0.041667in}}%
\pgfpathclose%
\pgfusepath{stroke,fill}%
}%
\begin{pgfscope}%
\pgfsys@transformshift{3.447217in}{2.413529in}%
\pgfsys@useobject{currentmarker}{}%
\end{pgfscope}%
\end{pgfscope}%
\begin{pgfscope}%
\definecolor{textcolor}{rgb}{0.150000,0.150000,0.150000}%
\pgfsetstrokecolor{textcolor}%
\pgfsetfillcolor{textcolor}%
\pgftext[x=3.647217in,y=2.374640in,left,base]{\color{textcolor}{\sffamily\fontsize{8.000000}{9.600000}\selectfont\catcode`\^=\active\def^{\ifmmode\sp\else\^{}\fi}\catcode`\%=\active\def%{\%}solid energy}}%
\end{pgfscope}%
\begin{pgfscope}%
\pgfsetbuttcap%
\pgfsetroundjoin%
\pgfsetlinewidth{1.505625pt}%
\definecolor{currentstroke}{rgb}{0.000000,0.000000,0.000000}%
\pgfsetstrokecolor{currentstroke}%
\pgfsetdash{{5.550000pt}{2.400000pt}}{0.000000pt}%
\pgfpathmoveto{\pgfqpoint{3.336106in}{2.253698in}}%
\pgfpathlineto{\pgfqpoint{3.447217in}{2.253698in}}%
\pgfpathlineto{\pgfqpoint{3.558328in}{2.253698in}}%
\pgfusepath{stroke}%
\end{pgfscope}%
\begin{pgfscope}%
\definecolor{textcolor}{rgb}{0.150000,0.150000,0.150000}%
\pgfsetstrokecolor{textcolor}%
\pgfsetfillcolor{textcolor}%
\pgftext[x=3.647217in,y=2.214809in,left,base]{\color{textcolor}{\sffamily\fontsize{8.000000}{9.600000}\selectfont\catcode`\^=\active\def^{\ifmmode\sp\else\^{}\fi}\catcode`\%=\active\def%{\%}acceptance limit}}%
\end{pgfscope}%
\begin{pgfscope}%
\pgfsetbuttcap%
\pgfsetmiterjoin%
\definecolor{currentfill}{rgb}{1.000000,1.000000,1.000000}%
\pgfsetfillcolor{currentfill}%
\pgfsetlinewidth{0.000000pt}%
\definecolor{currentstroke}{rgb}{0.000000,0.000000,0.000000}%
\pgfsetstrokecolor{currentstroke}%
\pgfsetstrokeopacity{0.000000}%
\pgfsetdash{}{0pt}%
\pgfpathmoveto{\pgfqpoint{5.483877in}{0.464855in}}%
\pgfpathlineto{\pgfqpoint{9.535013in}{0.464855in}}%
\pgfpathlineto{\pgfqpoint{9.535013in}{3.194855in}}%
\pgfpathlineto{\pgfqpoint{5.483877in}{3.194855in}}%
\pgfpathlineto{\pgfqpoint{5.483877in}{0.464855in}}%
\pgfpathclose%
\pgfusepath{fill}%
\end{pgfscope}%
\begin{pgfscope}%
\pgfpathrectangle{\pgfqpoint{5.483877in}{0.464855in}}{\pgfqpoint{4.051136in}{2.730000in}}%
\pgfusepath{clip}%
\pgfsetroundcap%
\pgfsetroundjoin%
\pgfsetlinewidth{0.803000pt}%
\definecolor{currentstroke}{rgb}{0.800000,0.800000,0.800000}%
\pgfsetstrokecolor{currentstroke}%
\pgfsetdash{}{0pt}%
\pgfpathmoveto{\pgfqpoint{6.240907in}{0.464855in}}%
\pgfpathlineto{\pgfqpoint{6.240907in}{3.194855in}}%
\pgfusepath{stroke}%
\end{pgfscope}%
\begin{pgfscope}%
\definecolor{textcolor}{rgb}{0.150000,0.150000,0.150000}%
\pgfsetstrokecolor{textcolor}%
\pgfsetfillcolor{textcolor}%
\pgftext[x=6.240907in,y=0.416243in,,top]{\color{textcolor}{\sffamily\fontsize{10.000000}{12.000000}\selectfont\catcode`\^=\active\def^{\ifmmode\sp\else\^{}\fi}\catcode`\%=\active\def%{\%}5}}%
\end{pgfscope}%
\begin{pgfscope}%
\pgfpathrectangle{\pgfqpoint{5.483877in}{0.464855in}}{\pgfqpoint{4.051136in}{2.730000in}}%
\pgfusepath{clip}%
\pgfsetroundcap%
\pgfsetroundjoin%
\pgfsetlinewidth{0.803000pt}%
\definecolor{currentstroke}{rgb}{0.800000,0.800000,0.800000}%
\pgfsetstrokecolor{currentstroke}%
\pgfsetdash{}{0pt}%
\pgfpathmoveto{\pgfqpoint{7.059319in}{0.464855in}}%
\pgfpathlineto{\pgfqpoint{7.059319in}{3.194855in}}%
\pgfusepath{stroke}%
\end{pgfscope}%
\begin{pgfscope}%
\definecolor{textcolor}{rgb}{0.150000,0.150000,0.150000}%
\pgfsetstrokecolor{textcolor}%
\pgfsetfillcolor{textcolor}%
\pgftext[x=7.059319in,y=0.416243in,,top]{\color{textcolor}{\sffamily\fontsize{10.000000}{12.000000}\selectfont\catcode`\^=\active\def^{\ifmmode\sp\else\^{}\fi}\catcode`\%=\active\def%{\%}10}}%
\end{pgfscope}%
\begin{pgfscope}%
\pgfpathrectangle{\pgfqpoint{5.483877in}{0.464855in}}{\pgfqpoint{4.051136in}{2.730000in}}%
\pgfusepath{clip}%
\pgfsetroundcap%
\pgfsetroundjoin%
\pgfsetlinewidth{0.803000pt}%
\definecolor{currentstroke}{rgb}{0.800000,0.800000,0.800000}%
\pgfsetstrokecolor{currentstroke}%
\pgfsetdash{}{0pt}%
\pgfpathmoveto{\pgfqpoint{7.877730in}{0.464855in}}%
\pgfpathlineto{\pgfqpoint{7.877730in}{3.194855in}}%
\pgfusepath{stroke}%
\end{pgfscope}%
\begin{pgfscope}%
\definecolor{textcolor}{rgb}{0.150000,0.150000,0.150000}%
\pgfsetstrokecolor{textcolor}%
\pgfsetfillcolor{textcolor}%
\pgftext[x=7.877730in,y=0.416243in,,top]{\color{textcolor}{\sffamily\fontsize{10.000000}{12.000000}\selectfont\catcode`\^=\active\def^{\ifmmode\sp\else\^{}\fi}\catcode`\%=\active\def%{\%}15}}%
\end{pgfscope}%
\begin{pgfscope}%
\pgfpathrectangle{\pgfqpoint{5.483877in}{0.464855in}}{\pgfqpoint{4.051136in}{2.730000in}}%
\pgfusepath{clip}%
\pgfsetroundcap%
\pgfsetroundjoin%
\pgfsetlinewidth{0.803000pt}%
\definecolor{currentstroke}{rgb}{0.800000,0.800000,0.800000}%
\pgfsetstrokecolor{currentstroke}%
\pgfsetdash{}{0pt}%
\pgfpathmoveto{\pgfqpoint{8.696141in}{0.464855in}}%
\pgfpathlineto{\pgfqpoint{8.696141in}{3.194855in}}%
\pgfusepath{stroke}%
\end{pgfscope}%
\begin{pgfscope}%
\definecolor{textcolor}{rgb}{0.150000,0.150000,0.150000}%
\pgfsetstrokecolor{textcolor}%
\pgfsetfillcolor{textcolor}%
\pgftext[x=8.696141in,y=0.416243in,,top]{\color{textcolor}{\sffamily\fontsize{10.000000}{12.000000}\selectfont\catcode`\^=\active\def^{\ifmmode\sp\else\^{}\fi}\catcode`\%=\active\def%{\%}20}}%
\end{pgfscope}%
\begin{pgfscope}%
\pgfpathrectangle{\pgfqpoint{5.483877in}{0.464855in}}{\pgfqpoint{4.051136in}{2.730000in}}%
\pgfusepath{clip}%
\pgfsetroundcap%
\pgfsetroundjoin%
\pgfsetlinewidth{0.803000pt}%
\definecolor{currentstroke}{rgb}{0.800000,0.800000,0.800000}%
\pgfsetstrokecolor{currentstroke}%
\pgfsetdash{}{0pt}%
\pgfpathmoveto{\pgfqpoint{9.514553in}{0.464855in}}%
\pgfpathlineto{\pgfqpoint{9.514553in}{3.194855in}}%
\pgfusepath{stroke}%
\end{pgfscope}%
\begin{pgfscope}%
\definecolor{textcolor}{rgb}{0.150000,0.150000,0.150000}%
\pgfsetstrokecolor{textcolor}%
\pgfsetfillcolor{textcolor}%
\pgftext[x=9.514553in,y=0.416243in,,top]{\color{textcolor}{\sffamily\fontsize{10.000000}{12.000000}\selectfont\catcode`\^=\active\def^{\ifmmode\sp\else\^{}\fi}\catcode`\%=\active\def%{\%}25}}%
\end{pgfscope}%
\begin{pgfscope}%
\definecolor{textcolor}{rgb}{0.150000,0.150000,0.150000}%
\pgfsetstrokecolor{textcolor}%
\pgfsetfillcolor{textcolor}%
\pgftext[x=7.509445in,y=0.230344in,,top]{\color{textcolor}{\sffamily\fontsize{10.000000}{12.000000}\selectfont\catcode`\^=\active\def^{\ifmmode\sp\else\^{}\fi}\catcode`\%=\active\def%{\%}time [h]}}%
\end{pgfscope}%
\begin{pgfscope}%
\pgfpathrectangle{\pgfqpoint{5.483877in}{0.464855in}}{\pgfqpoint{4.051136in}{2.730000in}}%
\pgfusepath{clip}%
\pgfsetroundcap%
\pgfsetroundjoin%
\pgfsetlinewidth{0.803000pt}%
\definecolor{currentstroke}{rgb}{0.800000,0.800000,0.800000}%
\pgfsetstrokecolor{currentstroke}%
\pgfsetdash{}{0pt}%
\pgfpathmoveto{\pgfqpoint{5.483877in}{0.586166in}}%
\pgfpathlineto{\pgfqpoint{9.535013in}{0.586166in}}%
\pgfusepath{stroke}%
\end{pgfscope}%
\begin{pgfscope}%
\definecolor{textcolor}{rgb}{0.150000,0.150000,0.150000}%
\pgfsetstrokecolor{textcolor}%
\pgfsetfillcolor{textcolor}%
\pgftext[x=5.358022in, y=0.535609in, left, base]{\color{textcolor}{\sffamily\fontsize{10.000000}{12.000000}\selectfont\catcode`\^=\active\def^{\ifmmode\sp\else\^{}\fi}\catcode`\%=\active\def%{\%}0}}%
\end{pgfscope}%
\begin{pgfscope}%
\pgfpathrectangle{\pgfqpoint{5.483877in}{0.464855in}}{\pgfqpoint{4.051136in}{2.730000in}}%
\pgfusepath{clip}%
\pgfsetroundcap%
\pgfsetroundjoin%
\pgfsetlinewidth{0.803000pt}%
\definecolor{currentstroke}{rgb}{0.800000,0.800000,0.800000}%
\pgfsetstrokecolor{currentstroke}%
\pgfsetdash{}{0pt}%
\pgfpathmoveto{\pgfqpoint{5.483877in}{0.963759in}}%
\pgfpathlineto{\pgfqpoint{9.535013in}{0.963759in}}%
\pgfusepath{stroke}%
\end{pgfscope}%
\begin{pgfscope}%
\definecolor{textcolor}{rgb}{0.150000,0.150000,0.150000}%
\pgfsetstrokecolor{textcolor}%
\pgfsetfillcolor{textcolor}%
\pgftext[x=5.358022in, y=0.913201in, left, base]{\color{textcolor}{\sffamily\fontsize{10.000000}{12.000000}\selectfont\catcode`\^=\active\def^{\ifmmode\sp\else\^{}\fi}\catcode`\%=\active\def%{\%}5}}%
\end{pgfscope}%
\begin{pgfscope}%
\pgfpathrectangle{\pgfqpoint{5.483877in}{0.464855in}}{\pgfqpoint{4.051136in}{2.730000in}}%
\pgfusepath{clip}%
\pgfsetroundcap%
\pgfsetroundjoin%
\pgfsetlinewidth{0.803000pt}%
\definecolor{currentstroke}{rgb}{0.800000,0.800000,0.800000}%
\pgfsetstrokecolor{currentstroke}%
\pgfsetdash{}{0pt}%
\pgfpathmoveto{\pgfqpoint{5.483877in}{1.341351in}}%
\pgfpathlineto{\pgfqpoint{9.535013in}{1.341351in}}%
\pgfusepath{stroke}%
\end{pgfscope}%
\begin{pgfscope}%
\definecolor{textcolor}{rgb}{0.150000,0.150000,0.150000}%
\pgfsetstrokecolor{textcolor}%
\pgfsetfillcolor{textcolor}%
\pgftext[x=5.280779in, y=1.290794in, left, base]{\color{textcolor}{\sffamily\fontsize{10.000000}{12.000000}\selectfont\catcode`\^=\active\def^{\ifmmode\sp\else\^{}\fi}\catcode`\%=\active\def%{\%}10}}%
\end{pgfscope}%
\begin{pgfscope}%
\pgfpathrectangle{\pgfqpoint{5.483877in}{0.464855in}}{\pgfqpoint{4.051136in}{2.730000in}}%
\pgfusepath{clip}%
\pgfsetroundcap%
\pgfsetroundjoin%
\pgfsetlinewidth{0.803000pt}%
\definecolor{currentstroke}{rgb}{0.800000,0.800000,0.800000}%
\pgfsetstrokecolor{currentstroke}%
\pgfsetdash{}{0pt}%
\pgfpathmoveto{\pgfqpoint{5.483877in}{1.718944in}}%
\pgfpathlineto{\pgfqpoint{9.535013in}{1.718944in}}%
\pgfusepath{stroke}%
\end{pgfscope}%
\begin{pgfscope}%
\definecolor{textcolor}{rgb}{0.150000,0.150000,0.150000}%
\pgfsetstrokecolor{textcolor}%
\pgfsetfillcolor{textcolor}%
\pgftext[x=5.280779in, y=1.668386in, left, base]{\color{textcolor}{\sffamily\fontsize{10.000000}{12.000000}\selectfont\catcode`\^=\active\def^{\ifmmode\sp\else\^{}\fi}\catcode`\%=\active\def%{\%}15}}%
\end{pgfscope}%
\begin{pgfscope}%
\pgfpathrectangle{\pgfqpoint{5.483877in}{0.464855in}}{\pgfqpoint{4.051136in}{2.730000in}}%
\pgfusepath{clip}%
\pgfsetroundcap%
\pgfsetroundjoin%
\pgfsetlinewidth{0.803000pt}%
\definecolor{currentstroke}{rgb}{0.800000,0.800000,0.800000}%
\pgfsetstrokecolor{currentstroke}%
\pgfsetdash{}{0pt}%
\pgfpathmoveto{\pgfqpoint{5.483877in}{2.096536in}}%
\pgfpathlineto{\pgfqpoint{9.535013in}{2.096536in}}%
\pgfusepath{stroke}%
\end{pgfscope}%
\begin{pgfscope}%
\definecolor{textcolor}{rgb}{0.150000,0.150000,0.150000}%
\pgfsetstrokecolor{textcolor}%
\pgfsetfillcolor{textcolor}%
\pgftext[x=5.280779in, y=2.045979in, left, base]{\color{textcolor}{\sffamily\fontsize{10.000000}{12.000000}\selectfont\catcode`\^=\active\def^{\ifmmode\sp\else\^{}\fi}\catcode`\%=\active\def%{\%}20}}%
\end{pgfscope}%
\begin{pgfscope}%
\pgfpathrectangle{\pgfqpoint{5.483877in}{0.464855in}}{\pgfqpoint{4.051136in}{2.730000in}}%
\pgfusepath{clip}%
\pgfsetroundcap%
\pgfsetroundjoin%
\pgfsetlinewidth{0.803000pt}%
\definecolor{currentstroke}{rgb}{0.800000,0.800000,0.800000}%
\pgfsetstrokecolor{currentstroke}%
\pgfsetdash{}{0pt}%
\pgfpathmoveto{\pgfqpoint{5.483877in}{2.474129in}}%
\pgfpathlineto{\pgfqpoint{9.535013in}{2.474129in}}%
\pgfusepath{stroke}%
\end{pgfscope}%
\begin{pgfscope}%
\definecolor{textcolor}{rgb}{0.150000,0.150000,0.150000}%
\pgfsetstrokecolor{textcolor}%
\pgfsetfillcolor{textcolor}%
\pgftext[x=5.280779in, y=2.423571in, left, base]{\color{textcolor}{\sffamily\fontsize{10.000000}{12.000000}\selectfont\catcode`\^=\active\def^{\ifmmode\sp\else\^{}\fi}\catcode`\%=\active\def%{\%}25}}%
\end{pgfscope}%
\begin{pgfscope}%
\pgfpathrectangle{\pgfqpoint{5.483877in}{0.464855in}}{\pgfqpoint{4.051136in}{2.730000in}}%
\pgfusepath{clip}%
\pgfsetroundcap%
\pgfsetroundjoin%
\pgfsetlinewidth{0.803000pt}%
\definecolor{currentstroke}{rgb}{0.800000,0.800000,0.800000}%
\pgfsetstrokecolor{currentstroke}%
\pgfsetdash{}{0pt}%
\pgfpathmoveto{\pgfqpoint{5.483877in}{2.851721in}}%
\pgfpathlineto{\pgfqpoint{9.535013in}{2.851721in}}%
\pgfusepath{stroke}%
\end{pgfscope}%
\begin{pgfscope}%
\definecolor{textcolor}{rgb}{0.150000,0.150000,0.150000}%
\pgfsetstrokecolor{textcolor}%
\pgfsetfillcolor{textcolor}%
\pgftext[x=5.280779in, y=2.801164in, left, base]{\color{textcolor}{\sffamily\fontsize{10.000000}{12.000000}\selectfont\catcode`\^=\active\def^{\ifmmode\sp\else\^{}\fi}\catcode`\%=\active\def%{\%}30}}%
\end{pgfscope}%
\begin{pgfscope}%
\definecolor{textcolor}{rgb}{0.150000,0.150000,0.150000}%
\pgfsetstrokecolor{textcolor}%
\pgfsetfillcolor{textcolor}%
\pgftext[x=5.225223in,y=1.829855in,,bottom,rotate=90.000000]{\color{textcolor}{\sffamily\fontsize{10.000000}{12.000000}\selectfont\catcode`\^=\active\def^{\ifmmode\sp\else\^{}\fi}\catcode`\%=\active\def%{\%}surface rate [m³/s]}}%
\end{pgfscope}%
\begin{pgfscope}%
\pgfpathrectangle{\pgfqpoint{5.483877in}{0.464855in}}{\pgfqpoint{4.051136in}{2.730000in}}%
\pgfusepath{clip}%
\pgfsetroundcap%
\pgfsetroundjoin%
\pgfsetlinewidth{1.505625pt}%
\definecolor{currentstroke}{rgb}{0.121569,0.466667,0.705882}%
\pgfsetstrokecolor{currentstroke}%
\pgfsetdash{}{0pt}%
\pgfpathmoveto{\pgfqpoint{5.668019in}{3.070764in}}%
\pgfpathlineto{\pgfqpoint{5.913543in}{3.029840in}}%
\pgfpathlineto{\pgfqpoint{6.159066in}{3.058763in}}%
\pgfpathlineto{\pgfqpoint{6.404590in}{3.034636in}}%
\pgfpathlineto{\pgfqpoint{6.650113in}{3.058983in}}%
\pgfpathlineto{\pgfqpoint{6.895636in}{3.038348in}}%
\pgfpathlineto{\pgfqpoint{7.141160in}{3.058103in}}%
\pgfpathlineto{\pgfqpoint{7.386683in}{3.057803in}}%
\pgfpathlineto{\pgfqpoint{7.632207in}{3.060753in}}%
\pgfpathlineto{\pgfqpoint{7.877730in}{3.058457in}}%
\pgfpathlineto{\pgfqpoint{8.123253in}{3.043184in}}%
\pgfpathlineto{\pgfqpoint{8.368777in}{3.043988in}}%
\pgfpathlineto{\pgfqpoint{8.614300in}{3.061076in}}%
\pgfpathlineto{\pgfqpoint{8.859824in}{3.060265in}}%
\pgfpathlineto{\pgfqpoint{9.105347in}{3.061051in}}%
\pgfpathlineto{\pgfqpoint{9.350871in}{3.061203in}}%
\pgfusepath{stroke}%
\end{pgfscope}%
\begin{pgfscope}%
\pgfpathrectangle{\pgfqpoint{5.483877in}{0.464855in}}{\pgfqpoint{4.051136in}{2.730000in}}%
\pgfusepath{clip}%
\pgfsetbuttcap%
\pgfsetroundjoin%
\definecolor{currentfill}{rgb}{0.121569,0.466667,0.705882}%
\pgfsetfillcolor{currentfill}%
\pgfsetlinewidth{1.003750pt}%
\definecolor{currentstroke}{rgb}{0.121569,0.466667,0.705882}%
\pgfsetstrokecolor{currentstroke}%
\pgfsetdash{}{0pt}%
\pgfsys@defobject{currentmarker}{\pgfqpoint{-0.041667in}{-0.041667in}}{\pgfqpoint{0.041667in}{0.041667in}}{%
\pgfpathmoveto{\pgfqpoint{0.000000in}{-0.041667in}}%
\pgfpathcurveto{\pgfqpoint{0.011050in}{-0.041667in}}{\pgfqpoint{0.021649in}{-0.037276in}}{\pgfqpoint{0.029463in}{-0.029463in}}%
\pgfpathcurveto{\pgfqpoint{0.037276in}{-0.021649in}}{\pgfqpoint{0.041667in}{-0.011050in}}{\pgfqpoint{0.041667in}{0.000000in}}%
\pgfpathcurveto{\pgfqpoint{0.041667in}{0.011050in}}{\pgfqpoint{0.037276in}{0.021649in}}{\pgfqpoint{0.029463in}{0.029463in}}%
\pgfpathcurveto{\pgfqpoint{0.021649in}{0.037276in}}{\pgfqpoint{0.011050in}{0.041667in}}{\pgfqpoint{0.000000in}{0.041667in}}%
\pgfpathcurveto{\pgfqpoint{-0.011050in}{0.041667in}}{\pgfqpoint{-0.021649in}{0.037276in}}{\pgfqpoint{-0.029463in}{0.029463in}}%
\pgfpathcurveto{\pgfqpoint{-0.037276in}{0.021649in}}{\pgfqpoint{-0.041667in}{0.011050in}}{\pgfqpoint{-0.041667in}{0.000000in}}%
\pgfpathcurveto{\pgfqpoint{-0.041667in}{-0.011050in}}{\pgfqpoint{-0.037276in}{-0.021649in}}{\pgfqpoint{-0.029463in}{-0.029463in}}%
\pgfpathcurveto{\pgfqpoint{-0.021649in}{-0.037276in}}{\pgfqpoint{-0.011050in}{-0.041667in}}{\pgfqpoint{0.000000in}{-0.041667in}}%
\pgfpathlineto{\pgfqpoint{0.000000in}{-0.041667in}}%
\pgfpathclose%
\pgfusepath{stroke,fill}%
}%
\begin{pgfscope}%
\pgfsys@transformshift{5.668019in}{3.070764in}%
\pgfsys@useobject{currentmarker}{}%
\end{pgfscope}%
\begin{pgfscope}%
\pgfsys@transformshift{5.913543in}{3.029840in}%
\pgfsys@useobject{currentmarker}{}%
\end{pgfscope}%
\begin{pgfscope}%
\pgfsys@transformshift{6.159066in}{3.058763in}%
\pgfsys@useobject{currentmarker}{}%
\end{pgfscope}%
\begin{pgfscope}%
\pgfsys@transformshift{6.404590in}{3.034636in}%
\pgfsys@useobject{currentmarker}{}%
\end{pgfscope}%
\begin{pgfscope}%
\pgfsys@transformshift{6.650113in}{3.058983in}%
\pgfsys@useobject{currentmarker}{}%
\end{pgfscope}%
\begin{pgfscope}%
\pgfsys@transformshift{6.895636in}{3.038348in}%
\pgfsys@useobject{currentmarker}{}%
\end{pgfscope}%
\begin{pgfscope}%
\pgfsys@transformshift{7.141160in}{3.058103in}%
\pgfsys@useobject{currentmarker}{}%
\end{pgfscope}%
\begin{pgfscope}%
\pgfsys@transformshift{7.386683in}{3.057803in}%
\pgfsys@useobject{currentmarker}{}%
\end{pgfscope}%
\begin{pgfscope}%
\pgfsys@transformshift{7.632207in}{3.060753in}%
\pgfsys@useobject{currentmarker}{}%
\end{pgfscope}%
\begin{pgfscope}%
\pgfsys@transformshift{7.877730in}{3.058457in}%
\pgfsys@useobject{currentmarker}{}%
\end{pgfscope}%
\begin{pgfscope}%
\pgfsys@transformshift{8.123253in}{3.043184in}%
\pgfsys@useobject{currentmarker}{}%
\end{pgfscope}%
\begin{pgfscope}%
\pgfsys@transformshift{8.368777in}{3.043988in}%
\pgfsys@useobject{currentmarker}{}%
\end{pgfscope}%
\begin{pgfscope}%
\pgfsys@transformshift{8.614300in}{3.061076in}%
\pgfsys@useobject{currentmarker}{}%
\end{pgfscope}%
\begin{pgfscope}%
\pgfsys@transformshift{8.859824in}{3.060265in}%
\pgfsys@useobject{currentmarker}{}%
\end{pgfscope}%
\begin{pgfscope}%
\pgfsys@transformshift{9.105347in}{3.061051in}%
\pgfsys@useobject{currentmarker}{}%
\end{pgfscope}%
\begin{pgfscope}%
\pgfsys@transformshift{9.350871in}{3.061203in}%
\pgfsys@useobject{currentmarker}{}%
\end{pgfscope}%
\end{pgfscope}%
\begin{pgfscope}%
\pgfpathrectangle{\pgfqpoint{5.483877in}{0.464855in}}{\pgfqpoint{4.051136in}{2.730000in}}%
\pgfusepath{clip}%
\pgfsetroundcap%
\pgfsetroundjoin%
\pgfsetlinewidth{1.505625pt}%
\definecolor{currentstroke}{rgb}{1.000000,0.498039,0.054902}%
\pgfsetstrokecolor{currentstroke}%
\pgfsetdash{}{0pt}%
\pgfpathmoveto{\pgfqpoint{5.668019in}{0.588946in}}%
\pgfpathlineto{\pgfqpoint{5.913543in}{0.588946in}}%
\pgfpathlineto{\pgfqpoint{6.159066in}{0.588946in}}%
\pgfpathlineto{\pgfqpoint{6.404590in}{0.588946in}}%
\pgfpathlineto{\pgfqpoint{6.650113in}{0.588946in}}%
\pgfpathlineto{\pgfqpoint{6.895636in}{0.588946in}}%
\pgfpathlineto{\pgfqpoint{7.141160in}{0.588946in}}%
\pgfpathlineto{\pgfqpoint{7.386683in}{0.588946in}}%
\pgfpathlineto{\pgfqpoint{7.632207in}{0.588946in}}%
\pgfpathlineto{\pgfqpoint{7.877730in}{0.588946in}}%
\pgfpathlineto{\pgfqpoint{8.123253in}{0.588946in}}%
\pgfpathlineto{\pgfqpoint{8.368777in}{0.588946in}}%
\pgfpathlineto{\pgfqpoint{8.614300in}{0.588946in}}%
\pgfpathlineto{\pgfqpoint{8.859824in}{0.588946in}}%
\pgfpathlineto{\pgfqpoint{9.105347in}{0.588946in}}%
\pgfpathlineto{\pgfqpoint{9.350871in}{0.588946in}}%
\pgfusepath{stroke}%
\end{pgfscope}%
\begin{pgfscope}%
\pgfpathrectangle{\pgfqpoint{5.483877in}{0.464855in}}{\pgfqpoint{4.051136in}{2.730000in}}%
\pgfusepath{clip}%
\pgfsetbuttcap%
\pgfsetroundjoin%
\definecolor{currentfill}{rgb}{1.000000,0.498039,0.054902}%
\pgfsetfillcolor{currentfill}%
\pgfsetlinewidth{1.003750pt}%
\definecolor{currentstroke}{rgb}{1.000000,0.498039,0.054902}%
\pgfsetstrokecolor{currentstroke}%
\pgfsetdash{}{0pt}%
\pgfsys@defobject{currentmarker}{\pgfqpoint{-0.041667in}{-0.041667in}}{\pgfqpoint{0.041667in}{0.041667in}}{%
\pgfpathmoveto{\pgfqpoint{0.000000in}{-0.041667in}}%
\pgfpathcurveto{\pgfqpoint{0.011050in}{-0.041667in}}{\pgfqpoint{0.021649in}{-0.037276in}}{\pgfqpoint{0.029463in}{-0.029463in}}%
\pgfpathcurveto{\pgfqpoint{0.037276in}{-0.021649in}}{\pgfqpoint{0.041667in}{-0.011050in}}{\pgfqpoint{0.041667in}{0.000000in}}%
\pgfpathcurveto{\pgfqpoint{0.041667in}{0.011050in}}{\pgfqpoint{0.037276in}{0.021649in}}{\pgfqpoint{0.029463in}{0.029463in}}%
\pgfpathcurveto{\pgfqpoint{0.021649in}{0.037276in}}{\pgfqpoint{0.011050in}{0.041667in}}{\pgfqpoint{0.000000in}{0.041667in}}%
\pgfpathcurveto{\pgfqpoint{-0.011050in}{0.041667in}}{\pgfqpoint{-0.021649in}{0.037276in}}{\pgfqpoint{-0.029463in}{0.029463in}}%
\pgfpathcurveto{\pgfqpoint{-0.037276in}{0.021649in}}{\pgfqpoint{-0.041667in}{0.011050in}}{\pgfqpoint{-0.041667in}{0.000000in}}%
\pgfpathcurveto{\pgfqpoint{-0.041667in}{-0.011050in}}{\pgfqpoint{-0.037276in}{-0.021649in}}{\pgfqpoint{-0.029463in}{-0.029463in}}%
\pgfpathcurveto{\pgfqpoint{-0.021649in}{-0.037276in}}{\pgfqpoint{-0.011050in}{-0.041667in}}{\pgfqpoint{0.000000in}{-0.041667in}}%
\pgfpathlineto{\pgfqpoint{0.000000in}{-0.041667in}}%
\pgfpathclose%
\pgfusepath{stroke,fill}%
}%
\begin{pgfscope}%
\pgfsys@transformshift{5.668019in}{0.588946in}%
\pgfsys@useobject{currentmarker}{}%
\end{pgfscope}%
\begin{pgfscope}%
\pgfsys@transformshift{5.913543in}{0.588946in}%
\pgfsys@useobject{currentmarker}{}%
\end{pgfscope}%
\begin{pgfscope}%
\pgfsys@transformshift{6.159066in}{0.588946in}%
\pgfsys@useobject{currentmarker}{}%
\end{pgfscope}%
\begin{pgfscope}%
\pgfsys@transformshift{6.404590in}{0.588946in}%
\pgfsys@useobject{currentmarker}{}%
\end{pgfscope}%
\begin{pgfscope}%
\pgfsys@transformshift{6.650113in}{0.588946in}%
\pgfsys@useobject{currentmarker}{}%
\end{pgfscope}%
\begin{pgfscope}%
\pgfsys@transformshift{6.895636in}{0.588946in}%
\pgfsys@useobject{currentmarker}{}%
\end{pgfscope}%
\begin{pgfscope}%
\pgfsys@transformshift{7.141160in}{0.588946in}%
\pgfsys@useobject{currentmarker}{}%
\end{pgfscope}%
\begin{pgfscope}%
\pgfsys@transformshift{7.386683in}{0.588946in}%
\pgfsys@useobject{currentmarker}{}%
\end{pgfscope}%
\begin{pgfscope}%
\pgfsys@transformshift{7.632207in}{0.588946in}%
\pgfsys@useobject{currentmarker}{}%
\end{pgfscope}%
\begin{pgfscope}%
\pgfsys@transformshift{7.877730in}{0.588946in}%
\pgfsys@useobject{currentmarker}{}%
\end{pgfscope}%
\begin{pgfscope}%
\pgfsys@transformshift{8.123253in}{0.588946in}%
\pgfsys@useobject{currentmarker}{}%
\end{pgfscope}%
\begin{pgfscope}%
\pgfsys@transformshift{8.368777in}{0.588946in}%
\pgfsys@useobject{currentmarker}{}%
\end{pgfscope}%
\begin{pgfscope}%
\pgfsys@transformshift{8.614300in}{0.588946in}%
\pgfsys@useobject{currentmarker}{}%
\end{pgfscope}%
\begin{pgfscope}%
\pgfsys@transformshift{8.859824in}{0.588946in}%
\pgfsys@useobject{currentmarker}{}%
\end{pgfscope}%
\begin{pgfscope}%
\pgfsys@transformshift{9.105347in}{0.588946in}%
\pgfsys@useobject{currentmarker}{}%
\end{pgfscope}%
\begin{pgfscope}%
\pgfsys@transformshift{9.350871in}{0.588946in}%
\pgfsys@useobject{currentmarker}{}%
\end{pgfscope}%
\end{pgfscope}%
\begin{pgfscope}%
\pgfsetrectcap%
\pgfsetmiterjoin%
\pgfsetlinewidth{1.003750pt}%
\definecolor{currentstroke}{rgb}{0.800000,0.800000,0.800000}%
\pgfsetstrokecolor{currentstroke}%
\pgfsetdash{}{0pt}%
\pgfpathmoveto{\pgfqpoint{5.483877in}{0.464855in}}%
\pgfpathlineto{\pgfqpoint{5.483877in}{3.194855in}}%
\pgfusepath{stroke}%
\end{pgfscope}%
\begin{pgfscope}%
\pgfsetrectcap%
\pgfsetmiterjoin%
\pgfsetlinewidth{1.003750pt}%
\definecolor{currentstroke}{rgb}{0.800000,0.800000,0.800000}%
\pgfsetstrokecolor{currentstroke}%
\pgfsetdash{}{0pt}%
\pgfpathmoveto{\pgfqpoint{9.535013in}{0.464855in}}%
\pgfpathlineto{\pgfqpoint{9.535013in}{3.194855in}}%
\pgfusepath{stroke}%
\end{pgfscope}%
\begin{pgfscope}%
\pgfsetrectcap%
\pgfsetmiterjoin%
\pgfsetlinewidth{1.003750pt}%
\definecolor{currentstroke}{rgb}{0.800000,0.800000,0.800000}%
\pgfsetstrokecolor{currentstroke}%
\pgfsetdash{}{0pt}%
\pgfpathmoveto{\pgfqpoint{5.483877in}{0.464855in}}%
\pgfpathlineto{\pgfqpoint{9.535013in}{0.464855in}}%
\pgfusepath{stroke}%
\end{pgfscope}%
\begin{pgfscope}%
\pgfsetrectcap%
\pgfsetmiterjoin%
\pgfsetlinewidth{1.003750pt}%
\definecolor{currentstroke}{rgb}{0.800000,0.800000,0.800000}%
\pgfsetstrokecolor{currentstroke}%
\pgfsetdash{}{0pt}%
\pgfpathmoveto{\pgfqpoint{5.483877in}{3.194855in}}%
\pgfpathlineto{\pgfqpoint{9.535013in}{3.194855in}}%
\pgfusepath{stroke}%
\end{pgfscope}%
\begin{pgfscope}%
\definecolor{textcolor}{rgb}{0.150000,0.150000,0.150000}%
\pgfsetstrokecolor{textcolor}%
\pgfsetfillcolor{textcolor}%
\pgftext[x=7.509445in,y=3.278188in,,base]{\color{textcolor}{\sffamily\fontsize{12.000000}{14.400000}\selectfont\catcode`\^=\active\def^{\ifmmode\sp\else\^{}\fi}\catcode`\%=\active\def%{\%}Active-well controls}}%
\end{pgfscope}%
\begin{pgfscope}%
\pgfsetroundcap%
\pgfsetroundjoin%
\pgfsetlinewidth{1.505625pt}%
\definecolor{currentstroke}{rgb}{0.121569,0.466667,0.705882}%
\pgfsetstrokecolor{currentstroke}%
\pgfsetdash{}{0pt}%
\pgfpathmoveto{\pgfqpoint{8.194980in}{1.942417in}}%
\pgfpathlineto{\pgfqpoint{8.333868in}{1.942417in}}%
\pgfpathlineto{\pgfqpoint{8.472757in}{1.942417in}}%
\pgfusepath{stroke}%
\end{pgfscope}%
\begin{pgfscope}%
\pgfsetbuttcap%
\pgfsetroundjoin%
\definecolor{currentfill}{rgb}{0.121569,0.466667,0.705882}%
\pgfsetfillcolor{currentfill}%
\pgfsetlinewidth{1.003750pt}%
\definecolor{currentstroke}{rgb}{0.121569,0.466667,0.705882}%
\pgfsetstrokecolor{currentstroke}%
\pgfsetdash{}{0pt}%
\pgfsys@defobject{currentmarker}{\pgfqpoint{-0.041667in}{-0.041667in}}{\pgfqpoint{0.041667in}{0.041667in}}{%
\pgfpathmoveto{\pgfqpoint{0.000000in}{-0.041667in}}%
\pgfpathcurveto{\pgfqpoint{0.011050in}{-0.041667in}}{\pgfqpoint{0.021649in}{-0.037276in}}{\pgfqpoint{0.029463in}{-0.029463in}}%
\pgfpathcurveto{\pgfqpoint{0.037276in}{-0.021649in}}{\pgfqpoint{0.041667in}{-0.011050in}}{\pgfqpoint{0.041667in}{0.000000in}}%
\pgfpathcurveto{\pgfqpoint{0.041667in}{0.011050in}}{\pgfqpoint{0.037276in}{0.021649in}}{\pgfqpoint{0.029463in}{0.029463in}}%
\pgfpathcurveto{\pgfqpoint{0.021649in}{0.037276in}}{\pgfqpoint{0.011050in}{0.041667in}}{\pgfqpoint{0.000000in}{0.041667in}}%
\pgfpathcurveto{\pgfqpoint{-0.011050in}{0.041667in}}{\pgfqpoint{-0.021649in}{0.037276in}}{\pgfqpoint{-0.029463in}{0.029463in}}%
\pgfpathcurveto{\pgfqpoint{-0.037276in}{0.021649in}}{\pgfqpoint{-0.041667in}{0.011050in}}{\pgfqpoint{-0.041667in}{0.000000in}}%
\pgfpathcurveto{\pgfqpoint{-0.041667in}{-0.011050in}}{\pgfqpoint{-0.037276in}{-0.021649in}}{\pgfqpoint{-0.029463in}{-0.029463in}}%
\pgfpathcurveto{\pgfqpoint{-0.021649in}{-0.037276in}}{\pgfqpoint{-0.011050in}{-0.041667in}}{\pgfqpoint{0.000000in}{-0.041667in}}%
\pgfpathlineto{\pgfqpoint{0.000000in}{-0.041667in}}%
\pgfpathclose%
\pgfusepath{stroke,fill}%
}%
\begin{pgfscope}%
\pgfsys@transformshift{8.333868in}{1.942417in}%
\pgfsys@useobject{currentmarker}{}%
\end{pgfscope}%
\end{pgfscope}%
\begin{pgfscope}%
\definecolor{textcolor}{rgb}{0.150000,0.150000,0.150000}%
\pgfsetstrokecolor{textcolor}%
\pgfsetfillcolor{textcolor}%
\pgftext[x=8.583868in,y=1.893806in,left,base]{\color{textcolor}{\sffamily\fontsize{10.000000}{12.000000}\selectfont\catcode`\^=\active\def^{\ifmmode\sp\else\^{}\fi}\catcode`\%=\active\def%{\%}gas injection}}%
\end{pgfscope}%
\begin{pgfscope}%
\pgfsetroundcap%
\pgfsetroundjoin%
\pgfsetlinewidth{1.505625pt}%
\definecolor{currentstroke}{rgb}{1.000000,0.498039,0.054902}%
\pgfsetstrokecolor{currentstroke}%
\pgfsetdash{}{0pt}%
\pgfpathmoveto{\pgfqpoint{8.194980in}{1.742629in}}%
\pgfpathlineto{\pgfqpoint{8.333868in}{1.742629in}}%
\pgfpathlineto{\pgfqpoint{8.472757in}{1.742629in}}%
\pgfusepath{stroke}%
\end{pgfscope}%
\begin{pgfscope}%
\pgfsetbuttcap%
\pgfsetroundjoin%
\definecolor{currentfill}{rgb}{1.000000,0.498039,0.054902}%
\pgfsetfillcolor{currentfill}%
\pgfsetlinewidth{1.003750pt}%
\definecolor{currentstroke}{rgb}{1.000000,0.498039,0.054902}%
\pgfsetstrokecolor{currentstroke}%
\pgfsetdash{}{0pt}%
\pgfsys@defobject{currentmarker}{\pgfqpoint{-0.041667in}{-0.041667in}}{\pgfqpoint{0.041667in}{0.041667in}}{%
\pgfpathmoveto{\pgfqpoint{0.000000in}{-0.041667in}}%
\pgfpathcurveto{\pgfqpoint{0.011050in}{-0.041667in}}{\pgfqpoint{0.021649in}{-0.037276in}}{\pgfqpoint{0.029463in}{-0.029463in}}%
\pgfpathcurveto{\pgfqpoint{0.037276in}{-0.021649in}}{\pgfqpoint{0.041667in}{-0.011050in}}{\pgfqpoint{0.041667in}{0.000000in}}%
\pgfpathcurveto{\pgfqpoint{0.041667in}{0.011050in}}{\pgfqpoint{0.037276in}{0.021649in}}{\pgfqpoint{0.029463in}{0.029463in}}%
\pgfpathcurveto{\pgfqpoint{0.021649in}{0.037276in}}{\pgfqpoint{0.011050in}{0.041667in}}{\pgfqpoint{0.000000in}{0.041667in}}%
\pgfpathcurveto{\pgfqpoint{-0.011050in}{0.041667in}}{\pgfqpoint{-0.021649in}{0.037276in}}{\pgfqpoint{-0.029463in}{0.029463in}}%
\pgfpathcurveto{\pgfqpoint{-0.037276in}{0.021649in}}{\pgfqpoint{-0.041667in}{0.011050in}}{\pgfqpoint{-0.041667in}{0.000000in}}%
\pgfpathcurveto{\pgfqpoint{-0.041667in}{-0.011050in}}{\pgfqpoint{-0.037276in}{-0.021649in}}{\pgfqpoint{-0.029463in}{-0.029463in}}%
\pgfpathcurveto{\pgfqpoint{-0.021649in}{-0.037276in}}{\pgfqpoint{-0.011050in}{-0.041667in}}{\pgfqpoint{0.000000in}{-0.041667in}}%
\pgfpathlineto{\pgfqpoint{0.000000in}{-0.041667in}}%
\pgfpathclose%
\pgfusepath{stroke,fill}%
}%
\begin{pgfscope}%
\pgfsys@transformshift{8.333868in}{1.742629in}%
\pgfsys@useobject{currentmarker}{}%
\end{pgfscope}%
\end{pgfscope}%
\begin{pgfscope}%
\definecolor{textcolor}{rgb}{0.150000,0.150000,0.150000}%
\pgfsetstrokecolor{textcolor}%
\pgfsetfillcolor{textcolor}%
\pgftext[x=8.583868in,y=1.694017in,left,base]{\color{textcolor}{\sffamily\fontsize{10.000000}{12.000000}\selectfont\catcode`\^=\active\def^{\ifmmode\sp\else\^{}\fi}\catcode`\%=\active\def%{\%}oil production}}%
\end{pgfscope}%
\begin{pgfscope}%
\definecolor{textcolor}{rgb}{0.150000,0.150000,0.150000}%
\pgfsetstrokecolor{textcolor}%
\pgfsetfillcolor{textcolor}%
\pgftext[x=4.935013in,y=7.250855in,,top]{\color{textcolor}{\sffamily\fontsize{12.000000}{14.400000}\selectfont\catcode`\^=\active\def^{\ifmmode\sp\else\^{}\fi}\catcode`\%=\active\def%{\%}SPE1 Case 1: one-day coupled CG/EG acceptance history}}%
\end{pgfscope}%
\end{pgfpicture}%
\makeatother%
\endgroup%
%
	}
	\caption{One-day SPE1 coupled-acceptance history.  Dashed lines mark the
		applicable balance and residual gates.}
	\label{fig:spe1_one_day_acceptance_history}
\end{figure}

\begin{figure}[ht]
	\centering
	\resizebox{\textwidth}{!}{%
		%% Creator: Matplotlib, PGF backend
%%
%% To include the figure in your LaTeX document, write
%%   \input{<filename>.pgf}
%%
%% Make sure the required packages are loaded in your preamble
%%   \usepackage{pgf}
%%
%% Also ensure that all the required font packages are loaded; for instance,
%% the lmodern package is sometimes necessary when using math font.
%%   \usepackage{lmodern}
%%
%% Figures using additional raster images can only be included by \input if
%% they are in the same directory as the main LaTeX file. For loading figures
%% from other directories you can use the `import` package
%%   \usepackage{import}
%%
%% and then include the figures with
%%   \import{<path to file>}{<filename>.pgf}
%%
%% Matplotlib used the following preamble
%%   \def\mathdefault#1{#1}
%%   \everymath=\expandafter{\the\everymath\displaystyle}
%%   \IfFileExists{scrextend.sty}{
%%     \usepackage[fontsize=10.000000pt]{scrextend}
%%   }{
%%     \renewcommand{\normalsize}{\fontsize{10.000000}{12.000000}\selectfont}
%%     \normalsize
%%   }
%%
%%   \ifdefined\pdftexversion\else  % non-pdftex case.
%%     \usepackage{fontspec}
%%     \setmainfont{DejaVuSerif.ttf}[Path=\detokenize{/home/john/miniconda3/envs/moose/lib/python3.14/site-packages/matplotlib/mpl-data/fonts/ttf/}]
%%     \setsansfont{Arial.ttf}[Path=\detokenize{/usr/share/fonts/truetype/msttcorefonts/}]
%%     \setmonofont{DejaVuSansMono.ttf}[Path=\detokenize{/home/john/miniconda3/envs/moose/lib/python3.14/site-packages/matplotlib/mpl-data/fonts/ttf/}]
%%   \fi
%%   \makeatletter\@ifpackageloaded{underscore}{}{\usepackage[strings]{underscore}}\makeatother
%%
\begingroup%
\makeatletter%
\begin{pgfpicture}%
\pgfpathrectangle{\pgfpointorigin}{\pgfqpoint{9.655984in}{4.305855in}}%
\pgfusepath{use as bounding box, clip}%
\begin{pgfscope}%
\pgfsetbuttcap%
\pgfsetmiterjoin%
\definecolor{currentfill}{rgb}{1.000000,1.000000,1.000000}%
\pgfsetfillcolor{currentfill}%
\pgfsetlinewidth{0.000000pt}%
\definecolor{currentstroke}{rgb}{1.000000,1.000000,1.000000}%
\pgfsetstrokecolor{currentstroke}%
\pgfsetdash{}{0pt}%
\pgfpathmoveto{\pgfqpoint{0.000000in}{0.000000in}}%
\pgfpathlineto{\pgfqpoint{9.655984in}{0.000000in}}%
\pgfpathlineto{\pgfqpoint{9.655984in}{4.305855in}}%
\pgfpathlineto{\pgfqpoint{0.000000in}{4.305855in}}%
\pgfpathlineto{\pgfqpoint{0.000000in}{0.000000in}}%
\pgfpathclose%
\pgfusepath{fill}%
\end{pgfscope}%
\begin{pgfscope}%
\pgfsetbuttcap%
\pgfsetmiterjoin%
\definecolor{currentfill}{rgb}{1.000000,1.000000,1.000000}%
\pgfsetfillcolor{currentfill}%
\pgfsetlinewidth{0.000000pt}%
\definecolor{currentstroke}{rgb}{0.000000,0.000000,0.000000}%
\pgfsetstrokecolor{currentstroke}%
\pgfsetstrokeopacity{0.000000}%
\pgfsetdash{}{0pt}%
\pgfpathmoveto{\pgfqpoint{0.643484in}{0.464855in}}%
\pgfpathlineto{\pgfqpoint{4.694621in}{0.464855in}}%
\pgfpathlineto{\pgfqpoint{4.694621in}{3.775855in}}%
\pgfpathlineto{\pgfqpoint{0.643484in}{3.775855in}}%
\pgfpathlineto{\pgfqpoint{0.643484in}{0.464855in}}%
\pgfpathclose%
\pgfusepath{fill}%
\end{pgfscope}%
\begin{pgfscope}%
\pgfpathrectangle{\pgfqpoint{0.643484in}{0.464855in}}{\pgfqpoint{4.051136in}{3.311000in}}%
\pgfusepath{clip}%
\pgfsetroundcap%
\pgfsetroundjoin%
\pgfsetlinewidth{0.803000pt}%
\definecolor{currentstroke}{rgb}{0.800000,0.800000,0.800000}%
\pgfsetstrokecolor{currentstroke}%
\pgfsetdash{}{0pt}%
\pgfpathmoveto{\pgfqpoint{0.826618in}{0.464855in}}%
\pgfpathlineto{\pgfqpoint{0.826618in}{3.775855in}}%
\pgfusepath{stroke}%
\end{pgfscope}%
\begin{pgfscope}%
\definecolor{textcolor}{rgb}{0.150000,0.150000,0.150000}%
\pgfsetstrokecolor{textcolor}%
\pgfsetfillcolor{textcolor}%
\pgftext[x=0.826618in,y=0.416243in,,top]{\color{textcolor}{\sffamily\fontsize{10.000000}{12.000000}\selectfont\catcode`\^=\active\def^{\ifmmode\sp\else\^{}\fi}\catcode`\%=\active\def%{\%}0}}%
\end{pgfscope}%
\begin{pgfscope}%
\pgfpathrectangle{\pgfqpoint{0.643484in}{0.464855in}}{\pgfqpoint{4.051136in}{3.311000in}}%
\pgfusepath{clip}%
\pgfsetroundcap%
\pgfsetroundjoin%
\pgfsetlinewidth{0.803000pt}%
\definecolor{currentstroke}{rgb}{0.800000,0.800000,0.800000}%
\pgfsetstrokecolor{currentstroke}%
\pgfsetdash{}{0pt}%
\pgfpathmoveto{\pgfqpoint{1.331256in}{0.464855in}}%
\pgfpathlineto{\pgfqpoint{1.331256in}{3.775855in}}%
\pgfusepath{stroke}%
\end{pgfscope}%
\begin{pgfscope}%
\definecolor{textcolor}{rgb}{0.150000,0.150000,0.150000}%
\pgfsetstrokecolor{textcolor}%
\pgfsetfillcolor{textcolor}%
\pgftext[x=1.331256in,y=0.416243in,,top]{\color{textcolor}{\sffamily\fontsize{10.000000}{12.000000}\selectfont\catcode`\^=\active\def^{\ifmmode\sp\else\^{}\fi}\catcode`\%=\active\def%{\%}500}}%
\end{pgfscope}%
\begin{pgfscope}%
\pgfpathrectangle{\pgfqpoint{0.643484in}{0.464855in}}{\pgfqpoint{4.051136in}{3.311000in}}%
\pgfusepath{clip}%
\pgfsetroundcap%
\pgfsetroundjoin%
\pgfsetlinewidth{0.803000pt}%
\definecolor{currentstroke}{rgb}{0.800000,0.800000,0.800000}%
\pgfsetstrokecolor{currentstroke}%
\pgfsetdash{}{0pt}%
\pgfpathmoveto{\pgfqpoint{1.835894in}{0.464855in}}%
\pgfpathlineto{\pgfqpoint{1.835894in}{3.775855in}}%
\pgfusepath{stroke}%
\end{pgfscope}%
\begin{pgfscope}%
\definecolor{textcolor}{rgb}{0.150000,0.150000,0.150000}%
\pgfsetstrokecolor{textcolor}%
\pgfsetfillcolor{textcolor}%
\pgftext[x=1.835894in,y=0.416243in,,top]{\color{textcolor}{\sffamily\fontsize{10.000000}{12.000000}\selectfont\catcode`\^=\active\def^{\ifmmode\sp\else\^{}\fi}\catcode`\%=\active\def%{\%}1000}}%
\end{pgfscope}%
\begin{pgfscope}%
\pgfpathrectangle{\pgfqpoint{0.643484in}{0.464855in}}{\pgfqpoint{4.051136in}{3.311000in}}%
\pgfusepath{clip}%
\pgfsetroundcap%
\pgfsetroundjoin%
\pgfsetlinewidth{0.803000pt}%
\definecolor{currentstroke}{rgb}{0.800000,0.800000,0.800000}%
\pgfsetstrokecolor{currentstroke}%
\pgfsetdash{}{0pt}%
\pgfpathmoveto{\pgfqpoint{2.340533in}{0.464855in}}%
\pgfpathlineto{\pgfqpoint{2.340533in}{3.775855in}}%
\pgfusepath{stroke}%
\end{pgfscope}%
\begin{pgfscope}%
\definecolor{textcolor}{rgb}{0.150000,0.150000,0.150000}%
\pgfsetstrokecolor{textcolor}%
\pgfsetfillcolor{textcolor}%
\pgftext[x=2.340533in,y=0.416243in,,top]{\color{textcolor}{\sffamily\fontsize{10.000000}{12.000000}\selectfont\catcode`\^=\active\def^{\ifmmode\sp\else\^{}\fi}\catcode`\%=\active\def%{\%}1500}}%
\end{pgfscope}%
\begin{pgfscope}%
\pgfpathrectangle{\pgfqpoint{0.643484in}{0.464855in}}{\pgfqpoint{4.051136in}{3.311000in}}%
\pgfusepath{clip}%
\pgfsetroundcap%
\pgfsetroundjoin%
\pgfsetlinewidth{0.803000pt}%
\definecolor{currentstroke}{rgb}{0.800000,0.800000,0.800000}%
\pgfsetstrokecolor{currentstroke}%
\pgfsetdash{}{0pt}%
\pgfpathmoveto{\pgfqpoint{2.845171in}{0.464855in}}%
\pgfpathlineto{\pgfqpoint{2.845171in}{3.775855in}}%
\pgfusepath{stroke}%
\end{pgfscope}%
\begin{pgfscope}%
\definecolor{textcolor}{rgb}{0.150000,0.150000,0.150000}%
\pgfsetstrokecolor{textcolor}%
\pgfsetfillcolor{textcolor}%
\pgftext[x=2.845171in,y=0.416243in,,top]{\color{textcolor}{\sffamily\fontsize{10.000000}{12.000000}\selectfont\catcode`\^=\active\def^{\ifmmode\sp\else\^{}\fi}\catcode`\%=\active\def%{\%}2000}}%
\end{pgfscope}%
\begin{pgfscope}%
\pgfpathrectangle{\pgfqpoint{0.643484in}{0.464855in}}{\pgfqpoint{4.051136in}{3.311000in}}%
\pgfusepath{clip}%
\pgfsetroundcap%
\pgfsetroundjoin%
\pgfsetlinewidth{0.803000pt}%
\definecolor{currentstroke}{rgb}{0.800000,0.800000,0.800000}%
\pgfsetstrokecolor{currentstroke}%
\pgfsetdash{}{0pt}%
\pgfpathmoveto{\pgfqpoint{3.349810in}{0.464855in}}%
\pgfpathlineto{\pgfqpoint{3.349810in}{3.775855in}}%
\pgfusepath{stroke}%
\end{pgfscope}%
\begin{pgfscope}%
\definecolor{textcolor}{rgb}{0.150000,0.150000,0.150000}%
\pgfsetstrokecolor{textcolor}%
\pgfsetfillcolor{textcolor}%
\pgftext[x=3.349810in,y=0.416243in,,top]{\color{textcolor}{\sffamily\fontsize{10.000000}{12.000000}\selectfont\catcode`\^=\active\def^{\ifmmode\sp\else\^{}\fi}\catcode`\%=\active\def%{\%}2500}}%
\end{pgfscope}%
\begin{pgfscope}%
\pgfpathrectangle{\pgfqpoint{0.643484in}{0.464855in}}{\pgfqpoint{4.051136in}{3.311000in}}%
\pgfusepath{clip}%
\pgfsetroundcap%
\pgfsetroundjoin%
\pgfsetlinewidth{0.803000pt}%
\definecolor{currentstroke}{rgb}{0.800000,0.800000,0.800000}%
\pgfsetstrokecolor{currentstroke}%
\pgfsetdash{}{0pt}%
\pgfpathmoveto{\pgfqpoint{3.854448in}{0.464855in}}%
\pgfpathlineto{\pgfqpoint{3.854448in}{3.775855in}}%
\pgfusepath{stroke}%
\end{pgfscope}%
\begin{pgfscope}%
\definecolor{textcolor}{rgb}{0.150000,0.150000,0.150000}%
\pgfsetstrokecolor{textcolor}%
\pgfsetfillcolor{textcolor}%
\pgftext[x=3.854448in,y=0.416243in,,top]{\color{textcolor}{\sffamily\fontsize{10.000000}{12.000000}\selectfont\catcode`\^=\active\def^{\ifmmode\sp\else\^{}\fi}\catcode`\%=\active\def%{\%}3000}}%
\end{pgfscope}%
\begin{pgfscope}%
\pgfpathrectangle{\pgfqpoint{0.643484in}{0.464855in}}{\pgfqpoint{4.051136in}{3.311000in}}%
\pgfusepath{clip}%
\pgfsetroundcap%
\pgfsetroundjoin%
\pgfsetlinewidth{0.803000pt}%
\definecolor{currentstroke}{rgb}{0.800000,0.800000,0.800000}%
\pgfsetstrokecolor{currentstroke}%
\pgfsetdash{}{0pt}%
\pgfpathmoveto{\pgfqpoint{4.359087in}{0.464855in}}%
\pgfpathlineto{\pgfqpoint{4.359087in}{3.775855in}}%
\pgfusepath{stroke}%
\end{pgfscope}%
\begin{pgfscope}%
\definecolor{textcolor}{rgb}{0.150000,0.150000,0.150000}%
\pgfsetstrokecolor{textcolor}%
\pgfsetfillcolor{textcolor}%
\pgftext[x=4.359087in,y=0.416243in,,top]{\color{textcolor}{\sffamily\fontsize{10.000000}{12.000000}\selectfont\catcode`\^=\active\def^{\ifmmode\sp\else\^{}\fi}\catcode`\%=\active\def%{\%}3500}}%
\end{pgfscope}%
\begin{pgfscope}%
\definecolor{textcolor}{rgb}{0.150000,0.150000,0.150000}%
\pgfsetstrokecolor{textcolor}%
\pgfsetfillcolor{textcolor}%
\pgftext[x=2.669052in,y=0.230344in,,top]{\color{textcolor}{\sffamily\fontsize{10.000000}{12.000000}\selectfont\catcode`\^=\active\def^{\ifmmode\sp\else\^{}\fi}\catcode`\%=\active\def%{\%}time [day]}}%
\end{pgfscope}%
\begin{pgfscope}%
\pgfpathrectangle{\pgfqpoint{0.643484in}{0.464855in}}{\pgfqpoint{4.051136in}{3.311000in}}%
\pgfusepath{clip}%
\pgfsetroundcap%
\pgfsetroundjoin%
\pgfsetlinewidth{0.803000pt}%
\definecolor{currentstroke}{rgb}{0.800000,0.800000,0.800000}%
\pgfsetstrokecolor{currentstroke}%
\pgfsetdash{}{0pt}%
\pgfpathmoveto{\pgfqpoint{0.643484in}{0.615355in}}%
\pgfpathlineto{\pgfqpoint{4.694621in}{0.615355in}}%
\pgfusepath{stroke}%
\end{pgfscope}%
\begin{pgfscope}%
\definecolor{textcolor}{rgb}{0.150000,0.150000,0.150000}%
\pgfsetstrokecolor{textcolor}%
\pgfsetfillcolor{textcolor}%
\pgftext[x=0.285900in, y=0.564797in, left, base]{\color{textcolor}{\sffamily\fontsize{10.000000}{12.000000}\selectfont\catcode`\^=\active\def^{\ifmmode\sp\else\^{}\fi}\catcode`\%=\active\def%{\%}1000}}%
\end{pgfscope}%
\begin{pgfscope}%
\pgfpathrectangle{\pgfqpoint{0.643484in}{0.464855in}}{\pgfqpoint{4.051136in}{3.311000in}}%
\pgfusepath{clip}%
\pgfsetroundcap%
\pgfsetroundjoin%
\pgfsetlinewidth{0.803000pt}%
\definecolor{currentstroke}{rgb}{0.800000,0.800000,0.800000}%
\pgfsetstrokecolor{currentstroke}%
\pgfsetdash{}{0pt}%
\pgfpathmoveto{\pgfqpoint{0.643484in}{1.112777in}}%
\pgfpathlineto{\pgfqpoint{4.694621in}{1.112777in}}%
\pgfusepath{stroke}%
\end{pgfscope}%
\begin{pgfscope}%
\definecolor{textcolor}{rgb}{0.150000,0.150000,0.150000}%
\pgfsetstrokecolor{textcolor}%
\pgfsetfillcolor{textcolor}%
\pgftext[x=0.285900in, y=1.062220in, left, base]{\color{textcolor}{\sffamily\fontsize{10.000000}{12.000000}\selectfont\catcode`\^=\active\def^{\ifmmode\sp\else\^{}\fi}\catcode`\%=\active\def%{\%}2000}}%
\end{pgfscope}%
\begin{pgfscope}%
\pgfpathrectangle{\pgfqpoint{0.643484in}{0.464855in}}{\pgfqpoint{4.051136in}{3.311000in}}%
\pgfusepath{clip}%
\pgfsetroundcap%
\pgfsetroundjoin%
\pgfsetlinewidth{0.803000pt}%
\definecolor{currentstroke}{rgb}{0.800000,0.800000,0.800000}%
\pgfsetstrokecolor{currentstroke}%
\pgfsetdash{}{0pt}%
\pgfpathmoveto{\pgfqpoint{0.643484in}{1.610200in}}%
\pgfpathlineto{\pgfqpoint{4.694621in}{1.610200in}}%
\pgfusepath{stroke}%
\end{pgfscope}%
\begin{pgfscope}%
\definecolor{textcolor}{rgb}{0.150000,0.150000,0.150000}%
\pgfsetstrokecolor{textcolor}%
\pgfsetfillcolor{textcolor}%
\pgftext[x=0.285900in, y=1.559643in, left, base]{\color{textcolor}{\sffamily\fontsize{10.000000}{12.000000}\selectfont\catcode`\^=\active\def^{\ifmmode\sp\else\^{}\fi}\catcode`\%=\active\def%{\%}3000}}%
\end{pgfscope}%
\begin{pgfscope}%
\pgfpathrectangle{\pgfqpoint{0.643484in}{0.464855in}}{\pgfqpoint{4.051136in}{3.311000in}}%
\pgfusepath{clip}%
\pgfsetroundcap%
\pgfsetroundjoin%
\pgfsetlinewidth{0.803000pt}%
\definecolor{currentstroke}{rgb}{0.800000,0.800000,0.800000}%
\pgfsetstrokecolor{currentstroke}%
\pgfsetdash{}{0pt}%
\pgfpathmoveto{\pgfqpoint{0.643484in}{2.107623in}}%
\pgfpathlineto{\pgfqpoint{4.694621in}{2.107623in}}%
\pgfusepath{stroke}%
\end{pgfscope}%
\begin{pgfscope}%
\definecolor{textcolor}{rgb}{0.150000,0.150000,0.150000}%
\pgfsetstrokecolor{textcolor}%
\pgfsetfillcolor{textcolor}%
\pgftext[x=0.285900in, y=2.057065in, left, base]{\color{textcolor}{\sffamily\fontsize{10.000000}{12.000000}\selectfont\catcode`\^=\active\def^{\ifmmode\sp\else\^{}\fi}\catcode`\%=\active\def%{\%}4000}}%
\end{pgfscope}%
\begin{pgfscope}%
\pgfpathrectangle{\pgfqpoint{0.643484in}{0.464855in}}{\pgfqpoint{4.051136in}{3.311000in}}%
\pgfusepath{clip}%
\pgfsetroundcap%
\pgfsetroundjoin%
\pgfsetlinewidth{0.803000pt}%
\definecolor{currentstroke}{rgb}{0.800000,0.800000,0.800000}%
\pgfsetstrokecolor{currentstroke}%
\pgfsetdash{}{0pt}%
\pgfpathmoveto{\pgfqpoint{0.643484in}{2.605045in}}%
\pgfpathlineto{\pgfqpoint{4.694621in}{2.605045in}}%
\pgfusepath{stroke}%
\end{pgfscope}%
\begin{pgfscope}%
\definecolor{textcolor}{rgb}{0.150000,0.150000,0.150000}%
\pgfsetstrokecolor{textcolor}%
\pgfsetfillcolor{textcolor}%
\pgftext[x=0.285900in, y=2.554488in, left, base]{\color{textcolor}{\sffamily\fontsize{10.000000}{12.000000}\selectfont\catcode`\^=\active\def^{\ifmmode\sp\else\^{}\fi}\catcode`\%=\active\def%{\%}5000}}%
\end{pgfscope}%
\begin{pgfscope}%
\pgfpathrectangle{\pgfqpoint{0.643484in}{0.464855in}}{\pgfqpoint{4.051136in}{3.311000in}}%
\pgfusepath{clip}%
\pgfsetroundcap%
\pgfsetroundjoin%
\pgfsetlinewidth{0.803000pt}%
\definecolor{currentstroke}{rgb}{0.800000,0.800000,0.800000}%
\pgfsetstrokecolor{currentstroke}%
\pgfsetdash{}{0pt}%
\pgfpathmoveto{\pgfqpoint{0.643484in}{3.102468in}}%
\pgfpathlineto{\pgfqpoint{4.694621in}{3.102468in}}%
\pgfusepath{stroke}%
\end{pgfscope}%
\begin{pgfscope}%
\definecolor{textcolor}{rgb}{0.150000,0.150000,0.150000}%
\pgfsetstrokecolor{textcolor}%
\pgfsetfillcolor{textcolor}%
\pgftext[x=0.285900in, y=3.051911in, left, base]{\color{textcolor}{\sffamily\fontsize{10.000000}{12.000000}\selectfont\catcode`\^=\active\def^{\ifmmode\sp\else\^{}\fi}\catcode`\%=\active\def%{\%}6000}}%
\end{pgfscope}%
\begin{pgfscope}%
\pgfpathrectangle{\pgfqpoint{0.643484in}{0.464855in}}{\pgfqpoint{4.051136in}{3.311000in}}%
\pgfusepath{clip}%
\pgfsetroundcap%
\pgfsetroundjoin%
\pgfsetlinewidth{0.803000pt}%
\definecolor{currentstroke}{rgb}{0.800000,0.800000,0.800000}%
\pgfsetstrokecolor{currentstroke}%
\pgfsetdash{}{0pt}%
\pgfpathmoveto{\pgfqpoint{0.643484in}{3.599891in}}%
\pgfpathlineto{\pgfqpoint{4.694621in}{3.599891in}}%
\pgfusepath{stroke}%
\end{pgfscope}%
\begin{pgfscope}%
\definecolor{textcolor}{rgb}{0.150000,0.150000,0.150000}%
\pgfsetstrokecolor{textcolor}%
\pgfsetfillcolor{textcolor}%
\pgftext[x=0.285900in, y=3.549333in, left, base]{\color{textcolor}{\sffamily\fontsize{10.000000}{12.000000}\selectfont\catcode`\^=\active\def^{\ifmmode\sp\else\^{}\fi}\catcode`\%=\active\def%{\%}7000}}%
\end{pgfscope}%
\begin{pgfscope}%
\definecolor{textcolor}{rgb}{0.150000,0.150000,0.150000}%
\pgfsetstrokecolor{textcolor}%
\pgfsetfillcolor{textcolor}%
\pgftext[x=0.230344in,y=2.120355in,,bottom,rotate=90.000000]{\color{textcolor}{\sffamily\fontsize{10.000000}{12.000000}\selectfont\catcode`\^=\active\def^{\ifmmode\sp\else\^{}\fi}\catcode`\%=\active\def%{\%}BHP [psia]}}%
\end{pgfscope}%
\begin{pgfscope}%
\pgfpathrectangle{\pgfqpoint{0.643484in}{0.464855in}}{\pgfqpoint{4.051136in}{3.311000in}}%
\pgfusepath{clip}%
\pgfsetbuttcap%
\pgfsetroundjoin%
\definecolor{currentfill}{rgb}{0.121569,0.466667,0.705882}%
\pgfsetfillcolor{currentfill}%
\pgfsetlinewidth{1.003750pt}%
\definecolor{currentstroke}{rgb}{0.121569,0.466667,0.705882}%
\pgfsetstrokecolor{currentstroke}%
\pgfsetdash{}{0pt}%
\pgfsys@defobject{currentmarker}{\pgfqpoint{-0.041667in}{-0.041667in}}{\pgfqpoint{0.041667in}{0.041667in}}{%
\pgfpathmoveto{\pgfqpoint{0.000000in}{0.041667in}}%
\pgfpathlineto{\pgfqpoint{-0.041667in}{-0.041667in}}%
\pgfpathlineto{\pgfqpoint{0.041667in}{-0.041667in}}%
\pgfpathlineto{\pgfqpoint{0.000000in}{0.041667in}}%
\pgfpathclose%
\pgfusepath{stroke,fill}%
}%
\begin{pgfscope}%
\pgfsys@transformshift{0.827627in}{2.721058in}%
\pgfsys@useobject{currentmarker}{}%
\end{pgfscope}%
\end{pgfscope}%
\begin{pgfscope}%
\pgfpathrectangle{\pgfqpoint{0.643484in}{0.464855in}}{\pgfqpoint{4.051136in}{3.311000in}}%
\pgfusepath{clip}%
\pgfsetbuttcap%
\pgfsetroundjoin%
\definecolor{currentfill}{rgb}{1.000000,0.498039,0.054902}%
\pgfsetfillcolor{currentfill}%
\pgfsetlinewidth{1.003750pt}%
\definecolor{currentstroke}{rgb}{1.000000,0.498039,0.054902}%
\pgfsetstrokecolor{currentstroke}%
\pgfsetdash{}{0pt}%
\pgfsys@defobject{currentmarker}{\pgfqpoint{-0.041667in}{-0.041667in}}{\pgfqpoint{0.041667in}{0.041667in}}{%
\pgfpathmoveto{\pgfqpoint{-0.000000in}{-0.041667in}}%
\pgfpathlineto{\pgfqpoint{0.041667in}{0.041667in}}%
\pgfpathlineto{\pgfqpoint{-0.041667in}{0.041667in}}%
\pgfpathlineto{\pgfqpoint{-0.000000in}{-0.041667in}}%
\pgfpathclose%
\pgfusepath{stroke,fill}%
}%
\begin{pgfscope}%
\pgfsys@transformshift{0.827627in}{1.594402in}%
\pgfsys@useobject{currentmarker}{}%
\end{pgfscope}%
\end{pgfscope}%
\begin{pgfscope}%
\pgfpathrectangle{\pgfqpoint{0.643484in}{0.464855in}}{\pgfqpoint{4.051136in}{3.311000in}}%
\pgfusepath{clip}%
\pgfsetroundcap%
\pgfsetroundjoin%
\pgfsetlinewidth{1.505625pt}%
\definecolor{currentstroke}{rgb}{0.121569,0.466667,0.705882}%
\pgfsetstrokecolor{currentstroke}%
\pgfsetdash{}{0pt}%
\pgfpathmoveto{\pgfqpoint{0.857905in}{3.544348in}}%
\pgfpathlineto{\pgfqpoint{0.886165in}{3.419384in}}%
\pgfpathlineto{\pgfqpoint{0.917452in}{3.375199in}}%
\pgfpathlineto{\pgfqpoint{0.947731in}{3.350124in}}%
\pgfpathlineto{\pgfqpoint{0.979018in}{3.328845in}}%
\pgfpathlineto{\pgfqpoint{1.009297in}{3.320511in}}%
\pgfpathlineto{\pgfqpoint{1.040584in}{3.320321in}}%
\pgfpathlineto{\pgfqpoint{1.071872in}{3.326407in}}%
\pgfpathlineto{\pgfqpoint{1.102150in}{3.339206in}}%
\pgfpathlineto{\pgfqpoint{1.133438in}{3.354452in}}%
\pgfpathlineto{\pgfqpoint{1.163716in}{3.370696in}}%
\pgfpathlineto{\pgfqpoint{1.195004in}{3.388673in}}%
\pgfpathlineto{\pgfqpoint{1.226291in}{3.407688in}}%
\pgfpathlineto{\pgfqpoint{1.254551in}{3.425419in}}%
\pgfpathlineto{\pgfqpoint{1.285838in}{3.445342in}}%
\pgfpathlineto{\pgfqpoint{1.316117in}{3.464941in}}%
\pgfpathlineto{\pgfqpoint{1.347404in}{3.485436in}}%
\pgfpathlineto{\pgfqpoint{1.377683in}{3.505311in}}%
\pgfpathlineto{\pgfqpoint{1.408970in}{3.525847in}}%
\pgfpathlineto{\pgfqpoint{1.440258in}{3.546334in}}%
\pgfpathlineto{\pgfqpoint{1.470536in}{3.566077in}}%
\pgfpathlineto{\pgfqpoint{1.501824in}{3.585700in}}%
\pgfpathlineto{\pgfqpoint{1.532102in}{3.602437in}}%
\pgfpathlineto{\pgfqpoint{1.563390in}{3.615472in}}%
\pgfpathlineto{\pgfqpoint{1.594677in}{3.623297in}}%
\pgfpathlineto{\pgfqpoint{1.622937in}{3.625355in}}%
\pgfpathlineto{\pgfqpoint{1.654225in}{3.620846in}}%
\pgfpathlineto{\pgfqpoint{1.684503in}{3.610329in}}%
\pgfpathlineto{\pgfqpoint{1.715790in}{3.593394in}}%
\pgfpathlineto{\pgfqpoint{1.746069in}{3.571799in}}%
\pgfpathlineto{\pgfqpoint{1.777356in}{3.544350in}}%
\pgfpathlineto{\pgfqpoint{1.808644in}{3.512217in}}%
\pgfpathlineto{\pgfqpoint{1.838922in}{3.478384in}}%
\pgfpathlineto{\pgfqpoint{1.870210in}{3.442248in}}%
\pgfpathlineto{\pgfqpoint{1.900488in}{3.407091in}}%
\pgfpathlineto{\pgfqpoint{1.931776in}{3.371153in}}%
\pgfpathlineto{\pgfqpoint{1.963063in}{3.335712in}}%
\pgfpathlineto{\pgfqpoint{1.991323in}{3.304231in}}%
\pgfpathlineto{\pgfqpoint{2.022611in}{3.270153in}}%
\pgfpathlineto{\pgfqpoint{2.052889in}{3.237923in}}%
\pgfpathlineto{\pgfqpoint{2.084176in}{3.205441in}}%
\pgfpathlineto{\pgfqpoint{2.114455in}{3.174852in}}%
\pgfpathlineto{\pgfqpoint{2.145742in}{3.144098in}}%
\pgfpathlineto{\pgfqpoint{2.177030in}{3.114262in}}%
\pgfpathlineto{\pgfqpoint{2.207308in}{3.086247in}}%
\pgfpathlineto{\pgfqpoint{2.238596in}{3.058208in}}%
\pgfpathlineto{\pgfqpoint{2.268874in}{3.031935in}}%
\pgfpathlineto{\pgfqpoint{2.300162in}{3.005702in}}%
\pgfpathlineto{\pgfqpoint{2.331449in}{2.980344in}}%
\pgfpathlineto{\pgfqpoint{2.359709in}{2.958187in}}%
\pgfpathlineto{\pgfqpoint{2.390997in}{2.934553in}}%
\pgfpathlineto{\pgfqpoint{2.421275in}{2.912533in}}%
\pgfpathlineto{\pgfqpoint{2.452563in}{2.890684in}}%
\pgfpathlineto{\pgfqpoint{2.482841in}{2.870354in}}%
\pgfpathlineto{\pgfqpoint{2.514128in}{2.850176in}}%
\pgfpathlineto{\pgfqpoint{2.545416in}{2.830616in}}%
\pgfpathlineto{\pgfqpoint{2.575694in}{2.812344in}}%
\pgfpathlineto{\pgfqpoint{2.606982in}{2.794120in}}%
\pgfpathlineto{\pgfqpoint{2.637260in}{2.777044in}}%
\pgfpathlineto{\pgfqpoint{2.668548in}{2.760053in}}%
\pgfpathlineto{\pgfqpoint{2.699835in}{2.743477in}}%
\pgfpathlineto{\pgfqpoint{2.728095in}{2.729032in}}%
\pgfpathlineto{\pgfqpoint{2.759383in}{2.713583in}}%
\pgfpathlineto{\pgfqpoint{2.789661in}{2.699090in}}%
\pgfpathlineto{\pgfqpoint{2.820949in}{2.684664in}}%
\pgfpathlineto{\pgfqpoint{2.851227in}{2.671182in}}%
\pgfpathlineto{\pgfqpoint{2.882514in}{2.657720in}}%
\pgfpathlineto{\pgfqpoint{2.913802in}{2.644673in}}%
\pgfpathlineto{\pgfqpoint{2.944080in}{2.632416in}}%
\pgfpathlineto{\pgfqpoint{2.975368in}{2.620055in}}%
\pgfpathlineto{\pgfqpoint{3.005646in}{2.608262in}}%
\pgfpathlineto{\pgfqpoint{3.036934in}{2.596150in}}%
\pgfpathlineto{\pgfqpoint{3.068221in}{2.584112in}}%
\pgfpathlineto{\pgfqpoint{3.096481in}{2.573412in}}%
\pgfpathlineto{\pgfqpoint{3.127769in}{2.561816in}}%
\pgfpathlineto{\pgfqpoint{3.158047in}{2.550885in}}%
\pgfpathlineto{\pgfqpoint{3.189335in}{2.539937in}}%
\pgfpathlineto{\pgfqpoint{3.219613in}{2.529699in}}%
\pgfpathlineto{\pgfqpoint{3.250901in}{2.519367in}}%
\pgfpathlineto{\pgfqpoint{3.282188in}{2.509292in}}%
\pgfpathlineto{\pgfqpoint{3.312466in}{2.499837in}}%
\pgfpathlineto{\pgfqpoint{3.343754in}{2.490325in}}%
\pgfpathlineto{\pgfqpoint{3.374032in}{2.481354in}}%
\pgfpathlineto{\pgfqpoint{3.405320in}{2.472235in}}%
\pgfpathlineto{\pgfqpoint{3.436607in}{2.463310in}}%
\pgfpathlineto{\pgfqpoint{3.464867in}{2.455434in}}%
\pgfpathlineto{\pgfqpoint{3.496155in}{2.446912in}}%
\pgfpathlineto{\pgfqpoint{3.526433in}{2.438864in}}%
\pgfpathlineto{\pgfqpoint{3.557721in}{2.430743in}}%
\pgfpathlineto{\pgfqpoint{3.587999in}{2.423083in}}%
\pgfpathlineto{\pgfqpoint{3.619287in}{2.415381in}}%
\pgfpathlineto{\pgfqpoint{3.650574in}{2.407894in}}%
\pgfpathlineto{\pgfqpoint{3.680852in}{2.400888in}}%
\pgfpathlineto{\pgfqpoint{3.712140in}{2.393905in}}%
\pgfpathlineto{\pgfqpoint{3.742418in}{2.387435in}}%
\pgfpathlineto{\pgfqpoint{3.773706in}{2.381037in}}%
\pgfpathlineto{\pgfqpoint{3.804994in}{2.374936in}}%
\pgfpathlineto{\pgfqpoint{3.833253in}{2.369652in}}%
\pgfpathlineto{\pgfqpoint{3.864541in}{2.363933in}}%
\pgfpathlineto{\pgfqpoint{3.894819in}{2.358542in}}%
\pgfpathlineto{\pgfqpoint{3.926107in}{2.352986in}}%
\pgfpathlineto{\pgfqpoint{3.956385in}{2.347726in}}%
\pgfpathlineto{\pgfqpoint{3.987673in}{2.342306in}}%
\pgfpathlineto{\pgfqpoint{4.018960in}{2.336915in}}%
\pgfpathlineto{\pgfqpoint{4.049239in}{2.331696in}}%
\pgfpathlineto{\pgfqpoint{4.080526in}{2.326269in}}%
\pgfpathlineto{\pgfqpoint{4.110804in}{2.321038in}}%
\pgfpathlineto{\pgfqpoint{4.142092in}{2.315561in}}%
\pgfpathlineto{\pgfqpoint{4.173380in}{2.310055in}}%
\pgfpathlineto{\pgfqpoint{4.201639in}{2.305047in}}%
\pgfpathlineto{\pgfqpoint{4.232927in}{2.299481in}}%
\pgfpathlineto{\pgfqpoint{4.263205in}{2.294045in}}%
\pgfpathlineto{\pgfqpoint{4.294493in}{2.288404in}}%
\pgfpathlineto{\pgfqpoint{4.324771in}{2.282885in}}%
\pgfpathlineto{\pgfqpoint{4.356059in}{2.277190in}}%
\pgfpathlineto{\pgfqpoint{4.387346in}{2.271553in}}%
\pgfpathlineto{\pgfqpoint{4.417625in}{2.266077in}}%
\pgfpathlineto{\pgfqpoint{4.448912in}{2.260426in}}%
\pgfpathlineto{\pgfqpoint{4.479190in}{2.254983in}}%
\pgfpathlineto{\pgfqpoint{4.510478in}{2.249383in}}%
\pgfusepath{stroke}%
\end{pgfscope}%
\begin{pgfscope}%
\pgfpathrectangle{\pgfqpoint{0.643484in}{0.464855in}}{\pgfqpoint{4.051136in}{3.311000in}}%
\pgfusepath{clip}%
\pgfsetroundcap%
\pgfsetroundjoin%
\pgfsetlinewidth{1.505625pt}%
\definecolor{currentstroke}{rgb}{1.000000,0.498039,0.054902}%
\pgfsetstrokecolor{currentstroke}%
\pgfsetdash{}{0pt}%
\pgfpathmoveto{\pgfqpoint{0.857905in}{1.252334in}}%
\pgfpathlineto{\pgfqpoint{0.886165in}{1.214923in}}%
\pgfpathlineto{\pgfqpoint{0.917452in}{1.232059in}}%
\pgfpathlineto{\pgfqpoint{0.947731in}{1.264717in}}%
\pgfpathlineto{\pgfqpoint{0.979018in}{1.303633in}}%
\pgfpathlineto{\pgfqpoint{1.009297in}{1.342619in}}%
\pgfpathlineto{\pgfqpoint{1.040584in}{1.382751in}}%
\pgfpathlineto{\pgfqpoint{1.071872in}{1.421932in}}%
\pgfpathlineto{\pgfqpoint{1.102150in}{1.458755in}}%
\pgfpathlineto{\pgfqpoint{1.133438in}{1.495585in}}%
\pgfpathlineto{\pgfqpoint{1.163716in}{1.530162in}}%
\pgfpathlineto{\pgfqpoint{1.195004in}{1.564939in}}%
\pgfpathlineto{\pgfqpoint{1.226291in}{1.598786in}}%
\pgfpathlineto{\pgfqpoint{1.254551in}{1.628674in}}%
\pgfpathlineto{\pgfqpoint{1.285838in}{1.661171in}}%
\pgfpathlineto{\pgfqpoint{1.316117in}{1.692344in}}%
\pgfpathlineto{\pgfqpoint{1.347404in}{1.724438in}}%
\pgfpathlineto{\pgfqpoint{1.377683in}{1.755504in}}%
\pgfpathlineto{\pgfqpoint{1.408970in}{1.787420in}}%
\pgfpathlineto{\pgfqpoint{1.440258in}{1.824767in}}%
\pgfpathlineto{\pgfqpoint{1.470536in}{1.861684in}}%
\pgfpathlineto{\pgfqpoint{1.501824in}{1.868951in}}%
\pgfpathlineto{\pgfqpoint{1.532102in}{1.792779in}}%
\pgfpathlineto{\pgfqpoint{1.563390in}{1.656054in}}%
\pgfpathlineto{\pgfqpoint{1.594677in}{1.499445in}}%
\pgfpathlineto{\pgfqpoint{1.622937in}{1.364904in}}%
\pgfpathlineto{\pgfqpoint{1.654225in}{1.231061in}}%
\pgfpathlineto{\pgfqpoint{1.684503in}{1.111351in}}%
\pgfpathlineto{\pgfqpoint{1.715790in}{0.992046in}}%
\pgfpathlineto{\pgfqpoint{1.746069in}{0.876596in}}%
\pgfpathlineto{\pgfqpoint{1.777356in}{0.757357in}}%
\pgfpathlineto{\pgfqpoint{1.808644in}{0.646377in}}%
\pgfpathlineto{\pgfqpoint{1.838922in}{0.615355in}}%
\pgfpathlineto{\pgfqpoint{1.870210in}{0.615355in}}%
\pgfpathlineto{\pgfqpoint{1.900488in}{0.615355in}}%
\pgfpathlineto{\pgfqpoint{1.931776in}{0.615355in}}%
\pgfpathlineto{\pgfqpoint{1.963063in}{0.615355in}}%
\pgfpathlineto{\pgfqpoint{1.991323in}{0.615355in}}%
\pgfpathlineto{\pgfqpoint{2.022611in}{0.615355in}}%
\pgfpathlineto{\pgfqpoint{2.052889in}{0.615355in}}%
\pgfpathlineto{\pgfqpoint{2.084176in}{0.615355in}}%
\pgfpathlineto{\pgfqpoint{2.114455in}{0.615355in}}%
\pgfpathlineto{\pgfqpoint{2.145742in}{0.615355in}}%
\pgfpathlineto{\pgfqpoint{2.177030in}{0.615355in}}%
\pgfpathlineto{\pgfqpoint{2.207308in}{0.615355in}}%
\pgfpathlineto{\pgfqpoint{2.238596in}{0.615355in}}%
\pgfpathlineto{\pgfqpoint{2.268874in}{0.615355in}}%
\pgfpathlineto{\pgfqpoint{2.300162in}{0.615355in}}%
\pgfpathlineto{\pgfqpoint{2.331449in}{0.615355in}}%
\pgfpathlineto{\pgfqpoint{2.359709in}{0.615355in}}%
\pgfpathlineto{\pgfqpoint{2.390997in}{0.615355in}}%
\pgfpathlineto{\pgfqpoint{2.421275in}{0.615355in}}%
\pgfpathlineto{\pgfqpoint{2.452563in}{0.615355in}}%
\pgfpathlineto{\pgfqpoint{2.482841in}{0.615355in}}%
\pgfpathlineto{\pgfqpoint{2.514128in}{0.615355in}}%
\pgfpathlineto{\pgfqpoint{2.545416in}{0.615355in}}%
\pgfpathlineto{\pgfqpoint{2.575694in}{0.615355in}}%
\pgfpathlineto{\pgfqpoint{2.606982in}{0.615355in}}%
\pgfpathlineto{\pgfqpoint{2.637260in}{0.615355in}}%
\pgfpathlineto{\pgfqpoint{2.668548in}{0.615355in}}%
\pgfpathlineto{\pgfqpoint{2.699835in}{0.615355in}}%
\pgfpathlineto{\pgfqpoint{2.728095in}{0.615355in}}%
\pgfpathlineto{\pgfqpoint{2.759383in}{0.615355in}}%
\pgfpathlineto{\pgfqpoint{2.789661in}{0.615355in}}%
\pgfpathlineto{\pgfqpoint{2.820949in}{0.615355in}}%
\pgfpathlineto{\pgfqpoint{2.851227in}{0.615355in}}%
\pgfpathlineto{\pgfqpoint{2.882514in}{0.615355in}}%
\pgfpathlineto{\pgfqpoint{2.913802in}{0.615355in}}%
\pgfpathlineto{\pgfqpoint{2.944080in}{0.615355in}}%
\pgfpathlineto{\pgfqpoint{2.975368in}{0.615355in}}%
\pgfpathlineto{\pgfqpoint{3.005646in}{0.615355in}}%
\pgfpathlineto{\pgfqpoint{3.036934in}{0.615355in}}%
\pgfpathlineto{\pgfqpoint{3.068221in}{0.615355in}}%
\pgfpathlineto{\pgfqpoint{3.096481in}{0.615355in}}%
\pgfpathlineto{\pgfqpoint{3.127769in}{0.615355in}}%
\pgfpathlineto{\pgfqpoint{3.158047in}{0.615355in}}%
\pgfpathlineto{\pgfqpoint{3.189335in}{0.615355in}}%
\pgfpathlineto{\pgfqpoint{3.219613in}{0.615355in}}%
\pgfpathlineto{\pgfqpoint{3.250901in}{0.615355in}}%
\pgfpathlineto{\pgfqpoint{3.282188in}{0.615355in}}%
\pgfpathlineto{\pgfqpoint{3.312466in}{0.615355in}}%
\pgfpathlineto{\pgfqpoint{3.343754in}{0.615355in}}%
\pgfpathlineto{\pgfqpoint{3.374032in}{0.615355in}}%
\pgfpathlineto{\pgfqpoint{3.405320in}{0.615355in}}%
\pgfpathlineto{\pgfqpoint{3.436607in}{0.615355in}}%
\pgfpathlineto{\pgfqpoint{3.464867in}{0.615355in}}%
\pgfpathlineto{\pgfqpoint{3.496155in}{0.615355in}}%
\pgfpathlineto{\pgfqpoint{3.526433in}{0.615355in}}%
\pgfpathlineto{\pgfqpoint{3.557721in}{0.615355in}}%
\pgfpathlineto{\pgfqpoint{3.587999in}{0.615355in}}%
\pgfpathlineto{\pgfqpoint{3.619287in}{0.615355in}}%
\pgfpathlineto{\pgfqpoint{3.650574in}{0.615355in}}%
\pgfpathlineto{\pgfqpoint{3.680852in}{0.615355in}}%
\pgfpathlineto{\pgfqpoint{3.712140in}{0.615355in}}%
\pgfpathlineto{\pgfqpoint{3.742418in}{0.615355in}}%
\pgfpathlineto{\pgfqpoint{3.773706in}{0.615355in}}%
\pgfpathlineto{\pgfqpoint{3.804994in}{0.615355in}}%
\pgfpathlineto{\pgfqpoint{3.833253in}{0.615355in}}%
\pgfpathlineto{\pgfqpoint{3.864541in}{0.615355in}}%
\pgfpathlineto{\pgfqpoint{3.894819in}{0.615355in}}%
\pgfpathlineto{\pgfqpoint{3.926107in}{0.615355in}}%
\pgfpathlineto{\pgfqpoint{3.956385in}{0.615355in}}%
\pgfpathlineto{\pgfqpoint{3.987673in}{0.615355in}}%
\pgfpathlineto{\pgfqpoint{4.018960in}{0.615355in}}%
\pgfpathlineto{\pgfqpoint{4.049239in}{0.615355in}}%
\pgfpathlineto{\pgfqpoint{4.080526in}{0.615355in}}%
\pgfpathlineto{\pgfqpoint{4.110804in}{0.615355in}}%
\pgfpathlineto{\pgfqpoint{4.142092in}{0.615355in}}%
\pgfpathlineto{\pgfqpoint{4.173380in}{0.615355in}}%
\pgfpathlineto{\pgfqpoint{4.201639in}{0.615355in}}%
\pgfpathlineto{\pgfqpoint{4.232927in}{0.615355in}}%
\pgfpathlineto{\pgfqpoint{4.263205in}{0.615355in}}%
\pgfpathlineto{\pgfqpoint{4.294493in}{0.615355in}}%
\pgfpathlineto{\pgfqpoint{4.324771in}{0.615355in}}%
\pgfpathlineto{\pgfqpoint{4.356059in}{0.615355in}}%
\pgfpathlineto{\pgfqpoint{4.387346in}{0.615355in}}%
\pgfpathlineto{\pgfqpoint{4.417625in}{0.615355in}}%
\pgfpathlineto{\pgfqpoint{4.448912in}{0.615355in}}%
\pgfpathlineto{\pgfqpoint{4.479190in}{0.615355in}}%
\pgfpathlineto{\pgfqpoint{4.510478in}{0.615355in}}%
\pgfusepath{stroke}%
\end{pgfscope}%
\begin{pgfscope}%
\pgfsetrectcap%
\pgfsetmiterjoin%
\pgfsetlinewidth{1.003750pt}%
\definecolor{currentstroke}{rgb}{0.800000,0.800000,0.800000}%
\pgfsetstrokecolor{currentstroke}%
\pgfsetdash{}{0pt}%
\pgfpathmoveto{\pgfqpoint{0.643484in}{0.464855in}}%
\pgfpathlineto{\pgfqpoint{0.643484in}{3.775855in}}%
\pgfusepath{stroke}%
\end{pgfscope}%
\begin{pgfscope}%
\pgfsetrectcap%
\pgfsetmiterjoin%
\pgfsetlinewidth{1.003750pt}%
\definecolor{currentstroke}{rgb}{0.800000,0.800000,0.800000}%
\pgfsetstrokecolor{currentstroke}%
\pgfsetdash{}{0pt}%
\pgfpathmoveto{\pgfqpoint{4.694621in}{0.464855in}}%
\pgfpathlineto{\pgfqpoint{4.694621in}{3.775855in}}%
\pgfusepath{stroke}%
\end{pgfscope}%
\begin{pgfscope}%
\pgfsetrectcap%
\pgfsetmiterjoin%
\pgfsetlinewidth{1.003750pt}%
\definecolor{currentstroke}{rgb}{0.800000,0.800000,0.800000}%
\pgfsetstrokecolor{currentstroke}%
\pgfsetdash{}{0pt}%
\pgfpathmoveto{\pgfqpoint{0.643484in}{0.464855in}}%
\pgfpathlineto{\pgfqpoint{4.694621in}{0.464855in}}%
\pgfusepath{stroke}%
\end{pgfscope}%
\begin{pgfscope}%
\pgfsetrectcap%
\pgfsetmiterjoin%
\pgfsetlinewidth{1.003750pt}%
\definecolor{currentstroke}{rgb}{0.800000,0.800000,0.800000}%
\pgfsetstrokecolor{currentstroke}%
\pgfsetdash{}{0pt}%
\pgfpathmoveto{\pgfqpoint{0.643484in}{3.775855in}}%
\pgfpathlineto{\pgfqpoint{4.694621in}{3.775855in}}%
\pgfusepath{stroke}%
\end{pgfscope}%
\begin{pgfscope}%
\definecolor{textcolor}{rgb}{0.150000,0.150000,0.150000}%
\pgfsetstrokecolor{textcolor}%
\pgfsetfillcolor{textcolor}%
\pgftext[x=2.669052in,y=3.859188in,,base]{\color{textcolor}{\sffamily\fontsize{12.000000}{14.400000}\selectfont\catcode`\^=\active\def^{\ifmmode\sp\else\^{}\fi}\catcode`\%=\active\def%{\%}Pressure schedule context}}%
\end{pgfscope}%
\begin{pgfscope}%
\pgfsetroundcap%
\pgfsetroundjoin%
\pgfsetlinewidth{1.505625pt}%
\definecolor{currentstroke}{rgb}{0.121569,0.466667,0.705882}%
\pgfsetstrokecolor{currentstroke}%
\pgfsetdash{}{0pt}%
\pgfpathmoveto{\pgfqpoint{3.153355in}{3.633797in}}%
\pgfpathlineto{\pgfqpoint{3.264466in}{3.633797in}}%
\pgfpathlineto{\pgfqpoint{3.375578in}{3.633797in}}%
\pgfusepath{stroke}%
\end{pgfscope}%
\begin{pgfscope}%
\definecolor{textcolor}{rgb}{0.150000,0.150000,0.150000}%
\pgfsetstrokecolor{textcolor}%
\pgfsetfillcolor{textcolor}%
\pgftext[x=3.464466in,y=3.594908in,left,base]{\color{textcolor}{\sffamily\fontsize{8.000000}{9.600000}\selectfont\catcode`\^=\active\def^{\ifmmode\sp\else\^{}\fi}\catcode`\%=\active\def%{\%}OPM injector BHP}}%
\end{pgfscope}%
\begin{pgfscope}%
\pgfsetroundcap%
\pgfsetroundjoin%
\pgfsetlinewidth{1.505625pt}%
\definecolor{currentstroke}{rgb}{1.000000,0.498039,0.054902}%
\pgfsetstrokecolor{currentstroke}%
\pgfsetdash{}{0pt}%
\pgfpathmoveto{\pgfqpoint{3.153355in}{3.473912in}}%
\pgfpathlineto{\pgfqpoint{3.264466in}{3.473912in}}%
\pgfpathlineto{\pgfqpoint{3.375578in}{3.473912in}}%
\pgfusepath{stroke}%
\end{pgfscope}%
\begin{pgfscope}%
\definecolor{textcolor}{rgb}{0.150000,0.150000,0.150000}%
\pgfsetstrokecolor{textcolor}%
\pgfsetfillcolor{textcolor}%
\pgftext[x=3.464466in,y=3.435023in,left,base]{\color{textcolor}{\sffamily\fontsize{8.000000}{9.600000}\selectfont\catcode`\^=\active\def^{\ifmmode\sp\else\^{}\fi}\catcode`\%=\active\def%{\%}OPM producer BHP}}%
\end{pgfscope}%
\begin{pgfscope}%
\pgfsetbuttcap%
\pgfsetroundjoin%
\definecolor{currentfill}{rgb}{0.121569,0.466667,0.705882}%
\pgfsetfillcolor{currentfill}%
\pgfsetlinewidth{1.003750pt}%
\definecolor{currentstroke}{rgb}{0.121569,0.466667,0.705882}%
\pgfsetstrokecolor{currentstroke}%
\pgfsetdash{}{0pt}%
\pgfsys@defobject{currentmarker}{\pgfqpoint{-0.041667in}{-0.041667in}}{\pgfqpoint{0.041667in}{0.041667in}}{%
\pgfpathmoveto{\pgfqpoint{0.000000in}{0.041667in}}%
\pgfpathlineto{\pgfqpoint{-0.041667in}{-0.041667in}}%
\pgfpathlineto{\pgfqpoint{0.041667in}{-0.041667in}}%
\pgfpathlineto{\pgfqpoint{0.000000in}{0.041667in}}%
\pgfpathclose%
\pgfusepath{stroke,fill}%
}%
\begin{pgfscope}%
\pgfsys@transformshift{3.264466in}{3.304359in}%
\pgfsys@useobject{currentmarker}{}%
\end{pgfscope}%
\end{pgfscope}%
\begin{pgfscope}%
\definecolor{textcolor}{rgb}{0.150000,0.150000,0.150000}%
\pgfsetstrokecolor{textcolor}%
\pgfsetfillcolor{textcolor}%
\pgftext[x=3.464466in,y=3.275193in,left,base]{\color{textcolor}{\sffamily\fontsize{8.000000}{9.600000}\selectfont\catcode`\^=\active\def^{\ifmmode\sp\else\^{}\fi}\catcode`\%=\active\def%{\%}CG/EG day 1 injector}}%
\end{pgfscope}%
\begin{pgfscope}%
\pgfsetbuttcap%
\pgfsetroundjoin%
\definecolor{currentfill}{rgb}{1.000000,0.498039,0.054902}%
\pgfsetfillcolor{currentfill}%
\pgfsetlinewidth{1.003750pt}%
\definecolor{currentstroke}{rgb}{1.000000,0.498039,0.054902}%
\pgfsetstrokecolor{currentstroke}%
\pgfsetdash{}{0pt}%
\pgfsys@defobject{currentmarker}{\pgfqpoint{-0.041667in}{-0.041667in}}{\pgfqpoint{0.041667in}{0.041667in}}{%
\pgfpathmoveto{\pgfqpoint{-0.000000in}{-0.041667in}}%
\pgfpathlineto{\pgfqpoint{0.041667in}{0.041667in}}%
\pgfpathlineto{\pgfqpoint{-0.041667in}{0.041667in}}%
\pgfpathlineto{\pgfqpoint{-0.000000in}{-0.041667in}}%
\pgfpathclose%
\pgfusepath{stroke,fill}%
}%
\begin{pgfscope}%
\pgfsys@transformshift{3.264466in}{3.144528in}%
\pgfsys@useobject{currentmarker}{}%
\end{pgfscope}%
\end{pgfscope}%
\begin{pgfscope}%
\definecolor{textcolor}{rgb}{0.150000,0.150000,0.150000}%
\pgfsetstrokecolor{textcolor}%
\pgfsetfillcolor{textcolor}%
\pgftext[x=3.464466in,y=3.115362in,left,base]{\color{textcolor}{\sffamily\fontsize{8.000000}{9.600000}\selectfont\catcode`\^=\active\def^{\ifmmode\sp\else\^{}\fi}\catcode`\%=\active\def%{\%}CG/EG day 1 producer}}%
\end{pgfscope}%
\begin{pgfscope}%
\pgfsetbuttcap%
\pgfsetmiterjoin%
\definecolor{currentfill}{rgb}{1.000000,1.000000,1.000000}%
\pgfsetfillcolor{currentfill}%
\pgfsetlinewidth{0.000000pt}%
\definecolor{currentstroke}{rgb}{0.000000,0.000000,0.000000}%
\pgfsetstrokecolor{currentstroke}%
\pgfsetstrokeopacity{0.000000}%
\pgfsetdash{}{0pt}%
\pgfpathmoveto{\pgfqpoint{5.504848in}{0.464855in}}%
\pgfpathlineto{\pgfqpoint{9.555984in}{0.464855in}}%
\pgfpathlineto{\pgfqpoint{9.555984in}{3.775855in}}%
\pgfpathlineto{\pgfqpoint{5.504848in}{3.775855in}}%
\pgfpathlineto{\pgfqpoint{5.504848in}{0.464855in}}%
\pgfpathclose%
\pgfusepath{fill}%
\end{pgfscope}%
\begin{pgfscope}%
\pgfpathrectangle{\pgfqpoint{5.504848in}{0.464855in}}{\pgfqpoint{4.051136in}{3.311000in}}%
\pgfusepath{clip}%
\pgfsetroundcap%
\pgfsetroundjoin%
\pgfsetlinewidth{0.803000pt}%
\definecolor{currentstroke}{rgb}{0.800000,0.800000,0.800000}%
\pgfsetstrokecolor{currentstroke}%
\pgfsetdash{}{0pt}%
\pgfpathmoveto{\pgfqpoint{6.507402in}{0.464855in}}%
\pgfpathlineto{\pgfqpoint{6.507402in}{3.775855in}}%
\pgfusepath{stroke}%
\end{pgfscope}%
\begin{pgfscope}%
\definecolor{textcolor}{rgb}{0.150000,0.150000,0.150000}%
\pgfsetstrokecolor{textcolor}%
\pgfsetfillcolor{textcolor}%
\pgftext[x=6.507402in,y=0.416243in,,top]{\color{textcolor}{\sffamily\fontsize{10.000000}{12.000000}\selectfont\catcode`\^=\active\def^{\ifmmode\sp\else\^{}\fi}\catcode`\%=\active\def%{\%}gas injection}}%
\end{pgfscope}%
\begin{pgfscope}%
\pgfpathrectangle{\pgfqpoint{5.504848in}{0.464855in}}{\pgfqpoint{4.051136in}{3.311000in}}%
\pgfusepath{clip}%
\pgfsetroundcap%
\pgfsetroundjoin%
\pgfsetlinewidth{0.803000pt}%
\definecolor{currentstroke}{rgb}{0.800000,0.800000,0.800000}%
\pgfsetstrokecolor{currentstroke}%
\pgfsetdash{}{0pt}%
\pgfpathmoveto{\pgfqpoint{8.553430in}{0.464855in}}%
\pgfpathlineto{\pgfqpoint{8.553430in}{3.775855in}}%
\pgfusepath{stroke}%
\end{pgfscope}%
\begin{pgfscope}%
\definecolor{textcolor}{rgb}{0.150000,0.150000,0.150000}%
\pgfsetstrokecolor{textcolor}%
\pgfsetfillcolor{textcolor}%
\pgftext[x=8.553430in,y=0.416243in,,top]{\color{textcolor}{\sffamily\fontsize{10.000000}{12.000000}\selectfont\catcode`\^=\active\def^{\ifmmode\sp\else\^{}\fi}\catcode`\%=\active\def%{\%}oil production}}%
\end{pgfscope}%
\begin{pgfscope}%
\pgfpathrectangle{\pgfqpoint{5.504848in}{0.464855in}}{\pgfqpoint{4.051136in}{3.311000in}}%
\pgfusepath{clip}%
\pgfsetroundcap%
\pgfsetroundjoin%
\pgfsetlinewidth{0.803000pt}%
\definecolor{currentstroke}{rgb}{0.800000,0.800000,0.800000}%
\pgfsetstrokecolor{currentstroke}%
\pgfsetdash{}{0pt}%
\pgfpathmoveto{\pgfqpoint{5.504848in}{0.464855in}}%
\pgfpathlineto{\pgfqpoint{9.555984in}{0.464855in}}%
\pgfusepath{stroke}%
\end{pgfscope}%
\begin{pgfscope}%
\definecolor{textcolor}{rgb}{0.150000,0.150000,0.150000}%
\pgfsetstrokecolor{textcolor}%
\pgfsetfillcolor{textcolor}%
\pgftext[x=5.263162in, y=0.414297in, left, base]{\color{textcolor}{\sffamily\fontsize{10.000000}{12.000000}\selectfont\catcode`\^=\active\def^{\ifmmode\sp\else\^{}\fi}\catcode`\%=\active\def%{\%}0.0}}%
\end{pgfscope}%
\begin{pgfscope}%
\pgfpathrectangle{\pgfqpoint{5.504848in}{0.464855in}}{\pgfqpoint{4.051136in}{3.311000in}}%
\pgfusepath{clip}%
\pgfsetroundcap%
\pgfsetroundjoin%
\pgfsetlinewidth{0.803000pt}%
\definecolor{currentstroke}{rgb}{0.800000,0.800000,0.800000}%
\pgfsetstrokecolor{currentstroke}%
\pgfsetdash{}{0pt}%
\pgfpathmoveto{\pgfqpoint{5.504848in}{1.095521in}}%
\pgfpathlineto{\pgfqpoint{9.555984in}{1.095521in}}%
\pgfusepath{stroke}%
\end{pgfscope}%
\begin{pgfscope}%
\definecolor{textcolor}{rgb}{0.150000,0.150000,0.150000}%
\pgfsetstrokecolor{textcolor}%
\pgfsetfillcolor{textcolor}%
\pgftext[x=5.263162in, y=1.044964in, left, base]{\color{textcolor}{\sffamily\fontsize{10.000000}{12.000000}\selectfont\catcode`\^=\active\def^{\ifmmode\sp\else\^{}\fi}\catcode`\%=\active\def%{\%}0.2}}%
\end{pgfscope}%
\begin{pgfscope}%
\pgfpathrectangle{\pgfqpoint{5.504848in}{0.464855in}}{\pgfqpoint{4.051136in}{3.311000in}}%
\pgfusepath{clip}%
\pgfsetroundcap%
\pgfsetroundjoin%
\pgfsetlinewidth{0.803000pt}%
\definecolor{currentstroke}{rgb}{0.800000,0.800000,0.800000}%
\pgfsetstrokecolor{currentstroke}%
\pgfsetdash{}{0pt}%
\pgfpathmoveto{\pgfqpoint{5.504848in}{1.726188in}}%
\pgfpathlineto{\pgfqpoint{9.555984in}{1.726188in}}%
\pgfusepath{stroke}%
\end{pgfscope}%
\begin{pgfscope}%
\definecolor{textcolor}{rgb}{0.150000,0.150000,0.150000}%
\pgfsetstrokecolor{textcolor}%
\pgfsetfillcolor{textcolor}%
\pgftext[x=5.263162in, y=1.675630in, left, base]{\color{textcolor}{\sffamily\fontsize{10.000000}{12.000000}\selectfont\catcode`\^=\active\def^{\ifmmode\sp\else\^{}\fi}\catcode`\%=\active\def%{\%}0.4}}%
\end{pgfscope}%
\begin{pgfscope}%
\pgfpathrectangle{\pgfqpoint{5.504848in}{0.464855in}}{\pgfqpoint{4.051136in}{3.311000in}}%
\pgfusepath{clip}%
\pgfsetroundcap%
\pgfsetroundjoin%
\pgfsetlinewidth{0.803000pt}%
\definecolor{currentstroke}{rgb}{0.800000,0.800000,0.800000}%
\pgfsetstrokecolor{currentstroke}%
\pgfsetdash{}{0pt}%
\pgfpathmoveto{\pgfqpoint{5.504848in}{2.356855in}}%
\pgfpathlineto{\pgfqpoint{9.555984in}{2.356855in}}%
\pgfusepath{stroke}%
\end{pgfscope}%
\begin{pgfscope}%
\definecolor{textcolor}{rgb}{0.150000,0.150000,0.150000}%
\pgfsetstrokecolor{textcolor}%
\pgfsetfillcolor{textcolor}%
\pgftext[x=5.263162in, y=2.306297in, left, base]{\color{textcolor}{\sffamily\fontsize{10.000000}{12.000000}\selectfont\catcode`\^=\active\def^{\ifmmode\sp\else\^{}\fi}\catcode`\%=\active\def%{\%}0.6}}%
\end{pgfscope}%
\begin{pgfscope}%
\pgfpathrectangle{\pgfqpoint{5.504848in}{0.464855in}}{\pgfqpoint{4.051136in}{3.311000in}}%
\pgfusepath{clip}%
\pgfsetroundcap%
\pgfsetroundjoin%
\pgfsetlinewidth{0.803000pt}%
\definecolor{currentstroke}{rgb}{0.800000,0.800000,0.800000}%
\pgfsetstrokecolor{currentstroke}%
\pgfsetdash{}{0pt}%
\pgfpathmoveto{\pgfqpoint{5.504848in}{2.987521in}}%
\pgfpathlineto{\pgfqpoint{9.555984in}{2.987521in}}%
\pgfusepath{stroke}%
\end{pgfscope}%
\begin{pgfscope}%
\definecolor{textcolor}{rgb}{0.150000,0.150000,0.150000}%
\pgfsetstrokecolor{textcolor}%
\pgfsetfillcolor{textcolor}%
\pgftext[x=5.263162in, y=2.936964in, left, base]{\color{textcolor}{\sffamily\fontsize{10.000000}{12.000000}\selectfont\catcode`\^=\active\def^{\ifmmode\sp\else\^{}\fi}\catcode`\%=\active\def%{\%}0.8}}%
\end{pgfscope}%
\begin{pgfscope}%
\pgfpathrectangle{\pgfqpoint{5.504848in}{0.464855in}}{\pgfqpoint{4.051136in}{3.311000in}}%
\pgfusepath{clip}%
\pgfsetroundcap%
\pgfsetroundjoin%
\pgfsetlinewidth{0.803000pt}%
\definecolor{currentstroke}{rgb}{0.800000,0.800000,0.800000}%
\pgfsetstrokecolor{currentstroke}%
\pgfsetdash{}{0pt}%
\pgfpathmoveto{\pgfqpoint{5.504848in}{3.618188in}}%
\pgfpathlineto{\pgfqpoint{9.555984in}{3.618188in}}%
\pgfusepath{stroke}%
\end{pgfscope}%
\begin{pgfscope}%
\definecolor{textcolor}{rgb}{0.150000,0.150000,0.150000}%
\pgfsetstrokecolor{textcolor}%
\pgfsetfillcolor{textcolor}%
\pgftext[x=5.263162in, y=3.567631in, left, base]{\color{textcolor}{\sffamily\fontsize{10.000000}{12.000000}\selectfont\catcode`\^=\active\def^{\ifmmode\sp\else\^{}\fi}\catcode`\%=\active\def%{\%}1.0}}%
\end{pgfscope}%
\begin{pgfscope}%
\definecolor{textcolor}{rgb}{0.150000,0.150000,0.150000}%
\pgfsetstrokecolor{textcolor}%
\pgfsetfillcolor{textcolor}%
\pgftext[x=5.207607in,y=2.120355in,,bottom,rotate=90.000000]{\color{textcolor}{\sffamily\fontsize{10.000000}{12.000000}\selectfont\catcode`\^=\active\def^{\ifmmode\sp\else\^{}\fi}\catcode`\%=\active\def%{\%}CG/EG / specified rate}}%
\end{pgfscope}%
\begin{pgfscope}%
\pgfpathrectangle{\pgfqpoint{5.504848in}{0.464855in}}{\pgfqpoint{4.051136in}{3.311000in}}%
\pgfusepath{clip}%
\pgfsetbuttcap%
\pgfsetmiterjoin%
\definecolor{currentfill}{rgb}{0.121569,0.466667,0.705882}%
\pgfsetfillcolor{currentfill}%
\pgfsetlinewidth{0.000000pt}%
\definecolor{currentstroke}{rgb}{0.000000,0.000000,0.000000}%
\pgfsetstrokecolor{currentstroke}%
\pgfsetstrokeopacity{0.000000}%
\pgfsetdash{}{0pt}%
\pgfpathmoveto{\pgfqpoint{5.688990in}{0.464855in}}%
\pgfpathlineto{\pgfqpoint{7.325813in}{0.464855in}}%
\pgfpathlineto{\pgfqpoint{7.325813in}{3.618168in}}%
\pgfpathlineto{\pgfqpoint{5.688990in}{3.618168in}}%
\pgfpathlineto{\pgfqpoint{5.688990in}{0.464855in}}%
\pgfpathclose%
\pgfusepath{fill}%
\end{pgfscope}%
\begin{pgfscope}%
\pgfpathrectangle{\pgfqpoint{5.504848in}{0.464855in}}{\pgfqpoint{4.051136in}{3.311000in}}%
\pgfusepath{clip}%
\pgfsetbuttcap%
\pgfsetmiterjoin%
\definecolor{currentfill}{rgb}{0.121569,0.466667,0.705882}%
\pgfsetfillcolor{currentfill}%
\pgfsetlinewidth{0.000000pt}%
\definecolor{currentstroke}{rgb}{0.000000,0.000000,0.000000}%
\pgfsetstrokecolor{currentstroke}%
\pgfsetstrokeopacity{0.000000}%
\pgfsetdash{}{0pt}%
\pgfpathmoveto{\pgfqpoint{7.735019in}{0.464855in}}%
\pgfpathlineto{\pgfqpoint{9.371842in}{0.464855in}}%
\pgfpathlineto{\pgfqpoint{9.371842in}{3.618188in}}%
\pgfpathlineto{\pgfqpoint{7.735019in}{3.618188in}}%
\pgfpathlineto{\pgfqpoint{7.735019in}{0.464855in}}%
\pgfpathclose%
\pgfusepath{fill}%
\end{pgfscope}%
\begin{pgfscope}%
\pgfpathrectangle{\pgfqpoint{5.504848in}{0.464855in}}{\pgfqpoint{4.051136in}{3.311000in}}%
\pgfusepath{clip}%
\pgfsetbuttcap%
\pgfsetroundjoin%
\pgfsetlinewidth{1.505625pt}%
\definecolor{currentstroke}{rgb}{0.000000,0.000000,0.000000}%
\pgfsetstrokecolor{currentstroke}%
\pgfsetdash{{5.550000pt}{2.400000pt}}{0.000000pt}%
\pgfpathmoveto{\pgfqpoint{5.504848in}{3.618188in}}%
\pgfpathlineto{\pgfqpoint{9.555984in}{3.618188in}}%
\pgfusepath{stroke}%
\end{pgfscope}%
\begin{pgfscope}%
\pgfsetrectcap%
\pgfsetmiterjoin%
\pgfsetlinewidth{1.003750pt}%
\definecolor{currentstroke}{rgb}{0.800000,0.800000,0.800000}%
\pgfsetstrokecolor{currentstroke}%
\pgfsetdash{}{0pt}%
\pgfpathmoveto{\pgfqpoint{5.504848in}{0.464855in}}%
\pgfpathlineto{\pgfqpoint{5.504848in}{3.775855in}}%
\pgfusepath{stroke}%
\end{pgfscope}%
\begin{pgfscope}%
\pgfsetrectcap%
\pgfsetmiterjoin%
\pgfsetlinewidth{1.003750pt}%
\definecolor{currentstroke}{rgb}{0.800000,0.800000,0.800000}%
\pgfsetstrokecolor{currentstroke}%
\pgfsetdash{}{0pt}%
\pgfpathmoveto{\pgfqpoint{9.555984in}{0.464855in}}%
\pgfpathlineto{\pgfqpoint{9.555984in}{3.775855in}}%
\pgfusepath{stroke}%
\end{pgfscope}%
\begin{pgfscope}%
\pgfsetrectcap%
\pgfsetmiterjoin%
\pgfsetlinewidth{1.003750pt}%
\definecolor{currentstroke}{rgb}{0.800000,0.800000,0.800000}%
\pgfsetstrokecolor{currentstroke}%
\pgfsetdash{}{0pt}%
\pgfpathmoveto{\pgfqpoint{5.504848in}{0.464855in}}%
\pgfpathlineto{\pgfqpoint{9.555984in}{0.464855in}}%
\pgfusepath{stroke}%
\end{pgfscope}%
\begin{pgfscope}%
\pgfsetrectcap%
\pgfsetmiterjoin%
\pgfsetlinewidth{1.003750pt}%
\definecolor{currentstroke}{rgb}{0.800000,0.800000,0.800000}%
\pgfsetstrokecolor{currentstroke}%
\pgfsetdash{}{0pt}%
\pgfpathmoveto{\pgfqpoint{5.504848in}{3.775855in}}%
\pgfpathlineto{\pgfqpoint{9.555984in}{3.775855in}}%
\pgfusepath{stroke}%
\end{pgfscope}%
\begin{pgfscope}%
\definecolor{textcolor}{rgb}{0.150000,0.150000,0.150000}%
\pgfsetstrokecolor{textcolor}%
\pgfsetfillcolor{textcolor}%
\pgftext[x=7.530416in,y=3.859188in,,base]{\color{textcolor}{\sffamily\fontsize{12.000000}{14.400000}\selectfont\catcode`\^=\active\def^{\ifmmode\sp\else\^{}\fi}\catcode`\%=\active\def%{\%}Day-1 control agreement}}%
\end{pgfscope}%
\begin{pgfscope}%
\pgfsetbuttcap%
\pgfsetroundjoin%
\pgfsetlinewidth{1.505625pt}%
\definecolor{currentstroke}{rgb}{0.000000,0.000000,0.000000}%
\pgfsetstrokecolor{currentstroke}%
\pgfsetdash{{5.550000pt}{2.400000pt}}{0.000000pt}%
\pgfpathmoveto{\pgfqpoint{5.629848in}{0.667695in}}%
\pgfpathlineto{\pgfqpoint{5.768737in}{0.667695in}}%
\pgfpathlineto{\pgfqpoint{5.907626in}{0.667695in}}%
\pgfusepath{stroke}%
\end{pgfscope}%
\begin{pgfscope}%
\definecolor{textcolor}{rgb}{0.150000,0.150000,0.150000}%
\pgfsetstrokecolor{textcolor}%
\pgfsetfillcolor{textcolor}%
\pgftext[x=6.018737in,y=0.619084in,left,base]{\color{textcolor}{\sffamily\fontsize{10.000000}{12.000000}\selectfont\catcode`\^=\active\def^{\ifmmode\sp\else\^{}\fi}\catcode`\%=\active\def%{\%}specified control}}%
\end{pgfscope}%
\begin{pgfscope}%
\definecolor{textcolor}{rgb}{0.150000,0.150000,0.150000}%
\pgfsetstrokecolor{textcolor}%
\pgfsetfillcolor{textcolor}%
\pgftext[x=4.955984in,y=4.205855in,,top]{\color{textcolor}{\sffamily\fontsize{12.000000}{14.400000}\selectfont\catcode`\^=\active\def^{\ifmmode\sp\else\^{}\fi}\catcode`\%=\active\def%{\%}SPE1 reference context only: the first like-for-like OPM report is day 31}}%
\end{pgfscope}%
\end{pgfpicture}%
\makeatother%
\endgroup%
%
	}
	\caption{SPE1 schedule context and day-one control agreement.  The pressure
		curves are the OPM schedule history; the day-one markers and normalized
		rate bars are the coupled Q2/CG--EG result.}
	\label{fig:spe1_one_day_opm_context}
\end{figure}

\FloatBarrier

\section{Open thermodynamic and AD derivative questions}
\label{sec:open_questions}

This section records issues that must be resolved before declaring the MOOSE
implementation thermodynamically faithful.

The automatic-differentiation strategy is no longer an open yes-or-no question:
MOOSE can use \texttt{ADDerivativeParsedMaterial} for scalar free-energy
derivatives, while tensor kinematics, phase closure algorithms, flux laws,
constraints, and weak forms still require explicit AD-aware implementation.
The open questions below specify where that division must be made precise.

\begin{enumerate}[leftmargin=*]
	\item \textbf{Primitive-variable Jacobian.}  Determine which primitive
	variables give robust AD derivatives for \(\bar p_f\),
	\(\mu_\xi^\alpha\), \(\bar p_E\), \(B\), and the phase fractions while
	respecting the volume and mass-fraction constraints.  This is open for the
	full theory, but the first closed problem uses solid-reference component
	storage variables, partition variables, and conversion potentials until a
	better primitive set is justified.

	\item \textbf{Held-fixed variables.}  For every derivative inherited from the
	Helmholtz free energy, record what is held fixed: entropy or temperature,
	component masses or mass fractions, intrinsic densities, solid distension,
	and skeleton deformation.  These choices control the stress and pressure
	terms in the Newton matrix.

	\item \textbf{Equivalent pressure and nonlinear Biot coefficient.}  Verify
	the AD path for \(\bar p_E\) and \(B\) against the theory manuscript,
	including the distinction between pore-volume filling and stress-generating
	skeleton deformation.

	\item \textbf{Reaction-rate derivatives.}  Decide whether kinetics depend on
	current volume fractions, Lagrangian component masses, chemical potentials,
	affinities, or activities.  The implementation should include the full AD
	derivative of the chosen closure.

	\item \textbf{Flux derivative completeness.}  Check derivatives of
	\(\mbf W_f\) with respect to \(\mbf F\), \(J\), phase fractions, intrinsic
	densities, potentials, pressure-like variables, reaction rates, and
	permeability.  The pulled-back permeability
	\(J\mbf F^{-1}\widetilde{\bs\kappa}_f\mbf F^{-T}\) is a likely source of
	missing terms, especially because
	\(\widetilde{\bs\kappa}_f\) itself contains
	\(\sum_\alpha\dot c_f^\alpha\mbf I\).

	\item \textbf{Fluid phase kinematics.}  Decide whether any later closure
	requires \(J_f\) or \(\mbf F_f\) as a finite-element variable, material
	internal variable, or user-object state.  The first review problem does not
	require those fields as residual unknowns; it discretizes solid-frame
	advection directly by \(\mbf W_f/(J\rho_f)\).  Separate \(J_f\) or
	\(\mbf F_f\) evolution should be added only if a later constitutive closure
	needs the full fluid phase map or phase-reference history variables.

	\item \textbf{Constraint enforcement.}  Compare variable elimination,
	multiplier residuals, and constrained material solves for the volume-fraction
	and mass-fraction constraints.  The selected method must preserve the
	thermodynamic interpretation of pressure and component potentials.

	\item \textbf{Consistent units.}  Audit whether each source term is per
	current mixture volume or per solid reference volume before multiplying by
	\(J\).  This is especially important for reaction rates, conversion
	momentum, and boundary fluxes.
\end{enumerate}

These questions should remain visible in the paper until each is answered by a
derivation, a source-code trace, and at least one verification case.


\bibliographystyle{plain}
\bibliography{../all}

\end{document}
